% F083｜12个函数逐式图解：一式一图，保留统一编号和一般条件。
\clearpage
\begin{explanationbox}
\textbf{\textcolor{CExplanation}{指数互换，差变号}}
% Formula-based samples; see F083 checks/build_tikz_diagrams.py.
\begin{center}\begin{tikzpicture}[font=\small,>=stealth,line join=round]
\draw[gray!40,rounded corners=3pt] (0,0) rectangle (14.2,8.4);
\node at (7.1,7.75) {奇1：$y=a^x-a^{-x}$};
\begin{scope}[shift={(4.400000,3.900000)},x=1.700000cm,y=0.450000cm]
\draw[->,gray] (-2,0)--(2,0) node[right] {$x$};
\draw[->,gray] (0,-5)--(0,5) node[above] {$y$};
\draw[gray] (-1,-.035/0.450000)--(-1,.035/0.450000);
\node[anchor=north] at (-1,-.07/0.450000) {$-1$};
\draw[gray] (1,-.035/0.450000)--(1,.035/0.450000);
\node[anchor=north] at (1,-.07/0.450000) {$1$};
\begin{scope}
\clip (-2,-5) rectangle (2,5);
\draw[blue,thick,solid] plot coordinates {(-2.0000000,-3.7500000) (-1.9933333,-3.7304008) (-1.9866667,-3.7108813) (-1.9800000,-3.6914409) (-1.9733333,-3.6720795) (-1.9666667,-3.6527964) (-1.9600000,-3.6335913) (-1.9533333,-3.6144639) (-1.9466667,-3.5954136) (-1.9400000,-3.5764400) (-1.9333333,-3.5575429) (-1.9266667,-3.5387217) (-1.9200000,-3.5199761) (-1.9133333,-3.5013056) (-1.9066667,-3.4827099) (-1.9000000,-3.4641886) (-1.8933333,-3.4457412) (-1.8866667,-3.4273675) (-1.8800000,-3.4090669) (-1.8733333,-3.3908391) (-1.8666667,-3.3726837) (-1.8600000,-3.3546003) (-1.8533333,-3.3365886) (-1.8466667,-3.3186481) (-1.8400000,-3.3007785) (-1.8333333,-3.2829794) (-1.8266667,-3.2652503) (-1.8200000,-3.2475910) (-1.8133333,-3.2300011) (-1.8066667,-3.2124801) (-1.8000000,-3.1950277) (-1.7933333,-3.1776435) (-1.7866667,-3.1603272) (-1.7800000,-3.1430783) (-1.7733333,-3.1258966) (-1.7666667,-3.1087817) (-1.7600000,-3.0917331) (-1.7533333,-3.0747505) (-1.7466667,-3.0578336) (-1.7400000,-3.0409820) (-1.7333333,-3.0241953) (-1.7266667,-3.0074732) (-1.7200000,-2.9908153) (-1.7133333,-2.9742213) (-1.7066667,-2.9576908) (-1.7000000,-2.9412235) (-1.6933333,-2.9248189) (-1.6866667,-2.9084769) (-1.6800000,-2.8921969) (-1.6733333,-2.8759787) (-1.6666667,-2.8598218) (-1.6600000,-2.8437261) (-1.6533333,-2.8276911) (-1.6466667,-2.8117164) (-1.6400000,-2.7958018) (-1.6333333,-2.7799469) (-1.6266667,-2.7641514) (-1.6200000,-2.7484149) (-1.6133333,-2.7327371) (-1.6066667,-2.7171176) (-1.6000000,-2.7015562) (-1.5933333,-2.6860524) (-1.5866667,-2.6706060) (-1.5800000,-2.6552166) (-1.5733333,-2.6398839) (-1.5666667,-2.6246076) (-1.5600000,-2.6093874) (-1.5533333,-2.5942228) (-1.5466667,-2.5791137) (-1.5400000,-2.5640596) (-1.5333333,-2.5490603) (-1.5266667,-2.5341154) (-1.5200000,-2.5192246) (-1.5133333,-2.5043876) (-1.5066667,-2.4896041) (-1.5000000,-2.4748737) (-1.4933333,-2.4601962) (-1.4866667,-2.4455713) (-1.4800000,-2.4309985) (-1.4733333,-2.4164777) (-1.4666667,-2.4020085) (-1.4600000,-2.3875905) (-1.4533333,-2.3732235) (-1.4466667,-2.3589073) (-1.4400000,-2.3446414) (-1.4333333,-2.3304255) (-1.4266667,-2.3162594) (-1.4200000,-2.3021428) (-1.4133333,-2.2880753) (-1.4066667,-2.2740567) (-1.4000000,-2.2600867) (-1.3933333,-2.2461649) (-1.3866667,-2.2322911) (-1.3800000,-2.2184649) (-1.3733333,-2.2046861) (-1.3666667,-2.1909544) (-1.3600000,-2.1772695) (-1.3533333,-2.1636311) (-1.3466667,-2.1500388) (-1.3400000,-2.1364925) (-1.3333333,-2.1229918) (-1.3266667,-2.1095365) (-1.3200000,-2.0961262) (-1.3133333,-2.0827606) (-1.3066667,-2.0694395) (-1.3000000,-2.0561626) (-1.2933333,-2.0429296) (-1.2866667,-2.0297403) (-1.2800000,-2.0165943) (-1.2733333,-2.0034913) (-1.2666667,-1.9904311) (-1.2600000,-1.9774134) (-1.2533333,-1.9644380) (-1.2466667,-1.9515045) (-1.2400000,-1.9386127) (-1.2333333,-1.9257622) (-1.2266667,-1.9129529) (-1.2200000,-1.9001845) (-1.2133333,-1.8874566) (-1.2066667,-1.8747690) (-1.2000000,-1.8621214) (-1.1933333,-1.8495136) (-1.1866667,-1.8369453) (-1.1800000,-1.8244163) (-1.1733333,-1.8119262) (-1.1666667,-1.7994747) (-1.1600000,-1.7870617) (-1.1533333,-1.7746869) (-1.1466667,-1.7623500) (-1.1400000,-1.7500507) (-1.1333333,-1.7377887) (-1.1266667,-1.7255639) (-1.1200000,-1.7133759) (-1.1133333,-1.7012245) (-1.1066667,-1.6891094) (-1.1000000,-1.6770304) (-1.0933333,-1.6649872) (-1.0866667,-1.6529796) (-1.0800000,-1.6410073) (-1.0733333,-1.6290700) (-1.0666667,-1.6171674) (-1.0600000,-1.6052995) (-1.0533333,-1.5934658) (-1.0466667,-1.5816661) (-1.0400000,-1.5699002) (-1.0333333,-1.5581678) (-1.0266667,-1.5464687) (-1.0200000,-1.5348026) (-1.0133333,-1.5231693) (-1.0066667,-1.5115685) (-1.0000000,-1.5000000) (-0.9933333,-1.4884635) (-0.9866667,-1.4769588) (-0.9800000,-1.4654857) (-0.9733333,-1.4540438) (-0.9666667,-1.4426330) (-0.9600000,-1.4312530) (-0.9533333,-1.4199035) (-0.9466667,-1.4085844) (-0.9400000,-1.3972954) (-0.9333333,-1.3860361) (-0.9266667,-1.3748065) (-0.9200000,-1.3636063) (-0.9133333,-1.3524351) (-0.9066667,-1.3412929) (-0.9000000,-1.3301793) (-0.8933333,-1.3190940) (-0.8866667,-1.3080370) (-0.8800000,-1.2970079) (-0.8733333,-1.2860064) (-0.8666667,-1.2750325) (-0.8600000,-1.2640858) (-0.8533333,-1.2531660) (-0.8466667,-1.2422730) (-0.8400000,-1.2314066) (-0.8333333,-1.2205664) (-0.8266667,-1.2097523) (-0.8200000,-1.1989640) (-0.8133333,-1.1882014) (-0.8066667,-1.1774641) (-0.8000000,-1.1667519) (-0.7933333,-1.1560647) (-0.7866667,-1.1454022) (-0.7800000,-1.1347641) (-0.7733333,-1.1241502) (-0.7666667,-1.1135604) (-0.7600000,-1.1029943) (-0.7533333,-1.0924518) (-0.7466667,-1.0819326) (-0.7400000,-1.0714365) (-0.7333333,-1.0609633) (-0.7266667,-1.0505127) (-0.7200000,-1.0400846) (-0.7133333,-1.0296787) (-0.7066667,-1.0192947) (-0.7000000,-1.0089326) (-0.6933333,-0.9985920) (-0.6866667,-0.9882727) (-0.6800000,-0.9779745) (-0.6733333,-0.9676972) (-0.6666667,-0.9574405) (-0.6600000,-0.9472043) (-0.6533333,-0.9369884) (-0.6466667,-0.9267924) (-0.6400000,-0.9166162) (-0.6333333,-0.9064596) (-0.6266667,-0.8963224) (-0.6200000,-0.8862043) (-0.6133333,-0.8761051) (-0.6066667,-0.8660246) (-0.6000000,-0.8559626) (-0.5933333,-0.8459189) (-0.5866667,-0.8358933) (-0.5800000,-0.8258855) (-0.5733333,-0.8158953) (-0.5666667,-0.8059226) (-0.5600000,-0.7959671) (-0.5533333,-0.7860285) (-0.5466667,-0.7761068) (-0.5400000,-0.7662016) (-0.5333333,-0.7563128) (-0.5266667,-0.7464401) (-0.5200000,-0.7365834) (-0.5133333,-0.7267424) (-0.5066667,-0.7169169) (-0.5000000,-0.7071068) (-0.4933333,-0.6973117) (-0.4866667,-0.6875315) (-0.4800000,-0.6777660) (-0.4733333,-0.6680150) (-0.4666667,-0.6582783) (-0.4600000,-0.6485556) (-0.4533333,-0.6388467) (-0.4466667,-0.6291515) (-0.4400000,-0.6194697) (-0.4333333,-0.6098012) (-0.4266667,-0.6001456) (-0.4200000,-0.5905029) (-0.4133333,-0.5808728) (-0.4066667,-0.5712551) (-0.4000000,-0.5616496) (-0.3933333,-0.5520561) (-0.3866667,-0.5424744) (-0.3800000,-0.5329043) (-0.3733333,-0.5233455) (-0.3666667,-0.5137979) (-0.3600000,-0.5042613) (-0.3533333,-0.4947355) (-0.3466667,-0.4852202) (-0.3400000,-0.4757153) (-0.3333333,-0.4662205) (-0.3266667,-0.4567357) (-0.3200000,-0.4472607) (-0.3133333,-0.4377952) (-0.3066667,-0.4283390) (-0.3000000,-0.4188920) (-0.2933333,-0.4094540) (-0.2866667,-0.4000246) (-0.2800000,-0.3906039) (-0.2733333,-0.3811914) (-0.2666667,-0.3717871) (-0.2600000,-0.3623908) (-0.2533333,-0.3530022) (-0.2466667,-0.3436211) (-0.2400000,-0.3342473) (-0.2333333,-0.3248807) (-0.2266667,-0.3155211) (-0.2200000,-0.3061682) (-0.2133333,-0.2968218) (-0.2066667,-0.2874817) (-0.2000000,-0.2781478) (-0.1933333,-0.2688198) (-0.1866667,-0.2594976) (-0.1800000,-0.2501809) (-0.1733333,-0.2408695) (-0.1666667,-0.2315633) (-0.1600000,-0.2222621) (-0.1533333,-0.2129656) (-0.1466667,-0.2036736) (-0.1400000,-0.1943860) (-0.1333333,-0.1851025) (-0.1266667,-0.1758230) (-0.1200000,-0.1665472) (-0.1133333,-0.1572750) (-0.1066667,-0.1480062) (-0.1000000,-0.1387405) (-0.0933333,-0.1294777) (-0.0866667,-0.1202178) (-0.0800000,-0.1109604) (-0.0733333,-0.1017054) (-0.0666667,-0.0924525) (-0.0600000,-0.0832016) (-0.0533333,-0.0739525) (-0.0466667,-0.0647050) (-0.0400000,-0.0554589) (-0.0333333,-0.0462139) (-0.0266667,-0.0369700) (-0.0200000,-0.0277268) (-0.0133333,-0.0184842) (-0.0066667,-0.0092420) (0.0000000,0.0000000) (0.0066667,0.0092420) (0.0133333,0.0184842) (0.0200000,0.0277268) (0.0266667,0.0369700) (0.0333333,0.0462139) (0.0400000,0.0554589) (0.0466667,0.0647050) (0.0533333,0.0739525) (0.0600000,0.0832016) (0.0666667,0.0924525) (0.0733333,0.1017054) (0.0800000,0.1109604) (0.0866667,0.1202178) (0.0933333,0.1294777) (0.1000000,0.1387405) (0.1066667,0.1480062) (0.1133333,0.1572750) (0.1200000,0.1665472) (0.1266667,0.1758230) (0.1333333,0.1851025) (0.1400000,0.1943860) (0.1466667,0.2036736) (0.1533333,0.2129656) (0.1600000,0.2222621) (0.1666667,0.2315633) (0.1733333,0.2408695) (0.1800000,0.2501809) (0.1866667,0.2594976) (0.1933333,0.2688198) (0.2000000,0.2781478) (0.2066667,0.2874817) (0.2133333,0.2968218) (0.2200000,0.3061682) (0.2266667,0.3155211) (0.2333333,0.3248807) (0.2400000,0.3342473) (0.2466667,0.3436211) (0.2533333,0.3530022) (0.2600000,0.3623908) (0.2666667,0.3717871) (0.2733333,0.3811914) (0.2800000,0.3906039) (0.2866667,0.4000246) (0.2933333,0.4094540) (0.3000000,0.4188920) (0.3066667,0.4283390) (0.3133333,0.4377952) (0.3200000,0.4472607) (0.3266667,0.4567357) (0.3333333,0.4662205) (0.3400000,0.4757153) (0.3466667,0.4852202) (0.3533333,0.4947355) (0.3600000,0.5042613) (0.3666667,0.5137979) (0.3733333,0.5233455) (0.3800000,0.5329043) (0.3866667,0.5424744) (0.3933333,0.5520561) (0.4000000,0.5616496) (0.4066667,0.5712551) (0.4133333,0.5808728) (0.4200000,0.5905029) (0.4266667,0.6001456) (0.4333333,0.6098012) (0.4400000,0.6194697) (0.4466667,0.6291515) (0.4533333,0.6388467) (0.4600000,0.6485556) (0.4666667,0.6582783) (0.4733333,0.6680150) (0.4800000,0.6777660) (0.4866667,0.6875315) (0.4933333,0.6973117) (0.5000000,0.7071068) (0.5066667,0.7169169) (0.5133333,0.7267424) (0.5200000,0.7365834) (0.5266667,0.7464401) (0.5333333,0.7563128) (0.5400000,0.7662016) (0.5466667,0.7761068) (0.5533333,0.7860285) (0.5600000,0.7959671) (0.5666667,0.8059226) (0.5733333,0.8158953) (0.5800000,0.8258855) (0.5866667,0.8358933) (0.5933333,0.8459189) (0.6000000,0.8559626) (0.6066667,0.8660246) (0.6133333,0.8761051) (0.6200000,0.8862043) (0.6266667,0.8963224) (0.6333333,0.9064596) (0.6400000,0.9166162) (0.6466667,0.9267924) (0.6533333,0.9369884) (0.6600000,0.9472043) (0.6666667,0.9574405) (0.6733333,0.9676972) (0.6800000,0.9779745) (0.6866667,0.9882727) (0.6933333,0.9985920) (0.7000000,1.0089326) (0.7066667,1.0192947) (0.7133333,1.0296787) (0.7200000,1.0400846) (0.7266667,1.0505127) (0.7333333,1.0609633) (0.7400000,1.0714365) (0.7466667,1.0819326) (0.7533333,1.0924518) (0.7600000,1.1029943) (0.7666667,1.1135604) (0.7733333,1.1241502) (0.7800000,1.1347641) (0.7866667,1.1454022) (0.7933333,1.1560647) (0.8000000,1.1667519) (0.8066667,1.1774641) (0.8133333,1.1882014) (0.8200000,1.1989640) (0.8266667,1.2097523) (0.8333333,1.2205664) (0.8400000,1.2314066) (0.8466667,1.2422730) (0.8533333,1.2531660) (0.8600000,1.2640858) (0.8666667,1.2750325) (0.8733333,1.2860064) (0.8800000,1.2970079) (0.8866667,1.3080370) (0.8933333,1.3190940) (0.9000000,1.3301793) (0.9066667,1.3412929) (0.9133333,1.3524351) (0.9200000,1.3636063) (0.9266667,1.3748065) (0.9333333,1.3860361) (0.9400000,1.3972954) (0.9466667,1.4085844) (0.9533333,1.4199035) (0.9600000,1.4312530) (0.9666667,1.4426330) (0.9733333,1.4540438) (0.9800000,1.4654857) (0.9866667,1.4769588) (0.9933333,1.4884635) (1.0000000,1.5000000) (1.0066667,1.5115685) (1.0133333,1.5231693) (1.0200000,1.5348026) (1.0266667,1.5464687) (1.0333333,1.5581678) (1.0400000,1.5699002) (1.0466667,1.5816661) (1.0533333,1.5934658) (1.0600000,1.6052995) (1.0666667,1.6171674) (1.0733333,1.6290700) (1.0800000,1.6410073) (1.0866667,1.6529796) (1.0933333,1.6649872) (1.1000000,1.6770304) (1.1066667,1.6891094) (1.1133333,1.7012245) (1.1200000,1.7133759) (1.1266667,1.7255639) (1.1333333,1.7377887) (1.1400000,1.7500507) (1.1466667,1.7623500) (1.1533333,1.7746869) (1.1600000,1.7870617) (1.1666667,1.7994747) (1.1733333,1.8119262) (1.1800000,1.8244163) (1.1866667,1.8369453) (1.1933333,1.8495136) (1.2000000,1.8621214) (1.2066667,1.8747690) (1.2133333,1.8874566) (1.2200000,1.9001845) (1.2266667,1.9129529) (1.2333333,1.9257622) (1.2400000,1.9386127) (1.2466667,1.9515045) (1.2533333,1.9644380) (1.2600000,1.9774134) (1.2666667,1.9904311) (1.2733333,2.0034913) (1.2800000,2.0165943) (1.2866667,2.0297403) (1.2933333,2.0429296) (1.3000000,2.0561626) (1.3066667,2.0694395) (1.3133333,2.0827606) (1.3200000,2.0961262) (1.3266667,2.1095365) (1.3333333,2.1229918) (1.3400000,2.1364925) (1.3466667,2.1500388) (1.3533333,2.1636311) (1.3600000,2.1772695) (1.3666667,2.1909544) (1.3733333,2.2046861) (1.3800000,2.2184649) (1.3866667,2.2322911) (1.3933333,2.2461649) (1.4000000,2.2600867) (1.4066667,2.2740567) (1.4133333,2.2880753) (1.4200000,2.3021428) (1.4266667,2.3162594) (1.4333333,2.3304255) (1.4400000,2.3446414) (1.4466667,2.3589073) (1.4533333,2.3732235) (1.4600000,2.3875905) (1.4666667,2.4020085) (1.4733333,2.4164777) (1.4800000,2.4309985) (1.4866667,2.4455713) (1.4933333,2.4601962) (1.5000000,2.4748737) (1.5066667,2.4896041) (1.5133333,2.5043876) (1.5200000,2.5192246) (1.5266667,2.5341154) (1.5333333,2.5490603) (1.5400000,2.5640596) (1.5466667,2.5791137) (1.5533333,2.5942228) (1.5600000,2.6093874) (1.5666667,2.6246076) (1.5733333,2.6398839) (1.5800000,2.6552166) (1.5866667,2.6706060) (1.5933333,2.6860524) (1.6000000,2.7015562) (1.6066667,2.7171176) (1.6133333,2.7327371) (1.6200000,2.7484149) (1.6266667,2.7641514) (1.6333333,2.7799469) (1.6400000,2.7958018) (1.6466667,2.8117164) (1.6533333,2.8276911) (1.6600000,2.8437261) (1.6666667,2.8598218) (1.6733333,2.8759787) (1.6800000,2.8921969) (1.6866667,2.9084769) (1.6933333,2.9248189) (1.7000000,2.9412235) (1.7066667,2.9576908) (1.7133333,2.9742213) (1.7200000,2.9908153) (1.7266667,3.0074732) (1.7333333,3.0241953) (1.7400000,3.0409820) (1.7466667,3.0578336) (1.7533333,3.0747505) (1.7600000,3.0917331) (1.7666667,3.1087817) (1.7733333,3.1258966) (1.7800000,3.1430783) (1.7866667,3.1603272) (1.7933333,3.1776435) (1.8000000,3.1950277) (1.8066667,3.2124801) (1.8133333,3.2300011) (1.8200000,3.2475910) (1.8266667,3.2652503) (1.8333333,3.2829794) (1.8400000,3.3007785) (1.8466667,3.3186481) (1.8533333,3.3365886) (1.8600000,3.3546003) (1.8666667,3.3726837) (1.8733333,3.3908391) (1.8800000,3.4090669) (1.8866667,3.4273675) (1.8933333,3.4457412) (1.9000000,3.4641886) (1.9066667,3.4827099) (1.9133333,3.5013056) (1.9200000,3.5199761) (1.9266667,3.5387217) (1.9333333,3.5575429) (1.9400000,3.5764400) (1.9466667,3.5954136) (1.9533333,3.6144639) (1.9600000,3.6335913) (1.9666667,3.6527964) (1.9733333,3.6720795) (1.9800000,3.6914409) (1.9866667,3.7108813) (1.9933333,3.7304008) (2.0000000,3.7500000)};
\draw[orange!85!black,thick,dashed] plot coordinates {(-2.0000000,3.7500000) (-1.9933333,3.7304008) (-1.9866667,3.7108813) (-1.9800000,3.6914409) (-1.9733333,3.6720795) (-1.9666667,3.6527964) (-1.9600000,3.6335913) (-1.9533333,3.6144639) (-1.9466667,3.5954136) (-1.9400000,3.5764400) (-1.9333333,3.5575429) (-1.9266667,3.5387217) (-1.9200000,3.5199761) (-1.9133333,3.5013056) (-1.9066667,3.4827099) (-1.9000000,3.4641886) (-1.8933333,3.4457412) (-1.8866667,3.4273675) (-1.8800000,3.4090669) (-1.8733333,3.3908391) (-1.8666667,3.3726837) (-1.8600000,3.3546003) (-1.8533333,3.3365886) (-1.8466667,3.3186481) (-1.8400000,3.3007785) (-1.8333333,3.2829794) (-1.8266667,3.2652503) (-1.8200000,3.2475910) (-1.8133333,3.2300011) (-1.8066667,3.2124801) (-1.8000000,3.1950277) (-1.7933333,3.1776435) (-1.7866667,3.1603272) (-1.7800000,3.1430783) (-1.7733333,3.1258966) (-1.7666667,3.1087817) (-1.7600000,3.0917331) (-1.7533333,3.0747505) (-1.7466667,3.0578336) (-1.7400000,3.0409820) (-1.7333333,3.0241953) (-1.7266667,3.0074732) (-1.7200000,2.9908153) (-1.7133333,2.9742213) (-1.7066667,2.9576908) (-1.7000000,2.9412235) (-1.6933333,2.9248189) (-1.6866667,2.9084769) (-1.6800000,2.8921969) (-1.6733333,2.8759787) (-1.6666667,2.8598218) (-1.6600000,2.8437261) (-1.6533333,2.8276911) (-1.6466667,2.8117164) (-1.6400000,2.7958018) (-1.6333333,2.7799469) (-1.6266667,2.7641514) (-1.6200000,2.7484149) (-1.6133333,2.7327371) (-1.6066667,2.7171176) (-1.6000000,2.7015562) (-1.5933333,2.6860524) (-1.5866667,2.6706060) (-1.5800000,2.6552166) (-1.5733333,2.6398839) (-1.5666667,2.6246076) (-1.5600000,2.6093874) (-1.5533333,2.5942228) (-1.5466667,2.5791137) (-1.5400000,2.5640596) (-1.5333333,2.5490603) (-1.5266667,2.5341154) (-1.5200000,2.5192246) (-1.5133333,2.5043876) (-1.5066667,2.4896041) (-1.5000000,2.4748737) (-1.4933333,2.4601962) (-1.4866667,2.4455713) (-1.4800000,2.4309985) (-1.4733333,2.4164777) (-1.4666667,2.4020085) (-1.4600000,2.3875905) (-1.4533333,2.3732235) (-1.4466667,2.3589073) (-1.4400000,2.3446414) (-1.4333333,2.3304255) (-1.4266667,2.3162594) (-1.4200000,2.3021428) (-1.4133333,2.2880753) (-1.4066667,2.2740567) (-1.4000000,2.2600867) (-1.3933333,2.2461649) (-1.3866667,2.2322911) (-1.3800000,2.2184649) (-1.3733333,2.2046861) (-1.3666667,2.1909544) (-1.3600000,2.1772695) (-1.3533333,2.1636311) (-1.3466667,2.1500388) (-1.3400000,2.1364925) (-1.3333333,2.1229918) (-1.3266667,2.1095365) (-1.3200000,2.0961262) (-1.3133333,2.0827606) (-1.3066667,2.0694395) (-1.3000000,2.0561626) (-1.2933333,2.0429296) (-1.2866667,2.0297403) (-1.2800000,2.0165943) (-1.2733333,2.0034913) (-1.2666667,1.9904311) (-1.2600000,1.9774134) (-1.2533333,1.9644380) (-1.2466667,1.9515045) (-1.2400000,1.9386127) (-1.2333333,1.9257622) (-1.2266667,1.9129529) (-1.2200000,1.9001845) (-1.2133333,1.8874566) (-1.2066667,1.8747690) (-1.2000000,1.8621214) (-1.1933333,1.8495136) (-1.1866667,1.8369453) (-1.1800000,1.8244163) (-1.1733333,1.8119262) (-1.1666667,1.7994747) (-1.1600000,1.7870617) (-1.1533333,1.7746869) (-1.1466667,1.7623500) (-1.1400000,1.7500507) (-1.1333333,1.7377887) (-1.1266667,1.7255639) (-1.1200000,1.7133759) (-1.1133333,1.7012245) (-1.1066667,1.6891094) (-1.1000000,1.6770304) (-1.0933333,1.6649872) (-1.0866667,1.6529796) (-1.0800000,1.6410073) (-1.0733333,1.6290700) (-1.0666667,1.6171674) (-1.0600000,1.6052995) (-1.0533333,1.5934658) (-1.0466667,1.5816661) (-1.0400000,1.5699002) (-1.0333333,1.5581678) (-1.0266667,1.5464687) (-1.0200000,1.5348026) (-1.0133333,1.5231693) (-1.0066667,1.5115685) (-1.0000000,1.5000000) (-0.9933333,1.4884635) (-0.9866667,1.4769588) (-0.9800000,1.4654857) (-0.9733333,1.4540438) (-0.9666667,1.4426330) (-0.9600000,1.4312530) (-0.9533333,1.4199035) (-0.9466667,1.4085844) (-0.9400000,1.3972954) (-0.9333333,1.3860361) (-0.9266667,1.3748065) (-0.9200000,1.3636063) (-0.9133333,1.3524351) (-0.9066667,1.3412929) (-0.9000000,1.3301793) (-0.8933333,1.3190940) (-0.8866667,1.3080370) (-0.8800000,1.2970079) (-0.8733333,1.2860064) (-0.8666667,1.2750325) (-0.8600000,1.2640858) (-0.8533333,1.2531660) (-0.8466667,1.2422730) (-0.8400000,1.2314066) (-0.8333333,1.2205664) (-0.8266667,1.2097523) (-0.8200000,1.1989640) (-0.8133333,1.1882014) (-0.8066667,1.1774641) (-0.8000000,1.1667519) (-0.7933333,1.1560647) (-0.7866667,1.1454022) (-0.7800000,1.1347641) (-0.7733333,1.1241502) (-0.7666667,1.1135604) (-0.7600000,1.1029943) (-0.7533333,1.0924518) (-0.7466667,1.0819326) (-0.7400000,1.0714365) (-0.7333333,1.0609633) (-0.7266667,1.0505127) (-0.7200000,1.0400846) (-0.7133333,1.0296787) (-0.7066667,1.0192947) (-0.7000000,1.0089326) (-0.6933333,0.9985920) (-0.6866667,0.9882727) (-0.6800000,0.9779745) (-0.6733333,0.9676972) (-0.6666667,0.9574405) (-0.6600000,0.9472043) (-0.6533333,0.9369884) (-0.6466667,0.9267924) (-0.6400000,0.9166162) (-0.6333333,0.9064596) (-0.6266667,0.8963224) (-0.6200000,0.8862043) (-0.6133333,0.8761051) (-0.6066667,0.8660246) (-0.6000000,0.8559626) (-0.5933333,0.8459189) (-0.5866667,0.8358933) (-0.5800000,0.8258855) (-0.5733333,0.8158953) (-0.5666667,0.8059226) (-0.5600000,0.7959671) (-0.5533333,0.7860285) (-0.5466667,0.7761068) (-0.5400000,0.7662016) (-0.5333333,0.7563128) (-0.5266667,0.7464401) (-0.5200000,0.7365834) (-0.5133333,0.7267424) (-0.5066667,0.7169169) (-0.5000000,0.7071068) (-0.4933333,0.6973117) (-0.4866667,0.6875315) (-0.4800000,0.6777660) (-0.4733333,0.6680150) (-0.4666667,0.6582783) (-0.4600000,0.6485556) (-0.4533333,0.6388467) (-0.4466667,0.6291515) (-0.4400000,0.6194697) (-0.4333333,0.6098012) (-0.4266667,0.6001456) (-0.4200000,0.5905029) (-0.4133333,0.5808728) (-0.4066667,0.5712551) (-0.4000000,0.5616496) (-0.3933333,0.5520561) (-0.3866667,0.5424744) (-0.3800000,0.5329043) (-0.3733333,0.5233455) (-0.3666667,0.5137979) (-0.3600000,0.5042613) (-0.3533333,0.4947355) (-0.3466667,0.4852202) (-0.3400000,0.4757153) (-0.3333333,0.4662205) (-0.3266667,0.4567357) (-0.3200000,0.4472607) (-0.3133333,0.4377952) (-0.3066667,0.4283390) (-0.3000000,0.4188920) (-0.2933333,0.4094540) (-0.2866667,0.4000246) (-0.2800000,0.3906039) (-0.2733333,0.3811914) (-0.2666667,0.3717871) (-0.2600000,0.3623908) (-0.2533333,0.3530022) (-0.2466667,0.3436211) (-0.2400000,0.3342473) (-0.2333333,0.3248807) (-0.2266667,0.3155211) (-0.2200000,0.3061682) (-0.2133333,0.2968218) (-0.2066667,0.2874817) (-0.2000000,0.2781478) (-0.1933333,0.2688198) (-0.1866667,0.2594976) (-0.1800000,0.2501809) (-0.1733333,0.2408695) (-0.1666667,0.2315633) (-0.1600000,0.2222621) (-0.1533333,0.2129656) (-0.1466667,0.2036736) (-0.1400000,0.1943860) (-0.1333333,0.1851025) (-0.1266667,0.1758230) (-0.1200000,0.1665472) (-0.1133333,0.1572750) (-0.1066667,0.1480062) (-0.1000000,0.1387405) (-0.0933333,0.1294777) (-0.0866667,0.1202178) (-0.0800000,0.1109604) (-0.0733333,0.1017054) (-0.0666667,0.0924525) (-0.0600000,0.0832016) (-0.0533333,0.0739525) (-0.0466667,0.0647050) (-0.0400000,0.0554589) (-0.0333333,0.0462139) (-0.0266667,0.0369700) (-0.0200000,0.0277268) (-0.0133333,0.0184842) (-0.0066667,0.0092420) (0.0000000,0.0000000) (0.0066667,-0.0092420) (0.0133333,-0.0184842) (0.0200000,-0.0277268) (0.0266667,-0.0369700) (0.0333333,-0.0462139) (0.0400000,-0.0554589) (0.0466667,-0.0647050) (0.0533333,-0.0739525) (0.0600000,-0.0832016) (0.0666667,-0.0924525) (0.0733333,-0.1017054) (0.0800000,-0.1109604) (0.0866667,-0.1202178) (0.0933333,-0.1294777) (0.1000000,-0.1387405) (0.1066667,-0.1480062) (0.1133333,-0.1572750) (0.1200000,-0.1665472) (0.1266667,-0.1758230) (0.1333333,-0.1851025) (0.1400000,-0.1943860) (0.1466667,-0.2036736) (0.1533333,-0.2129656) (0.1600000,-0.2222621) (0.1666667,-0.2315633) (0.1733333,-0.2408695) (0.1800000,-0.2501809) (0.1866667,-0.2594976) (0.1933333,-0.2688198) (0.2000000,-0.2781478) (0.2066667,-0.2874817) (0.2133333,-0.2968218) (0.2200000,-0.3061682) (0.2266667,-0.3155211) (0.2333333,-0.3248807) (0.2400000,-0.3342473) (0.2466667,-0.3436211) (0.2533333,-0.3530022) (0.2600000,-0.3623908) (0.2666667,-0.3717871) (0.2733333,-0.3811914) (0.2800000,-0.3906039) (0.2866667,-0.4000246) (0.2933333,-0.4094540) (0.3000000,-0.4188920) (0.3066667,-0.4283390) (0.3133333,-0.4377952) (0.3200000,-0.4472607) (0.3266667,-0.4567357) (0.3333333,-0.4662205) (0.3400000,-0.4757153) (0.3466667,-0.4852202) (0.3533333,-0.4947355) (0.3600000,-0.5042613) (0.3666667,-0.5137979) (0.3733333,-0.5233455) (0.3800000,-0.5329043) (0.3866667,-0.5424744) (0.3933333,-0.5520561) (0.4000000,-0.5616496) (0.4066667,-0.5712551) (0.4133333,-0.5808728) (0.4200000,-0.5905029) (0.4266667,-0.6001456) (0.4333333,-0.6098012) (0.4400000,-0.6194697) (0.4466667,-0.6291515) (0.4533333,-0.6388467) (0.4600000,-0.6485556) (0.4666667,-0.6582783) (0.4733333,-0.6680150) (0.4800000,-0.6777660) (0.4866667,-0.6875315) (0.4933333,-0.6973117) (0.5000000,-0.7071068) (0.5066667,-0.7169169) (0.5133333,-0.7267424) (0.5200000,-0.7365834) (0.5266667,-0.7464401) (0.5333333,-0.7563128) (0.5400000,-0.7662016) (0.5466667,-0.7761068) (0.5533333,-0.7860285) (0.5600000,-0.7959671) (0.5666667,-0.8059226) (0.5733333,-0.8158953) (0.5800000,-0.8258855) (0.5866667,-0.8358933) (0.5933333,-0.8459189) (0.6000000,-0.8559626) (0.6066667,-0.8660246) (0.6133333,-0.8761051) (0.6200000,-0.8862043) (0.6266667,-0.8963224) (0.6333333,-0.9064596) (0.6400000,-0.9166162) (0.6466667,-0.9267924) (0.6533333,-0.9369884) (0.6600000,-0.9472043) (0.6666667,-0.9574405) (0.6733333,-0.9676972) (0.6800000,-0.9779745) (0.6866667,-0.9882727) (0.6933333,-0.9985920) (0.7000000,-1.0089326) (0.7066667,-1.0192947) (0.7133333,-1.0296787) (0.7200000,-1.0400846) (0.7266667,-1.0505127) (0.7333333,-1.0609633) (0.7400000,-1.0714365) (0.7466667,-1.0819326) (0.7533333,-1.0924518) (0.7600000,-1.1029943) (0.7666667,-1.1135604) (0.7733333,-1.1241502) (0.7800000,-1.1347641) (0.7866667,-1.1454022) (0.7933333,-1.1560647) (0.8000000,-1.1667519) (0.8066667,-1.1774641) (0.8133333,-1.1882014) (0.8200000,-1.1989640) (0.8266667,-1.2097523) (0.8333333,-1.2205664) (0.8400000,-1.2314066) (0.8466667,-1.2422730) (0.8533333,-1.2531660) (0.8600000,-1.2640858) (0.8666667,-1.2750325) (0.8733333,-1.2860064) (0.8800000,-1.2970079) (0.8866667,-1.3080370) (0.8933333,-1.3190940) (0.9000000,-1.3301793) (0.9066667,-1.3412929) (0.9133333,-1.3524351) (0.9200000,-1.3636063) (0.9266667,-1.3748065) (0.9333333,-1.3860361) (0.9400000,-1.3972954) (0.9466667,-1.4085844) (0.9533333,-1.4199035) (0.9600000,-1.4312530) (0.9666667,-1.4426330) (0.9733333,-1.4540438) (0.9800000,-1.4654857) (0.9866667,-1.4769588) (0.9933333,-1.4884635) (1.0000000,-1.5000000) (1.0066667,-1.5115685) (1.0133333,-1.5231693) (1.0200000,-1.5348026) (1.0266667,-1.5464687) (1.0333333,-1.5581678) (1.0400000,-1.5699002) (1.0466667,-1.5816661) (1.0533333,-1.5934658) (1.0600000,-1.6052995) (1.0666667,-1.6171674) (1.0733333,-1.6290700) (1.0800000,-1.6410073) (1.0866667,-1.6529796) (1.0933333,-1.6649872) (1.1000000,-1.6770304) (1.1066667,-1.6891094) (1.1133333,-1.7012245) (1.1200000,-1.7133759) (1.1266667,-1.7255639) (1.1333333,-1.7377887) (1.1400000,-1.7500507) (1.1466667,-1.7623500) (1.1533333,-1.7746869) (1.1600000,-1.7870617) (1.1666667,-1.7994747) (1.1733333,-1.8119262) (1.1800000,-1.8244163) (1.1866667,-1.8369453) (1.1933333,-1.8495136) (1.2000000,-1.8621214) (1.2066667,-1.8747690) (1.2133333,-1.8874566) (1.2200000,-1.9001845) (1.2266667,-1.9129529) (1.2333333,-1.9257622) (1.2400000,-1.9386127) (1.2466667,-1.9515045) (1.2533333,-1.9644380) (1.2600000,-1.9774134) (1.2666667,-1.9904311) (1.2733333,-2.0034913) (1.2800000,-2.0165943) (1.2866667,-2.0297403) (1.2933333,-2.0429296) (1.3000000,-2.0561626) (1.3066667,-2.0694395) (1.3133333,-2.0827606) (1.3200000,-2.0961262) (1.3266667,-2.1095365) (1.3333333,-2.1229918) (1.3400000,-2.1364925) (1.3466667,-2.1500388) (1.3533333,-2.1636311) (1.3600000,-2.1772695) (1.3666667,-2.1909544) (1.3733333,-2.2046861) (1.3800000,-2.2184649) (1.3866667,-2.2322911) (1.3933333,-2.2461649) (1.4000000,-2.2600867) (1.4066667,-2.2740567) (1.4133333,-2.2880753) (1.4200000,-2.3021428) (1.4266667,-2.3162594) (1.4333333,-2.3304255) (1.4400000,-2.3446414) (1.4466667,-2.3589073) (1.4533333,-2.3732235) (1.4600000,-2.3875905) (1.4666667,-2.4020085) (1.4733333,-2.4164777) (1.4800000,-2.4309985) (1.4866667,-2.4455713) (1.4933333,-2.4601962) (1.5000000,-2.4748737) (1.5066667,-2.4896041) (1.5133333,-2.5043876) (1.5200000,-2.5192246) (1.5266667,-2.5341154) (1.5333333,-2.5490603) (1.5400000,-2.5640596) (1.5466667,-2.5791137) (1.5533333,-2.5942228) (1.5600000,-2.6093874) (1.5666667,-2.6246076) (1.5733333,-2.6398839) (1.5800000,-2.6552166) (1.5866667,-2.6706060) (1.5933333,-2.6860524) (1.6000000,-2.7015562) (1.6066667,-2.7171176) (1.6133333,-2.7327371) (1.6200000,-2.7484149) (1.6266667,-2.7641514) (1.6333333,-2.7799469) (1.6400000,-2.7958018) (1.6466667,-2.8117164) (1.6533333,-2.8276911) (1.6600000,-2.8437261) (1.6666667,-2.8598218) (1.6733333,-2.8759787) (1.6800000,-2.8921969) (1.6866667,-2.9084769) (1.6933333,-2.9248189) (1.7000000,-2.9412235) (1.7066667,-2.9576908) (1.7133333,-2.9742213) (1.7200000,-2.9908153) (1.7266667,-3.0074732) (1.7333333,-3.0241953) (1.7400000,-3.0409820) (1.7466667,-3.0578336) (1.7533333,-3.0747505) (1.7600000,-3.0917331) (1.7666667,-3.1087817) (1.7733333,-3.1258966) (1.7800000,-3.1430783) (1.7866667,-3.1603272) (1.7933333,-3.1776435) (1.8000000,-3.1950277) (1.8066667,-3.2124801) (1.8133333,-3.2300011) (1.8200000,-3.2475910) (1.8266667,-3.2652503) (1.8333333,-3.2829794) (1.8400000,-3.3007785) (1.8466667,-3.3186481) (1.8533333,-3.3365886) (1.8600000,-3.3546003) (1.8666667,-3.3726837) (1.8733333,-3.3908391) (1.8800000,-3.4090669) (1.8866667,-3.4273675) (1.8933333,-3.4457412) (1.9000000,-3.4641886) (1.9066667,-3.4827099) (1.9133333,-3.5013056) (1.9200000,-3.5199761) (1.9266667,-3.5387217) (1.9333333,-3.5575429) (1.9400000,-3.5764400) (1.9466667,-3.5954136) (1.9533333,-3.6144639) (1.9600000,-3.6335913) (1.9666667,-3.6527964) (1.9733333,-3.6720795) (1.9800000,-3.6914409) (1.9866667,-3.7108813) (1.9933333,-3.7304008) (2.0000000,-3.7500000)};
\draw[violet!75!black,densely dotted] (-1,-1.5000000)--(1,1.5000000);
\fill[violet!75!black] (-1,-1.5000000) circle(2.3pt);
\fill[violet!75!black] (1,1.5000000) circle(2.3pt);
\end{scope}
\end{scope}
\node[anchor=west,align=left,text width=4.35cm] at (8.9,4.4) {蓝实线：$a=2$\\[8pt]橙虚线：$a=\frac12$\\[14pt]两种底数下，\\点对均关于原点对称。};
\node[text width=12.8cm,align=center] at (7.1,.65) {$D=\mathbb R$；$f(0)=0$。指数互换后，差变号。};
\end{tikzpicture}\end{center}

\textbf{先看图：} 共同条件为 $a>0$ 且 $a\ne1$。先看蓝色曲线：$x=1$ 与 $x=-1$ 的值分别为 $\frac32$ 与 $-\frac32$；底数换为 $\frac12$ 后方向反转，但每条图像仍关于原点对称。

\textbf{再说明一般情形：}
\[
f(-x)=a^{-x}-a^x=-(a^x-a^{-x})=-f(x).
\]
紫色点对只帮助看见关系；上式才说明所有实数输入都满足奇函数定义。
\end{explanationbox}
\clearpage
\begin{explanationbox}
\textbf{\textcolor{CExplanation}{分母恒正，所以原点可以取}}
% Formula-based samples; see F083 checks/build_tikz_diagrams.py.
\begin{center}\begin{tikzpicture}[font=\small,>=stealth,line join=round]
\draw[gray!40,rounded corners=3pt] (0,0) rectangle (14.2,8.4);
\node at (7.1,7.75) {奇2：$y=\dfrac{a^x-1}{a^x+1}$};
\begin{scope}[shift={(4.400000,3.900000)},x=0.850000cm,y=1.730769cm]
\draw[gray,dotted] (-4,-1)--(4,-1);
\node[anchor=east] at (-4-.08/0.850000,-1) {$-1$};
\draw[gray,dotted] (-4,1)--(4,1);
\node[anchor=east] at (-4-.08/0.850000,1) {$1$};
\draw[->,gray] (-4,0)--(4,0) node[right] {$x$};
\draw[->,gray] (0,-1.3)--(0,1.3) node[above] {$y$};
\draw[gray] (-1,-.035/1.730769)--(-1,.035/1.730769);
\node[anchor=north] at (-1,-.07/1.730769) {$-1$};
\draw[gray] (1,-.035/1.730769)--(1,.035/1.730769);
\node[anchor=north] at (1,-.07/1.730769) {$1$};
\begin{scope}
\clip (-4,-1.3) rectangle (4,1.3);
\draw[blue,thick,solid] plot coordinates {(-4.0000000,-0.8823529) (-3.9866667,-0.8813254) (-3.9733333,-0.8802895) (-3.9600000,-0.8792451) (-3.9466667,-0.8781922) (-3.9333333,-0.8771307) (-3.9200000,-0.8760606) (-3.9066667,-0.8749818) (-3.8933333,-0.8738942) (-3.8800000,-0.8727978) (-3.8666667,-0.8716926) (-3.8533333,-0.8705783) (-3.8400000,-0.8694551) (-3.8266667,-0.8683228) (-3.8133333,-0.8671814) (-3.8000000,-0.8660308) (-3.7866667,-0.8648710) (-3.7733333,-0.8637019) (-3.7600000,-0.8625233) (-3.7466667,-0.8613354) (-3.7333333,-0.8601380) (-3.7200000,-0.8589310) (-3.7066667,-0.8577144) (-3.6933333,-0.8564881) (-3.6800000,-0.8552520) (-3.6666667,-0.8540062) (-3.6533333,-0.8527504) (-3.6400000,-0.8514848) (-3.6266667,-0.8502091) (-3.6133333,-0.8489234) (-3.6000000,-0.8476276) (-3.5866667,-0.8463215) (-3.5733333,-0.8450052) (-3.5600000,-0.8436786) (-3.5466667,-0.8423416) (-3.5333333,-0.8409942) (-3.5200000,-0.8396362) (-3.5066667,-0.8382677) (-3.4933333,-0.8368885) (-3.4800000,-0.8354986) (-3.4666667,-0.8340979) (-3.4533333,-0.8326864) (-3.4400000,-0.8312640) (-3.4266667,-0.8298306) (-3.4133333,-0.8283862) (-3.4000000,-0.8269307) (-3.3866667,-0.8254640) (-3.3733333,-0.8239861) (-3.3600000,-0.8224969) (-3.3466667,-0.8209963) (-3.3333333,-0.8194843) (-3.3200000,-0.8179608) (-3.3066667,-0.8164257) (-3.2933333,-0.8148791) (-3.2800000,-0.8133207) (-3.2666667,-0.8117505) (-3.2533333,-0.8101686) (-3.2400000,-0.8085747) (-3.2266667,-0.8069689) (-3.2133333,-0.8053511) (-3.2000000,-0.8037212) (-3.1866667,-0.8020791) (-3.1733333,-0.8004248) (-3.1600000,-0.7987583) (-3.1466667,-0.7970794) (-3.1333333,-0.7953880) (-3.1200000,-0.7936842) (-3.1066667,-0.7919679) (-3.0933333,-0.7902389) (-3.0800000,-0.7884973) (-3.0666667,-0.7867429) (-3.0533333,-0.7849758) (-3.0400000,-0.7831957) (-3.0266667,-0.7814028) (-3.0133333,-0.7795968) (-3.0000000,-0.7777778) (-2.9866667,-0.7759456) (-2.9733333,-0.7741003) (-2.9600000,-0.7722417) (-2.9466667,-0.7703698) (-2.9333333,-0.7684846) (-2.9200000,-0.7665859) (-2.9066667,-0.7646737) (-2.8933333,-0.7627479) (-2.8800000,-0.7608085) (-2.8666667,-0.7588554) (-2.8533333,-0.7568886) (-2.8400000,-0.7549080) (-2.8266667,-0.7529135) (-2.8133333,-0.7509051) (-2.8000000,-0.7488827) (-2.7866667,-0.7468462) (-2.7733333,-0.7447957) (-2.7600000,-0.7427309) (-2.7466667,-0.7406520) (-2.7333333,-0.7385588) (-2.7200000,-0.7364512) (-2.7066667,-0.7343293) (-2.6933333,-0.7321929) (-2.6800000,-0.7300420) (-2.6666667,-0.7278765) (-2.6533333,-0.7256964) (-2.6400000,-0.7235017) (-2.6266667,-0.7212922) (-2.6133333,-0.7190679) (-2.6000000,-0.7168288) (-2.5866667,-0.7145749) (-2.5733333,-0.7123060) (-2.5600000,-0.7100221) (-2.5466667,-0.7077232) (-2.5333333,-0.7054091) (-2.5200000,-0.7030800) (-2.5066667,-0.7007357) (-2.4933333,-0.6983761) (-2.4800000,-0.6960013) (-2.4666667,-0.6936111) (-2.4533333,-0.6912056) (-2.4400000,-0.6887846) (-2.4266667,-0.6863482) (-2.4133333,-0.6838963) (-2.4000000,-0.6814288) (-2.3866667,-0.6789458) (-2.3733333,-0.6764471) (-2.3600000,-0.6739327) (-2.3466667,-0.6714027) (-2.3333333,-0.6688569) (-2.3200000,-0.6662953) (-2.3066667,-0.6637178) (-2.2933333,-0.6611246) (-2.2800000,-0.6585154) (-2.2666667,-0.6558903) (-2.2533333,-0.6532492) (-2.2400000,-0.6505922) (-2.2266667,-0.6479191) (-2.2133333,-0.6452300) (-2.2000000,-0.6425248) (-2.1866667,-0.6398035) (-2.1733333,-0.6370660) (-2.1600000,-0.6343124) (-2.1466667,-0.6315425) (-2.1333333,-0.6287565) (-2.1200000,-0.6259542) (-2.1066667,-0.6231357) (-2.0933333,-0.6203009) (-2.0800000,-0.6174498) (-2.0666667,-0.6145824) (-2.0533333,-0.6116986) (-2.0400000,-0.6087985) (-2.0266667,-0.6058821) (-2.0133333,-0.6029492) (-2.0000000,-0.6000000) (-1.9866667,-0.5970344) (-1.9733333,-0.5940523) (-1.9600000,-0.5910539) (-1.9466667,-0.5880390) (-1.9333333,-0.5850077) (-1.9200000,-0.5819599) (-1.9066667,-0.5788958) (-1.8933333,-0.5758152) (-1.8800000,-0.5727181) (-1.8666667,-0.5696046) (-1.8533333,-0.5664747) (-1.8400000,-0.5633283) (-1.8266667,-0.5601656) (-1.8133333,-0.5569864) (-1.8000000,-0.5537908) (-1.7866667,-0.5505788) (-1.7733333,-0.5473504) (-1.7600000,-0.5441056) (-1.7466667,-0.5408445) (-1.7333333,-0.5375671) (-1.7200000,-0.5342733) (-1.7066667,-0.5309632) (-1.6933333,-0.5276369) (-1.6800000,-0.5242942) (-1.6666667,-0.5209354) (-1.6533333,-0.5175603) (-1.6400000,-0.5141691) (-1.6266667,-0.5107617) (-1.6133333,-0.5073381) (-1.6000000,-0.5038985) (-1.5866667,-0.5004428) (-1.5733333,-0.4969711) (-1.5600000,-0.4934835) (-1.5466667,-0.4899798) (-1.5333333,-0.4864603) (-1.5200000,-0.4829250) (-1.5066667,-0.4793738) (-1.4933333,-0.4758068) (-1.4800000,-0.4722241) (-1.4666667,-0.4686258) (-1.4533333,-0.4650118) (-1.4400000,-0.4613823) (-1.4266667,-0.4577372) (-1.4133333,-0.4540767) (-1.4000000,-0.4504009) (-1.3866667,-0.4467096) (-1.3733333,-0.4430031) (-1.3600000,-0.4392814) (-1.3466667,-0.4355446) (-1.3333333,-0.4317927) (-1.3200000,-0.4280258) (-1.3066667,-0.4242439) (-1.2933333,-0.4204472) (-1.2800000,-0.4166358) (-1.2666667,-0.4128096) (-1.2533333,-0.4089688) (-1.2400000,-0.4051134) (-1.2266667,-0.4012436) (-1.2133333,-0.3973594) (-1.2000000,-0.3934609) (-1.1866667,-0.3895482) (-1.1733333,-0.3856214) (-1.1600000,-0.3816806) (-1.1466667,-0.3777259) (-1.1333333,-0.3737573) (-1.1200000,-0.3697750) (-1.1066667,-0.3657790) (-1.0933333,-0.3617696) (-1.0800000,-0.3577467) (-1.0666667,-0.3537105) (-1.0533333,-0.3496610) (-1.0400000,-0.3455985) (-1.0266667,-0.3415229) (-1.0133333,-0.3374345) (-1.0000000,-0.3333333) (-0.9866667,-0.3292195) (-0.9733333,-0.3250931) (-0.9600000,-0.3209543) (-0.9466667,-0.3168032) (-0.9333333,-0.3126399) (-0.9200000,-0.3084646) (-0.9066667,-0.3042774) (-0.8933333,-0.3000784) (-0.8800000,-0.2958677) (-0.8666667,-0.2916455) (-0.8533333,-0.2874119) (-0.8400000,-0.2831670) (-0.8266667,-0.2789110) (-0.8133333,-0.2746440) (-0.8000000,-0.2703662) (-0.7866667,-0.2660777) (-0.7733333,-0.2617786) (-0.7600000,-0.2574691) (-0.7466667,-0.2531493) (-0.7333333,-0.2488195) (-0.7200000,-0.2444796) (-0.7066667,-0.2401299) (-0.6933333,-0.2357706) (-0.6800000,-0.2314018) (-0.6666667,-0.2270236) (-0.6533333,-0.2226362) (-0.6400000,-0.2182398) (-0.6266667,-0.2138345) (-0.6133333,-0.2094204) (-0.6000000,-0.2049979) (-0.5866667,-0.2005669) (-0.5733333,-0.1961278) (-0.5600000,-0.1916805) (-0.5466667,-0.1872254) (-0.5333333,-0.1827626) (-0.5200000,-0.1782922) (-0.5066667,-0.1738145) (-0.4933333,-0.1693295) (-0.4800000,-0.1648375) (-0.4666667,-0.1603387) (-0.4533333,-0.1558332) (-0.4400000,-0.1513213) (-0.4266667,-0.1468030) (-0.4133333,-0.1422785) (-0.4000000,-0.1377481) (-0.3866667,-0.1332120) (-0.3733333,-0.1286702) (-0.3600000,-0.1241231) (-0.3466667,-0.1195707) (-0.3333333,-0.1150133) (-0.3200000,-0.1104511) (-0.3066667,-0.1058842) (-0.2933333,-0.1013128) (-0.2800000,-0.0967371) (-0.2666667,-0.0921574) (-0.2533333,-0.0875737) (-0.2400000,-0.0829864) (-0.2266667,-0.0783955) (-0.2133333,-0.0738013) (-0.2000000,-0.0692039) (-0.1866667,-0.0646036) (-0.1733333,-0.0600006) (-0.1600000,-0.0553950) (-0.1466667,-0.0507871) (-0.1333333,-0.0461769) (-0.1200000,-0.0415649) (-0.1066667,-0.0369510) (-0.0933333,-0.0323356) (-0.0800000,-0.0277188) (-0.0666667,-0.0231008) (-0.0533333,-0.0184818) (-0.0400000,-0.0138621) (-0.0266667,-0.0092417) (-0.0133333,-0.0046209) (0.0000000,0.0000000) (0.0133333,0.0046209) (0.0266667,0.0092417) (0.0400000,0.0138621) (0.0533333,0.0184818) (0.0666667,0.0231008) (0.0800000,0.0277188) (0.0933333,0.0323356) (0.1066667,0.0369510) (0.1200000,0.0415649) (0.1333333,0.0461769) (0.1466667,0.0507871) (0.1600000,0.0553950) (0.1733333,0.0600006) (0.1866667,0.0646036) (0.2000000,0.0692039) (0.2133333,0.0738013) (0.2266667,0.0783955) (0.2400000,0.0829864) (0.2533333,0.0875737) (0.2666667,0.0921574) (0.2800000,0.0967371) (0.2933333,0.1013128) (0.3066667,0.1058842) (0.3200000,0.1104511) (0.3333333,0.1150133) (0.3466667,0.1195707) (0.3600000,0.1241231) (0.3733333,0.1286702) (0.3866667,0.1332120) (0.4000000,0.1377481) (0.4133333,0.1422785) (0.4266667,0.1468030) (0.4400000,0.1513213) (0.4533333,0.1558332) (0.4666667,0.1603387) (0.4800000,0.1648375) (0.4933333,0.1693295) (0.5066667,0.1738145) (0.5200000,0.1782922) (0.5333333,0.1827626) (0.5466667,0.1872254) (0.5600000,0.1916805) (0.5733333,0.1961278) (0.5866667,0.2005669) (0.6000000,0.2049979) (0.6133333,0.2094204) (0.6266667,0.2138345) (0.6400000,0.2182398) (0.6533333,0.2226362) (0.6666667,0.2270236) (0.6800000,0.2314018) (0.6933333,0.2357706) (0.7066667,0.2401299) (0.7200000,0.2444796) (0.7333333,0.2488195) (0.7466667,0.2531493) (0.7600000,0.2574691) (0.7733333,0.2617786) (0.7866667,0.2660777) (0.8000000,0.2703662) (0.8133333,0.2746440) (0.8266667,0.2789110) (0.8400000,0.2831670) (0.8533333,0.2874119) (0.8666667,0.2916455) (0.8800000,0.2958677) (0.8933333,0.3000784) (0.9066667,0.3042774) (0.9200000,0.3084646) (0.9333333,0.3126399) (0.9466667,0.3168032) (0.9600000,0.3209543) (0.9733333,0.3250931) (0.9866667,0.3292195) (1.0000000,0.3333333) (1.0133333,0.3374345) (1.0266667,0.3415229) (1.0400000,0.3455985) (1.0533333,0.3496610) (1.0666667,0.3537105) (1.0800000,0.3577467) (1.0933333,0.3617696) (1.1066667,0.3657790) (1.1200000,0.3697750) (1.1333333,0.3737573) (1.1466667,0.3777259) (1.1600000,0.3816806) (1.1733333,0.3856214) (1.1866667,0.3895482) (1.2000000,0.3934609) (1.2133333,0.3973594) (1.2266667,0.4012436) (1.2400000,0.4051134) (1.2533333,0.4089688) (1.2666667,0.4128096) (1.2800000,0.4166358) (1.2933333,0.4204472) (1.3066667,0.4242439) (1.3200000,0.4280258) (1.3333333,0.4317927) (1.3466667,0.4355446) (1.3600000,0.4392814) (1.3733333,0.4430031) (1.3866667,0.4467096) (1.4000000,0.4504009) (1.4133333,0.4540767) (1.4266667,0.4577372) (1.4400000,0.4613823) (1.4533333,0.4650118) (1.4666667,0.4686258) (1.4800000,0.4722241) (1.4933333,0.4758068) (1.5066667,0.4793738) (1.5200000,0.4829250) (1.5333333,0.4864603) (1.5466667,0.4899798) (1.5600000,0.4934835) (1.5733333,0.4969711) (1.5866667,0.5004428) (1.6000000,0.5038985) (1.6133333,0.5073381) (1.6266667,0.5107617) (1.6400000,0.5141691) (1.6533333,0.5175603) (1.6666667,0.5209354) (1.6800000,0.5242942) (1.6933333,0.5276369) (1.7066667,0.5309632) (1.7200000,0.5342733) (1.7333333,0.5375671) (1.7466667,0.5408445) (1.7600000,0.5441056) (1.7733333,0.5473504) (1.7866667,0.5505788) (1.8000000,0.5537908) (1.8133333,0.5569864) (1.8266667,0.5601656) (1.8400000,0.5633283) (1.8533333,0.5664747) (1.8666667,0.5696046) (1.8800000,0.5727181) (1.8933333,0.5758152) (1.9066667,0.5788958) (1.9200000,0.5819599) (1.9333333,0.5850077) (1.9466667,0.5880390) (1.9600000,0.5910539) (1.9733333,0.5940523) (1.9866667,0.5970344) (2.0000000,0.6000000) (2.0133333,0.6029492) (2.0266667,0.6058821) (2.0400000,0.6087985) (2.0533333,0.6116986) (2.0666667,0.6145824) (2.0800000,0.6174498) (2.0933333,0.6203009) (2.1066667,0.6231357) (2.1200000,0.6259542) (2.1333333,0.6287565) (2.1466667,0.6315425) (2.1600000,0.6343124) (2.1733333,0.6370660) (2.1866667,0.6398035) (2.2000000,0.6425248) (2.2133333,0.6452300) (2.2266667,0.6479191) (2.2400000,0.6505922) (2.2533333,0.6532492) (2.2666667,0.6558903) (2.2800000,0.6585154) (2.2933333,0.6611246) (2.3066667,0.6637178) (2.3200000,0.6662953) (2.3333333,0.6688569) (2.3466667,0.6714027) (2.3600000,0.6739327) (2.3733333,0.6764471) (2.3866667,0.6789458) (2.4000000,0.6814288) (2.4133333,0.6838963) (2.4266667,0.6863482) (2.4400000,0.6887846) (2.4533333,0.6912056) (2.4666667,0.6936111) (2.4800000,0.6960013) (2.4933333,0.6983761) (2.5066667,0.7007357) (2.5200000,0.7030800) (2.5333333,0.7054091) (2.5466667,0.7077232) (2.5600000,0.7100221) (2.5733333,0.7123060) (2.5866667,0.7145749) (2.6000000,0.7168288) (2.6133333,0.7190679) (2.6266667,0.7212922) (2.6400000,0.7235017) (2.6533333,0.7256964) (2.6666667,0.7278765) (2.6800000,0.7300420) (2.6933333,0.7321929) (2.7066667,0.7343293) (2.7200000,0.7364512) (2.7333333,0.7385588) (2.7466667,0.7406520) (2.7600000,0.7427309) (2.7733333,0.7447957) (2.7866667,0.7468462) (2.8000000,0.7488827) (2.8133333,0.7509051) (2.8266667,0.7529135) (2.8400000,0.7549080) (2.8533333,0.7568886) (2.8666667,0.7588554) (2.8800000,0.7608085) (2.8933333,0.7627479) (2.9066667,0.7646737) (2.9200000,0.7665859) (2.9333333,0.7684846) (2.9466667,0.7703698) (2.9600000,0.7722417) (2.9733333,0.7741003) (2.9866667,0.7759456) (3.0000000,0.7777778) (3.0133333,0.7795968) (3.0266667,0.7814028) (3.0400000,0.7831957) (3.0533333,0.7849758) (3.0666667,0.7867429) (3.0800000,0.7884973) (3.0933333,0.7902389) (3.1066667,0.7919679) (3.1200000,0.7936842) (3.1333333,0.7953880) (3.1466667,0.7970794) (3.1600000,0.7987583) (3.1733333,0.8004248) (3.1866667,0.8020791) (3.2000000,0.8037212) (3.2133333,0.8053511) (3.2266667,0.8069689) (3.2400000,0.8085747) (3.2533333,0.8101686) (3.2666667,0.8117505) (3.2800000,0.8133207) (3.2933333,0.8148791) (3.3066667,0.8164257) (3.3200000,0.8179608) (3.3333333,0.8194843) (3.3466667,0.8209963) (3.3600000,0.8224969) (3.3733333,0.8239861) (3.3866667,0.8254640) (3.4000000,0.8269307) (3.4133333,0.8283862) (3.4266667,0.8298306) (3.4400000,0.8312640) (3.4533333,0.8326864) (3.4666667,0.8340979) (3.4800000,0.8354986) (3.4933333,0.8368885) (3.5066667,0.8382677) (3.5200000,0.8396362) (3.5333333,0.8409942) (3.5466667,0.8423416) (3.5600000,0.8436786) (3.5733333,0.8450052) (3.5866667,0.8463215) (3.6000000,0.8476276) (3.6133333,0.8489234) (3.6266667,0.8502091) (3.6400000,0.8514848) (3.6533333,0.8527504) (3.6666667,0.8540062) (3.6800000,0.8552520) (3.6933333,0.8564881) (3.7066667,0.8577144) (3.7200000,0.8589310) (3.7333333,0.8601380) (3.7466667,0.8613354) (3.7600000,0.8625233) (3.7733333,0.8637019) (3.7866667,0.8648710) (3.8000000,0.8660308) (3.8133333,0.8671814) (3.8266667,0.8683228) (3.8400000,0.8694551) (3.8533333,0.8705783) (3.8666667,0.8716926) (3.8800000,0.8727978) (3.8933333,0.8738942) (3.9066667,0.8749818) (3.9200000,0.8760606) (3.9333333,0.8771307) (3.9466667,0.8781922) (3.9600000,0.8792451) (3.9733333,0.8802895) (3.9866667,0.8813254) (4.0000000,0.8823529)};
\draw[orange!85!black,thick,dashed] plot coordinates {(-4.0000000,0.8823529) (-3.9866667,0.8813254) (-3.9733333,0.8802895) (-3.9600000,0.8792451) (-3.9466667,0.8781922) (-3.9333333,0.8771307) (-3.9200000,0.8760606) (-3.9066667,0.8749818) (-3.8933333,0.8738942) (-3.8800000,0.8727978) (-3.8666667,0.8716926) (-3.8533333,0.8705783) (-3.8400000,0.8694551) (-3.8266667,0.8683228) (-3.8133333,0.8671814) (-3.8000000,0.8660308) (-3.7866667,0.8648710) (-3.7733333,0.8637019) (-3.7600000,0.8625233) (-3.7466667,0.8613354) (-3.7333333,0.8601380) (-3.7200000,0.8589310) (-3.7066667,0.8577144) (-3.6933333,0.8564881) (-3.6800000,0.8552520) (-3.6666667,0.8540062) (-3.6533333,0.8527504) (-3.6400000,0.8514848) (-3.6266667,0.8502091) (-3.6133333,0.8489234) (-3.6000000,0.8476276) (-3.5866667,0.8463215) (-3.5733333,0.8450052) (-3.5600000,0.8436786) (-3.5466667,0.8423416) (-3.5333333,0.8409942) (-3.5200000,0.8396362) (-3.5066667,0.8382677) (-3.4933333,0.8368885) (-3.4800000,0.8354986) (-3.4666667,0.8340979) (-3.4533333,0.8326864) (-3.4400000,0.8312640) (-3.4266667,0.8298306) (-3.4133333,0.8283862) (-3.4000000,0.8269307) (-3.3866667,0.8254640) (-3.3733333,0.8239861) (-3.3600000,0.8224969) (-3.3466667,0.8209963) (-3.3333333,0.8194843) (-3.3200000,0.8179608) (-3.3066667,0.8164257) (-3.2933333,0.8148791) (-3.2800000,0.8133207) (-3.2666667,0.8117505) (-3.2533333,0.8101686) (-3.2400000,0.8085747) (-3.2266667,0.8069689) (-3.2133333,0.8053511) (-3.2000000,0.8037212) (-3.1866667,0.8020791) (-3.1733333,0.8004248) (-3.1600000,0.7987583) (-3.1466667,0.7970794) (-3.1333333,0.7953880) (-3.1200000,0.7936842) (-3.1066667,0.7919679) (-3.0933333,0.7902389) (-3.0800000,0.7884973) (-3.0666667,0.7867429) (-3.0533333,0.7849758) (-3.0400000,0.7831957) (-3.0266667,0.7814028) (-3.0133333,0.7795968) (-3.0000000,0.7777778) (-2.9866667,0.7759456) (-2.9733333,0.7741003) (-2.9600000,0.7722417) (-2.9466667,0.7703698) (-2.9333333,0.7684846) (-2.9200000,0.7665859) (-2.9066667,0.7646737) (-2.8933333,0.7627479) (-2.8800000,0.7608085) (-2.8666667,0.7588554) (-2.8533333,0.7568886) (-2.8400000,0.7549080) (-2.8266667,0.7529135) (-2.8133333,0.7509051) (-2.8000000,0.7488827) (-2.7866667,0.7468462) (-2.7733333,0.7447957) (-2.7600000,0.7427309) (-2.7466667,0.7406520) (-2.7333333,0.7385588) (-2.7200000,0.7364512) (-2.7066667,0.7343293) (-2.6933333,0.7321929) (-2.6800000,0.7300420) (-2.6666667,0.7278765) (-2.6533333,0.7256964) (-2.6400000,0.7235017) (-2.6266667,0.7212922) (-2.6133333,0.7190679) (-2.6000000,0.7168288) (-2.5866667,0.7145749) (-2.5733333,0.7123060) (-2.5600000,0.7100221) (-2.5466667,0.7077232) (-2.5333333,0.7054091) (-2.5200000,0.7030800) (-2.5066667,0.7007357) (-2.4933333,0.6983761) (-2.4800000,0.6960013) (-2.4666667,0.6936111) (-2.4533333,0.6912056) (-2.4400000,0.6887846) (-2.4266667,0.6863482) (-2.4133333,0.6838963) (-2.4000000,0.6814288) (-2.3866667,0.6789458) (-2.3733333,0.6764471) (-2.3600000,0.6739327) (-2.3466667,0.6714027) (-2.3333333,0.6688569) (-2.3200000,0.6662953) (-2.3066667,0.6637178) (-2.2933333,0.6611246) (-2.2800000,0.6585154) (-2.2666667,0.6558903) (-2.2533333,0.6532492) (-2.2400000,0.6505922) (-2.2266667,0.6479191) (-2.2133333,0.6452300) (-2.2000000,0.6425248) (-2.1866667,0.6398035) (-2.1733333,0.6370660) (-2.1600000,0.6343124) (-2.1466667,0.6315425) (-2.1333333,0.6287565) (-2.1200000,0.6259542) (-2.1066667,0.6231357) (-2.0933333,0.6203009) (-2.0800000,0.6174498) (-2.0666667,0.6145824) (-2.0533333,0.6116986) (-2.0400000,0.6087985) (-2.0266667,0.6058821) (-2.0133333,0.6029492) (-2.0000000,0.6000000) (-1.9866667,0.5970344) (-1.9733333,0.5940523) (-1.9600000,0.5910539) (-1.9466667,0.5880390) (-1.9333333,0.5850077) (-1.9200000,0.5819599) (-1.9066667,0.5788958) (-1.8933333,0.5758152) (-1.8800000,0.5727181) (-1.8666667,0.5696046) (-1.8533333,0.5664747) (-1.8400000,0.5633283) (-1.8266667,0.5601656) (-1.8133333,0.5569864) (-1.8000000,0.5537908) (-1.7866667,0.5505788) (-1.7733333,0.5473504) (-1.7600000,0.5441056) (-1.7466667,0.5408445) (-1.7333333,0.5375671) (-1.7200000,0.5342733) (-1.7066667,0.5309632) (-1.6933333,0.5276369) (-1.6800000,0.5242942) (-1.6666667,0.5209354) (-1.6533333,0.5175603) (-1.6400000,0.5141691) (-1.6266667,0.5107617) (-1.6133333,0.5073381) (-1.6000000,0.5038985) (-1.5866667,0.5004428) (-1.5733333,0.4969711) (-1.5600000,0.4934835) (-1.5466667,0.4899798) (-1.5333333,0.4864603) (-1.5200000,0.4829250) (-1.5066667,0.4793738) (-1.4933333,0.4758068) (-1.4800000,0.4722241) (-1.4666667,0.4686258) (-1.4533333,0.4650118) (-1.4400000,0.4613823) (-1.4266667,0.4577372) (-1.4133333,0.4540767) (-1.4000000,0.4504009) (-1.3866667,0.4467096) (-1.3733333,0.4430031) (-1.3600000,0.4392814) (-1.3466667,0.4355446) (-1.3333333,0.4317927) (-1.3200000,0.4280258) (-1.3066667,0.4242439) (-1.2933333,0.4204472) (-1.2800000,0.4166358) (-1.2666667,0.4128096) (-1.2533333,0.4089688) (-1.2400000,0.4051134) (-1.2266667,0.4012436) (-1.2133333,0.3973594) (-1.2000000,0.3934609) (-1.1866667,0.3895482) (-1.1733333,0.3856214) (-1.1600000,0.3816806) (-1.1466667,0.3777259) (-1.1333333,0.3737573) (-1.1200000,0.3697750) (-1.1066667,0.3657790) (-1.0933333,0.3617696) (-1.0800000,0.3577467) (-1.0666667,0.3537105) (-1.0533333,0.3496610) (-1.0400000,0.3455985) (-1.0266667,0.3415229) (-1.0133333,0.3374345) (-1.0000000,0.3333333) (-0.9866667,0.3292195) (-0.9733333,0.3250931) (-0.9600000,0.3209543) (-0.9466667,0.3168032) (-0.9333333,0.3126399) (-0.9200000,0.3084646) (-0.9066667,0.3042774) (-0.8933333,0.3000784) (-0.8800000,0.2958677) (-0.8666667,0.2916455) (-0.8533333,0.2874119) (-0.8400000,0.2831670) (-0.8266667,0.2789110) (-0.8133333,0.2746440) (-0.8000000,0.2703662) (-0.7866667,0.2660777) (-0.7733333,0.2617786) (-0.7600000,0.2574691) (-0.7466667,0.2531493) (-0.7333333,0.2488195) (-0.7200000,0.2444796) (-0.7066667,0.2401299) (-0.6933333,0.2357706) (-0.6800000,0.2314018) (-0.6666667,0.2270236) (-0.6533333,0.2226362) (-0.6400000,0.2182398) (-0.6266667,0.2138345) (-0.6133333,0.2094204) (-0.6000000,0.2049979) (-0.5866667,0.2005669) (-0.5733333,0.1961278) (-0.5600000,0.1916805) (-0.5466667,0.1872254) (-0.5333333,0.1827626) (-0.5200000,0.1782922) (-0.5066667,0.1738145) (-0.4933333,0.1693295) (-0.4800000,0.1648375) (-0.4666667,0.1603387) (-0.4533333,0.1558332) (-0.4400000,0.1513213) (-0.4266667,0.1468030) (-0.4133333,0.1422785) (-0.4000000,0.1377481) (-0.3866667,0.1332120) (-0.3733333,0.1286702) (-0.3600000,0.1241231) (-0.3466667,0.1195707) (-0.3333333,0.1150133) (-0.3200000,0.1104511) (-0.3066667,0.1058842) (-0.2933333,0.1013128) (-0.2800000,0.0967371) (-0.2666667,0.0921574) (-0.2533333,0.0875737) (-0.2400000,0.0829864) (-0.2266667,0.0783955) (-0.2133333,0.0738013) (-0.2000000,0.0692039) (-0.1866667,0.0646036) (-0.1733333,0.0600006) (-0.1600000,0.0553950) (-0.1466667,0.0507871) (-0.1333333,0.0461769) (-0.1200000,0.0415649) (-0.1066667,0.0369510) (-0.0933333,0.0323356) (-0.0800000,0.0277188) (-0.0666667,0.0231008) (-0.0533333,0.0184818) (-0.0400000,0.0138621) (-0.0266667,0.0092417) (-0.0133333,0.0046209) (0.0000000,-0.0000000) (0.0133333,-0.0046209) (0.0266667,-0.0092417) (0.0400000,-0.0138621) (0.0533333,-0.0184818) (0.0666667,-0.0231008) (0.0800000,-0.0277188) (0.0933333,-0.0323356) (0.1066667,-0.0369510) (0.1200000,-0.0415649) (0.1333333,-0.0461769) (0.1466667,-0.0507871) (0.1600000,-0.0553950) (0.1733333,-0.0600006) (0.1866667,-0.0646036) (0.2000000,-0.0692039) (0.2133333,-0.0738013) (0.2266667,-0.0783955) (0.2400000,-0.0829864) (0.2533333,-0.0875737) (0.2666667,-0.0921574) (0.2800000,-0.0967371) (0.2933333,-0.1013128) (0.3066667,-0.1058842) (0.3200000,-0.1104511) (0.3333333,-0.1150133) (0.3466667,-0.1195707) (0.3600000,-0.1241231) (0.3733333,-0.1286702) (0.3866667,-0.1332120) (0.4000000,-0.1377481) (0.4133333,-0.1422785) (0.4266667,-0.1468030) (0.4400000,-0.1513213) (0.4533333,-0.1558332) (0.4666667,-0.1603387) (0.4800000,-0.1648375) (0.4933333,-0.1693295) (0.5066667,-0.1738145) (0.5200000,-0.1782922) (0.5333333,-0.1827626) (0.5466667,-0.1872254) (0.5600000,-0.1916805) (0.5733333,-0.1961278) (0.5866667,-0.2005669) (0.6000000,-0.2049979) (0.6133333,-0.2094204) (0.6266667,-0.2138345) (0.6400000,-0.2182398) (0.6533333,-0.2226362) (0.6666667,-0.2270236) (0.6800000,-0.2314018) (0.6933333,-0.2357706) (0.7066667,-0.2401299) (0.7200000,-0.2444796) (0.7333333,-0.2488195) (0.7466667,-0.2531493) (0.7600000,-0.2574691) (0.7733333,-0.2617786) (0.7866667,-0.2660777) (0.8000000,-0.2703662) (0.8133333,-0.2746440) (0.8266667,-0.2789110) (0.8400000,-0.2831670) (0.8533333,-0.2874119) (0.8666667,-0.2916455) (0.8800000,-0.2958677) (0.8933333,-0.3000784) (0.9066667,-0.3042774) (0.9200000,-0.3084646) (0.9333333,-0.3126399) (0.9466667,-0.3168032) (0.9600000,-0.3209543) (0.9733333,-0.3250931) (0.9866667,-0.3292195) (1.0000000,-0.3333333) (1.0133333,-0.3374345) (1.0266667,-0.3415229) (1.0400000,-0.3455985) (1.0533333,-0.3496610) (1.0666667,-0.3537105) (1.0800000,-0.3577467) (1.0933333,-0.3617696) (1.1066667,-0.3657790) (1.1200000,-0.3697750) (1.1333333,-0.3737573) (1.1466667,-0.3777259) (1.1600000,-0.3816806) (1.1733333,-0.3856214) (1.1866667,-0.3895482) (1.2000000,-0.3934609) (1.2133333,-0.3973594) (1.2266667,-0.4012436) (1.2400000,-0.4051134) (1.2533333,-0.4089688) (1.2666667,-0.4128096) (1.2800000,-0.4166358) (1.2933333,-0.4204472) (1.3066667,-0.4242439) (1.3200000,-0.4280258) (1.3333333,-0.4317927) (1.3466667,-0.4355446) (1.3600000,-0.4392814) (1.3733333,-0.4430031) (1.3866667,-0.4467096) (1.4000000,-0.4504009) (1.4133333,-0.4540767) (1.4266667,-0.4577372) (1.4400000,-0.4613823) (1.4533333,-0.4650118) (1.4666667,-0.4686258) (1.4800000,-0.4722241) (1.4933333,-0.4758068) (1.5066667,-0.4793738) (1.5200000,-0.4829250) (1.5333333,-0.4864603) (1.5466667,-0.4899798) (1.5600000,-0.4934835) (1.5733333,-0.4969711) (1.5866667,-0.5004428) (1.6000000,-0.5038985) (1.6133333,-0.5073381) (1.6266667,-0.5107617) (1.6400000,-0.5141691) (1.6533333,-0.5175603) (1.6666667,-0.5209354) (1.6800000,-0.5242942) (1.6933333,-0.5276369) (1.7066667,-0.5309632) (1.7200000,-0.5342733) (1.7333333,-0.5375671) (1.7466667,-0.5408445) (1.7600000,-0.5441056) (1.7733333,-0.5473504) (1.7866667,-0.5505788) (1.8000000,-0.5537908) (1.8133333,-0.5569864) (1.8266667,-0.5601656) (1.8400000,-0.5633283) (1.8533333,-0.5664747) (1.8666667,-0.5696046) (1.8800000,-0.5727181) (1.8933333,-0.5758152) (1.9066667,-0.5788958) (1.9200000,-0.5819599) (1.9333333,-0.5850077) (1.9466667,-0.5880390) (1.9600000,-0.5910539) (1.9733333,-0.5940523) (1.9866667,-0.5970344) (2.0000000,-0.6000000) (2.0133333,-0.6029492) (2.0266667,-0.6058821) (2.0400000,-0.6087985) (2.0533333,-0.6116986) (2.0666667,-0.6145824) (2.0800000,-0.6174498) (2.0933333,-0.6203009) (2.1066667,-0.6231357) (2.1200000,-0.6259542) (2.1333333,-0.6287565) (2.1466667,-0.6315425) (2.1600000,-0.6343124) (2.1733333,-0.6370660) (2.1866667,-0.6398035) (2.2000000,-0.6425248) (2.2133333,-0.6452300) (2.2266667,-0.6479191) (2.2400000,-0.6505922) (2.2533333,-0.6532492) (2.2666667,-0.6558903) (2.2800000,-0.6585154) (2.2933333,-0.6611246) (2.3066667,-0.6637178) (2.3200000,-0.6662953) (2.3333333,-0.6688569) (2.3466667,-0.6714027) (2.3600000,-0.6739327) (2.3733333,-0.6764471) (2.3866667,-0.6789458) (2.4000000,-0.6814288) (2.4133333,-0.6838963) (2.4266667,-0.6863482) (2.4400000,-0.6887846) (2.4533333,-0.6912056) (2.4666667,-0.6936111) (2.4800000,-0.6960013) (2.4933333,-0.6983761) (2.5066667,-0.7007357) (2.5200000,-0.7030800) (2.5333333,-0.7054091) (2.5466667,-0.7077232) (2.5600000,-0.7100221) (2.5733333,-0.7123060) (2.5866667,-0.7145749) (2.6000000,-0.7168288) (2.6133333,-0.7190679) (2.6266667,-0.7212922) (2.6400000,-0.7235017) (2.6533333,-0.7256964) (2.6666667,-0.7278765) (2.6800000,-0.7300420) (2.6933333,-0.7321929) (2.7066667,-0.7343293) (2.7200000,-0.7364512) (2.7333333,-0.7385588) (2.7466667,-0.7406520) (2.7600000,-0.7427309) (2.7733333,-0.7447957) (2.7866667,-0.7468462) (2.8000000,-0.7488827) (2.8133333,-0.7509051) (2.8266667,-0.7529135) (2.8400000,-0.7549080) (2.8533333,-0.7568886) (2.8666667,-0.7588554) (2.8800000,-0.7608085) (2.8933333,-0.7627479) (2.9066667,-0.7646737) (2.9200000,-0.7665859) (2.9333333,-0.7684846) (2.9466667,-0.7703698) (2.9600000,-0.7722417) (2.9733333,-0.7741003) (2.9866667,-0.7759456) (3.0000000,-0.7777778) (3.0133333,-0.7795968) (3.0266667,-0.7814028) (3.0400000,-0.7831957) (3.0533333,-0.7849758) (3.0666667,-0.7867429) (3.0800000,-0.7884973) (3.0933333,-0.7902389) (3.1066667,-0.7919679) (3.1200000,-0.7936842) (3.1333333,-0.7953880) (3.1466667,-0.7970794) (3.1600000,-0.7987583) (3.1733333,-0.8004248) (3.1866667,-0.8020791) (3.2000000,-0.8037212) (3.2133333,-0.8053511) (3.2266667,-0.8069689) (3.2400000,-0.8085747) (3.2533333,-0.8101686) (3.2666667,-0.8117505) (3.2800000,-0.8133207) (3.2933333,-0.8148791) (3.3066667,-0.8164257) (3.3200000,-0.8179608) (3.3333333,-0.8194843) (3.3466667,-0.8209963) (3.3600000,-0.8224969) (3.3733333,-0.8239861) (3.3866667,-0.8254640) (3.4000000,-0.8269307) (3.4133333,-0.8283862) (3.4266667,-0.8298306) (3.4400000,-0.8312640) (3.4533333,-0.8326864) (3.4666667,-0.8340979) (3.4800000,-0.8354986) (3.4933333,-0.8368885) (3.5066667,-0.8382677) (3.5200000,-0.8396362) (3.5333333,-0.8409942) (3.5466667,-0.8423416) (3.5600000,-0.8436786) (3.5733333,-0.8450052) (3.5866667,-0.8463215) (3.6000000,-0.8476276) (3.6133333,-0.8489234) (3.6266667,-0.8502091) (3.6400000,-0.8514848) (3.6533333,-0.8527504) (3.6666667,-0.8540062) (3.6800000,-0.8552520) (3.6933333,-0.8564881) (3.7066667,-0.8577144) (3.7200000,-0.8589310) (3.7333333,-0.8601380) (3.7466667,-0.8613354) (3.7600000,-0.8625233) (3.7733333,-0.8637019) (3.7866667,-0.8648710) (3.8000000,-0.8660308) (3.8133333,-0.8671814) (3.8266667,-0.8683228) (3.8400000,-0.8694551) (3.8533333,-0.8705783) (3.8666667,-0.8716926) (3.8800000,-0.8727978) (3.8933333,-0.8738942) (3.9066667,-0.8749818) (3.9200000,-0.8760606) (3.9333333,-0.8771307) (3.9466667,-0.8781922) (3.9600000,-0.8792451) (3.9733333,-0.8802895) (3.9866667,-0.8813254) (4.0000000,-0.8823529)};
\draw[violet!75!black,densely dotted] (-1,-0.3333333)--(1,0.3333333);
\fill[violet!75!black] (-1,-0.3333333) circle(2.3pt);
\fill[violet!75!black] (1,0.3333333) circle(2.3pt);
\end{scope}
\end{scope}
\node[anchor=west,align=left,text width=4.35cm] at (8.9,4.4) {蓝实线：$a=2$\\[8pt]橙虚线：$a=\frac12$\\[14pt]两图都经过原点。\\[8pt]靠近 $y=\pm1$，\\但取不到这两个值。};
\node[text width=12.8cm,align=center] at (7.1,.65) {$D=\mathbb R$；分母始终为正，$f(0)=0$。};
\end{tikzpicture}\end{center}

\textbf{先看图：} 图中两种底数的曲线都经过 $(0,0)$。因为 $a^x+1>0$，定义域为全体实数；分式的两个水平边界 $y=\pm1$ 取不到。

\textbf{再说明一般情形：}
\[
f(-x)=\frac{1-a^x}{1+a^x}=-f(x).
\]
先用 $a^{-x}=1/a^x$，再把分子、分母同乘 $a^x$。
\end{explanationbox}
\clearpage
\begin{explanationbox}
\textbf{\textcolor{CExplanation}{没有原点，仍是奇函数}}
% Formula-based samples; see F083 checks/build_tikz_diagrams.py.
\begin{center}\begin{tikzpicture}[font=\small,>=stealth,line join=round]
\draw[gray!40,rounded corners=3pt] (0,0) rectangle (14.2,8.4);
\node at (7.1,7.75) {奇3：$y=\dfrac{a^x+1}{a^x-1}$};
\begin{scope}[shift={(4.400000,3.900000)},x=0.850000cm,y=0.450000cm]
\draw[gray,dotted] (-4,-1)--(4,-1);
\node[anchor=east] at (-4-.08/0.850000,-1) {$-1$};
\draw[gray,dotted] (-4,1)--(4,1);
\node[anchor=east] at (-4-.08/0.850000,1) {$1$};
\draw[->,gray] (-4,0)--(4,0) node[right] {$x$};
\draw[->,gray] (0,-5)--(0,5) node[above] {$y$};
\draw[gray] (-1,-.035/0.450000)--(-1,.035/0.450000);
\node[anchor=north] at (-1,-.07/0.450000) {$-1$};
\draw[gray] (1,-.035/0.450000)--(1,.035/0.450000);
\node[anchor=north] at (1,-.07/0.450000) {$1$};
\begin{scope}
\clip (-4,-5) rectangle (4,5);
\draw[blue,thick,solid] plot coordinates {(-4.0000000,-1.1333333) (-3.9933333,-1.1339923) (-3.9866667,-1.1346547) (-3.9800000,-1.1353205) (-3.9733333,-1.1359899) (-3.9666667,-1.1366628) (-3.9600000,-1.1373393) (-3.9533333,-1.1380193) (-3.9466667,-1.1387029) (-3.9400000,-1.1393901) (-3.9333333,-1.1400809) (-3.9266667,-1.1407754) (-3.9200000,-1.1414735) (-3.9133333,-1.1421753) (-3.9066667,-1.1428809) (-3.9000000,-1.1435902) (-3.8933333,-1.1443032) (-3.8866667,-1.1450201) (-3.8800000,-1.1457407) (-3.8733333,-1.1464652) (-3.8666667,-1.1471935) (-3.8600000,-1.1479256) (-3.8533333,-1.1486617) (-3.8466667,-1.1494017) (-3.8400000,-1.1501456) (-3.8333333,-1.1508935) (-3.8266667,-1.1516454) (-3.8200000,-1.1524013) (-3.8133333,-1.1531612) (-3.8066667,-1.1539252) (-3.8000000,-1.1546933) (-3.7933333,-1.1554655) (-3.7866667,-1.1562418) (-3.7800000,-1.1570223) (-3.7733333,-1.1578069) (-3.7666667,-1.1585958) (-3.7600000,-1.1593889) (-3.7533333,-1.1601863) (-3.7466667,-1.1609879) (-3.7400000,-1.1617939) (-3.7333333,-1.1626042) (-3.7266667,-1.1634189) (-3.7200000,-1.1642379) (-3.7133333,-1.1650614) (-3.7066667,-1.1658893) (-3.7000000,-1.1667217) (-3.6933333,-1.1675586) (-3.6866667,-1.1684000) (-3.6800000,-1.1692460) (-3.6733333,-1.1700966) (-3.6666667,-1.1709517) (-3.6600000,-1.1718115) (-3.6533333,-1.1726760) (-3.6466667,-1.1735452) (-3.6400000,-1.1744191) (-3.6333333,-1.1752978) (-3.6266667,-1.1761812) (-3.6200000,-1.1770695) (-3.6133333,-1.1779626) (-3.6066667,-1.1788606) (-3.6000000,-1.1797634) (-3.5933333,-1.1806713) (-3.5866667,-1.1815841) (-3.5800000,-1.1825018) (-3.5733333,-1.1834246) (-3.5666667,-1.1843525) (-3.5600000,-1.1852855) (-3.5533333,-1.1862236) (-3.5466667,-1.1871668) (-3.5400000,-1.1881153) (-3.5333333,-1.1890689) (-3.5266667,-1.1900278) (-3.5200000,-1.1909920) (-3.5133333,-1.1919615) (-3.5066667,-1.1929364) (-3.5000000,-1.1939167) (-3.4933333,-1.1949024) (-3.4866667,-1.1958935) (-3.4800000,-1.1968901) (-3.4733333,-1.1978923) (-3.4666667,-1.1989000) (-3.4600000,-1.1999133) (-3.4533333,-1.2009323) (-3.4466667,-1.2019569) (-3.4400000,-1.2029873) (-3.4333333,-1.2040233) (-3.4266667,-1.2050652) (-3.4200000,-1.2061129) (-3.4133333,-1.2071664) (-3.4066667,-1.2082258) (-3.4000000,-1.2092912) (-3.3933333,-1.2103625) (-3.3866667,-1.2114399) (-3.3800000,-1.2125233) (-3.3733333,-1.2136127) (-3.3666667,-1.2147083) (-3.3600000,-1.2158101) (-3.3533333,-1.2169181) (-3.3466667,-1.2180323) (-3.3400000,-1.2191528) (-3.3333333,-1.2202796) (-3.3266667,-1.2214129) (-3.3200000,-1.2225525) (-3.3133333,-1.2236986) (-3.3066667,-1.2248511) (-3.3000000,-1.2260103) (-3.2933333,-1.2271760) (-3.2866667,-1.2283483) (-3.2800000,-1.2295273) (-3.2733333,-1.2307131) (-3.2666667,-1.2319056) (-3.2600000,-1.2331049) (-3.2533333,-1.2343110) (-3.2466667,-1.2355241) (-3.2400000,-1.2367441) (-3.2333333,-1.2379711) (-3.2266667,-1.2392051) (-3.2200000,-1.2404462) (-3.2133333,-1.2416945) (-3.2066667,-1.2429499) (-3.2000000,-1.2442126) (-3.1933333,-1.2454825) (-3.1866667,-1.2467598) (-3.1800000,-1.2480444) (-3.1733333,-1.2493365) (-3.1666667,-1.2506361) (-3.1600000,-1.2519432) (-3.1533333,-1.2532579) (-3.1466667,-1.2545802) (-3.1400000,-1.2559102) (-3.1333333,-1.2572480) (-3.1266667,-1.2585935) (-3.1200000,-1.2599469) (-3.1133333,-1.2613082) (-3.1066667,-1.2626775) (-3.1000000,-1.2640547) (-3.0933333,-1.2654401) (-3.0866667,-1.2668335) (-3.0800000,-1.2682352) (-3.0733333,-1.2696450) (-3.0666667,-1.2710632) (-3.0600000,-1.2724897) (-3.0533333,-1.2739247) (-3.0466667,-1.2753681) (-3.0400000,-1.2768200) (-3.0333333,-1.2782806) (-3.0266667,-1.2797498) (-3.0200000,-1.2812277) (-3.0133333,-1.2827143) (-3.0066667,-1.2842099) (-3.0000000,-1.2857143) (-2.9933333,-1.2872277) (-2.9866667,-1.2887501) (-2.9800000,-1.2902816) (-2.9733333,-1.2918223) (-2.9666667,-1.2933721) (-2.9600000,-1.2949313) (-2.9533333,-1.2964999) (-2.9466667,-1.2980778) (-2.9400000,-1.2996653) (-2.9333333,-1.3012623) (-2.9266667,-1.3028690) (-2.9200000,-1.3044853) (-2.9133333,-1.3061114) (-2.9066667,-1.3077474) (-2.9000000,-1.3093933) (-2.8933333,-1.3110492) (-2.8866667,-1.3127151) (-2.8800000,-1.3143912) (-2.8733333,-1.3160775) (-2.8666667,-1.3177741) (-2.8600000,-1.3194810) (-2.8533333,-1.3211984) (-2.8466667,-1.3229263) (-2.8400000,-1.3246648) (-2.8333333,-1.3264139) (-2.8266667,-1.3281738) (-2.8200000,-1.3299446) (-2.8133333,-1.3317263) (-2.8066667,-1.3335189) (-2.8000000,-1.3353227) (-2.7933333,-1.3371376) (-2.7866667,-1.3389637) (-2.7800000,-1.3408012) (-2.7733333,-1.3426501) (-2.7666667,-1.3445106) (-2.7600000,-1.3463826) (-2.7533333,-1.3482663) (-2.7466667,-1.3501617) (-2.7400000,-1.3520691) (-2.7333333,-1.3539884) (-2.7266667,-1.3559197) (-2.7200000,-1.3578632) (-2.7133333,-1.3598189) (-2.7066667,-1.3617869) (-2.7000000,-1.3637674) (-2.6933333,-1.3657604) (-2.6866667,-1.3677659) (-2.6800000,-1.3697843) (-2.6733333,-1.3718154) (-2.6666667,-1.3738594) (-2.6600000,-1.3759165) (-2.6533333,-1.3779867) (-2.6466667,-1.3800701) (-2.6400000,-1.3821668) (-2.6333333,-1.3842770) (-2.6266667,-1.3864007) (-2.6200000,-1.3885381) (-2.6133333,-1.3906892) (-2.6066667,-1.3928542) (-2.6000000,-1.3950331) (-2.5933333,-1.3972262) (-2.5866667,-1.3994335) (-2.5800000,-1.4016551) (-2.5733333,-1.4038911) (-2.5666667,-1.4061417) (-2.5600000,-1.4084069) (-2.5533333,-1.4106869) (-2.5466667,-1.4129819) (-2.5400000,-1.4152919) (-2.5333333,-1.4176170) (-2.5266667,-1.4199574) (-2.5200000,-1.4223133) (-2.5133333,-1.4246846) (-2.5066667,-1.4270716) (-2.5000000,-1.4294745) (-2.4933333,-1.4318932) (-2.4866667,-1.4343280) (-2.4800000,-1.4367790) (-2.4733333,-1.4392463) (-2.4666667,-1.4417301) (-2.4600000,-1.4442305) (-2.4533333,-1.4467476) (-2.4466667,-1.4492816) (-2.4400000,-1.4518327) (-2.4333333,-1.4544009) (-2.4266667,-1.4569864) (-2.4200000,-1.4595894) (-2.4133333,-1.4622100) (-2.4066667,-1.4648484) (-2.4000000,-1.4675047) (-2.3933333,-1.4701791) (-2.3866667,-1.4728717) (-2.3800000,-1.4755827) (-2.3733333,-1.4783122) (-2.3666667,-1.4810605) (-2.3600000,-1.4838276) (-2.3533333,-1.4866138) (-2.3466667,-1.4894192) (-2.3400000,-1.4922439) (-2.3333333,-1.4950882) (-2.3266667,-1.4979522) (-2.3200000,-1.5008361) (-2.3133333,-1.5037401) (-2.3066667,-1.5066643) (-2.3000000,-1.5096090) (-2.2933333,-1.5125742) (-2.2866667,-1.5155603) (-2.2800000,-1.5185674) (-2.2733333,-1.5215956) (-2.2666667,-1.5246452) (-2.2600000,-1.5277163) (-2.2533333,-1.5308093) (-2.2466667,-1.5339241) (-2.2400000,-1.5370612) (-2.2333333,-1.5402206) (-2.2266667,-1.5434025) (-2.2200000,-1.5466072) (-2.2133333,-1.5498350) (-2.2066667,-1.5530859) (-2.2000000,-1.5563602) (-2.1933333,-1.5596582) (-2.1866667,-1.5629800) (-2.1800000,-1.5663259) (-2.1733333,-1.5696961) (-2.1666667,-1.5730908) (-2.1600000,-1.5765103) (-2.1533333,-1.5799548) (-2.1466667,-1.5834246) (-2.1400000,-1.5869198) (-2.1333333,-1.5904408) (-2.1266667,-1.5939877) (-2.1200000,-1.5975609) (-2.1133333,-1.6011605) (-2.1066667,-1.6047868) (-2.1000000,-1.6084402) (-2.0933333,-1.6121208) (-2.0866667,-1.6158289) (-2.0800000,-1.6195648) (-2.0733333,-1.6233288) (-2.0666667,-1.6271211) (-2.0600000,-1.6309421) (-2.0533333,-1.6347919) (-2.0466667,-1.6386710) (-2.0400000,-1.6425795) (-2.0333333,-1.6465178) (-2.0266667,-1.6504862) (-2.0200000,-1.6544850) (-2.0133333,-1.6585144) (-2.0066667,-1.6625749) (-2.0000000,-1.6666667) (-1.9933333,-1.6707901) (-1.9866667,-1.6749454) (-1.9800000,-1.6791331) (-1.9733333,-1.6833534) (-1.9666667,-1.6876066) (-1.9600000,-1.6918931) (-1.9533333,-1.6962133) (-1.9466667,-1.7005675) (-1.9400000,-1.7049560) (-1.9333333,-1.7093792) (-1.9266667,-1.7138375) (-1.9200000,-1.7183313) (-1.9133333,-1.7228609) (-1.9066667,-1.7274267) (-1.9000000,-1.7320291) (-1.8933333,-1.7366684) (-1.8866667,-1.7413452) (-1.8800000,-1.7460597) (-1.8733333,-1.7508124) (-1.8666667,-1.7556038) (-1.8600000,-1.7604341) (-1.8533333,-1.7653039) (-1.8466667,-1.7702136) (-1.8400000,-1.7751637) (-1.8333333,-1.7801545) (-1.8266667,-1.7851865) (-1.8200000,-1.7902602) (-1.8133333,-1.7953761) (-1.8066667,-1.8005345) (-1.8000000,-1.8057361) (-1.7933333,-1.8109813) (-1.7866667,-1.8162705) (-1.7800000,-1.8216043) (-1.7733333,-1.8269832) (-1.7666667,-1.8324078) (-1.7600000,-1.8378784) (-1.7533333,-1.8433957) (-1.7466667,-1.8489602) (-1.7400000,-1.8545724) (-1.7333333,-1.8602330) (-1.7266667,-1.8659424) (-1.7200000,-1.8717012) (-1.7133333,-1.8775101) (-1.7066667,-1.8833696) (-1.7000000,-1.8892803) (-1.6933333,-1.8952429) (-1.6866667,-1.9012579) (-1.6800000,-1.9073259) (-1.6733333,-1.9134477) (-1.6666667,-1.9196239) (-1.6600000,-1.9258551) (-1.6533333,-1.9321420) (-1.6466667,-1.9384853) (-1.6400000,-1.9448856) (-1.6333333,-1.9513438) (-1.6266667,-1.9578604) (-1.6200000,-1.9644362) (-1.6133333,-1.9710721) (-1.6066667,-1.9777686) (-1.6000000,-1.9845266) (-1.5933333,-1.9913469) (-1.5866667,-1.9982302) (-1.5800000,-2.0051774) (-1.5733333,-2.0121893) (-1.5666667,-2.0192666) (-1.5600000,-2.0264103) (-1.5533333,-2.0336213) (-1.5466667,-2.0409003) (-1.5400000,-2.0482482) (-1.5333333,-2.0556661) (-1.5266667,-2.0631547) (-1.5200000,-2.0707151) (-1.5133333,-2.0783482) (-1.5066667,-2.0860549) (-1.5000000,-2.0938363) (-1.4933333,-2.1016934) (-1.4866667,-2.1096271) (-1.4800000,-2.1176385) (-1.4733333,-2.1257288) (-1.4666667,-2.1338988) (-1.4600000,-2.1421499) (-1.4533333,-2.1504830) (-1.4466667,-2.1588994) (-1.4400000,-2.1674001) (-1.4333333,-2.1759864) (-1.4266667,-2.1846595) (-1.4200000,-2.1934206) (-1.4133333,-2.2022709) (-1.4066667,-2.2112118) (-1.4000000,-2.2202445) (-1.3933333,-2.2293703) (-1.3866667,-2.2385906) (-1.3800000,-2.2479069) (-1.3733333,-2.2573204) (-1.3666667,-2.2668325) (-1.3600000,-2.2764449) (-1.3533333,-2.2861589) (-1.3466667,-2.2959761) (-1.3400000,-2.3058980) (-1.3333333,-2.3159262) (-1.3266667,-2.3260623) (-1.3200000,-2.3363079) (-1.3133333,-2.3466647) (-1.3066667,-2.3571344) (-1.3000000,-2.3677189) (-1.2933333,-2.3784197) (-1.2866667,-2.3892388) (-1.2800000,-2.4001780) (-1.2733333,-2.4112393) (-1.2666667,-2.4224244) (-1.2600000,-2.4337355) (-1.2533333,-2.4451745) (-1.2466667,-2.4567435) (-1.2400000,-2.4684446) (-1.2333333,-2.4802799) (-1.2266667,-2.4922517) (-1.2200000,-2.5043621) (-1.2133333,-2.5166135) (-1.2066667,-2.5290081) (-1.2000000,-2.5415485) (-1.1933333,-2.5542371) (-1.1866667,-2.5670763) (-1.1800000,-2.5800687) (-1.1733333,-2.5932169) (-1.1666667,-2.6065236) (-1.1600000,-2.6199916) (-1.1533333,-2.6336237) (-1.1466667,-2.6474226) (-1.1400000,-2.6613914) (-1.1333333,-2.6755331) (-1.1266667,-2.6898507) (-1.1200000,-2.7043473) (-1.1133333,-2.7190263) (-1.1066667,-2.7338909) (-1.1000000,-2.7489444) (-1.0933333,-2.7641904) (-1.0866667,-2.7796324) (-1.0800000,-2.7952740) (-1.0733333,-2.8111190) (-1.0666667,-2.8271711) (-1.0600000,-2.8434343) (-1.0533333,-2.8599126) (-1.0466667,-2.8766102) (-1.0400000,-2.8935312) (-1.0333333,-2.9106799) (-1.0266667,-2.9280609) (-1.0200000,-2.9456787) (-1.0133333,-2.9635379) (-1.0066667,-2.9816433) (-1.0000000,-3.0000000) (-0.9933333,-3.0186129) (-0.9866667,-3.0374873) (-0.9800000,-3.0566284) (-0.9733333,-3.0760419) (-0.9666667,-3.0957332) (-0.9600000,-3.1157083) (-0.9533333,-3.1359731) (-0.9466667,-3.1565336) (-0.9400000,-3.1773963) (-0.9333333,-3.1985676) (-0.9266667,-3.2200541) (-0.9200000,-3.2418627) (-0.9133333,-3.2640005) (-0.9066667,-3.2864747) (-0.9000000,-3.3092929) (-0.8933333,-3.3324626) (-0.8866667,-3.3559920) (-0.8800000,-3.3798891) (-0.8733333,-3.4041623) (-0.8666667,-3.4288205) (-0.8600000,-3.4538724) (-0.8533333,-3.4793274) (-0.8466667,-3.5051949) (-0.8400000,-3.5314849) (-0.8333333,-3.5582074) (-0.8266667,-3.5853730) (-0.8200000,-3.6129923) (-0.8133333,-3.6410767) (-0.8066667,-3.6696377) (-0.8000000,-3.6986870) (-0.7933333,-3.7282372) (-0.7866667,-3.7583009) (-0.7800000,-3.7888912) (-0.7733333,-3.8200219) (-0.7666667,-3.8517070) (-0.7600000,-3.8839611) (-0.7533333,-3.9167994) (-0.7466667,-3.9502374) (-0.7400000,-3.9842914) (-0.7333333,-4.0189782) (-0.7266667,-4.0543153) (-0.7200000,-4.0903206) (-0.7133333,-4.1270130) (-0.7066667,-4.1644119) (-0.7000000,-4.2025375) (-0.6933333,-4.2414107) (-0.6866667,-4.2810533) (-0.6800000,-4.3214880) (-0.6733333,-4.3627384) (-0.6666667,-4.4048288) (-0.6600000,-4.4477847) (-0.6533333,-4.4916327) (-0.6466667,-4.5364003) (-0.6400000,-4.5821163) (-0.6333333,-4.6288107) (-0.6266667,-4.6765146) (-0.6200000,-4.7252607) (-0.6133333,-4.7750829) (-0.6066667,-4.8260167) (-0.6000000,-4.8780992) (-0.5933333,-4.9313691) (-0.5866667,-4.9858669)};
\draw[blue,thick,solid] plot coordinates {(0.5866667,4.9858669) (0.5933333,4.9313691) (0.6000000,4.8780992) (0.6066667,4.8260167) (0.6133333,4.7750829) (0.6200000,4.7252607) (0.6266667,4.6765146) (0.6333333,4.6288107) (0.6400000,4.5821163) (0.6466667,4.5364003) (0.6533333,4.4916327) (0.6600000,4.4477847) (0.6666667,4.4048288) (0.6733333,4.3627384) (0.6800000,4.3214880) (0.6866667,4.2810533) (0.6933333,4.2414107) (0.7000000,4.2025375) (0.7066667,4.1644119) (0.7133333,4.1270130) (0.7200000,4.0903206) (0.7266667,4.0543153) (0.7333333,4.0189782) (0.7400000,3.9842914) (0.7466667,3.9502374) (0.7533333,3.9167994) (0.7600000,3.8839611) (0.7666667,3.8517070) (0.7733333,3.8200219) (0.7800000,3.7888912) (0.7866667,3.7583009) (0.7933333,3.7282372) (0.8000000,3.6986870) (0.8066667,3.6696377) (0.8133333,3.6410767) (0.8200000,3.6129923) (0.8266667,3.5853730) (0.8333333,3.5582074) (0.8400000,3.5314849) (0.8466667,3.5051949) (0.8533333,3.4793274) (0.8600000,3.4538724) (0.8666667,3.4288205) (0.8733333,3.4041623) (0.8800000,3.3798891) (0.8866667,3.3559920) (0.8933333,3.3324626) (0.9000000,3.3092929) (0.9066667,3.2864747) (0.9133333,3.2640005) (0.9200000,3.2418627) (0.9266667,3.2200541) (0.9333333,3.1985676) (0.9400000,3.1773963) (0.9466667,3.1565336) (0.9533333,3.1359731) (0.9600000,3.1157083) (0.9666667,3.0957332) (0.9733333,3.0760419) (0.9800000,3.0566284) (0.9866667,3.0374873) (0.9933333,3.0186129) (1.0000000,3.0000000) (1.0066667,2.9816433) (1.0133333,2.9635379) (1.0200000,2.9456787) (1.0266667,2.9280609) (1.0333333,2.9106799) (1.0400000,2.8935312) (1.0466667,2.8766102) (1.0533333,2.8599126) (1.0600000,2.8434343) (1.0666667,2.8271711) (1.0733333,2.8111190) (1.0800000,2.7952740) (1.0866667,2.7796324) (1.0933333,2.7641904) (1.1000000,2.7489444) (1.1066667,2.7338909) (1.1133333,2.7190263) (1.1200000,2.7043473) (1.1266667,2.6898507) (1.1333333,2.6755331) (1.1400000,2.6613914) (1.1466667,2.6474226) (1.1533333,2.6336237) (1.1600000,2.6199916) (1.1666667,2.6065236) (1.1733333,2.5932169) (1.1800000,2.5800687) (1.1866667,2.5670763) (1.1933333,2.5542371) (1.2000000,2.5415485) (1.2066667,2.5290081) (1.2133333,2.5166135) (1.2200000,2.5043621) (1.2266667,2.4922517) (1.2333333,2.4802799) (1.2400000,2.4684446) (1.2466667,2.4567435) (1.2533333,2.4451745) (1.2600000,2.4337355) (1.2666667,2.4224244) (1.2733333,2.4112393) (1.2800000,2.4001780) (1.2866667,2.3892388) (1.2933333,2.3784197) (1.3000000,2.3677189) (1.3066667,2.3571344) (1.3133333,2.3466647) (1.3200000,2.3363079) (1.3266667,2.3260623) (1.3333333,2.3159262) (1.3400000,2.3058980) (1.3466667,2.2959761) (1.3533333,2.2861589) (1.3600000,2.2764449) (1.3666667,2.2668325) (1.3733333,2.2573204) (1.3800000,2.2479069) (1.3866667,2.2385906) (1.3933333,2.2293703) (1.4000000,2.2202445) (1.4066667,2.2112118) (1.4133333,2.2022709) (1.4200000,2.1934206) (1.4266667,2.1846595) (1.4333333,2.1759864) (1.4400000,2.1674001) (1.4466667,2.1588994) (1.4533333,2.1504830) (1.4600000,2.1421499) (1.4666667,2.1338988) (1.4733333,2.1257288) (1.4800000,2.1176385) (1.4866667,2.1096271) (1.4933333,2.1016934) (1.5000000,2.0938363) (1.5066667,2.0860549) (1.5133333,2.0783482) (1.5200000,2.0707151) (1.5266667,2.0631547) (1.5333333,2.0556661) (1.5400000,2.0482482) (1.5466667,2.0409003) (1.5533333,2.0336213) (1.5600000,2.0264103) (1.5666667,2.0192666) (1.5733333,2.0121893) (1.5800000,2.0051774) (1.5866667,1.9982302) (1.5933333,1.9913469) (1.6000000,1.9845266) (1.6066667,1.9777686) (1.6133333,1.9710721) (1.6200000,1.9644362) (1.6266667,1.9578604) (1.6333333,1.9513438) (1.6400000,1.9448856) (1.6466667,1.9384853) (1.6533333,1.9321420) (1.6600000,1.9258551) (1.6666667,1.9196239) (1.6733333,1.9134477) (1.6800000,1.9073259) (1.6866667,1.9012579) (1.6933333,1.8952429) (1.7000000,1.8892803) (1.7066667,1.8833696) (1.7133333,1.8775101) (1.7200000,1.8717012) (1.7266667,1.8659424) (1.7333333,1.8602330) (1.7400000,1.8545724) (1.7466667,1.8489602) (1.7533333,1.8433957) (1.7600000,1.8378784) (1.7666667,1.8324078) (1.7733333,1.8269832) (1.7800000,1.8216043) (1.7866667,1.8162705) (1.7933333,1.8109813) (1.8000000,1.8057361) (1.8066667,1.8005345) (1.8133333,1.7953761) (1.8200000,1.7902602) (1.8266667,1.7851865) (1.8333333,1.7801545) (1.8400000,1.7751637) (1.8466667,1.7702136) (1.8533333,1.7653039) (1.8600000,1.7604341) (1.8666667,1.7556038) (1.8733333,1.7508124) (1.8800000,1.7460597) (1.8866667,1.7413452) (1.8933333,1.7366684) (1.9000000,1.7320291) (1.9066667,1.7274267) (1.9133333,1.7228609) (1.9200000,1.7183313) (1.9266667,1.7138375) (1.9333333,1.7093792) (1.9400000,1.7049560) (1.9466667,1.7005675) (1.9533333,1.6962133) (1.9600000,1.6918931) (1.9666667,1.6876066) (1.9733333,1.6833534) (1.9800000,1.6791331) (1.9866667,1.6749454) (1.9933333,1.6707901) (2.0000000,1.6666667) (2.0066667,1.6625749) (2.0133333,1.6585144) (2.0200000,1.6544850) (2.0266667,1.6504862) (2.0333333,1.6465178) (2.0400000,1.6425795) (2.0466667,1.6386710) (2.0533333,1.6347919) (2.0600000,1.6309421) (2.0666667,1.6271211) (2.0733333,1.6233288) (2.0800000,1.6195648) (2.0866667,1.6158289) (2.0933333,1.6121208) (2.1000000,1.6084402) (2.1066667,1.6047868) (2.1133333,1.6011605) (2.1200000,1.5975609) (2.1266667,1.5939877) (2.1333333,1.5904408) (2.1400000,1.5869198) (2.1466667,1.5834246) (2.1533333,1.5799548) (2.1600000,1.5765103) (2.1666667,1.5730908) (2.1733333,1.5696961) (2.1800000,1.5663259) (2.1866667,1.5629800) (2.1933333,1.5596582) (2.2000000,1.5563602) (2.2066667,1.5530859) (2.2133333,1.5498350) (2.2200000,1.5466072) (2.2266667,1.5434025) (2.2333333,1.5402206) (2.2400000,1.5370612) (2.2466667,1.5339241) (2.2533333,1.5308093) (2.2600000,1.5277163) (2.2666667,1.5246452) (2.2733333,1.5215956) (2.2800000,1.5185674) (2.2866667,1.5155603) (2.2933333,1.5125742) (2.3000000,1.5096090) (2.3066667,1.5066643) (2.3133333,1.5037401) (2.3200000,1.5008361) (2.3266667,1.4979522) (2.3333333,1.4950882) (2.3400000,1.4922439) (2.3466667,1.4894192) (2.3533333,1.4866138) (2.3600000,1.4838276) (2.3666667,1.4810605) (2.3733333,1.4783122) (2.3800000,1.4755827) (2.3866667,1.4728717) (2.3933333,1.4701791) (2.4000000,1.4675047) (2.4066667,1.4648484) (2.4133333,1.4622100) (2.4200000,1.4595894) (2.4266667,1.4569864) (2.4333333,1.4544009) (2.4400000,1.4518327) (2.4466667,1.4492816) (2.4533333,1.4467476) (2.4600000,1.4442305) (2.4666667,1.4417301) (2.4733333,1.4392463) (2.4800000,1.4367790) (2.4866667,1.4343280) (2.4933333,1.4318932) (2.5000000,1.4294745) (2.5066667,1.4270716) (2.5133333,1.4246846) (2.5200000,1.4223133) (2.5266667,1.4199574) (2.5333333,1.4176170) (2.5400000,1.4152919) (2.5466667,1.4129819) (2.5533333,1.4106869) (2.5600000,1.4084069) (2.5666667,1.4061417) (2.5733333,1.4038911) (2.5800000,1.4016551) (2.5866667,1.3994335) (2.5933333,1.3972262) (2.6000000,1.3950331) (2.6066667,1.3928542) (2.6133333,1.3906892) (2.6200000,1.3885381) (2.6266667,1.3864007) (2.6333333,1.3842770) (2.6400000,1.3821668) (2.6466667,1.3800701) (2.6533333,1.3779867) (2.6600000,1.3759165) (2.6666667,1.3738594) (2.6733333,1.3718154) (2.6800000,1.3697843) (2.6866667,1.3677659) (2.6933333,1.3657604) (2.7000000,1.3637674) (2.7066667,1.3617869) (2.7133333,1.3598189) (2.7200000,1.3578632) (2.7266667,1.3559197) (2.7333333,1.3539884) (2.7400000,1.3520691) (2.7466667,1.3501617) (2.7533333,1.3482663) (2.7600000,1.3463826) (2.7666667,1.3445106) (2.7733333,1.3426501) (2.7800000,1.3408012) (2.7866667,1.3389637) (2.7933333,1.3371376) (2.8000000,1.3353227) (2.8066667,1.3335189) (2.8133333,1.3317263) (2.8200000,1.3299446) (2.8266667,1.3281738) (2.8333333,1.3264139) (2.8400000,1.3246648) (2.8466667,1.3229263) (2.8533333,1.3211984) (2.8600000,1.3194810) (2.8666667,1.3177741) (2.8733333,1.3160775) (2.8800000,1.3143912) (2.8866667,1.3127151) (2.8933333,1.3110492) (2.9000000,1.3093933) (2.9066667,1.3077474) (2.9133333,1.3061114) (2.9200000,1.3044853) (2.9266667,1.3028690) (2.9333333,1.3012623) (2.9400000,1.2996653) (2.9466667,1.2980778) (2.9533333,1.2964999) (2.9600000,1.2949313) (2.9666667,1.2933721) (2.9733333,1.2918223) (2.9800000,1.2902816) (2.9866667,1.2887501) (2.9933333,1.2872277) (3.0000000,1.2857143) (3.0066667,1.2842099) (3.0133333,1.2827143) (3.0200000,1.2812277) (3.0266667,1.2797498) (3.0333333,1.2782806) (3.0400000,1.2768200) (3.0466667,1.2753681) (3.0533333,1.2739247) (3.0600000,1.2724897) (3.0666667,1.2710632) (3.0733333,1.2696450) (3.0800000,1.2682352) (3.0866667,1.2668335) (3.0933333,1.2654401) (3.1000000,1.2640547) (3.1066667,1.2626775) (3.1133333,1.2613082) (3.1200000,1.2599469) (3.1266667,1.2585935) (3.1333333,1.2572480) (3.1400000,1.2559102) (3.1466667,1.2545802) (3.1533333,1.2532579) (3.1600000,1.2519432) (3.1666667,1.2506361) (3.1733333,1.2493365) (3.1800000,1.2480444) (3.1866667,1.2467598) (3.1933333,1.2454825) (3.2000000,1.2442126) (3.2066667,1.2429499) (3.2133333,1.2416945) (3.2200000,1.2404462) (3.2266667,1.2392051) (3.2333333,1.2379711) (3.2400000,1.2367441) (3.2466667,1.2355241) (3.2533333,1.2343110) (3.2600000,1.2331049) (3.2666667,1.2319056) (3.2733333,1.2307131) (3.2800000,1.2295273) (3.2866667,1.2283483) (3.2933333,1.2271760) (3.3000000,1.2260103) (3.3066667,1.2248511) (3.3133333,1.2236986) (3.3200000,1.2225525) (3.3266667,1.2214129) (3.3333333,1.2202796) (3.3400000,1.2191528) (3.3466667,1.2180323) (3.3533333,1.2169181) (3.3600000,1.2158101) (3.3666667,1.2147083) (3.3733333,1.2136127) (3.3800000,1.2125233) (3.3866667,1.2114399) (3.3933333,1.2103625) (3.4000000,1.2092912) (3.4066667,1.2082258) (3.4133333,1.2071664) (3.4200000,1.2061129) (3.4266667,1.2050652) (3.4333333,1.2040233) (3.4400000,1.2029873) (3.4466667,1.2019569) (3.4533333,1.2009323) (3.4600000,1.1999133) (3.4666667,1.1989000) (3.4733333,1.1978923) (3.4800000,1.1968901) (3.4866667,1.1958935) (3.4933333,1.1949024) (3.5000000,1.1939167) (3.5066667,1.1929364) (3.5133333,1.1919615) (3.5200000,1.1909920) (3.5266667,1.1900278) (3.5333333,1.1890689) (3.5400000,1.1881153) (3.5466667,1.1871668) (3.5533333,1.1862236) (3.5600000,1.1852855) (3.5666667,1.1843525) (3.5733333,1.1834246) (3.5800000,1.1825018) (3.5866667,1.1815841) (3.5933333,1.1806713) (3.6000000,1.1797634) (3.6066667,1.1788606) (3.6133333,1.1779626) (3.6200000,1.1770695) (3.6266667,1.1761812) (3.6333333,1.1752978) (3.6400000,1.1744191) (3.6466667,1.1735452) (3.6533333,1.1726760) (3.6600000,1.1718115) (3.6666667,1.1709517) (3.6733333,1.1700966) (3.6800000,1.1692460) (3.6866667,1.1684000) (3.6933333,1.1675586) (3.7000000,1.1667217) (3.7066667,1.1658893) (3.7133333,1.1650614) (3.7200000,1.1642379) (3.7266667,1.1634189) (3.7333333,1.1626042) (3.7400000,1.1617939) (3.7466667,1.1609879) (3.7533333,1.1601863) (3.7600000,1.1593889) (3.7666667,1.1585958) (3.7733333,1.1578069) (3.7800000,1.1570223) (3.7866667,1.1562418) (3.7933333,1.1554655) (3.8000000,1.1546933) (3.8066667,1.1539252) (3.8133333,1.1531612) (3.8200000,1.1524013) (3.8266667,1.1516454) (3.8333333,1.1508935) (3.8400000,1.1501456) (3.8466667,1.1494017) (3.8533333,1.1486617) (3.8600000,1.1479256) (3.8666667,1.1471935) (3.8733333,1.1464652) (3.8800000,1.1457407) (3.8866667,1.1450201) (3.8933333,1.1443032) (3.9000000,1.1435902) (3.9066667,1.1428809) (3.9133333,1.1421753) (3.9200000,1.1414735) (3.9266667,1.1407754) (3.9333333,1.1400809) (3.9400000,1.1393901) (3.9466667,1.1387029) (3.9533333,1.1380193) (3.9600000,1.1373393) (3.9666667,1.1366628) (3.9733333,1.1359899) (3.9800000,1.1353205) (3.9866667,1.1346547) (3.9933333,1.1339923) (4.0000000,1.1333333)};
\draw[orange!85!black,thick,dashed] plot coordinates {(-4.0000000,1.1333333) (-3.9933333,1.1339923) (-3.9866667,1.1346547) (-3.9800000,1.1353205) (-3.9733333,1.1359899) (-3.9666667,1.1366628) (-3.9600000,1.1373393) (-3.9533333,1.1380193) (-3.9466667,1.1387029) (-3.9400000,1.1393901) (-3.9333333,1.1400809) (-3.9266667,1.1407754) (-3.9200000,1.1414735) (-3.9133333,1.1421753) (-3.9066667,1.1428809) (-3.9000000,1.1435902) (-3.8933333,1.1443032) (-3.8866667,1.1450201) (-3.8800000,1.1457407) (-3.8733333,1.1464652) (-3.8666667,1.1471935) (-3.8600000,1.1479256) (-3.8533333,1.1486617) (-3.8466667,1.1494017) (-3.8400000,1.1501456) (-3.8333333,1.1508935) (-3.8266667,1.1516454) (-3.8200000,1.1524013) (-3.8133333,1.1531612) (-3.8066667,1.1539252) (-3.8000000,1.1546933) (-3.7933333,1.1554655) (-3.7866667,1.1562418) (-3.7800000,1.1570223) (-3.7733333,1.1578069) (-3.7666667,1.1585958) (-3.7600000,1.1593889) (-3.7533333,1.1601863) (-3.7466667,1.1609879) (-3.7400000,1.1617939) (-3.7333333,1.1626042) (-3.7266667,1.1634189) (-3.7200000,1.1642379) (-3.7133333,1.1650614) (-3.7066667,1.1658893) (-3.7000000,1.1667217) (-3.6933333,1.1675586) (-3.6866667,1.1684000) (-3.6800000,1.1692460) (-3.6733333,1.1700966) (-3.6666667,1.1709517) (-3.6600000,1.1718115) (-3.6533333,1.1726760) (-3.6466667,1.1735452) (-3.6400000,1.1744191) (-3.6333333,1.1752978) (-3.6266667,1.1761812) (-3.6200000,1.1770695) (-3.6133333,1.1779626) (-3.6066667,1.1788606) (-3.6000000,1.1797634) (-3.5933333,1.1806713) (-3.5866667,1.1815841) (-3.5800000,1.1825018) (-3.5733333,1.1834246) (-3.5666667,1.1843525) (-3.5600000,1.1852855) (-3.5533333,1.1862236) (-3.5466667,1.1871668) (-3.5400000,1.1881153) (-3.5333333,1.1890689) (-3.5266667,1.1900278) (-3.5200000,1.1909920) (-3.5133333,1.1919615) (-3.5066667,1.1929364) (-3.5000000,1.1939167) (-3.4933333,1.1949024) (-3.4866667,1.1958935) (-3.4800000,1.1968901) (-3.4733333,1.1978923) (-3.4666667,1.1989000) (-3.4600000,1.1999133) (-3.4533333,1.2009323) (-3.4466667,1.2019569) (-3.4400000,1.2029873) (-3.4333333,1.2040233) (-3.4266667,1.2050652) (-3.4200000,1.2061129) (-3.4133333,1.2071664) (-3.4066667,1.2082258) (-3.4000000,1.2092912) (-3.3933333,1.2103625) (-3.3866667,1.2114399) (-3.3800000,1.2125233) (-3.3733333,1.2136127) (-3.3666667,1.2147083) (-3.3600000,1.2158101) (-3.3533333,1.2169181) (-3.3466667,1.2180323) (-3.3400000,1.2191528) (-3.3333333,1.2202796) (-3.3266667,1.2214129) (-3.3200000,1.2225525) (-3.3133333,1.2236986) (-3.3066667,1.2248511) (-3.3000000,1.2260103) (-3.2933333,1.2271760) (-3.2866667,1.2283483) (-3.2800000,1.2295273) (-3.2733333,1.2307131) (-3.2666667,1.2319056) (-3.2600000,1.2331049) (-3.2533333,1.2343110) (-3.2466667,1.2355241) (-3.2400000,1.2367441) (-3.2333333,1.2379711) (-3.2266667,1.2392051) (-3.2200000,1.2404462) (-3.2133333,1.2416945) (-3.2066667,1.2429499) (-3.2000000,1.2442126) (-3.1933333,1.2454825) (-3.1866667,1.2467598) (-3.1800000,1.2480444) (-3.1733333,1.2493365) (-3.1666667,1.2506361) (-3.1600000,1.2519432) (-3.1533333,1.2532579) (-3.1466667,1.2545802) (-3.1400000,1.2559102) (-3.1333333,1.2572480) (-3.1266667,1.2585935) (-3.1200000,1.2599469) (-3.1133333,1.2613082) (-3.1066667,1.2626775) (-3.1000000,1.2640547) (-3.0933333,1.2654401) (-3.0866667,1.2668335) (-3.0800000,1.2682352) (-3.0733333,1.2696450) (-3.0666667,1.2710632) (-3.0600000,1.2724897) (-3.0533333,1.2739247) (-3.0466667,1.2753681) (-3.0400000,1.2768200) (-3.0333333,1.2782806) (-3.0266667,1.2797498) (-3.0200000,1.2812277) (-3.0133333,1.2827143) (-3.0066667,1.2842099) (-3.0000000,1.2857143) (-2.9933333,1.2872277) (-2.9866667,1.2887501) (-2.9800000,1.2902816) (-2.9733333,1.2918223) (-2.9666667,1.2933721) (-2.9600000,1.2949313) (-2.9533333,1.2964999) (-2.9466667,1.2980778) (-2.9400000,1.2996653) (-2.9333333,1.3012623) (-2.9266667,1.3028690) (-2.9200000,1.3044853) (-2.9133333,1.3061114) (-2.9066667,1.3077474) (-2.9000000,1.3093933) (-2.8933333,1.3110492) (-2.8866667,1.3127151) (-2.8800000,1.3143912) (-2.8733333,1.3160775) (-2.8666667,1.3177741) (-2.8600000,1.3194810) (-2.8533333,1.3211984) (-2.8466667,1.3229263) (-2.8400000,1.3246648) (-2.8333333,1.3264139) (-2.8266667,1.3281738) (-2.8200000,1.3299446) (-2.8133333,1.3317263) (-2.8066667,1.3335189) (-2.8000000,1.3353227) (-2.7933333,1.3371376) (-2.7866667,1.3389637) (-2.7800000,1.3408012) (-2.7733333,1.3426501) (-2.7666667,1.3445106) (-2.7600000,1.3463826) (-2.7533333,1.3482663) (-2.7466667,1.3501617) (-2.7400000,1.3520691) (-2.7333333,1.3539884) (-2.7266667,1.3559197) (-2.7200000,1.3578632) (-2.7133333,1.3598189) (-2.7066667,1.3617869) (-2.7000000,1.3637674) (-2.6933333,1.3657604) (-2.6866667,1.3677659) (-2.6800000,1.3697843) (-2.6733333,1.3718154) (-2.6666667,1.3738594) (-2.6600000,1.3759165) (-2.6533333,1.3779867) (-2.6466667,1.3800701) (-2.6400000,1.3821668) (-2.6333333,1.3842770) (-2.6266667,1.3864007) (-2.6200000,1.3885381) (-2.6133333,1.3906892) (-2.6066667,1.3928542) (-2.6000000,1.3950331) (-2.5933333,1.3972262) (-2.5866667,1.3994335) (-2.5800000,1.4016551) (-2.5733333,1.4038911) (-2.5666667,1.4061417) (-2.5600000,1.4084069) (-2.5533333,1.4106869) (-2.5466667,1.4129819) (-2.5400000,1.4152919) (-2.5333333,1.4176170) (-2.5266667,1.4199574) (-2.5200000,1.4223133) (-2.5133333,1.4246846) (-2.5066667,1.4270716) (-2.5000000,1.4294745) (-2.4933333,1.4318932) (-2.4866667,1.4343280) (-2.4800000,1.4367790) (-2.4733333,1.4392463) (-2.4666667,1.4417301) (-2.4600000,1.4442305) (-2.4533333,1.4467476) (-2.4466667,1.4492816) (-2.4400000,1.4518327) (-2.4333333,1.4544009) (-2.4266667,1.4569864) (-2.4200000,1.4595894) (-2.4133333,1.4622100) (-2.4066667,1.4648484) (-2.4000000,1.4675047) (-2.3933333,1.4701791) (-2.3866667,1.4728717) (-2.3800000,1.4755827) (-2.3733333,1.4783122) (-2.3666667,1.4810605) (-2.3600000,1.4838276) (-2.3533333,1.4866138) (-2.3466667,1.4894192) (-2.3400000,1.4922439) (-2.3333333,1.4950882) (-2.3266667,1.4979522) (-2.3200000,1.5008361) (-2.3133333,1.5037401) (-2.3066667,1.5066643) (-2.3000000,1.5096090) (-2.2933333,1.5125742) (-2.2866667,1.5155603) (-2.2800000,1.5185674) (-2.2733333,1.5215956) (-2.2666667,1.5246452) (-2.2600000,1.5277163) (-2.2533333,1.5308093) (-2.2466667,1.5339241) (-2.2400000,1.5370612) (-2.2333333,1.5402206) (-2.2266667,1.5434025) (-2.2200000,1.5466072) (-2.2133333,1.5498350) (-2.2066667,1.5530859) (-2.2000000,1.5563602) (-2.1933333,1.5596582) (-2.1866667,1.5629800) (-2.1800000,1.5663259) (-2.1733333,1.5696961) (-2.1666667,1.5730908) (-2.1600000,1.5765103) (-2.1533333,1.5799548) (-2.1466667,1.5834246) (-2.1400000,1.5869198) (-2.1333333,1.5904408) (-2.1266667,1.5939877) (-2.1200000,1.5975609) (-2.1133333,1.6011605) (-2.1066667,1.6047868) (-2.1000000,1.6084402) (-2.0933333,1.6121208) (-2.0866667,1.6158289) (-2.0800000,1.6195648) (-2.0733333,1.6233288) (-2.0666667,1.6271211) (-2.0600000,1.6309421) (-2.0533333,1.6347919) (-2.0466667,1.6386710) (-2.0400000,1.6425795) (-2.0333333,1.6465178) (-2.0266667,1.6504862) (-2.0200000,1.6544850) (-2.0133333,1.6585144) (-2.0066667,1.6625749) (-2.0000000,1.6666667) (-1.9933333,1.6707901) (-1.9866667,1.6749454) (-1.9800000,1.6791331) (-1.9733333,1.6833534) (-1.9666667,1.6876066) (-1.9600000,1.6918931) (-1.9533333,1.6962133) (-1.9466667,1.7005675) (-1.9400000,1.7049560) (-1.9333333,1.7093792) (-1.9266667,1.7138375) (-1.9200000,1.7183313) (-1.9133333,1.7228609) (-1.9066667,1.7274267) (-1.9000000,1.7320291) (-1.8933333,1.7366684) (-1.8866667,1.7413452) (-1.8800000,1.7460597) (-1.8733333,1.7508124) (-1.8666667,1.7556038) (-1.8600000,1.7604341) (-1.8533333,1.7653039) (-1.8466667,1.7702136) (-1.8400000,1.7751637) (-1.8333333,1.7801545) (-1.8266667,1.7851865) (-1.8200000,1.7902602) (-1.8133333,1.7953761) (-1.8066667,1.8005345) (-1.8000000,1.8057361) (-1.7933333,1.8109813) (-1.7866667,1.8162705) (-1.7800000,1.8216043) (-1.7733333,1.8269832) (-1.7666667,1.8324078) (-1.7600000,1.8378784) (-1.7533333,1.8433957) (-1.7466667,1.8489602) (-1.7400000,1.8545724) (-1.7333333,1.8602330) (-1.7266667,1.8659424) (-1.7200000,1.8717012) (-1.7133333,1.8775101) (-1.7066667,1.8833696) (-1.7000000,1.8892803) (-1.6933333,1.8952429) (-1.6866667,1.9012579) (-1.6800000,1.9073259) (-1.6733333,1.9134477) (-1.6666667,1.9196239) (-1.6600000,1.9258551) (-1.6533333,1.9321420) (-1.6466667,1.9384853) (-1.6400000,1.9448856) (-1.6333333,1.9513438) (-1.6266667,1.9578604) (-1.6200000,1.9644362) (-1.6133333,1.9710721) (-1.6066667,1.9777686) (-1.6000000,1.9845266) (-1.5933333,1.9913469) (-1.5866667,1.9982302) (-1.5800000,2.0051774) (-1.5733333,2.0121893) (-1.5666667,2.0192666) (-1.5600000,2.0264103) (-1.5533333,2.0336213) (-1.5466667,2.0409003) (-1.5400000,2.0482482) (-1.5333333,2.0556661) (-1.5266667,2.0631547) (-1.5200000,2.0707151) (-1.5133333,2.0783482) (-1.5066667,2.0860549) (-1.5000000,2.0938363) (-1.4933333,2.1016934) (-1.4866667,2.1096271) (-1.4800000,2.1176385) (-1.4733333,2.1257288) (-1.4666667,2.1338988) (-1.4600000,2.1421499) (-1.4533333,2.1504830) (-1.4466667,2.1588994) (-1.4400000,2.1674001) (-1.4333333,2.1759864) (-1.4266667,2.1846595) (-1.4200000,2.1934206) (-1.4133333,2.2022709) (-1.4066667,2.2112118) (-1.4000000,2.2202445) (-1.3933333,2.2293703) (-1.3866667,2.2385906) (-1.3800000,2.2479069) (-1.3733333,2.2573204) (-1.3666667,2.2668325) (-1.3600000,2.2764449) (-1.3533333,2.2861589) (-1.3466667,2.2959761) (-1.3400000,2.3058980) (-1.3333333,2.3159262) (-1.3266667,2.3260623) (-1.3200000,2.3363079) (-1.3133333,2.3466647) (-1.3066667,2.3571344) (-1.3000000,2.3677189) (-1.2933333,2.3784197) (-1.2866667,2.3892388) (-1.2800000,2.4001780) (-1.2733333,2.4112393) (-1.2666667,2.4224244) (-1.2600000,2.4337355) (-1.2533333,2.4451745) (-1.2466667,2.4567435) (-1.2400000,2.4684446) (-1.2333333,2.4802799) (-1.2266667,2.4922517) (-1.2200000,2.5043621) (-1.2133333,2.5166135) (-1.2066667,2.5290081) (-1.2000000,2.5415485) (-1.1933333,2.5542371) (-1.1866667,2.5670763) (-1.1800000,2.5800687) (-1.1733333,2.5932169) (-1.1666667,2.6065236) (-1.1600000,2.6199916) (-1.1533333,2.6336237) (-1.1466667,2.6474226) (-1.1400000,2.6613914) (-1.1333333,2.6755331) (-1.1266667,2.6898507) (-1.1200000,2.7043473) (-1.1133333,2.7190263) (-1.1066667,2.7338909) (-1.1000000,2.7489444) (-1.0933333,2.7641904) (-1.0866667,2.7796324) (-1.0800000,2.7952740) (-1.0733333,2.8111190) (-1.0666667,2.8271711) (-1.0600000,2.8434343) (-1.0533333,2.8599126) (-1.0466667,2.8766102) (-1.0400000,2.8935312) (-1.0333333,2.9106799) (-1.0266667,2.9280609) (-1.0200000,2.9456787) (-1.0133333,2.9635379) (-1.0066667,2.9816433) (-1.0000000,3.0000000) (-0.9933333,3.0186129) (-0.9866667,3.0374873) (-0.9800000,3.0566284) (-0.9733333,3.0760419) (-0.9666667,3.0957332) (-0.9600000,3.1157083) (-0.9533333,3.1359731) (-0.9466667,3.1565336) (-0.9400000,3.1773963) (-0.9333333,3.1985676) (-0.9266667,3.2200541) (-0.9200000,3.2418627) (-0.9133333,3.2640005) (-0.9066667,3.2864747) (-0.9000000,3.3092929) (-0.8933333,3.3324626) (-0.8866667,3.3559920) (-0.8800000,3.3798891) (-0.8733333,3.4041623) (-0.8666667,3.4288205) (-0.8600000,3.4538724) (-0.8533333,3.4793274) (-0.8466667,3.5051949) (-0.8400000,3.5314849) (-0.8333333,3.5582074) (-0.8266667,3.5853730) (-0.8200000,3.6129923) (-0.8133333,3.6410767) (-0.8066667,3.6696377) (-0.8000000,3.6986870) (-0.7933333,3.7282372) (-0.7866667,3.7583009) (-0.7800000,3.7888912) (-0.7733333,3.8200219) (-0.7666667,3.8517070) (-0.7600000,3.8839611) (-0.7533333,3.9167994) (-0.7466667,3.9502374) (-0.7400000,3.9842914) (-0.7333333,4.0189782) (-0.7266667,4.0543153) (-0.7200000,4.0903206) (-0.7133333,4.1270130) (-0.7066667,4.1644119) (-0.7000000,4.2025375) (-0.6933333,4.2414107) (-0.6866667,4.2810533) (-0.6800000,4.3214880) (-0.6733333,4.3627384) (-0.6666667,4.4048288) (-0.6600000,4.4477847) (-0.6533333,4.4916327) (-0.6466667,4.5364003) (-0.6400000,4.5821163) (-0.6333333,4.6288107) (-0.6266667,4.6765146) (-0.6200000,4.7252607) (-0.6133333,4.7750829) (-0.6066667,4.8260167) (-0.6000000,4.8780992) (-0.5933333,4.9313691) (-0.5866667,4.9858669)};
\draw[orange!85!black,thick,dashed] plot coordinates {(0.5866667,-4.9858669) (0.5933333,-4.9313691) (0.6000000,-4.8780992) (0.6066667,-4.8260167) (0.6133333,-4.7750829) (0.6200000,-4.7252607) (0.6266667,-4.6765146) (0.6333333,-4.6288107) (0.6400000,-4.5821163) (0.6466667,-4.5364003) (0.6533333,-4.4916327) (0.6600000,-4.4477847) (0.6666667,-4.4048288) (0.6733333,-4.3627384) (0.6800000,-4.3214880) (0.6866667,-4.2810533) (0.6933333,-4.2414107) (0.7000000,-4.2025375) (0.7066667,-4.1644119) (0.7133333,-4.1270130) (0.7200000,-4.0903206) (0.7266667,-4.0543153) (0.7333333,-4.0189782) (0.7400000,-3.9842914) (0.7466667,-3.9502374) (0.7533333,-3.9167994) (0.7600000,-3.8839611) (0.7666667,-3.8517070) (0.7733333,-3.8200219) (0.7800000,-3.7888912) (0.7866667,-3.7583009) (0.7933333,-3.7282372) (0.8000000,-3.6986870) (0.8066667,-3.6696377) (0.8133333,-3.6410767) (0.8200000,-3.6129923) (0.8266667,-3.5853730) (0.8333333,-3.5582074) (0.8400000,-3.5314849) (0.8466667,-3.5051949) (0.8533333,-3.4793274) (0.8600000,-3.4538724) (0.8666667,-3.4288205) (0.8733333,-3.4041623) (0.8800000,-3.3798891) (0.8866667,-3.3559920) (0.8933333,-3.3324626) (0.9000000,-3.3092929) (0.9066667,-3.2864747) (0.9133333,-3.2640005) (0.9200000,-3.2418627) (0.9266667,-3.2200541) (0.9333333,-3.1985676) (0.9400000,-3.1773963) (0.9466667,-3.1565336) (0.9533333,-3.1359731) (0.9600000,-3.1157083) (0.9666667,-3.0957332) (0.9733333,-3.0760419) (0.9800000,-3.0566284) (0.9866667,-3.0374873) (0.9933333,-3.0186129) (1.0000000,-3.0000000) (1.0066667,-2.9816433) (1.0133333,-2.9635379) (1.0200000,-2.9456787) (1.0266667,-2.9280609) (1.0333333,-2.9106799) (1.0400000,-2.8935312) (1.0466667,-2.8766102) (1.0533333,-2.8599126) (1.0600000,-2.8434343) (1.0666667,-2.8271711) (1.0733333,-2.8111190) (1.0800000,-2.7952740) (1.0866667,-2.7796324) (1.0933333,-2.7641904) (1.1000000,-2.7489444) (1.1066667,-2.7338909) (1.1133333,-2.7190263) (1.1200000,-2.7043473) (1.1266667,-2.6898507) (1.1333333,-2.6755331) (1.1400000,-2.6613914) (1.1466667,-2.6474226) (1.1533333,-2.6336237) (1.1600000,-2.6199916) (1.1666667,-2.6065236) (1.1733333,-2.5932169) (1.1800000,-2.5800687) (1.1866667,-2.5670763) (1.1933333,-2.5542371) (1.2000000,-2.5415485) (1.2066667,-2.5290081) (1.2133333,-2.5166135) (1.2200000,-2.5043621) (1.2266667,-2.4922517) (1.2333333,-2.4802799) (1.2400000,-2.4684446) (1.2466667,-2.4567435) (1.2533333,-2.4451745) (1.2600000,-2.4337355) (1.2666667,-2.4224244) (1.2733333,-2.4112393) (1.2800000,-2.4001780) (1.2866667,-2.3892388) (1.2933333,-2.3784197) (1.3000000,-2.3677189) (1.3066667,-2.3571344) (1.3133333,-2.3466647) (1.3200000,-2.3363079) (1.3266667,-2.3260623) (1.3333333,-2.3159262) (1.3400000,-2.3058980) (1.3466667,-2.2959761) (1.3533333,-2.2861589) (1.3600000,-2.2764449) (1.3666667,-2.2668325) (1.3733333,-2.2573204) (1.3800000,-2.2479069) (1.3866667,-2.2385906) (1.3933333,-2.2293703) (1.4000000,-2.2202445) (1.4066667,-2.2112118) (1.4133333,-2.2022709) (1.4200000,-2.1934206) (1.4266667,-2.1846595) (1.4333333,-2.1759864) (1.4400000,-2.1674001) (1.4466667,-2.1588994) (1.4533333,-2.1504830) (1.4600000,-2.1421499) (1.4666667,-2.1338988) (1.4733333,-2.1257288) (1.4800000,-2.1176385) (1.4866667,-2.1096271) (1.4933333,-2.1016934) (1.5000000,-2.0938363) (1.5066667,-2.0860549) (1.5133333,-2.0783482) (1.5200000,-2.0707151) (1.5266667,-2.0631547) (1.5333333,-2.0556661) (1.5400000,-2.0482482) (1.5466667,-2.0409003) (1.5533333,-2.0336213) (1.5600000,-2.0264103) (1.5666667,-2.0192666) (1.5733333,-2.0121893) (1.5800000,-2.0051774) (1.5866667,-1.9982302) (1.5933333,-1.9913469) (1.6000000,-1.9845266) (1.6066667,-1.9777686) (1.6133333,-1.9710721) (1.6200000,-1.9644362) (1.6266667,-1.9578604) (1.6333333,-1.9513438) (1.6400000,-1.9448856) (1.6466667,-1.9384853) (1.6533333,-1.9321420) (1.6600000,-1.9258551) (1.6666667,-1.9196239) (1.6733333,-1.9134477) (1.6800000,-1.9073259) (1.6866667,-1.9012579) (1.6933333,-1.8952429) (1.7000000,-1.8892803) (1.7066667,-1.8833696) (1.7133333,-1.8775101) (1.7200000,-1.8717012) (1.7266667,-1.8659424) (1.7333333,-1.8602330) (1.7400000,-1.8545724) (1.7466667,-1.8489602) (1.7533333,-1.8433957) (1.7600000,-1.8378784) (1.7666667,-1.8324078) (1.7733333,-1.8269832) (1.7800000,-1.8216043) (1.7866667,-1.8162705) (1.7933333,-1.8109813) (1.8000000,-1.8057361) (1.8066667,-1.8005345) (1.8133333,-1.7953761) (1.8200000,-1.7902602) (1.8266667,-1.7851865) (1.8333333,-1.7801545) (1.8400000,-1.7751637) (1.8466667,-1.7702136) (1.8533333,-1.7653039) (1.8600000,-1.7604341) (1.8666667,-1.7556038) (1.8733333,-1.7508124) (1.8800000,-1.7460597) (1.8866667,-1.7413452) (1.8933333,-1.7366684) (1.9000000,-1.7320291) (1.9066667,-1.7274267) (1.9133333,-1.7228609) (1.9200000,-1.7183313) (1.9266667,-1.7138375) (1.9333333,-1.7093792) (1.9400000,-1.7049560) (1.9466667,-1.7005675) (1.9533333,-1.6962133) (1.9600000,-1.6918931) (1.9666667,-1.6876066) (1.9733333,-1.6833534) (1.9800000,-1.6791331) (1.9866667,-1.6749454) (1.9933333,-1.6707901) (2.0000000,-1.6666667) (2.0066667,-1.6625749) (2.0133333,-1.6585144) (2.0200000,-1.6544850) (2.0266667,-1.6504862) (2.0333333,-1.6465178) (2.0400000,-1.6425795) (2.0466667,-1.6386710) (2.0533333,-1.6347919) (2.0600000,-1.6309421) (2.0666667,-1.6271211) (2.0733333,-1.6233288) (2.0800000,-1.6195648) (2.0866667,-1.6158289) (2.0933333,-1.6121208) (2.1000000,-1.6084402) (2.1066667,-1.6047868) (2.1133333,-1.6011605) (2.1200000,-1.5975609) (2.1266667,-1.5939877) (2.1333333,-1.5904408) (2.1400000,-1.5869198) (2.1466667,-1.5834246) (2.1533333,-1.5799548) (2.1600000,-1.5765103) (2.1666667,-1.5730908) (2.1733333,-1.5696961) (2.1800000,-1.5663259) (2.1866667,-1.5629800) (2.1933333,-1.5596582) (2.2000000,-1.5563602) (2.2066667,-1.5530859) (2.2133333,-1.5498350) (2.2200000,-1.5466072) (2.2266667,-1.5434025) (2.2333333,-1.5402206) (2.2400000,-1.5370612) (2.2466667,-1.5339241) (2.2533333,-1.5308093) (2.2600000,-1.5277163) (2.2666667,-1.5246452) (2.2733333,-1.5215956) (2.2800000,-1.5185674) (2.2866667,-1.5155603) (2.2933333,-1.5125742) (2.3000000,-1.5096090) (2.3066667,-1.5066643) (2.3133333,-1.5037401) (2.3200000,-1.5008361) (2.3266667,-1.4979522) (2.3333333,-1.4950882) (2.3400000,-1.4922439) (2.3466667,-1.4894192) (2.3533333,-1.4866138) (2.3600000,-1.4838276) (2.3666667,-1.4810605) (2.3733333,-1.4783122) (2.3800000,-1.4755827) (2.3866667,-1.4728717) (2.3933333,-1.4701791) (2.4000000,-1.4675047) (2.4066667,-1.4648484) (2.4133333,-1.4622100) (2.4200000,-1.4595894) (2.4266667,-1.4569864) (2.4333333,-1.4544009) (2.4400000,-1.4518327) (2.4466667,-1.4492816) (2.4533333,-1.4467476) (2.4600000,-1.4442305) (2.4666667,-1.4417301) (2.4733333,-1.4392463) (2.4800000,-1.4367790) (2.4866667,-1.4343280) (2.4933333,-1.4318932) (2.5000000,-1.4294745) (2.5066667,-1.4270716) (2.5133333,-1.4246846) (2.5200000,-1.4223133) (2.5266667,-1.4199574) (2.5333333,-1.4176170) (2.5400000,-1.4152919) (2.5466667,-1.4129819) (2.5533333,-1.4106869) (2.5600000,-1.4084069) (2.5666667,-1.4061417) (2.5733333,-1.4038911) (2.5800000,-1.4016551) (2.5866667,-1.3994335) (2.5933333,-1.3972262) (2.6000000,-1.3950331) (2.6066667,-1.3928542) (2.6133333,-1.3906892) (2.6200000,-1.3885381) (2.6266667,-1.3864007) (2.6333333,-1.3842770) (2.6400000,-1.3821668) (2.6466667,-1.3800701) (2.6533333,-1.3779867) (2.6600000,-1.3759165) (2.6666667,-1.3738594) (2.6733333,-1.3718154) (2.6800000,-1.3697843) (2.6866667,-1.3677659) (2.6933333,-1.3657604) (2.7000000,-1.3637674) (2.7066667,-1.3617869) (2.7133333,-1.3598189) (2.7200000,-1.3578632) (2.7266667,-1.3559197) (2.7333333,-1.3539884) (2.7400000,-1.3520691) (2.7466667,-1.3501617) (2.7533333,-1.3482663) (2.7600000,-1.3463826) (2.7666667,-1.3445106) (2.7733333,-1.3426501) (2.7800000,-1.3408012) (2.7866667,-1.3389637) (2.7933333,-1.3371376) (2.8000000,-1.3353227) (2.8066667,-1.3335189) (2.8133333,-1.3317263) (2.8200000,-1.3299446) (2.8266667,-1.3281738) (2.8333333,-1.3264139) (2.8400000,-1.3246648) (2.8466667,-1.3229263) (2.8533333,-1.3211984) (2.8600000,-1.3194810) (2.8666667,-1.3177741) (2.8733333,-1.3160775) (2.8800000,-1.3143912) (2.8866667,-1.3127151) (2.8933333,-1.3110492) (2.9000000,-1.3093933) (2.9066667,-1.3077474) (2.9133333,-1.3061114) (2.9200000,-1.3044853) (2.9266667,-1.3028690) (2.9333333,-1.3012623) (2.9400000,-1.2996653) (2.9466667,-1.2980778) (2.9533333,-1.2964999) (2.9600000,-1.2949313) (2.9666667,-1.2933721) (2.9733333,-1.2918223) (2.9800000,-1.2902816) (2.9866667,-1.2887501) (2.9933333,-1.2872277) (3.0000000,-1.2857143) (3.0066667,-1.2842099) (3.0133333,-1.2827143) (3.0200000,-1.2812277) (3.0266667,-1.2797498) (3.0333333,-1.2782806) (3.0400000,-1.2768200) (3.0466667,-1.2753681) (3.0533333,-1.2739247) (3.0600000,-1.2724897) (3.0666667,-1.2710632) (3.0733333,-1.2696450) (3.0800000,-1.2682352) (3.0866667,-1.2668335) (3.0933333,-1.2654401) (3.1000000,-1.2640547) (3.1066667,-1.2626775) (3.1133333,-1.2613082) (3.1200000,-1.2599469) (3.1266667,-1.2585935) (3.1333333,-1.2572480) (3.1400000,-1.2559102) (3.1466667,-1.2545802) (3.1533333,-1.2532579) (3.1600000,-1.2519432) (3.1666667,-1.2506361) (3.1733333,-1.2493365) (3.1800000,-1.2480444) (3.1866667,-1.2467598) (3.1933333,-1.2454825) (3.2000000,-1.2442126) (3.2066667,-1.2429499) (3.2133333,-1.2416945) (3.2200000,-1.2404462) (3.2266667,-1.2392051) (3.2333333,-1.2379711) (3.2400000,-1.2367441) (3.2466667,-1.2355241) (3.2533333,-1.2343110) (3.2600000,-1.2331049) (3.2666667,-1.2319056) (3.2733333,-1.2307131) (3.2800000,-1.2295273) (3.2866667,-1.2283483) (3.2933333,-1.2271760) (3.3000000,-1.2260103) (3.3066667,-1.2248511) (3.3133333,-1.2236986) (3.3200000,-1.2225525) (3.3266667,-1.2214129) (3.3333333,-1.2202796) (3.3400000,-1.2191528) (3.3466667,-1.2180323) (3.3533333,-1.2169181) (3.3600000,-1.2158101) (3.3666667,-1.2147083) (3.3733333,-1.2136127) (3.3800000,-1.2125233) (3.3866667,-1.2114399) (3.3933333,-1.2103625) (3.4000000,-1.2092912) (3.4066667,-1.2082258) (3.4133333,-1.2071664) (3.4200000,-1.2061129) (3.4266667,-1.2050652) (3.4333333,-1.2040233) (3.4400000,-1.2029873) (3.4466667,-1.2019569) (3.4533333,-1.2009323) (3.4600000,-1.1999133) (3.4666667,-1.1989000) (3.4733333,-1.1978923) (3.4800000,-1.1968901) (3.4866667,-1.1958935) (3.4933333,-1.1949024) (3.5000000,-1.1939167) (3.5066667,-1.1929364) (3.5133333,-1.1919615) (3.5200000,-1.1909920) (3.5266667,-1.1900278) (3.5333333,-1.1890689) (3.5400000,-1.1881153) (3.5466667,-1.1871668) (3.5533333,-1.1862236) (3.5600000,-1.1852855) (3.5666667,-1.1843525) (3.5733333,-1.1834246) (3.5800000,-1.1825018) (3.5866667,-1.1815841) (3.5933333,-1.1806713) (3.6000000,-1.1797634) (3.6066667,-1.1788606) (3.6133333,-1.1779626) (3.6200000,-1.1770695) (3.6266667,-1.1761812) (3.6333333,-1.1752978) (3.6400000,-1.1744191) (3.6466667,-1.1735452) (3.6533333,-1.1726760) (3.6600000,-1.1718115) (3.6666667,-1.1709517) (3.6733333,-1.1700966) (3.6800000,-1.1692460) (3.6866667,-1.1684000) (3.6933333,-1.1675586) (3.7000000,-1.1667217) (3.7066667,-1.1658893) (3.7133333,-1.1650614) (3.7200000,-1.1642379) (3.7266667,-1.1634189) (3.7333333,-1.1626042) (3.7400000,-1.1617939) (3.7466667,-1.1609879) (3.7533333,-1.1601863) (3.7600000,-1.1593889) (3.7666667,-1.1585958) (3.7733333,-1.1578069) (3.7800000,-1.1570223) (3.7866667,-1.1562418) (3.7933333,-1.1554655) (3.8000000,-1.1546933) (3.8066667,-1.1539252) (3.8133333,-1.1531612) (3.8200000,-1.1524013) (3.8266667,-1.1516454) (3.8333333,-1.1508935) (3.8400000,-1.1501456) (3.8466667,-1.1494017) (3.8533333,-1.1486617) (3.8600000,-1.1479256) (3.8666667,-1.1471935) (3.8733333,-1.1464652) (3.8800000,-1.1457407) (3.8866667,-1.1450201) (3.8933333,-1.1443032) (3.9000000,-1.1435902) (3.9066667,-1.1428809) (3.9133333,-1.1421753) (3.9200000,-1.1414735) (3.9266667,-1.1407754) (3.9333333,-1.1400809) (3.9400000,-1.1393901) (3.9466667,-1.1387029) (3.9533333,-1.1380193) (3.9600000,-1.1373393) (3.9666667,-1.1366628) (3.9733333,-1.1359899) (3.9800000,-1.1353205) (3.9866667,-1.1346547) (3.9933333,-1.1339923) (4.0000000,-1.1333333)};
\draw[violet!75!black,densely dotted] (-1,-3.0000000)--(1,3.0000000);
\fill[violet!75!black] (-1,-3.0000000) circle(2.3pt);
\fill[violet!75!black] (1,3.0000000) circle(2.3pt);
\end{scope}
\end{scope}
\node[anchor=west,align=left,text width=4.35cm] at (8.9,4.4) {蓝实线：$a=2$\\[8pt]橙虚线：$a=\frac12$\\[14pt]左右分支关于原点对称。\\[8pt]$x=0$ 为竖直渐近线；\\图像不经过原点。};
\node[text width=12.8cm,align=center] at (7.1,.65) {$D=\mathbb R\setminus\{0\}$；$x=0$ 无定义，仍然是奇函数。};
\end{tikzpicture}\end{center}

\textbf{先看图：} 图像分成两侧，靠近 $x=0$ 时向上或向下无界延伸；左右分支关于原点对称。这里 $a^x-1=0$ 当且仅当 $x=0$，因此定义域是 $\mathbb R\setminus\{0\}$。

\textbf{再说明一般情形：}
\[
f(-x)=\frac{1+a^x}{1-a^x}=-f(x)\qquad(x\ne0).
\]
与奇2比较：奇2过原点，本式以纵轴为渐近线；两者都为奇函数。
\end{explanationbox}
\clearpage
\begin{explanationbox}
\textbf{\textcolor{CExplanation}{先看真数允许在哪些位置取值}}
% Formula-based samples; see F083 checks/build_tikz_diagrams.py.
\begin{center}\begin{tikzpicture}[font=\small,>=stealth,line join=round]
\draw[gray!40,rounded corners=3pt] (0,0) rectangle (14.2,8.4);
\node at (7.1,7.75) {奇4：$y=\ln\dfrac{kx-b}{kx+b}$};
\begin{scope}[shift={(4.400000,3.900000)},x=0.850000cm,y=0.562500cm]
\fill[gray!12] (-1,-4) rectangle (1,4);
\draw[gray,dashed] (-1,-4)--(-1,4);
\draw[gray,dashed] (1,-4)--(1,4);
\draw[->,gray] (-4,0)--(4,0) node[right] {$x$};
\draw[->,gray] (0,-4)--(0,4) node[above] {$y$};
\node[anchor=north] at (-1,-4-.07/0.562500) {$-1$};
\node[anchor=north] at (1,-4-.07/0.562500) {$1$};
\draw[gray] (-2,-.035/0.562500)--(-2,.035/0.562500);
\node[anchor=north] at (-2,-.07/0.562500) {$-2$};
\draw[gray] (2,-.035/0.562500)--(2,.035/0.562500);
\node[anchor=north] at (2,-.07/0.562500) {$2$};
\begin{scope}
\clip (-4,-4) rectangle (4,4);
\draw[blue,thick,solid] plot coordinates {(-4.0000000,0.5108256) (-3.9950000,0.5114932) (-3.9900000,0.5121625) (-3.9850000,0.5128337) (-3.9800000,0.5135066) (-3.9750000,0.5141813) (-3.9700000,0.5148579) (-3.9650000,0.5155363) (-3.9600000,0.5162165) (-3.9550000,0.5168985) (-3.9500000,0.5175824) (-3.9450000,0.5182681) (-3.9400000,0.5189557) (-3.9350000,0.5196452) (-3.9300000,0.5203366) (-3.9250000,0.5210298) (-3.9200000,0.5217249) (-3.9150000,0.5224219) (-3.9100000,0.5231209) (-3.9050000,0.5238217) (-3.9000000,0.5245245) (-3.8950000,0.5252292) (-3.8900000,0.5259358) (-3.8850000,0.5266444) (-3.8800000,0.5273549) (-3.8750000,0.5280674) (-3.8700000,0.5287819) (-3.8650000,0.5294984) (-3.8600000,0.5302168) (-3.8550000,0.5309373) (-3.8500000,0.5316597) (-3.8450000,0.5323842) (-3.8400000,0.5331107) (-3.8350000,0.5338392) (-3.8300000,0.5345698) (-3.8250000,0.5353024) (-3.8200000,0.5360370) (-3.8150000,0.5367738) (-3.8100000,0.5375126) (-3.8050000,0.5382535) (-3.8000000,0.5389965) (-3.7950000,0.5397416) (-3.7900000,0.5404888) (-3.7850000,0.5412382) (-3.7800000,0.5419896) (-3.7750000,0.5427432) (-3.7700000,0.5434990) (-3.7650000,0.5442569) (-3.7600000,0.5450170) (-3.7550000,0.5457793) (-3.7500000,0.5465437) (-3.7450000,0.5473104) (-3.7400000,0.5480792) (-3.7350000,0.5488503) (-3.7300000,0.5496236) (-3.7250000,0.5503991) (-3.7200000,0.5511769) (-3.7150000,0.5519570) (-3.7100000,0.5527393) (-3.7050000,0.5535239) (-3.7000000,0.5543107) (-3.6950000,0.5550999) (-3.6900000,0.5558914) (-3.6850000,0.5566852) (-3.6800000,0.5574813) (-3.6750000,0.5582798) (-3.6700000,0.5590806) (-3.6650000,0.5598838) (-3.6600000,0.5606893) (-3.6550000,0.5614973) (-3.6500000,0.5623076) (-3.6450000,0.5631203) (-3.6400000,0.5639354) (-3.6350000,0.5647530) (-3.6300000,0.5655730) (-3.6250000,0.5663955) (-3.6200000,0.5672204) (-3.6150000,0.5680478) (-3.6100000,0.5688776) (-3.6050000,0.5697100) (-3.6000000,0.5705449) (-3.5950000,0.5713822) (-3.5900000,0.5722221) (-3.5850000,0.5730646) (-3.5800000,0.5739096) (-3.5750000,0.5747572) (-3.5700000,0.5756073) (-3.5650000,0.5764600) (-3.5600000,0.5773154) (-3.5550000,0.5781733) (-3.5500000,0.5790339) (-3.5450000,0.5798971) (-3.5400000,0.5807629) (-3.5350000,0.5816314) (-3.5300000,0.5825026) (-3.5250000,0.5833765) (-3.5200000,0.5842531) (-3.5150000,0.5851324) (-3.5100000,0.5860144) (-3.5050000,0.5868992) (-3.5000000,0.5877867) (-3.4950000,0.5886769) (-3.4900000,0.5895700) (-3.4850000,0.5904658) (-3.4800000,0.5913645) (-3.4750000,0.5922660) (-3.4700000,0.5931703) (-3.4650000,0.5940774) (-3.4600000,0.5949874) (-3.4550000,0.5959003) (-3.4500000,0.5968161) (-3.4450000,0.5977347) (-3.4400000,0.5986563) (-3.4350000,0.5995809) (-3.4300000,0.6005083) (-3.4250000,0.6014388) (-3.4200000,0.6023722) (-3.4150000,0.6033085) (-3.4100000,0.6042479) (-3.4050000,0.6051904) (-3.4000000,0.6061358) (-3.3950000,0.6070843) (-3.3900000,0.6080359) (-3.3850000,0.6089905) (-3.3800000,0.6099482) (-3.3750000,0.6109091) (-3.3700000,0.6118731) (-3.3650000,0.6128402) (-3.3600000,0.6138104) (-3.3550000,0.6147839) (-3.3500000,0.6157605) (-3.3450000,0.6167404) (-3.3400000,0.6177234) (-3.3350000,0.6187097) (-3.3300000,0.6196993) (-3.3250000,0.6206921) (-3.3200000,0.6216882) (-3.3150000,0.6226876) (-3.3100000,0.6236904) (-3.3050000,0.6246965) (-3.3000000,0.6257059) (-3.2950000,0.6267187) (-3.2900000,0.6277349) (-3.2850000,0.6287545) (-3.2800000,0.6297776) (-3.2750000,0.6308040) (-3.2700000,0.6318340) (-3.2650000,0.6328674) (-3.2600000,0.6339043) (-3.2550000,0.6349448) (-3.2500000,0.6359888) (-3.2450000,0.6370363) (-3.2400000,0.6380874) (-3.2350000,0.6391421) (-3.2300000,0.6402004) (-3.2250000,0.6412623) (-3.2200000,0.6423279) (-3.2150000,0.6433972) (-3.2100000,0.6444701) (-3.2050000,0.6455468) (-3.2000000,0.6466272) (-3.1950000,0.6477113) (-3.1900000,0.6487992) (-3.1850000,0.6498909) (-3.1800000,0.6509864) (-3.1750000,0.6520857) (-3.1700000,0.6531889) (-3.1650000,0.6542959) (-3.1600000,0.6554069) (-3.1550000,0.6565217) (-3.1500000,0.6576405) (-3.1450000,0.6587632) (-3.1400000,0.6598900) (-3.1350000,0.6610207) (-3.1300000,0.6621554) (-3.1250000,0.6632942) (-3.1200000,0.6644371) (-3.1150000,0.6655840) (-3.1100000,0.6667351) (-3.1050000,0.6678903) (-3.1000000,0.6690496) (-3.0950000,0.6702132) (-3.0900000,0.6713809) (-3.0850000,0.6725529) (-3.0800000,0.6737291) (-3.0750000,0.6749096) (-3.0700000,0.6760944) (-3.0650000,0.6772835) (-3.0600000,0.6784770) (-3.0550000,0.6796748) (-3.0500000,0.6808771) (-3.0450000,0.6820838) (-3.0400000,0.6832949) (-3.0350000,0.6845105) (-3.0300000,0.6857306) (-3.0250000,0.6869552) (-3.0200000,0.6881844) (-3.0150000,0.6894182) (-3.0100000,0.6906565) (-3.0050000,0.6918995) (-3.0000000,0.6931472) (-2.9950000,0.6943995) (-2.9900000,0.6956566) (-2.9850000,0.6969184) (-2.9800000,0.6981850) (-2.9750000,0.6994563) (-2.9700000,0.7007326) (-2.9650000,0.7020136) (-2.9600000,0.7032996) (-2.9550000,0.7045904) (-2.9500000,0.7058862) (-2.9450000,0.7071870) (-2.9400000,0.7084928) (-2.9350000,0.7098036) (-2.9300000,0.7111194) (-2.9250000,0.7124404) (-2.9200000,0.7137665) (-2.9150000,0.7150977) (-2.9100000,0.7164341) (-2.9050000,0.7177758) (-2.9000000,0.7191227) (-2.8950000,0.7204748) (-2.8900000,0.7218323) (-2.8850000,0.7231952) (-2.8800000,0.7245634) (-2.8750000,0.7259370) (-2.8700000,0.7273161) (-2.8650000,0.7287006) (-2.8600000,0.7300907) (-2.8550000,0.7314863) (-2.8500000,0.7328875) (-2.8450000,0.7342943) (-2.8400000,0.7357068) (-2.8350000,0.7371250) (-2.8300000,0.7385488) (-2.8250000,0.7399785) (-2.8200000,0.7414139) (-2.8150000,0.7428552) (-2.8100000,0.7443023) (-2.8050000,0.7457554) (-2.8000000,0.7472144) (-2.7950000,0.7486794) (-2.7900000,0.7501504) (-2.7850000,0.7516275) (-2.7800000,0.7531106) (-2.7750000,0.7546000) (-2.7700000,0.7560955) (-2.7650000,0.7575972) (-2.7600000,0.7591051) (-2.7550000,0.7606194) (-2.7500000,0.7621401) (-2.7450000,0.7636671) (-2.7400000,0.7652005) (-2.7350000,0.7667404) (-2.7300000,0.7682868) (-2.7250000,0.7698398) (-2.7200000,0.7713994) (-2.7150000,0.7729656) (-2.7100000,0.7745385) (-2.7050000,0.7761181) (-2.7000000,0.7777046) (-2.6950000,0.7792978) (-2.6900000,0.7808979) (-2.6850000,0.7825050) (-2.6800000,0.7841190) (-2.6750000,0.7857400) (-2.6700000,0.7873680) (-2.6650000,0.7890032) (-2.6600000,0.7906455) (-2.6550000,0.7922951) (-2.6500000,0.7939519) (-2.6450000,0.7956160) (-2.6400000,0.7972874) (-2.6350000,0.7989663) (-2.6300000,0.8006526) (-2.6250000,0.8023465) (-2.6200000,0.8040479) (-2.6150000,0.8057569) (-2.6100000,0.8074736) (-2.6050000,0.8091980) (-2.6000000,0.8109302) (-2.5950000,0.8126703) (-2.5900000,0.8144182) (-2.5850000,0.8161741) (-2.5800000,0.8179380) (-2.5750000,0.8197099) (-2.5700000,0.8214900) (-2.5650000,0.8232782) (-2.5600000,0.8250747) (-2.5550000,0.8268795) (-2.5500000,0.8286927) (-2.5450000,0.8305142) (-2.5400000,0.8323443) (-2.5350000,0.8341829) (-2.5300000,0.8360301) (-2.5250000,0.8378860) (-2.5200000,0.8397507) (-2.5150000,0.8416241) (-2.5100000,0.8435064) (-2.5050000,0.8453976) (-2.5000000,0.8472979) (-2.4950000,0.8492072) (-2.4900000,0.8511256) (-2.4850000,0.8530533) (-2.4800000,0.8549902) (-2.4750000,0.8569365) (-2.4700000,0.8588922) (-2.4650000,0.8608574) (-2.4600000,0.8628322) (-2.4550000,0.8648166) (-2.4500000,0.8668107) (-2.4450000,0.8688146) (-2.4400000,0.8708284) (-2.4350000,0.8728521) (-2.4300000,0.8748858) (-2.4250000,0.8769297) (-2.4200000,0.8789837) (-2.4150000,0.8810480) (-2.4100000,0.8831226) (-2.4050000,0.8852076) (-2.4000000,0.8873032) (-2.3950000,0.8894093) (-2.3900000,0.8915262) (-2.3850000,0.8936538) (-2.3800000,0.8957922) (-2.3750000,0.8979416) (-2.3700000,0.9001020) (-2.3650000,0.9022735) (-2.3600000,0.9044563) (-2.3550000,0.9066503) (-2.3500000,0.9088558) (-2.3450000,0.9110727) (-2.3400000,0.9133012) (-2.3350000,0.9155414) (-2.3300000,0.9177934) (-2.3250000,0.9200572) (-2.3200000,0.9223330) (-2.3150000,0.9246210) (-2.3100000,0.9269211) (-2.3050000,0.9292334) (-2.3000000,0.9315582) (-2.2950000,0.9338955) (-2.2900000,0.9362453) (-2.2850000,0.9386079) (-2.2800000,0.9409833) (-2.2750000,0.9433717) (-2.2700000,0.9457731) (-2.2650000,0.9481876) (-2.2600000,0.9506155) (-2.2550000,0.9530567) (-2.2500000,0.9555114) (-2.2450000,0.9579798) (-2.2400000,0.9604620) (-2.2350000,0.9629580) (-2.2300000,0.9654680) (-2.2250000,0.9679921) (-2.2200000,0.9705305) (-2.2150000,0.9730833) (-2.2100000,0.9756506) (-2.2050000,0.9782325) (-2.2000000,0.9808293) (-2.1950000,0.9834409) (-2.1900000,0.9860676) (-2.1850000,0.9887095) (-2.1800000,0.9913668) (-2.1750000,0.9940395) (-2.1700000,0.9967278) (-2.1650000,0.9994320) (-2.1600000,1.0021520) (-2.1550000,1.0048882) (-2.1500000,1.0076405) (-2.1450000,1.0104093) (-2.1400000,1.0131945) (-2.1350000,1.0159965) (-2.1300000,1.0188154) (-2.1250000,1.0216512) (-2.1200000,1.0245043) (-2.1150000,1.0273747) (-2.1100000,1.0302627) (-2.1050000,1.0331684) (-2.1000000,1.0360919) (-2.0950000,1.0390335) (-2.0900000,1.0419934) (-2.0850000,1.0449717) (-2.0800000,1.0479686) (-2.0750000,1.0509842) (-2.0700000,1.0540189) (-2.0650000,1.0570728) (-2.0600000,1.0601460) (-2.0550000,1.0632388) (-2.0500000,1.0663514) (-2.0450000,1.0694840) (-2.0400000,1.0726368) (-2.0350000,1.0758100) (-2.0300000,1.0790038) (-2.0250000,1.0822185) (-2.0200000,1.0854542) (-2.0150000,1.0887112) (-2.0100000,1.0919897) (-2.0050000,1.0952900) (-2.0000000,1.0986123) (-1.9950000,1.1019568) (-1.9900000,1.1053237) (-1.9850000,1.1087134) (-1.9800000,1.1121260) (-1.9750000,1.1155618) (-1.9700000,1.1190212) (-1.9650000,1.1225042) (-1.9600000,1.1260113) (-1.9550000,1.1295426) (-1.9500000,1.1330985) (-1.9450000,1.1366792) (-1.9400000,1.1402850) (-1.9350000,1.1439162) (-1.9300000,1.1475731) (-1.9250000,1.1512560) (-1.9200000,1.1549652) (-1.9150000,1.1587010) (-1.9100000,1.1624638) (-1.9050000,1.1662537) (-1.9000000,1.1700713) (-1.8950000,1.1739167) (-1.8900000,1.1777903) (-1.8850000,1.1816925) (-1.8800000,1.1856237) (-1.8750000,1.1895841) (-1.8700000,1.1935741) (-1.8650000,1.1975941) (-1.8600000,1.2016445) (-1.8550000,1.2057257) (-1.8500000,1.2098379) (-1.8450000,1.2139817) (-1.8400000,1.2181574) (-1.8350000,1.2223655) (-1.8300000,1.2266063) (-1.8250000,1.2308803) (-1.8200000,1.2351878) (-1.8150000,1.2395294) (-1.8100000,1.2439055) (-1.8050000,1.2483165) (-1.8000000,1.2527630) (-1.7950000,1.2572453) (-1.7900000,1.2617639) (-1.7850000,1.2663194) (-1.7800000,1.2709123) (-1.7750000,1.2755430) (-1.7700000,1.2802121) (-1.7650000,1.2849201) (-1.7600000,1.2896675) (-1.7550000,1.2944550) (-1.7500000,1.2992830) (-1.7450000,1.3041521) (-1.7400000,1.3090630) (-1.7350000,1.3140162) (-1.7300000,1.3190124) (-1.7250000,1.3240521) (-1.7200000,1.3291359) (-1.7150000,1.3342647) (-1.7100000,1.3394389) (-1.7050000,1.3446594) (-1.7000000,1.3499267) (-1.6950000,1.3552416) (-1.6900000,1.3606049) (-1.6850000,1.3660172) (-1.6800000,1.3714793) (-1.6750000,1.3769920) (-1.6700000,1.3825560) (-1.6650000,1.3881723) (-1.6600000,1.3938416) (-1.6550000,1.3995647) (-1.6500000,1.4053426) (-1.6450000,1.4111760) (-1.6400000,1.4170660) (-1.6350000,1.4230135) (-1.6300000,1.4290193) (-1.6250000,1.4350845) (-1.6200000,1.4412101) (-1.6150000,1.4473971) (-1.6100000,1.4536465) (-1.6050000,1.4599595) (-1.6000000,1.4663371) (-1.5950000,1.4727804) (-1.5900000,1.4792906) (-1.5850000,1.4858689) (-1.5800000,1.4925166) (-1.5750000,1.4992348) (-1.5700000,1.5060248) (-1.5650000,1.5128880) (-1.5600000,1.5198258) (-1.5550000,1.5268394) (-1.5500000,1.5339304) (-1.5450000,1.5411001) (-1.5400000,1.5483502) (-1.5350000,1.5556822) (-1.5300000,1.5630976) (-1.5250000,1.5705981) (-1.5200000,1.5781854) (-1.5150000,1.5858612) (-1.5100000,1.5936273) (-1.5050000,1.6014856) (-1.5000000,1.6094379) (-1.4950000,1.6174862) (-1.4900000,1.6256326) (-1.4850000,1.6338790) (-1.4800000,1.6422277) (-1.4750000,1.6506809) (-1.4700000,1.6592407) (-1.4650000,1.6679097) (-1.4600000,1.6766901) (-1.4550000,1.6855846) (-1.4500000,1.6945957) (-1.4450000,1.7037261) (-1.4400000,1.7129786) (-1.4350000,1.7223560) (-1.4300000,1.7318613) (-1.4250000,1.7414976) (-1.4200000,1.7512681) (-1.4150000,1.7611760) (-1.4100000,1.7712249) (-1.4050000,1.7814181) (-1.4000000,1.7917595) (-1.3950000,1.8022527) (-1.3900000,1.8129019) (-1.3850000,1.8237111) (-1.3800000,1.8346845) (-1.3750000,1.8458267) (-1.3700000,1.8571422) (-1.3650000,1.8686359) (-1.3600000,1.8803129) (-1.3550000,1.8921782) (-1.3500000,1.9042375) (-1.3450000,1.9164963) (-1.3400000,1.9289606) (-1.3350000,1.9416366) (-1.3300000,1.9545309) (-1.3250000,1.9676501) (-1.3200000,1.9810015) (-1.3150000,1.9945923) (-1.3100000,2.0084305) (-1.3050000,2.0225242) (-1.3000000,2.0368819) (-1.2950000,2.0515128) (-1.2900000,2.0664262) (-1.2850000,2.0816321) (-1.2800000,2.0971411) (-1.2750000,2.1129642) (-1.2700000,2.1291132) (-1.2650000,2.1456002) (-1.2600000,2.1624385) (-1.2550000,2.1796417) (-1.2500000,2.1972246) (-1.2450000,2.2152026) (-1.2400000,2.2335922) (-1.2350000,2.2524110) (-1.2300000,2.2716776) (-1.2250000,2.2914118) (-1.2200000,2.3116349) (-1.2150000,2.3323697) (-1.2100000,2.3536403) (-1.2050000,2.3754728) (-1.2000000,2.3978953) (-1.1950000,2.4209378) (-1.1900000,2.4446328) (-1.1850000,2.4690153) (-1.1800000,2.4941233) (-1.1750000,2.5199980) (-1.1700000,2.5466840) (-1.1650000,2.5742302) (-1.1600000,2.6026897) (-1.1550000,2.6321209) (-1.1500000,2.6625878) (-1.1450000,2.6941611) (-1.1400000,2.7269187) (-1.1350000,2.7609471) (-1.1300000,2.7963428) (-1.1250000,2.8332133) (-1.1200000,2.8716796) (-1.1150000,2.9118780) (-1.1100000,2.9539629) (-1.1050000,2.9981104) (-1.1000000,3.0445224) (-1.0950000,3.0934319) (-1.0900000,3.1451097) (-1.0850000,3.1998729) (-1.0800000,3.2580965) (-1.0750000,3.3202283) (-1.0700000,3.3868086) (-1.0650000,3.4584982) (-1.0600000,3.5361167) (-1.0550000,3.6206979) (-1.0500000,3.7135721) (-1.0450000,3.8164906) (-1.0400000,3.9318256)};
\draw[blue,thick,solid] plot coordinates {(1.0400000,-3.9318256) (1.0450000,-3.8164906) (1.0500000,-3.7135721) (1.0550000,-3.6206979) (1.0600000,-3.5361167) (1.0650000,-3.4584982) (1.0700000,-3.3868086) (1.0750000,-3.3202283) (1.0800000,-3.2580965) (1.0850000,-3.1998729) (1.0900000,-3.1451097) (1.0950000,-3.0934319) (1.1000000,-3.0445224) (1.1050000,-2.9981104) (1.1100000,-2.9539629) (1.1150000,-2.9118780) (1.1200000,-2.8716796) (1.1250000,-2.8332133) (1.1300000,-2.7963428) (1.1350000,-2.7609471) (1.1400000,-2.7269187) (1.1450000,-2.6941611) (1.1500000,-2.6625878) (1.1550000,-2.6321209) (1.1600000,-2.6026897) (1.1650000,-2.5742302) (1.1700000,-2.5466840) (1.1750000,-2.5199980) (1.1800000,-2.4941233) (1.1850000,-2.4690153) (1.1900000,-2.4446328) (1.1950000,-2.4209378) (1.2000000,-2.3978953) (1.2050000,-2.3754728) (1.2100000,-2.3536403) (1.2150000,-2.3323697) (1.2200000,-2.3116349) (1.2250000,-2.2914118) (1.2300000,-2.2716776) (1.2350000,-2.2524110) (1.2400000,-2.2335922) (1.2450000,-2.2152026) (1.2500000,-2.1972246) (1.2550000,-2.1796417) (1.2600000,-2.1624385) (1.2650000,-2.1456002) (1.2700000,-2.1291132) (1.2750000,-2.1129642) (1.2800000,-2.0971411) (1.2850000,-2.0816321) (1.2900000,-2.0664262) (1.2950000,-2.0515128) (1.3000000,-2.0368819) (1.3050000,-2.0225242) (1.3100000,-2.0084305) (1.3150000,-1.9945923) (1.3200000,-1.9810015) (1.3250000,-1.9676501) (1.3300000,-1.9545309) (1.3350000,-1.9416366) (1.3400000,-1.9289606) (1.3450000,-1.9164963) (1.3500000,-1.9042375) (1.3550000,-1.8921782) (1.3600000,-1.8803129) (1.3650000,-1.8686359) (1.3700000,-1.8571422) (1.3750000,-1.8458267) (1.3800000,-1.8346845) (1.3850000,-1.8237111) (1.3900000,-1.8129019) (1.3950000,-1.8022527) (1.4000000,-1.7917595) (1.4050000,-1.7814181) (1.4100000,-1.7712249) (1.4150000,-1.7611760) (1.4200000,-1.7512681) (1.4250000,-1.7414976) (1.4300000,-1.7318613) (1.4350000,-1.7223560) (1.4400000,-1.7129786) (1.4450000,-1.7037261) (1.4500000,-1.6945957) (1.4550000,-1.6855846) (1.4600000,-1.6766901) (1.4650000,-1.6679097) (1.4700000,-1.6592407) (1.4750000,-1.6506809) (1.4800000,-1.6422277) (1.4850000,-1.6338790) (1.4900000,-1.6256326) (1.4950000,-1.6174862) (1.5000000,-1.6094379) (1.5050000,-1.6014856) (1.5100000,-1.5936273) (1.5150000,-1.5858612) (1.5200000,-1.5781854) (1.5250000,-1.5705981) (1.5300000,-1.5630976) (1.5350000,-1.5556822) (1.5400000,-1.5483502) (1.5450000,-1.5411001) (1.5500000,-1.5339304) (1.5550000,-1.5268394) (1.5600000,-1.5198258) (1.5650000,-1.5128880) (1.5700000,-1.5060248) (1.5750000,-1.4992348) (1.5800000,-1.4925166) (1.5850000,-1.4858689) (1.5900000,-1.4792906) (1.5950000,-1.4727804) (1.6000000,-1.4663371) (1.6050000,-1.4599595) (1.6100000,-1.4536465) (1.6150000,-1.4473971) (1.6200000,-1.4412101) (1.6250000,-1.4350845) (1.6300000,-1.4290193) (1.6350000,-1.4230135) (1.6400000,-1.4170660) (1.6450000,-1.4111760) (1.6500000,-1.4053426) (1.6550000,-1.3995647) (1.6600000,-1.3938416) (1.6650000,-1.3881723) (1.6700000,-1.3825560) (1.6750000,-1.3769920) (1.6800000,-1.3714793) (1.6850000,-1.3660172) (1.6900000,-1.3606049) (1.6950000,-1.3552416) (1.7000000,-1.3499267) (1.7050000,-1.3446594) (1.7100000,-1.3394389) (1.7150000,-1.3342647) (1.7200000,-1.3291359) (1.7250000,-1.3240521) (1.7300000,-1.3190124) (1.7350000,-1.3140162) (1.7400000,-1.3090630) (1.7450000,-1.3041521) (1.7500000,-1.2992830) (1.7550000,-1.2944550) (1.7600000,-1.2896675) (1.7650000,-1.2849201) (1.7700000,-1.2802121) (1.7750000,-1.2755430) (1.7800000,-1.2709123) (1.7850000,-1.2663194) (1.7900000,-1.2617639) (1.7950000,-1.2572453) (1.8000000,-1.2527630) (1.8050000,-1.2483165) (1.8100000,-1.2439055) (1.8150000,-1.2395294) (1.8200000,-1.2351878) (1.8250000,-1.2308803) (1.8300000,-1.2266063) (1.8350000,-1.2223655) (1.8400000,-1.2181574) (1.8450000,-1.2139817) (1.8500000,-1.2098379) (1.8550000,-1.2057257) (1.8600000,-1.2016445) (1.8650000,-1.1975941) (1.8700000,-1.1935741) (1.8750000,-1.1895841) (1.8800000,-1.1856237) (1.8850000,-1.1816925) (1.8900000,-1.1777903) (1.8950000,-1.1739167) (1.9000000,-1.1700713) (1.9050000,-1.1662537) (1.9100000,-1.1624638) (1.9150000,-1.1587010) (1.9200000,-1.1549652) (1.9250000,-1.1512560) (1.9300000,-1.1475731) (1.9350000,-1.1439162) (1.9400000,-1.1402850) (1.9450000,-1.1366792) (1.9500000,-1.1330985) (1.9550000,-1.1295426) (1.9600000,-1.1260113) (1.9650000,-1.1225042) (1.9700000,-1.1190212) (1.9750000,-1.1155618) (1.9800000,-1.1121260) (1.9850000,-1.1087134) (1.9900000,-1.1053237) (1.9950000,-1.1019568) (2.0000000,-1.0986123) (2.0050000,-1.0952900) (2.0100000,-1.0919897) (2.0150000,-1.0887112) (2.0200000,-1.0854542) (2.0250000,-1.0822185) (2.0300000,-1.0790038) (2.0350000,-1.0758100) (2.0400000,-1.0726368) (2.0450000,-1.0694840) (2.0500000,-1.0663514) (2.0550000,-1.0632388) (2.0600000,-1.0601460) (2.0650000,-1.0570728) (2.0700000,-1.0540189) (2.0750000,-1.0509842) (2.0800000,-1.0479686) (2.0850000,-1.0449717) (2.0900000,-1.0419934) (2.0950000,-1.0390335) (2.1000000,-1.0360919) (2.1050000,-1.0331684) (2.1100000,-1.0302627) (2.1150000,-1.0273747) (2.1200000,-1.0245043) (2.1250000,-1.0216512) (2.1300000,-1.0188154) (2.1350000,-1.0159965) (2.1400000,-1.0131945) (2.1450000,-1.0104093) (2.1500000,-1.0076405) (2.1550000,-1.0048882) (2.1600000,-1.0021520) (2.1650000,-0.9994320) (2.1700000,-0.9967278) (2.1750000,-0.9940395) (2.1800000,-0.9913668) (2.1850000,-0.9887095) (2.1900000,-0.9860676) (2.1950000,-0.9834409) (2.2000000,-0.9808293) (2.2050000,-0.9782325) (2.2100000,-0.9756506) (2.2150000,-0.9730833) (2.2200000,-0.9705305) (2.2250000,-0.9679921) (2.2300000,-0.9654680) (2.2350000,-0.9629580) (2.2400000,-0.9604620) (2.2450000,-0.9579798) (2.2500000,-0.9555114) (2.2550000,-0.9530567) (2.2600000,-0.9506155) (2.2650000,-0.9481876) (2.2700000,-0.9457731) (2.2750000,-0.9433717) (2.2800000,-0.9409833) (2.2850000,-0.9386079) (2.2900000,-0.9362453) (2.2950000,-0.9338955) (2.3000000,-0.9315582) (2.3050000,-0.9292334) (2.3100000,-0.9269211) (2.3150000,-0.9246210) (2.3200000,-0.9223330) (2.3250000,-0.9200572) (2.3300000,-0.9177934) (2.3350000,-0.9155414) (2.3400000,-0.9133012) (2.3450000,-0.9110727) (2.3500000,-0.9088558) (2.3550000,-0.9066503) (2.3600000,-0.9044563) (2.3650000,-0.9022735) (2.3700000,-0.9001020) (2.3750000,-0.8979416) (2.3800000,-0.8957922) (2.3850000,-0.8936538) (2.3900000,-0.8915262) (2.3950000,-0.8894093) (2.4000000,-0.8873032) (2.4050000,-0.8852076) (2.4100000,-0.8831226) (2.4150000,-0.8810480) (2.4200000,-0.8789837) (2.4250000,-0.8769297) (2.4300000,-0.8748858) (2.4350000,-0.8728521) (2.4400000,-0.8708284) (2.4450000,-0.8688146) (2.4500000,-0.8668107) (2.4550000,-0.8648166) (2.4600000,-0.8628322) (2.4650000,-0.8608574) (2.4700000,-0.8588922) (2.4750000,-0.8569365) (2.4800000,-0.8549902) (2.4850000,-0.8530533) (2.4900000,-0.8511256) (2.4950000,-0.8492072) (2.5000000,-0.8472979) (2.5050000,-0.8453976) (2.5100000,-0.8435064) (2.5150000,-0.8416241) (2.5200000,-0.8397507) (2.5250000,-0.8378860) (2.5300000,-0.8360301) (2.5350000,-0.8341829) (2.5400000,-0.8323443) (2.5450000,-0.8305142) (2.5500000,-0.8286927) (2.5550000,-0.8268795) (2.5600000,-0.8250747) (2.5650000,-0.8232782) (2.5700000,-0.8214900) (2.5750000,-0.8197099) (2.5800000,-0.8179380) (2.5850000,-0.8161741) (2.5900000,-0.8144182) (2.5950000,-0.8126703) (2.6000000,-0.8109302) (2.6050000,-0.8091980) (2.6100000,-0.8074736) (2.6150000,-0.8057569) (2.6200000,-0.8040479) (2.6250000,-0.8023465) (2.6300000,-0.8006526) (2.6350000,-0.7989663) (2.6400000,-0.7972874) (2.6450000,-0.7956160) (2.6500000,-0.7939519) (2.6550000,-0.7922951) (2.6600000,-0.7906455) (2.6650000,-0.7890032) (2.6700000,-0.7873680) (2.6750000,-0.7857400) (2.6800000,-0.7841190) (2.6850000,-0.7825050) (2.6900000,-0.7808979) (2.6950000,-0.7792978) (2.7000000,-0.7777046) (2.7050000,-0.7761181) (2.7100000,-0.7745385) (2.7150000,-0.7729656) (2.7200000,-0.7713994) (2.7250000,-0.7698398) (2.7300000,-0.7682868) (2.7350000,-0.7667404) (2.7400000,-0.7652005) (2.7450000,-0.7636671) (2.7500000,-0.7621401) (2.7550000,-0.7606194) (2.7600000,-0.7591051) (2.7650000,-0.7575972) (2.7700000,-0.7560955) (2.7750000,-0.7546000) (2.7800000,-0.7531106) (2.7850000,-0.7516275) (2.7900000,-0.7501504) (2.7950000,-0.7486794) (2.8000000,-0.7472144) (2.8050000,-0.7457554) (2.8100000,-0.7443023) (2.8150000,-0.7428552) (2.8200000,-0.7414139) (2.8250000,-0.7399785) (2.8300000,-0.7385488) (2.8350000,-0.7371250) (2.8400000,-0.7357068) (2.8450000,-0.7342943) (2.8500000,-0.7328875) (2.8550000,-0.7314863) (2.8600000,-0.7300907) (2.8650000,-0.7287006) (2.8700000,-0.7273161) (2.8750000,-0.7259370) (2.8800000,-0.7245634) (2.8850000,-0.7231952) (2.8900000,-0.7218323) (2.8950000,-0.7204748) (2.9000000,-0.7191227) (2.9050000,-0.7177758) (2.9100000,-0.7164341) (2.9150000,-0.7150977) (2.9200000,-0.7137665) (2.9250000,-0.7124404) (2.9300000,-0.7111194) (2.9350000,-0.7098036) (2.9400000,-0.7084928) (2.9450000,-0.7071870) (2.9500000,-0.7058862) (2.9550000,-0.7045904) (2.9600000,-0.7032996) (2.9650000,-0.7020136) (2.9700000,-0.7007326) (2.9750000,-0.6994563) (2.9800000,-0.6981850) (2.9850000,-0.6969184) (2.9900000,-0.6956566) (2.9950000,-0.6943995) (3.0000000,-0.6931472) (3.0050000,-0.6918995) (3.0100000,-0.6906565) (3.0150000,-0.6894182) (3.0200000,-0.6881844) (3.0250000,-0.6869552) (3.0300000,-0.6857306) (3.0350000,-0.6845105) (3.0400000,-0.6832949) (3.0450000,-0.6820838) (3.0500000,-0.6808771) (3.0550000,-0.6796748) (3.0600000,-0.6784770) (3.0650000,-0.6772835) (3.0700000,-0.6760944) (3.0750000,-0.6749096) (3.0800000,-0.6737291) (3.0850000,-0.6725529) (3.0900000,-0.6713809) (3.0950000,-0.6702132) (3.1000000,-0.6690496) (3.1050000,-0.6678903) (3.1100000,-0.6667351) (3.1150000,-0.6655840) (3.1200000,-0.6644371) (3.1250000,-0.6632942) (3.1300000,-0.6621554) (3.1350000,-0.6610207) (3.1400000,-0.6598900) (3.1450000,-0.6587632) (3.1500000,-0.6576405) (3.1550000,-0.6565217) (3.1600000,-0.6554069) (3.1650000,-0.6542959) (3.1700000,-0.6531889) (3.1750000,-0.6520857) (3.1800000,-0.6509864) (3.1850000,-0.6498909) (3.1900000,-0.6487992) (3.1950000,-0.6477113) (3.2000000,-0.6466272) (3.2050000,-0.6455468) (3.2100000,-0.6444701) (3.2150000,-0.6433972) (3.2200000,-0.6423279) (3.2250000,-0.6412623) (3.2300000,-0.6402004) (3.2350000,-0.6391421) (3.2400000,-0.6380874) (3.2450000,-0.6370363) (3.2500000,-0.6359888) (3.2550000,-0.6349448) (3.2600000,-0.6339043) (3.2650000,-0.6328674) (3.2700000,-0.6318340) (3.2750000,-0.6308040) (3.2800000,-0.6297776) (3.2850000,-0.6287545) (3.2900000,-0.6277349) (3.2950000,-0.6267187) (3.3000000,-0.6257059) (3.3050000,-0.6246965) (3.3100000,-0.6236904) (3.3150000,-0.6226876) (3.3200000,-0.6216882) (3.3250000,-0.6206921) (3.3300000,-0.6196993) (3.3350000,-0.6187097) (3.3400000,-0.6177234) (3.3450000,-0.6167404) (3.3500000,-0.6157605) (3.3550000,-0.6147839) (3.3600000,-0.6138104) (3.3650000,-0.6128402) (3.3700000,-0.6118731) (3.3750000,-0.6109091) (3.3800000,-0.6099482) (3.3850000,-0.6089905) (3.3900000,-0.6080359) (3.3950000,-0.6070843) (3.4000000,-0.6061358) (3.4050000,-0.6051904) (3.4100000,-0.6042479) (3.4150000,-0.6033085) (3.4200000,-0.6023722) (3.4250000,-0.6014388) (3.4300000,-0.6005083) (3.4350000,-0.5995809) (3.4400000,-0.5986563) (3.4450000,-0.5977347) (3.4500000,-0.5968161) (3.4550000,-0.5959003) (3.4600000,-0.5949874) (3.4650000,-0.5940774) (3.4700000,-0.5931703) (3.4750000,-0.5922660) (3.4800000,-0.5913645) (3.4850000,-0.5904658) (3.4900000,-0.5895700) (3.4950000,-0.5886769) (3.5000000,-0.5877867) (3.5050000,-0.5868992) (3.5100000,-0.5860144) (3.5150000,-0.5851324) (3.5200000,-0.5842531) (3.5250000,-0.5833765) (3.5300000,-0.5825026) (3.5350000,-0.5816314) (3.5400000,-0.5807629) (3.5450000,-0.5798971) (3.5500000,-0.5790339) (3.5550000,-0.5781733) (3.5600000,-0.5773154) (3.5650000,-0.5764600) (3.5700000,-0.5756073) (3.5750000,-0.5747572) (3.5800000,-0.5739096) (3.5850000,-0.5730646) (3.5900000,-0.5722221) (3.5950000,-0.5713822) (3.6000000,-0.5705449) (3.6050000,-0.5697100) (3.6100000,-0.5688776) (3.6150000,-0.5680478) (3.6200000,-0.5672204) (3.6250000,-0.5663955) (3.6300000,-0.5655730) (3.6350000,-0.5647530) (3.6400000,-0.5639354) (3.6450000,-0.5631203) (3.6500000,-0.5623076) (3.6550000,-0.5614973) (3.6600000,-0.5606893) (3.6650000,-0.5598838) (3.6700000,-0.5590806) (3.6750000,-0.5582798) (3.6800000,-0.5574813) (3.6850000,-0.5566852) (3.6900000,-0.5558914) (3.6950000,-0.5550999) (3.7000000,-0.5543107) (3.7050000,-0.5535239) (3.7100000,-0.5527393) (3.7150000,-0.5519570) (3.7200000,-0.5511769) (3.7250000,-0.5503991) (3.7300000,-0.5496236) (3.7350000,-0.5488503) (3.7400000,-0.5480792) (3.7450000,-0.5473104) (3.7500000,-0.5465437) (3.7550000,-0.5457793) (3.7600000,-0.5450170) (3.7650000,-0.5442569) (3.7700000,-0.5434990) (3.7750000,-0.5427432) (3.7800000,-0.5419896) (3.7850000,-0.5412382) (3.7900000,-0.5404888) (3.7950000,-0.5397416) (3.8000000,-0.5389965) (3.8050000,-0.5382535) (3.8100000,-0.5375126) (3.8150000,-0.5367738) (3.8200000,-0.5360370) (3.8250000,-0.5353024) (3.8300000,-0.5345698) (3.8350000,-0.5338392) (3.8400000,-0.5331107) (3.8450000,-0.5323842) (3.8500000,-0.5316597) (3.8550000,-0.5309373) (3.8600000,-0.5302168) (3.8650000,-0.5294984) (3.8700000,-0.5287819) (3.8750000,-0.5280674) (3.8800000,-0.5273549) (3.8850000,-0.5266444) (3.8900000,-0.5259358) (3.8950000,-0.5252292) (3.9000000,-0.5245245) (3.9050000,-0.5238217) (3.9100000,-0.5231209) (3.9150000,-0.5224219) (3.9200000,-0.5217249) (3.9250000,-0.5210298) (3.9300000,-0.5203366) (3.9350000,-0.5196452) (3.9400000,-0.5189557) (3.9450000,-0.5182681) (3.9500000,-0.5175824) (3.9550000,-0.5168985) (3.9600000,-0.5162165) (3.9650000,-0.5155363) (3.9700000,-0.5148579) (3.9750000,-0.5141813) (3.9800000,-0.5135066) (3.9850000,-0.5128337) (3.9900000,-0.5121625) (3.9950000,-0.5114932) (4.0000000,-0.5108256)};
\draw[orange!85!black,thick,dashed] plot coordinates {(-4.0000000,-0.5108256) (-3.9950000,-0.5114932) (-3.9900000,-0.5121625) (-3.9850000,-0.5128337) (-3.9800000,-0.5135066) (-3.9750000,-0.5141813) (-3.9700000,-0.5148579) (-3.9650000,-0.5155363) (-3.9600000,-0.5162165) (-3.9550000,-0.5168985) (-3.9500000,-0.5175824) (-3.9450000,-0.5182681) (-3.9400000,-0.5189557) (-3.9350000,-0.5196452) (-3.9300000,-0.5203366) (-3.9250000,-0.5210298) (-3.9200000,-0.5217249) (-3.9150000,-0.5224219) (-3.9100000,-0.5231209) (-3.9050000,-0.5238217) (-3.9000000,-0.5245245) (-3.8950000,-0.5252292) (-3.8900000,-0.5259358) (-3.8850000,-0.5266444) (-3.8800000,-0.5273549) (-3.8750000,-0.5280674) (-3.8700000,-0.5287819) (-3.8650000,-0.5294984) (-3.8600000,-0.5302168) (-3.8550000,-0.5309373) (-3.8500000,-0.5316597) (-3.8450000,-0.5323842) (-3.8400000,-0.5331107) (-3.8350000,-0.5338392) (-3.8300000,-0.5345698) (-3.8250000,-0.5353024) (-3.8200000,-0.5360370) (-3.8150000,-0.5367738) (-3.8100000,-0.5375126) (-3.8050000,-0.5382535) (-3.8000000,-0.5389965) (-3.7950000,-0.5397416) (-3.7900000,-0.5404888) (-3.7850000,-0.5412382) (-3.7800000,-0.5419896) (-3.7750000,-0.5427432) (-3.7700000,-0.5434990) (-3.7650000,-0.5442569) (-3.7600000,-0.5450170) (-3.7550000,-0.5457793) (-3.7500000,-0.5465437) (-3.7450000,-0.5473104) (-3.7400000,-0.5480792) (-3.7350000,-0.5488503) (-3.7300000,-0.5496236) (-3.7250000,-0.5503991) (-3.7200000,-0.5511769) (-3.7150000,-0.5519570) (-3.7100000,-0.5527393) (-3.7050000,-0.5535239) (-3.7000000,-0.5543107) (-3.6950000,-0.5550999) (-3.6900000,-0.5558914) (-3.6850000,-0.5566852) (-3.6800000,-0.5574813) (-3.6750000,-0.5582798) (-3.6700000,-0.5590806) (-3.6650000,-0.5598838) (-3.6600000,-0.5606893) (-3.6550000,-0.5614973) (-3.6500000,-0.5623076) (-3.6450000,-0.5631203) (-3.6400000,-0.5639354) (-3.6350000,-0.5647530) (-3.6300000,-0.5655730) (-3.6250000,-0.5663955) (-3.6200000,-0.5672204) (-3.6150000,-0.5680478) (-3.6100000,-0.5688776) (-3.6050000,-0.5697100) (-3.6000000,-0.5705449) (-3.5950000,-0.5713822) (-3.5900000,-0.5722221) (-3.5850000,-0.5730646) (-3.5800000,-0.5739096) (-3.5750000,-0.5747572) (-3.5700000,-0.5756073) (-3.5650000,-0.5764600) (-3.5600000,-0.5773154) (-3.5550000,-0.5781733) (-3.5500000,-0.5790339) (-3.5450000,-0.5798971) (-3.5400000,-0.5807629) (-3.5350000,-0.5816314) (-3.5300000,-0.5825026) (-3.5250000,-0.5833765) (-3.5200000,-0.5842531) (-3.5150000,-0.5851324) (-3.5100000,-0.5860144) (-3.5050000,-0.5868992) (-3.5000000,-0.5877867) (-3.4950000,-0.5886769) (-3.4900000,-0.5895700) (-3.4850000,-0.5904658) (-3.4800000,-0.5913645) (-3.4750000,-0.5922660) (-3.4700000,-0.5931703) (-3.4650000,-0.5940774) (-3.4600000,-0.5949874) (-3.4550000,-0.5959003) (-3.4500000,-0.5968161) (-3.4450000,-0.5977347) (-3.4400000,-0.5986563) (-3.4350000,-0.5995809) (-3.4300000,-0.6005083) (-3.4250000,-0.6014388) (-3.4200000,-0.6023722) (-3.4150000,-0.6033085) (-3.4100000,-0.6042479) (-3.4050000,-0.6051904) (-3.4000000,-0.6061358) (-3.3950000,-0.6070843) (-3.3900000,-0.6080359) (-3.3850000,-0.6089905) (-3.3800000,-0.6099482) (-3.3750000,-0.6109091) (-3.3700000,-0.6118731) (-3.3650000,-0.6128402) (-3.3600000,-0.6138104) (-3.3550000,-0.6147839) (-3.3500000,-0.6157605) (-3.3450000,-0.6167404) (-3.3400000,-0.6177234) (-3.3350000,-0.6187097) (-3.3300000,-0.6196993) (-3.3250000,-0.6206921) (-3.3200000,-0.6216882) (-3.3150000,-0.6226876) (-3.3100000,-0.6236904) (-3.3050000,-0.6246965) (-3.3000000,-0.6257059) (-3.2950000,-0.6267187) (-3.2900000,-0.6277349) (-3.2850000,-0.6287545) (-3.2800000,-0.6297776) (-3.2750000,-0.6308040) (-3.2700000,-0.6318340) (-3.2650000,-0.6328674) (-3.2600000,-0.6339043) (-3.2550000,-0.6349448) (-3.2500000,-0.6359888) (-3.2450000,-0.6370363) (-3.2400000,-0.6380874) (-3.2350000,-0.6391421) (-3.2300000,-0.6402004) (-3.2250000,-0.6412623) (-3.2200000,-0.6423279) (-3.2150000,-0.6433972) (-3.2100000,-0.6444701) (-3.2050000,-0.6455468) (-3.2000000,-0.6466272) (-3.1950000,-0.6477113) (-3.1900000,-0.6487992) (-3.1850000,-0.6498909) (-3.1800000,-0.6509864) (-3.1750000,-0.6520857) (-3.1700000,-0.6531889) (-3.1650000,-0.6542959) (-3.1600000,-0.6554069) (-3.1550000,-0.6565217) (-3.1500000,-0.6576405) (-3.1450000,-0.6587632) (-3.1400000,-0.6598900) (-3.1350000,-0.6610207) (-3.1300000,-0.6621554) (-3.1250000,-0.6632942) (-3.1200000,-0.6644371) (-3.1150000,-0.6655840) (-3.1100000,-0.6667351) (-3.1050000,-0.6678903) (-3.1000000,-0.6690496) (-3.0950000,-0.6702132) (-3.0900000,-0.6713809) (-3.0850000,-0.6725529) (-3.0800000,-0.6737291) (-3.0750000,-0.6749096) (-3.0700000,-0.6760944) (-3.0650000,-0.6772835) (-3.0600000,-0.6784770) (-3.0550000,-0.6796748) (-3.0500000,-0.6808771) (-3.0450000,-0.6820838) (-3.0400000,-0.6832949) (-3.0350000,-0.6845105) (-3.0300000,-0.6857306) (-3.0250000,-0.6869552) (-3.0200000,-0.6881844) (-3.0150000,-0.6894182) (-3.0100000,-0.6906565) (-3.0050000,-0.6918995) (-3.0000000,-0.6931472) (-2.9950000,-0.6943995) (-2.9900000,-0.6956566) (-2.9850000,-0.6969184) (-2.9800000,-0.6981850) (-2.9750000,-0.6994563) (-2.9700000,-0.7007326) (-2.9650000,-0.7020136) (-2.9600000,-0.7032996) (-2.9550000,-0.7045904) (-2.9500000,-0.7058862) (-2.9450000,-0.7071870) (-2.9400000,-0.7084928) (-2.9350000,-0.7098036) (-2.9300000,-0.7111194) (-2.9250000,-0.7124404) (-2.9200000,-0.7137665) (-2.9150000,-0.7150977) (-2.9100000,-0.7164341) (-2.9050000,-0.7177758) (-2.9000000,-0.7191227) (-2.8950000,-0.7204748) (-2.8900000,-0.7218323) (-2.8850000,-0.7231952) (-2.8800000,-0.7245634) (-2.8750000,-0.7259370) (-2.8700000,-0.7273161) (-2.8650000,-0.7287006) (-2.8600000,-0.7300907) (-2.8550000,-0.7314863) (-2.8500000,-0.7328875) (-2.8450000,-0.7342943) (-2.8400000,-0.7357068) (-2.8350000,-0.7371250) (-2.8300000,-0.7385488) (-2.8250000,-0.7399785) (-2.8200000,-0.7414139) (-2.8150000,-0.7428552) (-2.8100000,-0.7443023) (-2.8050000,-0.7457554) (-2.8000000,-0.7472144) (-2.7950000,-0.7486794) (-2.7900000,-0.7501504) (-2.7850000,-0.7516275) (-2.7800000,-0.7531106) (-2.7750000,-0.7546000) (-2.7700000,-0.7560955) (-2.7650000,-0.7575972) (-2.7600000,-0.7591051) (-2.7550000,-0.7606194) (-2.7500000,-0.7621401) (-2.7450000,-0.7636671) (-2.7400000,-0.7652005) (-2.7350000,-0.7667404) (-2.7300000,-0.7682868) (-2.7250000,-0.7698398) (-2.7200000,-0.7713994) (-2.7150000,-0.7729656) (-2.7100000,-0.7745385) (-2.7050000,-0.7761181) (-2.7000000,-0.7777046) (-2.6950000,-0.7792978) (-2.6900000,-0.7808979) (-2.6850000,-0.7825050) (-2.6800000,-0.7841190) (-2.6750000,-0.7857400) (-2.6700000,-0.7873680) (-2.6650000,-0.7890032) (-2.6600000,-0.7906455) (-2.6550000,-0.7922951) (-2.6500000,-0.7939519) (-2.6450000,-0.7956160) (-2.6400000,-0.7972874) (-2.6350000,-0.7989663) (-2.6300000,-0.8006526) (-2.6250000,-0.8023465) (-2.6200000,-0.8040479) (-2.6150000,-0.8057569) (-2.6100000,-0.8074736) (-2.6050000,-0.8091980) (-2.6000000,-0.8109302) (-2.5950000,-0.8126703) (-2.5900000,-0.8144182) (-2.5850000,-0.8161741) (-2.5800000,-0.8179380) (-2.5750000,-0.8197099) (-2.5700000,-0.8214900) (-2.5650000,-0.8232782) (-2.5600000,-0.8250747) (-2.5550000,-0.8268795) (-2.5500000,-0.8286927) (-2.5450000,-0.8305142) (-2.5400000,-0.8323443) (-2.5350000,-0.8341829) (-2.5300000,-0.8360301) (-2.5250000,-0.8378860) (-2.5200000,-0.8397507) (-2.5150000,-0.8416241) (-2.5100000,-0.8435064) (-2.5050000,-0.8453976) (-2.5000000,-0.8472979) (-2.4950000,-0.8492072) (-2.4900000,-0.8511256) (-2.4850000,-0.8530533) (-2.4800000,-0.8549902) (-2.4750000,-0.8569365) (-2.4700000,-0.8588922) (-2.4650000,-0.8608574) (-2.4600000,-0.8628322) (-2.4550000,-0.8648166) (-2.4500000,-0.8668107) (-2.4450000,-0.8688146) (-2.4400000,-0.8708284) (-2.4350000,-0.8728521) (-2.4300000,-0.8748858) (-2.4250000,-0.8769297) (-2.4200000,-0.8789837) (-2.4150000,-0.8810480) (-2.4100000,-0.8831226) (-2.4050000,-0.8852076) (-2.4000000,-0.8873032) (-2.3950000,-0.8894093) (-2.3900000,-0.8915262) (-2.3850000,-0.8936538) (-2.3800000,-0.8957922) (-2.3750000,-0.8979416) (-2.3700000,-0.9001020) (-2.3650000,-0.9022735) (-2.3600000,-0.9044563) (-2.3550000,-0.9066503) (-2.3500000,-0.9088558) (-2.3450000,-0.9110727) (-2.3400000,-0.9133012) (-2.3350000,-0.9155414) (-2.3300000,-0.9177934) (-2.3250000,-0.9200572) (-2.3200000,-0.9223330) (-2.3150000,-0.9246210) (-2.3100000,-0.9269211) (-2.3050000,-0.9292334) (-2.3000000,-0.9315582) (-2.2950000,-0.9338955) (-2.2900000,-0.9362453) (-2.2850000,-0.9386079) (-2.2800000,-0.9409833) (-2.2750000,-0.9433717) (-2.2700000,-0.9457731) (-2.2650000,-0.9481876) (-2.2600000,-0.9506155) (-2.2550000,-0.9530567) (-2.2500000,-0.9555114) (-2.2450000,-0.9579798) (-2.2400000,-0.9604620) (-2.2350000,-0.9629580) (-2.2300000,-0.9654680) (-2.2250000,-0.9679921) (-2.2200000,-0.9705305) (-2.2150000,-0.9730833) (-2.2100000,-0.9756506) (-2.2050000,-0.9782325) (-2.2000000,-0.9808293) (-2.1950000,-0.9834409) (-2.1900000,-0.9860676) (-2.1850000,-0.9887095) (-2.1800000,-0.9913668) (-2.1750000,-0.9940395) (-2.1700000,-0.9967278) (-2.1650000,-0.9994320) (-2.1600000,-1.0021520) (-2.1550000,-1.0048882) (-2.1500000,-1.0076405) (-2.1450000,-1.0104093) (-2.1400000,-1.0131945) (-2.1350000,-1.0159965) (-2.1300000,-1.0188154) (-2.1250000,-1.0216512) (-2.1200000,-1.0245043) (-2.1150000,-1.0273747) (-2.1100000,-1.0302627) (-2.1050000,-1.0331684) (-2.1000000,-1.0360919) (-2.0950000,-1.0390335) (-2.0900000,-1.0419934) (-2.0850000,-1.0449717) (-2.0800000,-1.0479686) (-2.0750000,-1.0509842) (-2.0700000,-1.0540189) (-2.0650000,-1.0570728) (-2.0600000,-1.0601460) (-2.0550000,-1.0632388) (-2.0500000,-1.0663514) (-2.0450000,-1.0694840) (-2.0400000,-1.0726368) (-2.0350000,-1.0758100) (-2.0300000,-1.0790038) (-2.0250000,-1.0822185) (-2.0200000,-1.0854542) (-2.0150000,-1.0887112) (-2.0100000,-1.0919897) (-2.0050000,-1.0952900) (-2.0000000,-1.0986123) (-1.9950000,-1.1019568) (-1.9900000,-1.1053237) (-1.9850000,-1.1087134) (-1.9800000,-1.1121260) (-1.9750000,-1.1155618) (-1.9700000,-1.1190212) (-1.9650000,-1.1225042) (-1.9600000,-1.1260113) (-1.9550000,-1.1295426) (-1.9500000,-1.1330985) (-1.9450000,-1.1366792) (-1.9400000,-1.1402850) (-1.9350000,-1.1439162) (-1.9300000,-1.1475731) (-1.9250000,-1.1512560) (-1.9200000,-1.1549652) (-1.9150000,-1.1587010) (-1.9100000,-1.1624638) (-1.9050000,-1.1662537) (-1.9000000,-1.1700713) (-1.8950000,-1.1739167) (-1.8900000,-1.1777903) (-1.8850000,-1.1816925) (-1.8800000,-1.1856237) (-1.8750000,-1.1895841) (-1.8700000,-1.1935741) (-1.8650000,-1.1975941) (-1.8600000,-1.2016445) (-1.8550000,-1.2057257) (-1.8500000,-1.2098379) (-1.8450000,-1.2139817) (-1.8400000,-1.2181574) (-1.8350000,-1.2223655) (-1.8300000,-1.2266063) (-1.8250000,-1.2308803) (-1.8200000,-1.2351878) (-1.8150000,-1.2395294) (-1.8100000,-1.2439055) (-1.8050000,-1.2483165) (-1.8000000,-1.2527630) (-1.7950000,-1.2572453) (-1.7900000,-1.2617639) (-1.7850000,-1.2663194) (-1.7800000,-1.2709123) (-1.7750000,-1.2755430) (-1.7700000,-1.2802121) (-1.7650000,-1.2849201) (-1.7600000,-1.2896675) (-1.7550000,-1.2944550) (-1.7500000,-1.2992830) (-1.7450000,-1.3041521) (-1.7400000,-1.3090630) (-1.7350000,-1.3140162) (-1.7300000,-1.3190124) (-1.7250000,-1.3240521) (-1.7200000,-1.3291359) (-1.7150000,-1.3342647) (-1.7100000,-1.3394389) (-1.7050000,-1.3446594) (-1.7000000,-1.3499267) (-1.6950000,-1.3552416) (-1.6900000,-1.3606049) (-1.6850000,-1.3660172) (-1.6800000,-1.3714793) (-1.6750000,-1.3769920) (-1.6700000,-1.3825560) (-1.6650000,-1.3881723) (-1.6600000,-1.3938416) (-1.6550000,-1.3995647) (-1.6500000,-1.4053426) (-1.6450000,-1.4111760) (-1.6400000,-1.4170660) (-1.6350000,-1.4230135) (-1.6300000,-1.4290193) (-1.6250000,-1.4350845) (-1.6200000,-1.4412101) (-1.6150000,-1.4473971) (-1.6100000,-1.4536465) (-1.6050000,-1.4599595) (-1.6000000,-1.4663371) (-1.5950000,-1.4727804) (-1.5900000,-1.4792906) (-1.5850000,-1.4858689) (-1.5800000,-1.4925166) (-1.5750000,-1.4992348) (-1.5700000,-1.5060248) (-1.5650000,-1.5128880) (-1.5600000,-1.5198258) (-1.5550000,-1.5268394) (-1.5500000,-1.5339304) (-1.5450000,-1.5411001) (-1.5400000,-1.5483502) (-1.5350000,-1.5556822) (-1.5300000,-1.5630976) (-1.5250000,-1.5705981) (-1.5200000,-1.5781854) (-1.5150000,-1.5858612) (-1.5100000,-1.5936273) (-1.5050000,-1.6014856) (-1.5000000,-1.6094379) (-1.4950000,-1.6174862) (-1.4900000,-1.6256326) (-1.4850000,-1.6338790) (-1.4800000,-1.6422277) (-1.4750000,-1.6506809) (-1.4700000,-1.6592407) (-1.4650000,-1.6679097) (-1.4600000,-1.6766901) (-1.4550000,-1.6855846) (-1.4500000,-1.6945957) (-1.4450000,-1.7037261) (-1.4400000,-1.7129786) (-1.4350000,-1.7223560) (-1.4300000,-1.7318613) (-1.4250000,-1.7414976) (-1.4200000,-1.7512681) (-1.4150000,-1.7611760) (-1.4100000,-1.7712249) (-1.4050000,-1.7814181) (-1.4000000,-1.7917595) (-1.3950000,-1.8022527) (-1.3900000,-1.8129019) (-1.3850000,-1.8237111) (-1.3800000,-1.8346845) (-1.3750000,-1.8458267) (-1.3700000,-1.8571422) (-1.3650000,-1.8686359) (-1.3600000,-1.8803129) (-1.3550000,-1.8921782) (-1.3500000,-1.9042375) (-1.3450000,-1.9164963) (-1.3400000,-1.9289606) (-1.3350000,-1.9416366) (-1.3300000,-1.9545309) (-1.3250000,-1.9676501) (-1.3200000,-1.9810015) (-1.3150000,-1.9945923) (-1.3100000,-2.0084305) (-1.3050000,-2.0225242) (-1.3000000,-2.0368819) (-1.2950000,-2.0515128) (-1.2900000,-2.0664262) (-1.2850000,-2.0816321) (-1.2800000,-2.0971411) (-1.2750000,-2.1129642) (-1.2700000,-2.1291132) (-1.2650000,-2.1456002) (-1.2600000,-2.1624385) (-1.2550000,-2.1796417) (-1.2500000,-2.1972246) (-1.2450000,-2.2152026) (-1.2400000,-2.2335922) (-1.2350000,-2.2524110) (-1.2300000,-2.2716776) (-1.2250000,-2.2914118) (-1.2200000,-2.3116349) (-1.2150000,-2.3323697) (-1.2100000,-2.3536403) (-1.2050000,-2.3754728) (-1.2000000,-2.3978953) (-1.1950000,-2.4209378) (-1.1900000,-2.4446328) (-1.1850000,-2.4690153) (-1.1800000,-2.4941233) (-1.1750000,-2.5199980) (-1.1700000,-2.5466840) (-1.1650000,-2.5742302) (-1.1600000,-2.6026897) (-1.1550000,-2.6321209) (-1.1500000,-2.6625878) (-1.1450000,-2.6941611) (-1.1400000,-2.7269187) (-1.1350000,-2.7609471) (-1.1300000,-2.7963428) (-1.1250000,-2.8332133) (-1.1200000,-2.8716796) (-1.1150000,-2.9118780) (-1.1100000,-2.9539629) (-1.1050000,-2.9981104) (-1.1000000,-3.0445224) (-1.0950000,-3.0934319) (-1.0900000,-3.1451097) (-1.0850000,-3.1998729) (-1.0800000,-3.2580965) (-1.0750000,-3.3202283) (-1.0700000,-3.3868086) (-1.0650000,-3.4584982) (-1.0600000,-3.5361167) (-1.0550000,-3.6206979) (-1.0500000,-3.7135721) (-1.0450000,-3.8164906) (-1.0400000,-3.9318256)};
\draw[orange!85!black,thick,dashed] plot coordinates {(1.0400000,3.9318256) (1.0450000,3.8164906) (1.0500000,3.7135721) (1.0550000,3.6206979) (1.0600000,3.5361167) (1.0650000,3.4584982) (1.0700000,3.3868086) (1.0750000,3.3202283) (1.0800000,3.2580965) (1.0850000,3.1998729) (1.0900000,3.1451097) (1.0950000,3.0934319) (1.1000000,3.0445224) (1.1050000,2.9981104) (1.1100000,2.9539629) (1.1150000,2.9118780) (1.1200000,2.8716796) (1.1250000,2.8332133) (1.1300000,2.7963428) (1.1350000,2.7609471) (1.1400000,2.7269187) (1.1450000,2.6941611) (1.1500000,2.6625878) (1.1550000,2.6321209) (1.1600000,2.6026897) (1.1650000,2.5742302) (1.1700000,2.5466840) (1.1750000,2.5199980) (1.1800000,2.4941233) (1.1850000,2.4690153) (1.1900000,2.4446328) (1.1950000,2.4209378) (1.2000000,2.3978953) (1.2050000,2.3754728) (1.2100000,2.3536403) (1.2150000,2.3323697) (1.2200000,2.3116349) (1.2250000,2.2914118) (1.2300000,2.2716776) (1.2350000,2.2524110) (1.2400000,2.2335922) (1.2450000,2.2152026) (1.2500000,2.1972246) (1.2550000,2.1796417) (1.2600000,2.1624385) (1.2650000,2.1456002) (1.2700000,2.1291132) (1.2750000,2.1129642) (1.2800000,2.0971411) (1.2850000,2.0816321) (1.2900000,2.0664262) (1.2950000,2.0515128) (1.3000000,2.0368819) (1.3050000,2.0225242) (1.3100000,2.0084305) (1.3150000,1.9945923) (1.3200000,1.9810015) (1.3250000,1.9676501) (1.3300000,1.9545309) (1.3350000,1.9416366) (1.3400000,1.9289606) (1.3450000,1.9164963) (1.3500000,1.9042375) (1.3550000,1.8921782) (1.3600000,1.8803129) (1.3650000,1.8686359) (1.3700000,1.8571422) (1.3750000,1.8458267) (1.3800000,1.8346845) (1.3850000,1.8237111) (1.3900000,1.8129019) (1.3950000,1.8022527) (1.4000000,1.7917595) (1.4050000,1.7814181) (1.4100000,1.7712249) (1.4150000,1.7611760) (1.4200000,1.7512681) (1.4250000,1.7414976) (1.4300000,1.7318613) (1.4350000,1.7223560) (1.4400000,1.7129786) (1.4450000,1.7037261) (1.4500000,1.6945957) (1.4550000,1.6855846) (1.4600000,1.6766901) (1.4650000,1.6679097) (1.4700000,1.6592407) (1.4750000,1.6506809) (1.4800000,1.6422277) (1.4850000,1.6338790) (1.4900000,1.6256326) (1.4950000,1.6174862) (1.5000000,1.6094379) (1.5050000,1.6014856) (1.5100000,1.5936273) (1.5150000,1.5858612) (1.5200000,1.5781854) (1.5250000,1.5705981) (1.5300000,1.5630976) (1.5350000,1.5556822) (1.5400000,1.5483502) (1.5450000,1.5411001) (1.5500000,1.5339304) (1.5550000,1.5268394) (1.5600000,1.5198258) (1.5650000,1.5128880) (1.5700000,1.5060248) (1.5750000,1.4992348) (1.5800000,1.4925166) (1.5850000,1.4858689) (1.5900000,1.4792906) (1.5950000,1.4727804) (1.6000000,1.4663371) (1.6050000,1.4599595) (1.6100000,1.4536465) (1.6150000,1.4473971) (1.6200000,1.4412101) (1.6250000,1.4350845) (1.6300000,1.4290193) (1.6350000,1.4230135) (1.6400000,1.4170660) (1.6450000,1.4111760) (1.6500000,1.4053426) (1.6550000,1.3995647) (1.6600000,1.3938416) (1.6650000,1.3881723) (1.6700000,1.3825560) (1.6750000,1.3769920) (1.6800000,1.3714793) (1.6850000,1.3660172) (1.6900000,1.3606049) (1.6950000,1.3552416) (1.7000000,1.3499267) (1.7050000,1.3446594) (1.7100000,1.3394389) (1.7150000,1.3342647) (1.7200000,1.3291359) (1.7250000,1.3240521) (1.7300000,1.3190124) (1.7350000,1.3140162) (1.7400000,1.3090630) (1.7450000,1.3041521) (1.7500000,1.2992830) (1.7550000,1.2944550) (1.7600000,1.2896675) (1.7650000,1.2849201) (1.7700000,1.2802121) (1.7750000,1.2755430) (1.7800000,1.2709123) (1.7850000,1.2663194) (1.7900000,1.2617639) (1.7950000,1.2572453) (1.8000000,1.2527630) (1.8050000,1.2483165) (1.8100000,1.2439055) (1.8150000,1.2395294) (1.8200000,1.2351878) (1.8250000,1.2308803) (1.8300000,1.2266063) (1.8350000,1.2223655) (1.8400000,1.2181574) (1.8450000,1.2139817) (1.8500000,1.2098379) (1.8550000,1.2057257) (1.8600000,1.2016445) (1.8650000,1.1975941) (1.8700000,1.1935741) (1.8750000,1.1895841) (1.8800000,1.1856237) (1.8850000,1.1816925) (1.8900000,1.1777903) (1.8950000,1.1739167) (1.9000000,1.1700713) (1.9050000,1.1662537) (1.9100000,1.1624638) (1.9150000,1.1587010) (1.9200000,1.1549652) (1.9250000,1.1512560) (1.9300000,1.1475731) (1.9350000,1.1439162) (1.9400000,1.1402850) (1.9450000,1.1366792) (1.9500000,1.1330985) (1.9550000,1.1295426) (1.9600000,1.1260113) (1.9650000,1.1225042) (1.9700000,1.1190212) (1.9750000,1.1155618) (1.9800000,1.1121260) (1.9850000,1.1087134) (1.9900000,1.1053237) (1.9950000,1.1019568) (2.0000000,1.0986123) (2.0050000,1.0952900) (2.0100000,1.0919897) (2.0150000,1.0887112) (2.0200000,1.0854542) (2.0250000,1.0822185) (2.0300000,1.0790038) (2.0350000,1.0758100) (2.0400000,1.0726368) (2.0450000,1.0694840) (2.0500000,1.0663514) (2.0550000,1.0632388) (2.0600000,1.0601460) (2.0650000,1.0570728) (2.0700000,1.0540189) (2.0750000,1.0509842) (2.0800000,1.0479686) (2.0850000,1.0449717) (2.0900000,1.0419934) (2.0950000,1.0390335) (2.1000000,1.0360919) (2.1050000,1.0331684) (2.1100000,1.0302627) (2.1150000,1.0273747) (2.1200000,1.0245043) (2.1250000,1.0216512) (2.1300000,1.0188154) (2.1350000,1.0159965) (2.1400000,1.0131945) (2.1450000,1.0104093) (2.1500000,1.0076405) (2.1550000,1.0048882) (2.1600000,1.0021520) (2.1650000,0.9994320) (2.1700000,0.9967278) (2.1750000,0.9940395) (2.1800000,0.9913668) (2.1850000,0.9887095) (2.1900000,0.9860676) (2.1950000,0.9834409) (2.2000000,0.9808293) (2.2050000,0.9782325) (2.2100000,0.9756506) (2.2150000,0.9730833) (2.2200000,0.9705305) (2.2250000,0.9679921) (2.2300000,0.9654680) (2.2350000,0.9629580) (2.2400000,0.9604620) (2.2450000,0.9579798) (2.2500000,0.9555114) (2.2550000,0.9530567) (2.2600000,0.9506155) (2.2650000,0.9481876) (2.2700000,0.9457731) (2.2750000,0.9433717) (2.2800000,0.9409833) (2.2850000,0.9386079) (2.2900000,0.9362453) (2.2950000,0.9338955) (2.3000000,0.9315582) (2.3050000,0.9292334) (2.3100000,0.9269211) (2.3150000,0.9246210) (2.3200000,0.9223330) (2.3250000,0.9200572) (2.3300000,0.9177934) (2.3350000,0.9155414) (2.3400000,0.9133012) (2.3450000,0.9110727) (2.3500000,0.9088558) (2.3550000,0.9066503) (2.3600000,0.9044563) (2.3650000,0.9022735) (2.3700000,0.9001020) (2.3750000,0.8979416) (2.3800000,0.8957922) (2.3850000,0.8936538) (2.3900000,0.8915262) (2.3950000,0.8894093) (2.4000000,0.8873032) (2.4050000,0.8852076) (2.4100000,0.8831226) (2.4150000,0.8810480) (2.4200000,0.8789837) (2.4250000,0.8769297) (2.4300000,0.8748858) (2.4350000,0.8728521) (2.4400000,0.8708284) (2.4450000,0.8688146) (2.4500000,0.8668107) (2.4550000,0.8648166) (2.4600000,0.8628322) (2.4650000,0.8608574) (2.4700000,0.8588922) (2.4750000,0.8569365) (2.4800000,0.8549902) (2.4850000,0.8530533) (2.4900000,0.8511256) (2.4950000,0.8492072) (2.5000000,0.8472979) (2.5050000,0.8453976) (2.5100000,0.8435064) (2.5150000,0.8416241) (2.5200000,0.8397507) (2.5250000,0.8378860) (2.5300000,0.8360301) (2.5350000,0.8341829) (2.5400000,0.8323443) (2.5450000,0.8305142) (2.5500000,0.8286927) (2.5550000,0.8268795) (2.5600000,0.8250747) (2.5650000,0.8232782) (2.5700000,0.8214900) (2.5750000,0.8197099) (2.5800000,0.8179380) (2.5850000,0.8161741) (2.5900000,0.8144182) (2.5950000,0.8126703) (2.6000000,0.8109302) (2.6050000,0.8091980) (2.6100000,0.8074736) (2.6150000,0.8057569) (2.6200000,0.8040479) (2.6250000,0.8023465) (2.6300000,0.8006526) (2.6350000,0.7989663) (2.6400000,0.7972874) (2.6450000,0.7956160) (2.6500000,0.7939519) (2.6550000,0.7922951) (2.6600000,0.7906455) (2.6650000,0.7890032) (2.6700000,0.7873680) (2.6750000,0.7857400) (2.6800000,0.7841190) (2.6850000,0.7825050) (2.6900000,0.7808979) (2.6950000,0.7792978) (2.7000000,0.7777046) (2.7050000,0.7761181) (2.7100000,0.7745385) (2.7150000,0.7729656) (2.7200000,0.7713994) (2.7250000,0.7698398) (2.7300000,0.7682868) (2.7350000,0.7667404) (2.7400000,0.7652005) (2.7450000,0.7636671) (2.7500000,0.7621401) (2.7550000,0.7606194) (2.7600000,0.7591051) (2.7650000,0.7575972) (2.7700000,0.7560955) (2.7750000,0.7546000) (2.7800000,0.7531106) (2.7850000,0.7516275) (2.7900000,0.7501504) (2.7950000,0.7486794) (2.8000000,0.7472144) (2.8050000,0.7457554) (2.8100000,0.7443023) (2.8150000,0.7428552) (2.8200000,0.7414139) (2.8250000,0.7399785) (2.8300000,0.7385488) (2.8350000,0.7371250) (2.8400000,0.7357068) (2.8450000,0.7342943) (2.8500000,0.7328875) (2.8550000,0.7314863) (2.8600000,0.7300907) (2.8650000,0.7287006) (2.8700000,0.7273161) (2.8750000,0.7259370) (2.8800000,0.7245634) (2.8850000,0.7231952) (2.8900000,0.7218323) (2.8950000,0.7204748) (2.9000000,0.7191227) (2.9050000,0.7177758) (2.9100000,0.7164341) (2.9150000,0.7150977) (2.9200000,0.7137665) (2.9250000,0.7124404) (2.9300000,0.7111194) (2.9350000,0.7098036) (2.9400000,0.7084928) (2.9450000,0.7071870) (2.9500000,0.7058862) (2.9550000,0.7045904) (2.9600000,0.7032996) (2.9650000,0.7020136) (2.9700000,0.7007326) (2.9750000,0.6994563) (2.9800000,0.6981850) (2.9850000,0.6969184) (2.9900000,0.6956566) (2.9950000,0.6943995) (3.0000000,0.6931472) (3.0050000,0.6918995) (3.0100000,0.6906565) (3.0150000,0.6894182) (3.0200000,0.6881844) (3.0250000,0.6869552) (3.0300000,0.6857306) (3.0350000,0.6845105) (3.0400000,0.6832949) (3.0450000,0.6820838) (3.0500000,0.6808771) (3.0550000,0.6796748) (3.0600000,0.6784770) (3.0650000,0.6772835) (3.0700000,0.6760944) (3.0750000,0.6749096) (3.0800000,0.6737291) (3.0850000,0.6725529) (3.0900000,0.6713809) (3.0950000,0.6702132) (3.1000000,0.6690496) (3.1050000,0.6678903) (3.1100000,0.6667351) (3.1150000,0.6655840) (3.1200000,0.6644371) (3.1250000,0.6632942) (3.1300000,0.6621554) (3.1350000,0.6610207) (3.1400000,0.6598900) (3.1450000,0.6587632) (3.1500000,0.6576405) (3.1550000,0.6565217) (3.1600000,0.6554069) (3.1650000,0.6542959) (3.1700000,0.6531889) (3.1750000,0.6520857) (3.1800000,0.6509864) (3.1850000,0.6498909) (3.1900000,0.6487992) (3.1950000,0.6477113) (3.2000000,0.6466272) (3.2050000,0.6455468) (3.2100000,0.6444701) (3.2150000,0.6433972) (3.2200000,0.6423279) (3.2250000,0.6412623) (3.2300000,0.6402004) (3.2350000,0.6391421) (3.2400000,0.6380874) (3.2450000,0.6370363) (3.2500000,0.6359888) (3.2550000,0.6349448) (3.2600000,0.6339043) (3.2650000,0.6328674) (3.2700000,0.6318340) (3.2750000,0.6308040) (3.2800000,0.6297776) (3.2850000,0.6287545) (3.2900000,0.6277349) (3.2950000,0.6267187) (3.3000000,0.6257059) (3.3050000,0.6246965) (3.3100000,0.6236904) (3.3150000,0.6226876) (3.3200000,0.6216882) (3.3250000,0.6206921) (3.3300000,0.6196993) (3.3350000,0.6187097) (3.3400000,0.6177234) (3.3450000,0.6167404) (3.3500000,0.6157605) (3.3550000,0.6147839) (3.3600000,0.6138104) (3.3650000,0.6128402) (3.3700000,0.6118731) (3.3750000,0.6109091) (3.3800000,0.6099482) (3.3850000,0.6089905) (3.3900000,0.6080359) (3.3950000,0.6070843) (3.4000000,0.6061358) (3.4050000,0.6051904) (3.4100000,0.6042479) (3.4150000,0.6033085) (3.4200000,0.6023722) (3.4250000,0.6014388) (3.4300000,0.6005083) (3.4350000,0.5995809) (3.4400000,0.5986563) (3.4450000,0.5977347) (3.4500000,0.5968161) (3.4550000,0.5959003) (3.4600000,0.5949874) (3.4650000,0.5940774) (3.4700000,0.5931703) (3.4750000,0.5922660) (3.4800000,0.5913645) (3.4850000,0.5904658) (3.4900000,0.5895700) (3.4950000,0.5886769) (3.5000000,0.5877867) (3.5050000,0.5868992) (3.5100000,0.5860144) (3.5150000,0.5851324) (3.5200000,0.5842531) (3.5250000,0.5833765) (3.5300000,0.5825026) (3.5350000,0.5816314) (3.5400000,0.5807629) (3.5450000,0.5798971) (3.5500000,0.5790339) (3.5550000,0.5781733) (3.5600000,0.5773154) (3.5650000,0.5764600) (3.5700000,0.5756073) (3.5750000,0.5747572) (3.5800000,0.5739096) (3.5850000,0.5730646) (3.5900000,0.5722221) (3.5950000,0.5713822) (3.6000000,0.5705449) (3.6050000,0.5697100) (3.6100000,0.5688776) (3.6150000,0.5680478) (3.6200000,0.5672204) (3.6250000,0.5663955) (3.6300000,0.5655730) (3.6350000,0.5647530) (3.6400000,0.5639354) (3.6450000,0.5631203) (3.6500000,0.5623076) (3.6550000,0.5614973) (3.6600000,0.5606893) (3.6650000,0.5598838) (3.6700000,0.5590806) (3.6750000,0.5582798) (3.6800000,0.5574813) (3.6850000,0.5566852) (3.6900000,0.5558914) (3.6950000,0.5550999) (3.7000000,0.5543107) (3.7050000,0.5535239) (3.7100000,0.5527393) (3.7150000,0.5519570) (3.7200000,0.5511769) (3.7250000,0.5503991) (3.7300000,0.5496236) (3.7350000,0.5488503) (3.7400000,0.5480792) (3.7450000,0.5473104) (3.7500000,0.5465437) (3.7550000,0.5457793) (3.7600000,0.5450170) (3.7650000,0.5442569) (3.7700000,0.5434990) (3.7750000,0.5427432) (3.7800000,0.5419896) (3.7850000,0.5412382) (3.7900000,0.5404888) (3.7950000,0.5397416) (3.8000000,0.5389965) (3.8050000,0.5382535) (3.8100000,0.5375126) (3.8150000,0.5367738) (3.8200000,0.5360370) (3.8250000,0.5353024) (3.8300000,0.5345698) (3.8350000,0.5338392) (3.8400000,0.5331107) (3.8450000,0.5323842) (3.8500000,0.5316597) (3.8550000,0.5309373) (3.8600000,0.5302168) (3.8650000,0.5294984) (3.8700000,0.5287819) (3.8750000,0.5280674) (3.8800000,0.5273549) (3.8850000,0.5266444) (3.8900000,0.5259358) (3.8950000,0.5252292) (3.9000000,0.5245245) (3.9050000,0.5238217) (3.9100000,0.5231209) (3.9150000,0.5224219) (3.9200000,0.5217249) (3.9250000,0.5210298) (3.9300000,0.5203366) (3.9350000,0.5196452) (3.9400000,0.5189557) (3.9450000,0.5182681) (3.9500000,0.5175824) (3.9550000,0.5168985) (3.9600000,0.5162165) (3.9650000,0.5155363) (3.9700000,0.5148579) (3.9750000,0.5141813) (3.9800000,0.5135066) (3.9850000,0.5128337) (3.9900000,0.5121625) (3.9950000,0.5114932) (4.0000000,0.5108256)};
\draw[violet!75!black,densely dotted] (-2,1.0986123)--(2,-1.0986123);
\fill[violet!75!black] (-2,1.0986123) circle(2.3pt);
\fill[violet!75!black] (2,-1.0986123) circle(2.3pt);
\end{scope}
\end{scope}
\node[anchor=west,align=left,text width=4.35cm] at (8.9,4.4) {固定 $k=1$。\\[8pt]蓝实线：$b=1$\\[8pt]橙虚线：$b=-1$\\[14pt]阴影内没有函数值。\\两条边界均不能取。};
\node[text width=12.8cm,align=center] at (7.1,.65) {$k,b\ne0$；$D=\{x:|x|>|b/k|\}$。图中禁区为 $[-1,1]$。};
\end{tikzpicture}\end{center}

\textbf{先看图：} 图中固定 $k=1$，分别取 $b=1,-1$，阴影区间 $[-1,1]$ 内没有函数值。一般 $k,b\ne0$ 时，真数为正等价于 $|kx|>|b|$，故必须 $|x|>|b/k|$。

\textbf{再说明一般情形：}
\[
R(x)=\frac{kx-b}{kx+b}>0,\qquad R(-x)=\frac1{R(x)}.
\]
由 $\ln(1/u)=-\ln u$（$u>0$），可得 $f(-x)=-f(x)$。不能只检查分母非零。
\end{explanationbox}
\clearpage
\begin{explanationbox}
\textbf{\textcolor{CExplanation}{分子分母互换，对数整体反号}}
% Formula-based samples; see F083 checks/build_tikz_diagrams.py.
\begin{center}\begin{tikzpicture}[font=\small,>=stealth,line join=round]
\draw[gray!40,rounded corners=3pt] (0,0) rectangle (14.2,8.4);
\node at (7.1,7.75) {奇5：$y=\ln\dfrac{kx+b}{kx-b}$};
\begin{scope}[shift={(4.400000,3.900000)},x=0.850000cm,y=0.562500cm]
\fill[gray!12] (-1,-4) rectangle (1,4);
\draw[gray,dashed] (-1,-4)--(-1,4);
\draw[gray,dashed] (1,-4)--(1,4);
\draw[->,gray] (-4,0)--(4,0) node[right] {$x$};
\draw[->,gray] (0,-4)--(0,4) node[above] {$y$};
\node[anchor=north] at (-1,-4-.07/0.562500) {$-1$};
\node[anchor=north] at (1,-4-.07/0.562500) {$1$};
\draw[gray] (-2,-.035/0.562500)--(-2,.035/0.562500);
\node[anchor=north] at (-2,-.07/0.562500) {$-2$};
\draw[gray] (2,-.035/0.562500)--(2,.035/0.562500);
\node[anchor=north] at (2,-.07/0.562500) {$2$};
\begin{scope}
\clip (-4,-4) rectangle (4,4);
\draw[blue,thick,solid] plot coordinates {(-4.0000000,-0.5108256) (-3.9950000,-0.5114932) (-3.9900000,-0.5121625) (-3.9850000,-0.5128337) (-3.9800000,-0.5135066) (-3.9750000,-0.5141813) (-3.9700000,-0.5148579) (-3.9650000,-0.5155363) (-3.9600000,-0.5162165) (-3.9550000,-0.5168985) (-3.9500000,-0.5175824) (-3.9450000,-0.5182681) (-3.9400000,-0.5189557) (-3.9350000,-0.5196452) (-3.9300000,-0.5203366) (-3.9250000,-0.5210298) (-3.9200000,-0.5217249) (-3.9150000,-0.5224219) (-3.9100000,-0.5231209) (-3.9050000,-0.5238217) (-3.9000000,-0.5245245) (-3.8950000,-0.5252292) (-3.8900000,-0.5259358) (-3.8850000,-0.5266444) (-3.8800000,-0.5273549) (-3.8750000,-0.5280674) (-3.8700000,-0.5287819) (-3.8650000,-0.5294984) (-3.8600000,-0.5302168) (-3.8550000,-0.5309373) (-3.8500000,-0.5316597) (-3.8450000,-0.5323842) (-3.8400000,-0.5331107) (-3.8350000,-0.5338392) (-3.8300000,-0.5345698) (-3.8250000,-0.5353024) (-3.8200000,-0.5360370) (-3.8150000,-0.5367738) (-3.8100000,-0.5375126) (-3.8050000,-0.5382535) (-3.8000000,-0.5389965) (-3.7950000,-0.5397416) (-3.7900000,-0.5404888) (-3.7850000,-0.5412382) (-3.7800000,-0.5419896) (-3.7750000,-0.5427432) (-3.7700000,-0.5434990) (-3.7650000,-0.5442569) (-3.7600000,-0.5450170) (-3.7550000,-0.5457793) (-3.7500000,-0.5465437) (-3.7450000,-0.5473104) (-3.7400000,-0.5480792) (-3.7350000,-0.5488503) (-3.7300000,-0.5496236) (-3.7250000,-0.5503991) (-3.7200000,-0.5511769) (-3.7150000,-0.5519570) (-3.7100000,-0.5527393) (-3.7050000,-0.5535239) (-3.7000000,-0.5543107) (-3.6950000,-0.5550999) (-3.6900000,-0.5558914) (-3.6850000,-0.5566852) (-3.6800000,-0.5574813) (-3.6750000,-0.5582798) (-3.6700000,-0.5590806) (-3.6650000,-0.5598838) (-3.6600000,-0.5606893) (-3.6550000,-0.5614973) (-3.6500000,-0.5623076) (-3.6450000,-0.5631203) (-3.6400000,-0.5639354) (-3.6350000,-0.5647530) (-3.6300000,-0.5655730) (-3.6250000,-0.5663955) (-3.6200000,-0.5672204) (-3.6150000,-0.5680478) (-3.6100000,-0.5688776) (-3.6050000,-0.5697100) (-3.6000000,-0.5705449) (-3.5950000,-0.5713822) (-3.5900000,-0.5722221) (-3.5850000,-0.5730646) (-3.5800000,-0.5739096) (-3.5750000,-0.5747572) (-3.5700000,-0.5756073) (-3.5650000,-0.5764600) (-3.5600000,-0.5773154) (-3.5550000,-0.5781733) (-3.5500000,-0.5790339) (-3.5450000,-0.5798971) (-3.5400000,-0.5807629) (-3.5350000,-0.5816314) (-3.5300000,-0.5825026) (-3.5250000,-0.5833765) (-3.5200000,-0.5842531) (-3.5150000,-0.5851324) (-3.5100000,-0.5860144) (-3.5050000,-0.5868992) (-3.5000000,-0.5877867) (-3.4950000,-0.5886769) (-3.4900000,-0.5895700) (-3.4850000,-0.5904658) (-3.4800000,-0.5913645) (-3.4750000,-0.5922660) (-3.4700000,-0.5931703) (-3.4650000,-0.5940774) (-3.4600000,-0.5949874) (-3.4550000,-0.5959003) (-3.4500000,-0.5968161) (-3.4450000,-0.5977347) (-3.4400000,-0.5986563) (-3.4350000,-0.5995809) (-3.4300000,-0.6005083) (-3.4250000,-0.6014388) (-3.4200000,-0.6023722) (-3.4150000,-0.6033085) (-3.4100000,-0.6042479) (-3.4050000,-0.6051904) (-3.4000000,-0.6061358) (-3.3950000,-0.6070843) (-3.3900000,-0.6080359) (-3.3850000,-0.6089905) (-3.3800000,-0.6099482) (-3.3750000,-0.6109091) (-3.3700000,-0.6118731) (-3.3650000,-0.6128402) (-3.3600000,-0.6138104) (-3.3550000,-0.6147839) (-3.3500000,-0.6157605) (-3.3450000,-0.6167404) (-3.3400000,-0.6177234) (-3.3350000,-0.6187097) (-3.3300000,-0.6196993) (-3.3250000,-0.6206921) (-3.3200000,-0.6216882) (-3.3150000,-0.6226876) (-3.3100000,-0.6236904) (-3.3050000,-0.6246965) (-3.3000000,-0.6257059) (-3.2950000,-0.6267187) (-3.2900000,-0.6277349) (-3.2850000,-0.6287545) (-3.2800000,-0.6297776) (-3.2750000,-0.6308040) (-3.2700000,-0.6318340) (-3.2650000,-0.6328674) (-3.2600000,-0.6339043) (-3.2550000,-0.6349448) (-3.2500000,-0.6359888) (-3.2450000,-0.6370363) (-3.2400000,-0.6380874) (-3.2350000,-0.6391421) (-3.2300000,-0.6402004) (-3.2250000,-0.6412623) (-3.2200000,-0.6423279) (-3.2150000,-0.6433972) (-3.2100000,-0.6444701) (-3.2050000,-0.6455468) (-3.2000000,-0.6466272) (-3.1950000,-0.6477113) (-3.1900000,-0.6487992) (-3.1850000,-0.6498909) (-3.1800000,-0.6509864) (-3.1750000,-0.6520857) (-3.1700000,-0.6531889) (-3.1650000,-0.6542959) (-3.1600000,-0.6554069) (-3.1550000,-0.6565217) (-3.1500000,-0.6576405) (-3.1450000,-0.6587632) (-3.1400000,-0.6598900) (-3.1350000,-0.6610207) (-3.1300000,-0.6621554) (-3.1250000,-0.6632942) (-3.1200000,-0.6644371) (-3.1150000,-0.6655840) (-3.1100000,-0.6667351) (-3.1050000,-0.6678903) (-3.1000000,-0.6690496) (-3.0950000,-0.6702132) (-3.0900000,-0.6713809) (-3.0850000,-0.6725529) (-3.0800000,-0.6737291) (-3.0750000,-0.6749096) (-3.0700000,-0.6760944) (-3.0650000,-0.6772835) (-3.0600000,-0.6784770) (-3.0550000,-0.6796748) (-3.0500000,-0.6808771) (-3.0450000,-0.6820838) (-3.0400000,-0.6832949) (-3.0350000,-0.6845105) (-3.0300000,-0.6857306) (-3.0250000,-0.6869552) (-3.0200000,-0.6881844) (-3.0150000,-0.6894182) (-3.0100000,-0.6906565) (-3.0050000,-0.6918995) (-3.0000000,-0.6931472) (-2.9950000,-0.6943995) (-2.9900000,-0.6956566) (-2.9850000,-0.6969184) (-2.9800000,-0.6981850) (-2.9750000,-0.6994563) (-2.9700000,-0.7007326) (-2.9650000,-0.7020136) (-2.9600000,-0.7032996) (-2.9550000,-0.7045904) (-2.9500000,-0.7058862) (-2.9450000,-0.7071870) (-2.9400000,-0.7084928) (-2.9350000,-0.7098036) (-2.9300000,-0.7111194) (-2.9250000,-0.7124404) (-2.9200000,-0.7137665) (-2.9150000,-0.7150977) (-2.9100000,-0.7164341) (-2.9050000,-0.7177758) (-2.9000000,-0.7191227) (-2.8950000,-0.7204748) (-2.8900000,-0.7218323) (-2.8850000,-0.7231952) (-2.8800000,-0.7245634) (-2.8750000,-0.7259370) (-2.8700000,-0.7273161) (-2.8650000,-0.7287006) (-2.8600000,-0.7300907) (-2.8550000,-0.7314863) (-2.8500000,-0.7328875) (-2.8450000,-0.7342943) (-2.8400000,-0.7357068) (-2.8350000,-0.7371250) (-2.8300000,-0.7385488) (-2.8250000,-0.7399785) (-2.8200000,-0.7414139) (-2.8150000,-0.7428552) (-2.8100000,-0.7443023) (-2.8050000,-0.7457554) (-2.8000000,-0.7472144) (-2.7950000,-0.7486794) (-2.7900000,-0.7501504) (-2.7850000,-0.7516275) (-2.7800000,-0.7531106) (-2.7750000,-0.7546000) (-2.7700000,-0.7560955) (-2.7650000,-0.7575972) (-2.7600000,-0.7591051) (-2.7550000,-0.7606194) (-2.7500000,-0.7621401) (-2.7450000,-0.7636671) (-2.7400000,-0.7652005) (-2.7350000,-0.7667404) (-2.7300000,-0.7682868) (-2.7250000,-0.7698398) (-2.7200000,-0.7713994) (-2.7150000,-0.7729656) (-2.7100000,-0.7745385) (-2.7050000,-0.7761181) (-2.7000000,-0.7777046) (-2.6950000,-0.7792978) (-2.6900000,-0.7808979) (-2.6850000,-0.7825050) (-2.6800000,-0.7841190) (-2.6750000,-0.7857400) (-2.6700000,-0.7873680) (-2.6650000,-0.7890032) (-2.6600000,-0.7906455) (-2.6550000,-0.7922951) (-2.6500000,-0.7939519) (-2.6450000,-0.7956160) (-2.6400000,-0.7972874) (-2.6350000,-0.7989663) (-2.6300000,-0.8006526) (-2.6250000,-0.8023465) (-2.6200000,-0.8040479) (-2.6150000,-0.8057569) (-2.6100000,-0.8074736) (-2.6050000,-0.8091980) (-2.6000000,-0.8109302) (-2.5950000,-0.8126703) (-2.5900000,-0.8144182) (-2.5850000,-0.8161741) (-2.5800000,-0.8179380) (-2.5750000,-0.8197099) (-2.5700000,-0.8214900) (-2.5650000,-0.8232782) (-2.5600000,-0.8250747) (-2.5550000,-0.8268795) (-2.5500000,-0.8286927) (-2.5450000,-0.8305142) (-2.5400000,-0.8323443) (-2.5350000,-0.8341829) (-2.5300000,-0.8360301) (-2.5250000,-0.8378860) (-2.5200000,-0.8397507) (-2.5150000,-0.8416241) (-2.5100000,-0.8435064) (-2.5050000,-0.8453976) (-2.5000000,-0.8472979) (-2.4950000,-0.8492072) (-2.4900000,-0.8511256) (-2.4850000,-0.8530533) (-2.4800000,-0.8549902) (-2.4750000,-0.8569365) (-2.4700000,-0.8588922) (-2.4650000,-0.8608574) (-2.4600000,-0.8628322) (-2.4550000,-0.8648166) (-2.4500000,-0.8668107) (-2.4450000,-0.8688146) (-2.4400000,-0.8708284) (-2.4350000,-0.8728521) (-2.4300000,-0.8748858) (-2.4250000,-0.8769297) (-2.4200000,-0.8789837) (-2.4150000,-0.8810480) (-2.4100000,-0.8831226) (-2.4050000,-0.8852076) (-2.4000000,-0.8873032) (-2.3950000,-0.8894093) (-2.3900000,-0.8915262) (-2.3850000,-0.8936538) (-2.3800000,-0.8957922) (-2.3750000,-0.8979416) (-2.3700000,-0.9001020) (-2.3650000,-0.9022735) (-2.3600000,-0.9044563) (-2.3550000,-0.9066503) (-2.3500000,-0.9088558) (-2.3450000,-0.9110727) (-2.3400000,-0.9133012) (-2.3350000,-0.9155414) (-2.3300000,-0.9177934) (-2.3250000,-0.9200572) (-2.3200000,-0.9223330) (-2.3150000,-0.9246210) (-2.3100000,-0.9269211) (-2.3050000,-0.9292334) (-2.3000000,-0.9315582) (-2.2950000,-0.9338955) (-2.2900000,-0.9362453) (-2.2850000,-0.9386079) (-2.2800000,-0.9409833) (-2.2750000,-0.9433717) (-2.2700000,-0.9457731) (-2.2650000,-0.9481876) (-2.2600000,-0.9506155) (-2.2550000,-0.9530567) (-2.2500000,-0.9555114) (-2.2450000,-0.9579798) (-2.2400000,-0.9604620) (-2.2350000,-0.9629580) (-2.2300000,-0.9654680) (-2.2250000,-0.9679921) (-2.2200000,-0.9705305) (-2.2150000,-0.9730833) (-2.2100000,-0.9756506) (-2.2050000,-0.9782325) (-2.2000000,-0.9808293) (-2.1950000,-0.9834409) (-2.1900000,-0.9860676) (-2.1850000,-0.9887095) (-2.1800000,-0.9913668) (-2.1750000,-0.9940395) (-2.1700000,-0.9967278) (-2.1650000,-0.9994320) (-2.1600000,-1.0021520) (-2.1550000,-1.0048882) (-2.1500000,-1.0076405) (-2.1450000,-1.0104093) (-2.1400000,-1.0131945) (-2.1350000,-1.0159965) (-2.1300000,-1.0188154) (-2.1250000,-1.0216512) (-2.1200000,-1.0245043) (-2.1150000,-1.0273747) (-2.1100000,-1.0302627) (-2.1050000,-1.0331684) (-2.1000000,-1.0360919) (-2.0950000,-1.0390335) (-2.0900000,-1.0419934) (-2.0850000,-1.0449717) (-2.0800000,-1.0479686) (-2.0750000,-1.0509842) (-2.0700000,-1.0540189) (-2.0650000,-1.0570728) (-2.0600000,-1.0601460) (-2.0550000,-1.0632388) (-2.0500000,-1.0663514) (-2.0450000,-1.0694840) (-2.0400000,-1.0726368) (-2.0350000,-1.0758100) (-2.0300000,-1.0790038) (-2.0250000,-1.0822185) (-2.0200000,-1.0854542) (-2.0150000,-1.0887112) (-2.0100000,-1.0919897) (-2.0050000,-1.0952900) (-2.0000000,-1.0986123) (-1.9950000,-1.1019568) (-1.9900000,-1.1053237) (-1.9850000,-1.1087134) (-1.9800000,-1.1121260) (-1.9750000,-1.1155618) (-1.9700000,-1.1190212) (-1.9650000,-1.1225042) (-1.9600000,-1.1260113) (-1.9550000,-1.1295426) (-1.9500000,-1.1330985) (-1.9450000,-1.1366792) (-1.9400000,-1.1402850) (-1.9350000,-1.1439162) (-1.9300000,-1.1475731) (-1.9250000,-1.1512560) (-1.9200000,-1.1549652) (-1.9150000,-1.1587010) (-1.9100000,-1.1624638) (-1.9050000,-1.1662537) (-1.9000000,-1.1700713) (-1.8950000,-1.1739167) (-1.8900000,-1.1777903) (-1.8850000,-1.1816925) (-1.8800000,-1.1856237) (-1.8750000,-1.1895841) (-1.8700000,-1.1935741) (-1.8650000,-1.1975941) (-1.8600000,-1.2016445) (-1.8550000,-1.2057257) (-1.8500000,-1.2098379) (-1.8450000,-1.2139817) (-1.8400000,-1.2181574) (-1.8350000,-1.2223655) (-1.8300000,-1.2266063) (-1.8250000,-1.2308803) (-1.8200000,-1.2351878) (-1.8150000,-1.2395294) (-1.8100000,-1.2439055) (-1.8050000,-1.2483165) (-1.8000000,-1.2527630) (-1.7950000,-1.2572453) (-1.7900000,-1.2617639) (-1.7850000,-1.2663194) (-1.7800000,-1.2709123) (-1.7750000,-1.2755430) (-1.7700000,-1.2802121) (-1.7650000,-1.2849201) (-1.7600000,-1.2896675) (-1.7550000,-1.2944550) (-1.7500000,-1.2992830) (-1.7450000,-1.3041521) (-1.7400000,-1.3090630) (-1.7350000,-1.3140162) (-1.7300000,-1.3190124) (-1.7250000,-1.3240521) (-1.7200000,-1.3291359) (-1.7150000,-1.3342647) (-1.7100000,-1.3394389) (-1.7050000,-1.3446594) (-1.7000000,-1.3499267) (-1.6950000,-1.3552416) (-1.6900000,-1.3606049) (-1.6850000,-1.3660172) (-1.6800000,-1.3714793) (-1.6750000,-1.3769920) (-1.6700000,-1.3825560) (-1.6650000,-1.3881723) (-1.6600000,-1.3938416) (-1.6550000,-1.3995647) (-1.6500000,-1.4053426) (-1.6450000,-1.4111760) (-1.6400000,-1.4170660) (-1.6350000,-1.4230135) (-1.6300000,-1.4290193) (-1.6250000,-1.4350845) (-1.6200000,-1.4412101) (-1.6150000,-1.4473971) (-1.6100000,-1.4536465) (-1.6050000,-1.4599595) (-1.6000000,-1.4663371) (-1.5950000,-1.4727804) (-1.5900000,-1.4792906) (-1.5850000,-1.4858689) (-1.5800000,-1.4925166) (-1.5750000,-1.4992348) (-1.5700000,-1.5060248) (-1.5650000,-1.5128880) (-1.5600000,-1.5198258) (-1.5550000,-1.5268394) (-1.5500000,-1.5339304) (-1.5450000,-1.5411001) (-1.5400000,-1.5483502) (-1.5350000,-1.5556822) (-1.5300000,-1.5630976) (-1.5250000,-1.5705981) (-1.5200000,-1.5781854) (-1.5150000,-1.5858612) (-1.5100000,-1.5936273) (-1.5050000,-1.6014856) (-1.5000000,-1.6094379) (-1.4950000,-1.6174862) (-1.4900000,-1.6256326) (-1.4850000,-1.6338790) (-1.4800000,-1.6422277) (-1.4750000,-1.6506809) (-1.4700000,-1.6592407) (-1.4650000,-1.6679097) (-1.4600000,-1.6766901) (-1.4550000,-1.6855846) (-1.4500000,-1.6945957) (-1.4450000,-1.7037261) (-1.4400000,-1.7129786) (-1.4350000,-1.7223560) (-1.4300000,-1.7318613) (-1.4250000,-1.7414976) (-1.4200000,-1.7512681) (-1.4150000,-1.7611760) (-1.4100000,-1.7712249) (-1.4050000,-1.7814181) (-1.4000000,-1.7917595) (-1.3950000,-1.8022527) (-1.3900000,-1.8129019) (-1.3850000,-1.8237111) (-1.3800000,-1.8346845) (-1.3750000,-1.8458267) (-1.3700000,-1.8571422) (-1.3650000,-1.8686359) (-1.3600000,-1.8803129) (-1.3550000,-1.8921782) (-1.3500000,-1.9042375) (-1.3450000,-1.9164963) (-1.3400000,-1.9289606) (-1.3350000,-1.9416366) (-1.3300000,-1.9545309) (-1.3250000,-1.9676501) (-1.3200000,-1.9810015) (-1.3150000,-1.9945923) (-1.3100000,-2.0084305) (-1.3050000,-2.0225242) (-1.3000000,-2.0368819) (-1.2950000,-2.0515128) (-1.2900000,-2.0664262) (-1.2850000,-2.0816321) (-1.2800000,-2.0971411) (-1.2750000,-2.1129642) (-1.2700000,-2.1291132) (-1.2650000,-2.1456002) (-1.2600000,-2.1624385) (-1.2550000,-2.1796417) (-1.2500000,-2.1972246) (-1.2450000,-2.2152026) (-1.2400000,-2.2335922) (-1.2350000,-2.2524110) (-1.2300000,-2.2716776) (-1.2250000,-2.2914118) (-1.2200000,-2.3116349) (-1.2150000,-2.3323697) (-1.2100000,-2.3536403) (-1.2050000,-2.3754728) (-1.2000000,-2.3978953) (-1.1950000,-2.4209378) (-1.1900000,-2.4446328) (-1.1850000,-2.4690153) (-1.1800000,-2.4941233) (-1.1750000,-2.5199980) (-1.1700000,-2.5466840) (-1.1650000,-2.5742302) (-1.1600000,-2.6026897) (-1.1550000,-2.6321209) (-1.1500000,-2.6625878) (-1.1450000,-2.6941611) (-1.1400000,-2.7269187) (-1.1350000,-2.7609471) (-1.1300000,-2.7963428) (-1.1250000,-2.8332133) (-1.1200000,-2.8716796) (-1.1150000,-2.9118780) (-1.1100000,-2.9539629) (-1.1050000,-2.9981104) (-1.1000000,-3.0445224) (-1.0950000,-3.0934319) (-1.0900000,-3.1451097) (-1.0850000,-3.1998729) (-1.0800000,-3.2580965) (-1.0750000,-3.3202283) (-1.0700000,-3.3868086) (-1.0650000,-3.4584982) (-1.0600000,-3.5361167) (-1.0550000,-3.6206979) (-1.0500000,-3.7135721) (-1.0450000,-3.8164906) (-1.0400000,-3.9318256)};
\draw[blue,thick,solid] plot coordinates {(1.0400000,3.9318256) (1.0450000,3.8164906) (1.0500000,3.7135721) (1.0550000,3.6206979) (1.0600000,3.5361167) (1.0650000,3.4584982) (1.0700000,3.3868086) (1.0750000,3.3202283) (1.0800000,3.2580965) (1.0850000,3.1998729) (1.0900000,3.1451097) (1.0950000,3.0934319) (1.1000000,3.0445224) (1.1050000,2.9981104) (1.1100000,2.9539629) (1.1150000,2.9118780) (1.1200000,2.8716796) (1.1250000,2.8332133) (1.1300000,2.7963428) (1.1350000,2.7609471) (1.1400000,2.7269187) (1.1450000,2.6941611) (1.1500000,2.6625878) (1.1550000,2.6321209) (1.1600000,2.6026897) (1.1650000,2.5742302) (1.1700000,2.5466840) (1.1750000,2.5199980) (1.1800000,2.4941233) (1.1850000,2.4690153) (1.1900000,2.4446328) (1.1950000,2.4209378) (1.2000000,2.3978953) (1.2050000,2.3754728) (1.2100000,2.3536403) (1.2150000,2.3323697) (1.2200000,2.3116349) (1.2250000,2.2914118) (1.2300000,2.2716776) (1.2350000,2.2524110) (1.2400000,2.2335922) (1.2450000,2.2152026) (1.2500000,2.1972246) (1.2550000,2.1796417) (1.2600000,2.1624385) (1.2650000,2.1456002) (1.2700000,2.1291132) (1.2750000,2.1129642) (1.2800000,2.0971411) (1.2850000,2.0816321) (1.2900000,2.0664262) (1.2950000,2.0515128) (1.3000000,2.0368819) (1.3050000,2.0225242) (1.3100000,2.0084305) (1.3150000,1.9945923) (1.3200000,1.9810015) (1.3250000,1.9676501) (1.3300000,1.9545309) (1.3350000,1.9416366) (1.3400000,1.9289606) (1.3450000,1.9164963) (1.3500000,1.9042375) (1.3550000,1.8921782) (1.3600000,1.8803129) (1.3650000,1.8686359) (1.3700000,1.8571422) (1.3750000,1.8458267) (1.3800000,1.8346845) (1.3850000,1.8237111) (1.3900000,1.8129019) (1.3950000,1.8022527) (1.4000000,1.7917595) (1.4050000,1.7814181) (1.4100000,1.7712249) (1.4150000,1.7611760) (1.4200000,1.7512681) (1.4250000,1.7414976) (1.4300000,1.7318613) (1.4350000,1.7223560) (1.4400000,1.7129786) (1.4450000,1.7037261) (1.4500000,1.6945957) (1.4550000,1.6855846) (1.4600000,1.6766901) (1.4650000,1.6679097) (1.4700000,1.6592407) (1.4750000,1.6506809) (1.4800000,1.6422277) (1.4850000,1.6338790) (1.4900000,1.6256326) (1.4950000,1.6174862) (1.5000000,1.6094379) (1.5050000,1.6014856) (1.5100000,1.5936273) (1.5150000,1.5858612) (1.5200000,1.5781854) (1.5250000,1.5705981) (1.5300000,1.5630976) (1.5350000,1.5556822) (1.5400000,1.5483502) (1.5450000,1.5411001) (1.5500000,1.5339304) (1.5550000,1.5268394) (1.5600000,1.5198258) (1.5650000,1.5128880) (1.5700000,1.5060248) (1.5750000,1.4992348) (1.5800000,1.4925166) (1.5850000,1.4858689) (1.5900000,1.4792906) (1.5950000,1.4727804) (1.6000000,1.4663371) (1.6050000,1.4599595) (1.6100000,1.4536465) (1.6150000,1.4473971) (1.6200000,1.4412101) (1.6250000,1.4350845) (1.6300000,1.4290193) (1.6350000,1.4230135) (1.6400000,1.4170660) (1.6450000,1.4111760) (1.6500000,1.4053426) (1.6550000,1.3995647) (1.6600000,1.3938416) (1.6650000,1.3881723) (1.6700000,1.3825560) (1.6750000,1.3769920) (1.6800000,1.3714793) (1.6850000,1.3660172) (1.6900000,1.3606049) (1.6950000,1.3552416) (1.7000000,1.3499267) (1.7050000,1.3446594) (1.7100000,1.3394389) (1.7150000,1.3342647) (1.7200000,1.3291359) (1.7250000,1.3240521) (1.7300000,1.3190124) (1.7350000,1.3140162) (1.7400000,1.3090630) (1.7450000,1.3041521) (1.7500000,1.2992830) (1.7550000,1.2944550) (1.7600000,1.2896675) (1.7650000,1.2849201) (1.7700000,1.2802121) (1.7750000,1.2755430) (1.7800000,1.2709123) (1.7850000,1.2663194) (1.7900000,1.2617639) (1.7950000,1.2572453) (1.8000000,1.2527630) (1.8050000,1.2483165) (1.8100000,1.2439055) (1.8150000,1.2395294) (1.8200000,1.2351878) (1.8250000,1.2308803) (1.8300000,1.2266063) (1.8350000,1.2223655) (1.8400000,1.2181574) (1.8450000,1.2139817) (1.8500000,1.2098379) (1.8550000,1.2057257) (1.8600000,1.2016445) (1.8650000,1.1975941) (1.8700000,1.1935741) (1.8750000,1.1895841) (1.8800000,1.1856237) (1.8850000,1.1816925) (1.8900000,1.1777903) (1.8950000,1.1739167) (1.9000000,1.1700713) (1.9050000,1.1662537) (1.9100000,1.1624638) (1.9150000,1.1587010) (1.9200000,1.1549652) (1.9250000,1.1512560) (1.9300000,1.1475731) (1.9350000,1.1439162) (1.9400000,1.1402850) (1.9450000,1.1366792) (1.9500000,1.1330985) (1.9550000,1.1295426) (1.9600000,1.1260113) (1.9650000,1.1225042) (1.9700000,1.1190212) (1.9750000,1.1155618) (1.9800000,1.1121260) (1.9850000,1.1087134) (1.9900000,1.1053237) (1.9950000,1.1019568) (2.0000000,1.0986123) (2.0050000,1.0952900) (2.0100000,1.0919897) (2.0150000,1.0887112) (2.0200000,1.0854542) (2.0250000,1.0822185) (2.0300000,1.0790038) (2.0350000,1.0758100) (2.0400000,1.0726368) (2.0450000,1.0694840) (2.0500000,1.0663514) (2.0550000,1.0632388) (2.0600000,1.0601460) (2.0650000,1.0570728) (2.0700000,1.0540189) (2.0750000,1.0509842) (2.0800000,1.0479686) (2.0850000,1.0449717) (2.0900000,1.0419934) (2.0950000,1.0390335) (2.1000000,1.0360919) (2.1050000,1.0331684) (2.1100000,1.0302627) (2.1150000,1.0273747) (2.1200000,1.0245043) (2.1250000,1.0216512) (2.1300000,1.0188154) (2.1350000,1.0159965) (2.1400000,1.0131945) (2.1450000,1.0104093) (2.1500000,1.0076405) (2.1550000,1.0048882) (2.1600000,1.0021520) (2.1650000,0.9994320) (2.1700000,0.9967278) (2.1750000,0.9940395) (2.1800000,0.9913668) (2.1850000,0.9887095) (2.1900000,0.9860676) (2.1950000,0.9834409) (2.2000000,0.9808293) (2.2050000,0.9782325) (2.2100000,0.9756506) (2.2150000,0.9730833) (2.2200000,0.9705305) (2.2250000,0.9679921) (2.2300000,0.9654680) (2.2350000,0.9629580) (2.2400000,0.9604620) (2.2450000,0.9579798) (2.2500000,0.9555114) (2.2550000,0.9530567) (2.2600000,0.9506155) (2.2650000,0.9481876) (2.2700000,0.9457731) (2.2750000,0.9433717) (2.2800000,0.9409833) (2.2850000,0.9386079) (2.2900000,0.9362453) (2.2950000,0.9338955) (2.3000000,0.9315582) (2.3050000,0.9292334) (2.3100000,0.9269211) (2.3150000,0.9246210) (2.3200000,0.9223330) (2.3250000,0.9200572) (2.3300000,0.9177934) (2.3350000,0.9155414) (2.3400000,0.9133012) (2.3450000,0.9110727) (2.3500000,0.9088558) (2.3550000,0.9066503) (2.3600000,0.9044563) (2.3650000,0.9022735) (2.3700000,0.9001020) (2.3750000,0.8979416) (2.3800000,0.8957922) (2.3850000,0.8936538) (2.3900000,0.8915262) (2.3950000,0.8894093) (2.4000000,0.8873032) (2.4050000,0.8852076) (2.4100000,0.8831226) (2.4150000,0.8810480) (2.4200000,0.8789837) (2.4250000,0.8769297) (2.4300000,0.8748858) (2.4350000,0.8728521) (2.4400000,0.8708284) (2.4450000,0.8688146) (2.4500000,0.8668107) (2.4550000,0.8648166) (2.4600000,0.8628322) (2.4650000,0.8608574) (2.4700000,0.8588922) (2.4750000,0.8569365) (2.4800000,0.8549902) (2.4850000,0.8530533) (2.4900000,0.8511256) (2.4950000,0.8492072) (2.5000000,0.8472979) (2.5050000,0.8453976) (2.5100000,0.8435064) (2.5150000,0.8416241) (2.5200000,0.8397507) (2.5250000,0.8378860) (2.5300000,0.8360301) (2.5350000,0.8341829) (2.5400000,0.8323443) (2.5450000,0.8305142) (2.5500000,0.8286927) (2.5550000,0.8268795) (2.5600000,0.8250747) (2.5650000,0.8232782) (2.5700000,0.8214900) (2.5750000,0.8197099) (2.5800000,0.8179380) (2.5850000,0.8161741) (2.5900000,0.8144182) (2.5950000,0.8126703) (2.6000000,0.8109302) (2.6050000,0.8091980) (2.6100000,0.8074736) (2.6150000,0.8057569) (2.6200000,0.8040479) (2.6250000,0.8023465) (2.6300000,0.8006526) (2.6350000,0.7989663) (2.6400000,0.7972874) (2.6450000,0.7956160) (2.6500000,0.7939519) (2.6550000,0.7922951) (2.6600000,0.7906455) (2.6650000,0.7890032) (2.6700000,0.7873680) (2.6750000,0.7857400) (2.6800000,0.7841190) (2.6850000,0.7825050) (2.6900000,0.7808979) (2.6950000,0.7792978) (2.7000000,0.7777046) (2.7050000,0.7761181) (2.7100000,0.7745385) (2.7150000,0.7729656) (2.7200000,0.7713994) (2.7250000,0.7698398) (2.7300000,0.7682868) (2.7350000,0.7667404) (2.7400000,0.7652005) (2.7450000,0.7636671) (2.7500000,0.7621401) (2.7550000,0.7606194) (2.7600000,0.7591051) (2.7650000,0.7575972) (2.7700000,0.7560955) (2.7750000,0.7546000) (2.7800000,0.7531106) (2.7850000,0.7516275) (2.7900000,0.7501504) (2.7950000,0.7486794) (2.8000000,0.7472144) (2.8050000,0.7457554) (2.8100000,0.7443023) (2.8150000,0.7428552) (2.8200000,0.7414139) (2.8250000,0.7399785) (2.8300000,0.7385488) (2.8350000,0.7371250) (2.8400000,0.7357068) (2.8450000,0.7342943) (2.8500000,0.7328875) (2.8550000,0.7314863) (2.8600000,0.7300907) (2.8650000,0.7287006) (2.8700000,0.7273161) (2.8750000,0.7259370) (2.8800000,0.7245634) (2.8850000,0.7231952) (2.8900000,0.7218323) (2.8950000,0.7204748) (2.9000000,0.7191227) (2.9050000,0.7177758) (2.9100000,0.7164341) (2.9150000,0.7150977) (2.9200000,0.7137665) (2.9250000,0.7124404) (2.9300000,0.7111194) (2.9350000,0.7098036) (2.9400000,0.7084928) (2.9450000,0.7071870) (2.9500000,0.7058862) (2.9550000,0.7045904) (2.9600000,0.7032996) (2.9650000,0.7020136) (2.9700000,0.7007326) (2.9750000,0.6994563) (2.9800000,0.6981850) (2.9850000,0.6969184) (2.9900000,0.6956566) (2.9950000,0.6943995) (3.0000000,0.6931472) (3.0050000,0.6918995) (3.0100000,0.6906565) (3.0150000,0.6894182) (3.0200000,0.6881844) (3.0250000,0.6869552) (3.0300000,0.6857306) (3.0350000,0.6845105) (3.0400000,0.6832949) (3.0450000,0.6820838) (3.0500000,0.6808771) (3.0550000,0.6796748) (3.0600000,0.6784770) (3.0650000,0.6772835) (3.0700000,0.6760944) (3.0750000,0.6749096) (3.0800000,0.6737291) (3.0850000,0.6725529) (3.0900000,0.6713809) (3.0950000,0.6702132) (3.1000000,0.6690496) (3.1050000,0.6678903) (3.1100000,0.6667351) (3.1150000,0.6655840) (3.1200000,0.6644371) (3.1250000,0.6632942) (3.1300000,0.6621554) (3.1350000,0.6610207) (3.1400000,0.6598900) (3.1450000,0.6587632) (3.1500000,0.6576405) (3.1550000,0.6565217) (3.1600000,0.6554069) (3.1650000,0.6542959) (3.1700000,0.6531889) (3.1750000,0.6520857) (3.1800000,0.6509864) (3.1850000,0.6498909) (3.1900000,0.6487992) (3.1950000,0.6477113) (3.2000000,0.6466272) (3.2050000,0.6455468) (3.2100000,0.6444701) (3.2150000,0.6433972) (3.2200000,0.6423279) (3.2250000,0.6412623) (3.2300000,0.6402004) (3.2350000,0.6391421) (3.2400000,0.6380874) (3.2450000,0.6370363) (3.2500000,0.6359888) (3.2550000,0.6349448) (3.2600000,0.6339043) (3.2650000,0.6328674) (3.2700000,0.6318340) (3.2750000,0.6308040) (3.2800000,0.6297776) (3.2850000,0.6287545) (3.2900000,0.6277349) (3.2950000,0.6267187) (3.3000000,0.6257059) (3.3050000,0.6246965) (3.3100000,0.6236904) (3.3150000,0.6226876) (3.3200000,0.6216882) (3.3250000,0.6206921) (3.3300000,0.6196993) (3.3350000,0.6187097) (3.3400000,0.6177234) (3.3450000,0.6167404) (3.3500000,0.6157605) (3.3550000,0.6147839) (3.3600000,0.6138104) (3.3650000,0.6128402) (3.3700000,0.6118731) (3.3750000,0.6109091) (3.3800000,0.6099482) (3.3850000,0.6089905) (3.3900000,0.6080359) (3.3950000,0.6070843) (3.4000000,0.6061358) (3.4050000,0.6051904) (3.4100000,0.6042479) (3.4150000,0.6033085) (3.4200000,0.6023722) (3.4250000,0.6014388) (3.4300000,0.6005083) (3.4350000,0.5995809) (3.4400000,0.5986563) (3.4450000,0.5977347) (3.4500000,0.5968161) (3.4550000,0.5959003) (3.4600000,0.5949874) (3.4650000,0.5940774) (3.4700000,0.5931703) (3.4750000,0.5922660) (3.4800000,0.5913645) (3.4850000,0.5904658) (3.4900000,0.5895700) (3.4950000,0.5886769) (3.5000000,0.5877867) (3.5050000,0.5868992) (3.5100000,0.5860144) (3.5150000,0.5851324) (3.5200000,0.5842531) (3.5250000,0.5833765) (3.5300000,0.5825026) (3.5350000,0.5816314) (3.5400000,0.5807629) (3.5450000,0.5798971) (3.5500000,0.5790339) (3.5550000,0.5781733) (3.5600000,0.5773154) (3.5650000,0.5764600) (3.5700000,0.5756073) (3.5750000,0.5747572) (3.5800000,0.5739096) (3.5850000,0.5730646) (3.5900000,0.5722221) (3.5950000,0.5713822) (3.6000000,0.5705449) (3.6050000,0.5697100) (3.6100000,0.5688776) (3.6150000,0.5680478) (3.6200000,0.5672204) (3.6250000,0.5663955) (3.6300000,0.5655730) (3.6350000,0.5647530) (3.6400000,0.5639354) (3.6450000,0.5631203) (3.6500000,0.5623076) (3.6550000,0.5614973) (3.6600000,0.5606893) (3.6650000,0.5598838) (3.6700000,0.5590806) (3.6750000,0.5582798) (3.6800000,0.5574813) (3.6850000,0.5566852) (3.6900000,0.5558914) (3.6950000,0.5550999) (3.7000000,0.5543107) (3.7050000,0.5535239) (3.7100000,0.5527393) (3.7150000,0.5519570) (3.7200000,0.5511769) (3.7250000,0.5503991) (3.7300000,0.5496236) (3.7350000,0.5488503) (3.7400000,0.5480792) (3.7450000,0.5473104) (3.7500000,0.5465437) (3.7550000,0.5457793) (3.7600000,0.5450170) (3.7650000,0.5442569) (3.7700000,0.5434990) (3.7750000,0.5427432) (3.7800000,0.5419896) (3.7850000,0.5412382) (3.7900000,0.5404888) (3.7950000,0.5397416) (3.8000000,0.5389965) (3.8050000,0.5382535) (3.8100000,0.5375126) (3.8150000,0.5367738) (3.8200000,0.5360370) (3.8250000,0.5353024) (3.8300000,0.5345698) (3.8350000,0.5338392) (3.8400000,0.5331107) (3.8450000,0.5323842) (3.8500000,0.5316597) (3.8550000,0.5309373) (3.8600000,0.5302168) (3.8650000,0.5294984) (3.8700000,0.5287819) (3.8750000,0.5280674) (3.8800000,0.5273549) (3.8850000,0.5266444) (3.8900000,0.5259358) (3.8950000,0.5252292) (3.9000000,0.5245245) (3.9050000,0.5238217) (3.9100000,0.5231209) (3.9150000,0.5224219) (3.9200000,0.5217249) (3.9250000,0.5210298) (3.9300000,0.5203366) (3.9350000,0.5196452) (3.9400000,0.5189557) (3.9450000,0.5182681) (3.9500000,0.5175824) (3.9550000,0.5168985) (3.9600000,0.5162165) (3.9650000,0.5155363) (3.9700000,0.5148579) (3.9750000,0.5141813) (3.9800000,0.5135066) (3.9850000,0.5128337) (3.9900000,0.5121625) (3.9950000,0.5114932) (4.0000000,0.5108256)};
\draw[orange!85!black,thick,dashed] plot coordinates {(-4.0000000,0.5108256) (-3.9950000,0.5114932) (-3.9900000,0.5121625) (-3.9850000,0.5128337) (-3.9800000,0.5135066) (-3.9750000,0.5141813) (-3.9700000,0.5148579) (-3.9650000,0.5155363) (-3.9600000,0.5162165) (-3.9550000,0.5168985) (-3.9500000,0.5175824) (-3.9450000,0.5182681) (-3.9400000,0.5189557) (-3.9350000,0.5196452) (-3.9300000,0.5203366) (-3.9250000,0.5210298) (-3.9200000,0.5217249) (-3.9150000,0.5224219) (-3.9100000,0.5231209) (-3.9050000,0.5238217) (-3.9000000,0.5245245) (-3.8950000,0.5252292) (-3.8900000,0.5259358) (-3.8850000,0.5266444) (-3.8800000,0.5273549) (-3.8750000,0.5280674) (-3.8700000,0.5287819) (-3.8650000,0.5294984) (-3.8600000,0.5302168) (-3.8550000,0.5309373) (-3.8500000,0.5316597) (-3.8450000,0.5323842) (-3.8400000,0.5331107) (-3.8350000,0.5338392) (-3.8300000,0.5345698) (-3.8250000,0.5353024) (-3.8200000,0.5360370) (-3.8150000,0.5367738) (-3.8100000,0.5375126) (-3.8050000,0.5382535) (-3.8000000,0.5389965) (-3.7950000,0.5397416) (-3.7900000,0.5404888) (-3.7850000,0.5412382) (-3.7800000,0.5419896) (-3.7750000,0.5427432) (-3.7700000,0.5434990) (-3.7650000,0.5442569) (-3.7600000,0.5450170) (-3.7550000,0.5457793) (-3.7500000,0.5465437) (-3.7450000,0.5473104) (-3.7400000,0.5480792) (-3.7350000,0.5488503) (-3.7300000,0.5496236) (-3.7250000,0.5503991) (-3.7200000,0.5511769) (-3.7150000,0.5519570) (-3.7100000,0.5527393) (-3.7050000,0.5535239) (-3.7000000,0.5543107) (-3.6950000,0.5550999) (-3.6900000,0.5558914) (-3.6850000,0.5566852) (-3.6800000,0.5574813) (-3.6750000,0.5582798) (-3.6700000,0.5590806) (-3.6650000,0.5598838) (-3.6600000,0.5606893) (-3.6550000,0.5614973) (-3.6500000,0.5623076) (-3.6450000,0.5631203) (-3.6400000,0.5639354) (-3.6350000,0.5647530) (-3.6300000,0.5655730) (-3.6250000,0.5663955) (-3.6200000,0.5672204) (-3.6150000,0.5680478) (-3.6100000,0.5688776) (-3.6050000,0.5697100) (-3.6000000,0.5705449) (-3.5950000,0.5713822) (-3.5900000,0.5722221) (-3.5850000,0.5730646) (-3.5800000,0.5739096) (-3.5750000,0.5747572) (-3.5700000,0.5756073) (-3.5650000,0.5764600) (-3.5600000,0.5773154) (-3.5550000,0.5781733) (-3.5500000,0.5790339) (-3.5450000,0.5798971) (-3.5400000,0.5807629) (-3.5350000,0.5816314) (-3.5300000,0.5825026) (-3.5250000,0.5833765) (-3.5200000,0.5842531) (-3.5150000,0.5851324) (-3.5100000,0.5860144) (-3.5050000,0.5868992) (-3.5000000,0.5877867) (-3.4950000,0.5886769) (-3.4900000,0.5895700) (-3.4850000,0.5904658) (-3.4800000,0.5913645) (-3.4750000,0.5922660) (-3.4700000,0.5931703) (-3.4650000,0.5940774) (-3.4600000,0.5949874) (-3.4550000,0.5959003) (-3.4500000,0.5968161) (-3.4450000,0.5977347) (-3.4400000,0.5986563) (-3.4350000,0.5995809) (-3.4300000,0.6005083) (-3.4250000,0.6014388) (-3.4200000,0.6023722) (-3.4150000,0.6033085) (-3.4100000,0.6042479) (-3.4050000,0.6051904) (-3.4000000,0.6061358) (-3.3950000,0.6070843) (-3.3900000,0.6080359) (-3.3850000,0.6089905) (-3.3800000,0.6099482) (-3.3750000,0.6109091) (-3.3700000,0.6118731) (-3.3650000,0.6128402) (-3.3600000,0.6138104) (-3.3550000,0.6147839) (-3.3500000,0.6157605) (-3.3450000,0.6167404) (-3.3400000,0.6177234) (-3.3350000,0.6187097) (-3.3300000,0.6196993) (-3.3250000,0.6206921) (-3.3200000,0.6216882) (-3.3150000,0.6226876) (-3.3100000,0.6236904) (-3.3050000,0.6246965) (-3.3000000,0.6257059) (-3.2950000,0.6267187) (-3.2900000,0.6277349) (-3.2850000,0.6287545) (-3.2800000,0.6297776) (-3.2750000,0.6308040) (-3.2700000,0.6318340) (-3.2650000,0.6328674) (-3.2600000,0.6339043) (-3.2550000,0.6349448) (-3.2500000,0.6359888) (-3.2450000,0.6370363) (-3.2400000,0.6380874) (-3.2350000,0.6391421) (-3.2300000,0.6402004) (-3.2250000,0.6412623) (-3.2200000,0.6423279) (-3.2150000,0.6433972) (-3.2100000,0.6444701) (-3.2050000,0.6455468) (-3.2000000,0.6466272) (-3.1950000,0.6477113) (-3.1900000,0.6487992) (-3.1850000,0.6498909) (-3.1800000,0.6509864) (-3.1750000,0.6520857) (-3.1700000,0.6531889) (-3.1650000,0.6542959) (-3.1600000,0.6554069) (-3.1550000,0.6565217) (-3.1500000,0.6576405) (-3.1450000,0.6587632) (-3.1400000,0.6598900) (-3.1350000,0.6610207) (-3.1300000,0.6621554) (-3.1250000,0.6632942) (-3.1200000,0.6644371) (-3.1150000,0.6655840) (-3.1100000,0.6667351) (-3.1050000,0.6678903) (-3.1000000,0.6690496) (-3.0950000,0.6702132) (-3.0900000,0.6713809) (-3.0850000,0.6725529) (-3.0800000,0.6737291) (-3.0750000,0.6749096) (-3.0700000,0.6760944) (-3.0650000,0.6772835) (-3.0600000,0.6784770) (-3.0550000,0.6796748) (-3.0500000,0.6808771) (-3.0450000,0.6820838) (-3.0400000,0.6832949) (-3.0350000,0.6845105) (-3.0300000,0.6857306) (-3.0250000,0.6869552) (-3.0200000,0.6881844) (-3.0150000,0.6894182) (-3.0100000,0.6906565) (-3.0050000,0.6918995) (-3.0000000,0.6931472) (-2.9950000,0.6943995) (-2.9900000,0.6956566) (-2.9850000,0.6969184) (-2.9800000,0.6981850) (-2.9750000,0.6994563) (-2.9700000,0.7007326) (-2.9650000,0.7020136) (-2.9600000,0.7032996) (-2.9550000,0.7045904) (-2.9500000,0.7058862) (-2.9450000,0.7071870) (-2.9400000,0.7084928) (-2.9350000,0.7098036) (-2.9300000,0.7111194) (-2.9250000,0.7124404) (-2.9200000,0.7137665) (-2.9150000,0.7150977) (-2.9100000,0.7164341) (-2.9050000,0.7177758) (-2.9000000,0.7191227) (-2.8950000,0.7204748) (-2.8900000,0.7218323) (-2.8850000,0.7231952) (-2.8800000,0.7245634) (-2.8750000,0.7259370) (-2.8700000,0.7273161) (-2.8650000,0.7287006) (-2.8600000,0.7300907) (-2.8550000,0.7314863) (-2.8500000,0.7328875) (-2.8450000,0.7342943) (-2.8400000,0.7357068) (-2.8350000,0.7371250) (-2.8300000,0.7385488) (-2.8250000,0.7399785) (-2.8200000,0.7414139) (-2.8150000,0.7428552) (-2.8100000,0.7443023) (-2.8050000,0.7457554) (-2.8000000,0.7472144) (-2.7950000,0.7486794) (-2.7900000,0.7501504) (-2.7850000,0.7516275) (-2.7800000,0.7531106) (-2.7750000,0.7546000) (-2.7700000,0.7560955) (-2.7650000,0.7575972) (-2.7600000,0.7591051) (-2.7550000,0.7606194) (-2.7500000,0.7621401) (-2.7450000,0.7636671) (-2.7400000,0.7652005) (-2.7350000,0.7667404) (-2.7300000,0.7682868) (-2.7250000,0.7698398) (-2.7200000,0.7713994) (-2.7150000,0.7729656) (-2.7100000,0.7745385) (-2.7050000,0.7761181) (-2.7000000,0.7777046) (-2.6950000,0.7792978) (-2.6900000,0.7808979) (-2.6850000,0.7825050) (-2.6800000,0.7841190) (-2.6750000,0.7857400) (-2.6700000,0.7873680) (-2.6650000,0.7890032) (-2.6600000,0.7906455) (-2.6550000,0.7922951) (-2.6500000,0.7939519) (-2.6450000,0.7956160) (-2.6400000,0.7972874) (-2.6350000,0.7989663) (-2.6300000,0.8006526) (-2.6250000,0.8023465) (-2.6200000,0.8040479) (-2.6150000,0.8057569) (-2.6100000,0.8074736) (-2.6050000,0.8091980) (-2.6000000,0.8109302) (-2.5950000,0.8126703) (-2.5900000,0.8144182) (-2.5850000,0.8161741) (-2.5800000,0.8179380) (-2.5750000,0.8197099) (-2.5700000,0.8214900) (-2.5650000,0.8232782) (-2.5600000,0.8250747) (-2.5550000,0.8268795) (-2.5500000,0.8286927) (-2.5450000,0.8305142) (-2.5400000,0.8323443) (-2.5350000,0.8341829) (-2.5300000,0.8360301) (-2.5250000,0.8378860) (-2.5200000,0.8397507) (-2.5150000,0.8416241) (-2.5100000,0.8435064) (-2.5050000,0.8453976) (-2.5000000,0.8472979) (-2.4950000,0.8492072) (-2.4900000,0.8511256) (-2.4850000,0.8530533) (-2.4800000,0.8549902) (-2.4750000,0.8569365) (-2.4700000,0.8588922) (-2.4650000,0.8608574) (-2.4600000,0.8628322) (-2.4550000,0.8648166) (-2.4500000,0.8668107) (-2.4450000,0.8688146) (-2.4400000,0.8708284) (-2.4350000,0.8728521) (-2.4300000,0.8748858) (-2.4250000,0.8769297) (-2.4200000,0.8789837) (-2.4150000,0.8810480) (-2.4100000,0.8831226) (-2.4050000,0.8852076) (-2.4000000,0.8873032) (-2.3950000,0.8894093) (-2.3900000,0.8915262) (-2.3850000,0.8936538) (-2.3800000,0.8957922) (-2.3750000,0.8979416) (-2.3700000,0.9001020) (-2.3650000,0.9022735) (-2.3600000,0.9044563) (-2.3550000,0.9066503) (-2.3500000,0.9088558) (-2.3450000,0.9110727) (-2.3400000,0.9133012) (-2.3350000,0.9155414) (-2.3300000,0.9177934) (-2.3250000,0.9200572) (-2.3200000,0.9223330) (-2.3150000,0.9246210) (-2.3100000,0.9269211) (-2.3050000,0.9292334) (-2.3000000,0.9315582) (-2.2950000,0.9338955) (-2.2900000,0.9362453) (-2.2850000,0.9386079) (-2.2800000,0.9409833) (-2.2750000,0.9433717) (-2.2700000,0.9457731) (-2.2650000,0.9481876) (-2.2600000,0.9506155) (-2.2550000,0.9530567) (-2.2500000,0.9555114) (-2.2450000,0.9579798) (-2.2400000,0.9604620) (-2.2350000,0.9629580) (-2.2300000,0.9654680) (-2.2250000,0.9679921) (-2.2200000,0.9705305) (-2.2150000,0.9730833) (-2.2100000,0.9756506) (-2.2050000,0.9782325) (-2.2000000,0.9808293) (-2.1950000,0.9834409) (-2.1900000,0.9860676) (-2.1850000,0.9887095) (-2.1800000,0.9913668) (-2.1750000,0.9940395) (-2.1700000,0.9967278) (-2.1650000,0.9994320) (-2.1600000,1.0021520) (-2.1550000,1.0048882) (-2.1500000,1.0076405) (-2.1450000,1.0104093) (-2.1400000,1.0131945) (-2.1350000,1.0159965) (-2.1300000,1.0188154) (-2.1250000,1.0216512) (-2.1200000,1.0245043) (-2.1150000,1.0273747) (-2.1100000,1.0302627) (-2.1050000,1.0331684) (-2.1000000,1.0360919) (-2.0950000,1.0390335) (-2.0900000,1.0419934) (-2.0850000,1.0449717) (-2.0800000,1.0479686) (-2.0750000,1.0509842) (-2.0700000,1.0540189) (-2.0650000,1.0570728) (-2.0600000,1.0601460) (-2.0550000,1.0632388) (-2.0500000,1.0663514) (-2.0450000,1.0694840) (-2.0400000,1.0726368) (-2.0350000,1.0758100) (-2.0300000,1.0790038) (-2.0250000,1.0822185) (-2.0200000,1.0854542) (-2.0150000,1.0887112) (-2.0100000,1.0919897) (-2.0050000,1.0952900) (-2.0000000,1.0986123) (-1.9950000,1.1019568) (-1.9900000,1.1053237) (-1.9850000,1.1087134) (-1.9800000,1.1121260) (-1.9750000,1.1155618) (-1.9700000,1.1190212) (-1.9650000,1.1225042) (-1.9600000,1.1260113) (-1.9550000,1.1295426) (-1.9500000,1.1330985) (-1.9450000,1.1366792) (-1.9400000,1.1402850) (-1.9350000,1.1439162) (-1.9300000,1.1475731) (-1.9250000,1.1512560) (-1.9200000,1.1549652) (-1.9150000,1.1587010) (-1.9100000,1.1624638) (-1.9050000,1.1662537) (-1.9000000,1.1700713) (-1.8950000,1.1739167) (-1.8900000,1.1777903) (-1.8850000,1.1816925) (-1.8800000,1.1856237) (-1.8750000,1.1895841) (-1.8700000,1.1935741) (-1.8650000,1.1975941) (-1.8600000,1.2016445) (-1.8550000,1.2057257) (-1.8500000,1.2098379) (-1.8450000,1.2139817) (-1.8400000,1.2181574) (-1.8350000,1.2223655) (-1.8300000,1.2266063) (-1.8250000,1.2308803) (-1.8200000,1.2351878) (-1.8150000,1.2395294) (-1.8100000,1.2439055) (-1.8050000,1.2483165) (-1.8000000,1.2527630) (-1.7950000,1.2572453) (-1.7900000,1.2617639) (-1.7850000,1.2663194) (-1.7800000,1.2709123) (-1.7750000,1.2755430) (-1.7700000,1.2802121) (-1.7650000,1.2849201) (-1.7600000,1.2896675) (-1.7550000,1.2944550) (-1.7500000,1.2992830) (-1.7450000,1.3041521) (-1.7400000,1.3090630) (-1.7350000,1.3140162) (-1.7300000,1.3190124) (-1.7250000,1.3240521) (-1.7200000,1.3291359) (-1.7150000,1.3342647) (-1.7100000,1.3394389) (-1.7050000,1.3446594) (-1.7000000,1.3499267) (-1.6950000,1.3552416) (-1.6900000,1.3606049) (-1.6850000,1.3660172) (-1.6800000,1.3714793) (-1.6750000,1.3769920) (-1.6700000,1.3825560) (-1.6650000,1.3881723) (-1.6600000,1.3938416) (-1.6550000,1.3995647) (-1.6500000,1.4053426) (-1.6450000,1.4111760) (-1.6400000,1.4170660) (-1.6350000,1.4230135) (-1.6300000,1.4290193) (-1.6250000,1.4350845) (-1.6200000,1.4412101) (-1.6150000,1.4473971) (-1.6100000,1.4536465) (-1.6050000,1.4599595) (-1.6000000,1.4663371) (-1.5950000,1.4727804) (-1.5900000,1.4792906) (-1.5850000,1.4858689) (-1.5800000,1.4925166) (-1.5750000,1.4992348) (-1.5700000,1.5060248) (-1.5650000,1.5128880) (-1.5600000,1.5198258) (-1.5550000,1.5268394) (-1.5500000,1.5339304) (-1.5450000,1.5411001) (-1.5400000,1.5483502) (-1.5350000,1.5556822) (-1.5300000,1.5630976) (-1.5250000,1.5705981) (-1.5200000,1.5781854) (-1.5150000,1.5858612) (-1.5100000,1.5936273) (-1.5050000,1.6014856) (-1.5000000,1.6094379) (-1.4950000,1.6174862) (-1.4900000,1.6256326) (-1.4850000,1.6338790) (-1.4800000,1.6422277) (-1.4750000,1.6506809) (-1.4700000,1.6592407) (-1.4650000,1.6679097) (-1.4600000,1.6766901) (-1.4550000,1.6855846) (-1.4500000,1.6945957) (-1.4450000,1.7037261) (-1.4400000,1.7129786) (-1.4350000,1.7223560) (-1.4300000,1.7318613) (-1.4250000,1.7414976) (-1.4200000,1.7512681) (-1.4150000,1.7611760) (-1.4100000,1.7712249) (-1.4050000,1.7814181) (-1.4000000,1.7917595) (-1.3950000,1.8022527) (-1.3900000,1.8129019) (-1.3850000,1.8237111) (-1.3800000,1.8346845) (-1.3750000,1.8458267) (-1.3700000,1.8571422) (-1.3650000,1.8686359) (-1.3600000,1.8803129) (-1.3550000,1.8921782) (-1.3500000,1.9042375) (-1.3450000,1.9164963) (-1.3400000,1.9289606) (-1.3350000,1.9416366) (-1.3300000,1.9545309) (-1.3250000,1.9676501) (-1.3200000,1.9810015) (-1.3150000,1.9945923) (-1.3100000,2.0084305) (-1.3050000,2.0225242) (-1.3000000,2.0368819) (-1.2950000,2.0515128) (-1.2900000,2.0664262) (-1.2850000,2.0816321) (-1.2800000,2.0971411) (-1.2750000,2.1129642) (-1.2700000,2.1291132) (-1.2650000,2.1456002) (-1.2600000,2.1624385) (-1.2550000,2.1796417) (-1.2500000,2.1972246) (-1.2450000,2.2152026) (-1.2400000,2.2335922) (-1.2350000,2.2524110) (-1.2300000,2.2716776) (-1.2250000,2.2914118) (-1.2200000,2.3116349) (-1.2150000,2.3323697) (-1.2100000,2.3536403) (-1.2050000,2.3754728) (-1.2000000,2.3978953) (-1.1950000,2.4209378) (-1.1900000,2.4446328) (-1.1850000,2.4690153) (-1.1800000,2.4941233) (-1.1750000,2.5199980) (-1.1700000,2.5466840) (-1.1650000,2.5742302) (-1.1600000,2.6026897) (-1.1550000,2.6321209) (-1.1500000,2.6625878) (-1.1450000,2.6941611) (-1.1400000,2.7269187) (-1.1350000,2.7609471) (-1.1300000,2.7963428) (-1.1250000,2.8332133) (-1.1200000,2.8716796) (-1.1150000,2.9118780) (-1.1100000,2.9539629) (-1.1050000,2.9981104) (-1.1000000,3.0445224) (-1.0950000,3.0934319) (-1.0900000,3.1451097) (-1.0850000,3.1998729) (-1.0800000,3.2580965) (-1.0750000,3.3202283) (-1.0700000,3.3868086) (-1.0650000,3.4584982) (-1.0600000,3.5361167) (-1.0550000,3.6206979) (-1.0500000,3.7135721) (-1.0450000,3.8164906) (-1.0400000,3.9318256)};
\draw[orange!85!black,thick,dashed] plot coordinates {(1.0400000,-3.9318256) (1.0450000,-3.8164906) (1.0500000,-3.7135721) (1.0550000,-3.6206979) (1.0600000,-3.5361167) (1.0650000,-3.4584982) (1.0700000,-3.3868086) (1.0750000,-3.3202283) (1.0800000,-3.2580965) (1.0850000,-3.1998729) (1.0900000,-3.1451097) (1.0950000,-3.0934319) (1.1000000,-3.0445224) (1.1050000,-2.9981104) (1.1100000,-2.9539629) (1.1150000,-2.9118780) (1.1200000,-2.8716796) (1.1250000,-2.8332133) (1.1300000,-2.7963428) (1.1350000,-2.7609471) (1.1400000,-2.7269187) (1.1450000,-2.6941611) (1.1500000,-2.6625878) (1.1550000,-2.6321209) (1.1600000,-2.6026897) (1.1650000,-2.5742302) (1.1700000,-2.5466840) (1.1750000,-2.5199980) (1.1800000,-2.4941233) (1.1850000,-2.4690153) (1.1900000,-2.4446328) (1.1950000,-2.4209378) (1.2000000,-2.3978953) (1.2050000,-2.3754728) (1.2100000,-2.3536403) (1.2150000,-2.3323697) (1.2200000,-2.3116349) (1.2250000,-2.2914118) (1.2300000,-2.2716776) (1.2350000,-2.2524110) (1.2400000,-2.2335922) (1.2450000,-2.2152026) (1.2500000,-2.1972246) (1.2550000,-2.1796417) (1.2600000,-2.1624385) (1.2650000,-2.1456002) (1.2700000,-2.1291132) (1.2750000,-2.1129642) (1.2800000,-2.0971411) (1.2850000,-2.0816321) (1.2900000,-2.0664262) (1.2950000,-2.0515128) (1.3000000,-2.0368819) (1.3050000,-2.0225242) (1.3100000,-2.0084305) (1.3150000,-1.9945923) (1.3200000,-1.9810015) (1.3250000,-1.9676501) (1.3300000,-1.9545309) (1.3350000,-1.9416366) (1.3400000,-1.9289606) (1.3450000,-1.9164963) (1.3500000,-1.9042375) (1.3550000,-1.8921782) (1.3600000,-1.8803129) (1.3650000,-1.8686359) (1.3700000,-1.8571422) (1.3750000,-1.8458267) (1.3800000,-1.8346845) (1.3850000,-1.8237111) (1.3900000,-1.8129019) (1.3950000,-1.8022527) (1.4000000,-1.7917595) (1.4050000,-1.7814181) (1.4100000,-1.7712249) (1.4150000,-1.7611760) (1.4200000,-1.7512681) (1.4250000,-1.7414976) (1.4300000,-1.7318613) (1.4350000,-1.7223560) (1.4400000,-1.7129786) (1.4450000,-1.7037261) (1.4500000,-1.6945957) (1.4550000,-1.6855846) (1.4600000,-1.6766901) (1.4650000,-1.6679097) (1.4700000,-1.6592407) (1.4750000,-1.6506809) (1.4800000,-1.6422277) (1.4850000,-1.6338790) (1.4900000,-1.6256326) (1.4950000,-1.6174862) (1.5000000,-1.6094379) (1.5050000,-1.6014856) (1.5100000,-1.5936273) (1.5150000,-1.5858612) (1.5200000,-1.5781854) (1.5250000,-1.5705981) (1.5300000,-1.5630976) (1.5350000,-1.5556822) (1.5400000,-1.5483502) (1.5450000,-1.5411001) (1.5500000,-1.5339304) (1.5550000,-1.5268394) (1.5600000,-1.5198258) (1.5650000,-1.5128880) (1.5700000,-1.5060248) (1.5750000,-1.4992348) (1.5800000,-1.4925166) (1.5850000,-1.4858689) (1.5900000,-1.4792906) (1.5950000,-1.4727804) (1.6000000,-1.4663371) (1.6050000,-1.4599595) (1.6100000,-1.4536465) (1.6150000,-1.4473971) (1.6200000,-1.4412101) (1.6250000,-1.4350845) (1.6300000,-1.4290193) (1.6350000,-1.4230135) (1.6400000,-1.4170660) (1.6450000,-1.4111760) (1.6500000,-1.4053426) (1.6550000,-1.3995647) (1.6600000,-1.3938416) (1.6650000,-1.3881723) (1.6700000,-1.3825560) (1.6750000,-1.3769920) (1.6800000,-1.3714793) (1.6850000,-1.3660172) (1.6900000,-1.3606049) (1.6950000,-1.3552416) (1.7000000,-1.3499267) (1.7050000,-1.3446594) (1.7100000,-1.3394389) (1.7150000,-1.3342647) (1.7200000,-1.3291359) (1.7250000,-1.3240521) (1.7300000,-1.3190124) (1.7350000,-1.3140162) (1.7400000,-1.3090630) (1.7450000,-1.3041521) (1.7500000,-1.2992830) (1.7550000,-1.2944550) (1.7600000,-1.2896675) (1.7650000,-1.2849201) (1.7700000,-1.2802121) (1.7750000,-1.2755430) (1.7800000,-1.2709123) (1.7850000,-1.2663194) (1.7900000,-1.2617639) (1.7950000,-1.2572453) (1.8000000,-1.2527630) (1.8050000,-1.2483165) (1.8100000,-1.2439055) (1.8150000,-1.2395294) (1.8200000,-1.2351878) (1.8250000,-1.2308803) (1.8300000,-1.2266063) (1.8350000,-1.2223655) (1.8400000,-1.2181574) (1.8450000,-1.2139817) (1.8500000,-1.2098379) (1.8550000,-1.2057257) (1.8600000,-1.2016445) (1.8650000,-1.1975941) (1.8700000,-1.1935741) (1.8750000,-1.1895841) (1.8800000,-1.1856237) (1.8850000,-1.1816925) (1.8900000,-1.1777903) (1.8950000,-1.1739167) (1.9000000,-1.1700713) (1.9050000,-1.1662537) (1.9100000,-1.1624638) (1.9150000,-1.1587010) (1.9200000,-1.1549652) (1.9250000,-1.1512560) (1.9300000,-1.1475731) (1.9350000,-1.1439162) (1.9400000,-1.1402850) (1.9450000,-1.1366792) (1.9500000,-1.1330985) (1.9550000,-1.1295426) (1.9600000,-1.1260113) (1.9650000,-1.1225042) (1.9700000,-1.1190212) (1.9750000,-1.1155618) (1.9800000,-1.1121260) (1.9850000,-1.1087134) (1.9900000,-1.1053237) (1.9950000,-1.1019568) (2.0000000,-1.0986123) (2.0050000,-1.0952900) (2.0100000,-1.0919897) (2.0150000,-1.0887112) (2.0200000,-1.0854542) (2.0250000,-1.0822185) (2.0300000,-1.0790038) (2.0350000,-1.0758100) (2.0400000,-1.0726368) (2.0450000,-1.0694840) (2.0500000,-1.0663514) (2.0550000,-1.0632388) (2.0600000,-1.0601460) (2.0650000,-1.0570728) (2.0700000,-1.0540189) (2.0750000,-1.0509842) (2.0800000,-1.0479686) (2.0850000,-1.0449717) (2.0900000,-1.0419934) (2.0950000,-1.0390335) (2.1000000,-1.0360919) (2.1050000,-1.0331684) (2.1100000,-1.0302627) (2.1150000,-1.0273747) (2.1200000,-1.0245043) (2.1250000,-1.0216512) (2.1300000,-1.0188154) (2.1350000,-1.0159965) (2.1400000,-1.0131945) (2.1450000,-1.0104093) (2.1500000,-1.0076405) (2.1550000,-1.0048882) (2.1600000,-1.0021520) (2.1650000,-0.9994320) (2.1700000,-0.9967278) (2.1750000,-0.9940395) (2.1800000,-0.9913668) (2.1850000,-0.9887095) (2.1900000,-0.9860676) (2.1950000,-0.9834409) (2.2000000,-0.9808293) (2.2050000,-0.9782325) (2.2100000,-0.9756506) (2.2150000,-0.9730833) (2.2200000,-0.9705305) (2.2250000,-0.9679921) (2.2300000,-0.9654680) (2.2350000,-0.9629580) (2.2400000,-0.9604620) (2.2450000,-0.9579798) (2.2500000,-0.9555114) (2.2550000,-0.9530567) (2.2600000,-0.9506155) (2.2650000,-0.9481876) (2.2700000,-0.9457731) (2.2750000,-0.9433717) (2.2800000,-0.9409833) (2.2850000,-0.9386079) (2.2900000,-0.9362453) (2.2950000,-0.9338955) (2.3000000,-0.9315582) (2.3050000,-0.9292334) (2.3100000,-0.9269211) (2.3150000,-0.9246210) (2.3200000,-0.9223330) (2.3250000,-0.9200572) (2.3300000,-0.9177934) (2.3350000,-0.9155414) (2.3400000,-0.9133012) (2.3450000,-0.9110727) (2.3500000,-0.9088558) (2.3550000,-0.9066503) (2.3600000,-0.9044563) (2.3650000,-0.9022735) (2.3700000,-0.9001020) (2.3750000,-0.8979416) (2.3800000,-0.8957922) (2.3850000,-0.8936538) (2.3900000,-0.8915262) (2.3950000,-0.8894093) (2.4000000,-0.8873032) (2.4050000,-0.8852076) (2.4100000,-0.8831226) (2.4150000,-0.8810480) (2.4200000,-0.8789837) (2.4250000,-0.8769297) (2.4300000,-0.8748858) (2.4350000,-0.8728521) (2.4400000,-0.8708284) (2.4450000,-0.8688146) (2.4500000,-0.8668107) (2.4550000,-0.8648166) (2.4600000,-0.8628322) (2.4650000,-0.8608574) (2.4700000,-0.8588922) (2.4750000,-0.8569365) (2.4800000,-0.8549902) (2.4850000,-0.8530533) (2.4900000,-0.8511256) (2.4950000,-0.8492072) (2.5000000,-0.8472979) (2.5050000,-0.8453976) (2.5100000,-0.8435064) (2.5150000,-0.8416241) (2.5200000,-0.8397507) (2.5250000,-0.8378860) (2.5300000,-0.8360301) (2.5350000,-0.8341829) (2.5400000,-0.8323443) (2.5450000,-0.8305142) (2.5500000,-0.8286927) (2.5550000,-0.8268795) (2.5600000,-0.8250747) (2.5650000,-0.8232782) (2.5700000,-0.8214900) (2.5750000,-0.8197099) (2.5800000,-0.8179380) (2.5850000,-0.8161741) (2.5900000,-0.8144182) (2.5950000,-0.8126703) (2.6000000,-0.8109302) (2.6050000,-0.8091980) (2.6100000,-0.8074736) (2.6150000,-0.8057569) (2.6200000,-0.8040479) (2.6250000,-0.8023465) (2.6300000,-0.8006526) (2.6350000,-0.7989663) (2.6400000,-0.7972874) (2.6450000,-0.7956160) (2.6500000,-0.7939519) (2.6550000,-0.7922951) (2.6600000,-0.7906455) (2.6650000,-0.7890032) (2.6700000,-0.7873680) (2.6750000,-0.7857400) (2.6800000,-0.7841190) (2.6850000,-0.7825050) (2.6900000,-0.7808979) (2.6950000,-0.7792978) (2.7000000,-0.7777046) (2.7050000,-0.7761181) (2.7100000,-0.7745385) (2.7150000,-0.7729656) (2.7200000,-0.7713994) (2.7250000,-0.7698398) (2.7300000,-0.7682868) (2.7350000,-0.7667404) (2.7400000,-0.7652005) (2.7450000,-0.7636671) (2.7500000,-0.7621401) (2.7550000,-0.7606194) (2.7600000,-0.7591051) (2.7650000,-0.7575972) (2.7700000,-0.7560955) (2.7750000,-0.7546000) (2.7800000,-0.7531106) (2.7850000,-0.7516275) (2.7900000,-0.7501504) (2.7950000,-0.7486794) (2.8000000,-0.7472144) (2.8050000,-0.7457554) (2.8100000,-0.7443023) (2.8150000,-0.7428552) (2.8200000,-0.7414139) (2.8250000,-0.7399785) (2.8300000,-0.7385488) (2.8350000,-0.7371250) (2.8400000,-0.7357068) (2.8450000,-0.7342943) (2.8500000,-0.7328875) (2.8550000,-0.7314863) (2.8600000,-0.7300907) (2.8650000,-0.7287006) (2.8700000,-0.7273161) (2.8750000,-0.7259370) (2.8800000,-0.7245634) (2.8850000,-0.7231952) (2.8900000,-0.7218323) (2.8950000,-0.7204748) (2.9000000,-0.7191227) (2.9050000,-0.7177758) (2.9100000,-0.7164341) (2.9150000,-0.7150977) (2.9200000,-0.7137665) (2.9250000,-0.7124404) (2.9300000,-0.7111194) (2.9350000,-0.7098036) (2.9400000,-0.7084928) (2.9450000,-0.7071870) (2.9500000,-0.7058862) (2.9550000,-0.7045904) (2.9600000,-0.7032996) (2.9650000,-0.7020136) (2.9700000,-0.7007326) (2.9750000,-0.6994563) (2.9800000,-0.6981850) (2.9850000,-0.6969184) (2.9900000,-0.6956566) (2.9950000,-0.6943995) (3.0000000,-0.6931472) (3.0050000,-0.6918995) (3.0100000,-0.6906565) (3.0150000,-0.6894182) (3.0200000,-0.6881844) (3.0250000,-0.6869552) (3.0300000,-0.6857306) (3.0350000,-0.6845105) (3.0400000,-0.6832949) (3.0450000,-0.6820838) (3.0500000,-0.6808771) (3.0550000,-0.6796748) (3.0600000,-0.6784770) (3.0650000,-0.6772835) (3.0700000,-0.6760944) (3.0750000,-0.6749096) (3.0800000,-0.6737291) (3.0850000,-0.6725529) (3.0900000,-0.6713809) (3.0950000,-0.6702132) (3.1000000,-0.6690496) (3.1050000,-0.6678903) (3.1100000,-0.6667351) (3.1150000,-0.6655840) (3.1200000,-0.6644371) (3.1250000,-0.6632942) (3.1300000,-0.6621554) (3.1350000,-0.6610207) (3.1400000,-0.6598900) (3.1450000,-0.6587632) (3.1500000,-0.6576405) (3.1550000,-0.6565217) (3.1600000,-0.6554069) (3.1650000,-0.6542959) (3.1700000,-0.6531889) (3.1750000,-0.6520857) (3.1800000,-0.6509864) (3.1850000,-0.6498909) (3.1900000,-0.6487992) (3.1950000,-0.6477113) (3.2000000,-0.6466272) (3.2050000,-0.6455468) (3.2100000,-0.6444701) (3.2150000,-0.6433972) (3.2200000,-0.6423279) (3.2250000,-0.6412623) (3.2300000,-0.6402004) (3.2350000,-0.6391421) (3.2400000,-0.6380874) (3.2450000,-0.6370363) (3.2500000,-0.6359888) (3.2550000,-0.6349448) (3.2600000,-0.6339043) (3.2650000,-0.6328674) (3.2700000,-0.6318340) (3.2750000,-0.6308040) (3.2800000,-0.6297776) (3.2850000,-0.6287545) (3.2900000,-0.6277349) (3.2950000,-0.6267187) (3.3000000,-0.6257059) (3.3050000,-0.6246965) (3.3100000,-0.6236904) (3.3150000,-0.6226876) (3.3200000,-0.6216882) (3.3250000,-0.6206921) (3.3300000,-0.6196993) (3.3350000,-0.6187097) (3.3400000,-0.6177234) (3.3450000,-0.6167404) (3.3500000,-0.6157605) (3.3550000,-0.6147839) (3.3600000,-0.6138104) (3.3650000,-0.6128402) (3.3700000,-0.6118731) (3.3750000,-0.6109091) (3.3800000,-0.6099482) (3.3850000,-0.6089905) (3.3900000,-0.6080359) (3.3950000,-0.6070843) (3.4000000,-0.6061358) (3.4050000,-0.6051904) (3.4100000,-0.6042479) (3.4150000,-0.6033085) (3.4200000,-0.6023722) (3.4250000,-0.6014388) (3.4300000,-0.6005083) (3.4350000,-0.5995809) (3.4400000,-0.5986563) (3.4450000,-0.5977347) (3.4500000,-0.5968161) (3.4550000,-0.5959003) (3.4600000,-0.5949874) (3.4650000,-0.5940774) (3.4700000,-0.5931703) (3.4750000,-0.5922660) (3.4800000,-0.5913645) (3.4850000,-0.5904658) (3.4900000,-0.5895700) (3.4950000,-0.5886769) (3.5000000,-0.5877867) (3.5050000,-0.5868992) (3.5100000,-0.5860144) (3.5150000,-0.5851324) (3.5200000,-0.5842531) (3.5250000,-0.5833765) (3.5300000,-0.5825026) (3.5350000,-0.5816314) (3.5400000,-0.5807629) (3.5450000,-0.5798971) (3.5500000,-0.5790339) (3.5550000,-0.5781733) (3.5600000,-0.5773154) (3.5650000,-0.5764600) (3.5700000,-0.5756073) (3.5750000,-0.5747572) (3.5800000,-0.5739096) (3.5850000,-0.5730646) (3.5900000,-0.5722221) (3.5950000,-0.5713822) (3.6000000,-0.5705449) (3.6050000,-0.5697100) (3.6100000,-0.5688776) (3.6150000,-0.5680478) (3.6200000,-0.5672204) (3.6250000,-0.5663955) (3.6300000,-0.5655730) (3.6350000,-0.5647530) (3.6400000,-0.5639354) (3.6450000,-0.5631203) (3.6500000,-0.5623076) (3.6550000,-0.5614973) (3.6600000,-0.5606893) (3.6650000,-0.5598838) (3.6700000,-0.5590806) (3.6750000,-0.5582798) (3.6800000,-0.5574813) (3.6850000,-0.5566852) (3.6900000,-0.5558914) (3.6950000,-0.5550999) (3.7000000,-0.5543107) (3.7050000,-0.5535239) (3.7100000,-0.5527393) (3.7150000,-0.5519570) (3.7200000,-0.5511769) (3.7250000,-0.5503991) (3.7300000,-0.5496236) (3.7350000,-0.5488503) (3.7400000,-0.5480792) (3.7450000,-0.5473104) (3.7500000,-0.5465437) (3.7550000,-0.5457793) (3.7600000,-0.5450170) (3.7650000,-0.5442569) (3.7700000,-0.5434990) (3.7750000,-0.5427432) (3.7800000,-0.5419896) (3.7850000,-0.5412382) (3.7900000,-0.5404888) (3.7950000,-0.5397416) (3.8000000,-0.5389965) (3.8050000,-0.5382535) (3.8100000,-0.5375126) (3.8150000,-0.5367738) (3.8200000,-0.5360370) (3.8250000,-0.5353024) (3.8300000,-0.5345698) (3.8350000,-0.5338392) (3.8400000,-0.5331107) (3.8450000,-0.5323842) (3.8500000,-0.5316597) (3.8550000,-0.5309373) (3.8600000,-0.5302168) (3.8650000,-0.5294984) (3.8700000,-0.5287819) (3.8750000,-0.5280674) (3.8800000,-0.5273549) (3.8850000,-0.5266444) (3.8900000,-0.5259358) (3.8950000,-0.5252292) (3.9000000,-0.5245245) (3.9050000,-0.5238217) (3.9100000,-0.5231209) (3.9150000,-0.5224219) (3.9200000,-0.5217249) (3.9250000,-0.5210298) (3.9300000,-0.5203366) (3.9350000,-0.5196452) (3.9400000,-0.5189557) (3.9450000,-0.5182681) (3.9500000,-0.5175824) (3.9550000,-0.5168985) (3.9600000,-0.5162165) (3.9650000,-0.5155363) (3.9700000,-0.5148579) (3.9750000,-0.5141813) (3.9800000,-0.5135066) (3.9850000,-0.5128337) (3.9900000,-0.5121625) (3.9950000,-0.5114932) (4.0000000,-0.5108256)};
\draw[violet!75!black,densely dotted] (-2,-1.0986123)--(2,1.0986123);
\fill[violet!75!black] (-2,-1.0986123) circle(2.3pt);
\fill[violet!75!black] (2,1.0986123) circle(2.3pt);
\end{scope}
\end{scope}
\node[anchor=west,align=left,text width=4.35cm] at (8.9,4.4) {固定 $k=1$。\\[8pt]蓝实线：$b=1$\\[8pt]橙虚线：$b=-1$\\[14pt]同一组 $k,b$ 下，\\本式与奇4互为相反数。};
\node[text width=12.8cm,align=center] at (7.1,.65) {$k,b\ne0$；$D=\{x:|x|>|b/k|\}$。取倒数后，对数反号。};
\end{tikzpicture}\end{center}

\textbf{先看图：} 仍取 $k,b\ne0$，定义域与奇4相同。对同一组 $k,b$，本式的真数恰好是奇4真数的倒数，因此两条对应曲线上下互为相反数。

\textbf{再说明一般情形：}
\[
\ln\frac{kx+b}{kx-b}=-\ln\frac{kx-b}{kx+b}.
\]
奇4为奇函数，乘以 $-1$ 后仍是奇函数。蓝实线与橙虚线各自均关于原点对称。
\end{explanationbox}
\clearpage
\begin{explanationbox}
\textbf{\textcolor{CExplanation}{根号加上 $x$，真数始终为正}}
% Formula-based samples; see F083 checks/build_tikz_diagrams.py.
\begin{center}\begin{tikzpicture}[font=\small,>=stealth,line join=round]
\draw[gray!40,rounded corners=3pt] (0,0) rectangle (14.2,8.4);
\node at (7.1,7.75) {奇6：$y=\log_a\!\bigl(x+\sqrt{x^2+1}\bigr)$};
\begin{scope}[shift={(4.400000,3.900000)},x=1.133333cm,y=0.703125cm]
\draw[->,gray] (-3,0)--(3,0) node[right] {$x$};
\draw[->,gray] (0,-3.2)--(0,3.2) node[above] {$y$};
\draw[gray] (-1,-.035/0.703125)--(-1,.035/0.703125);
\node[anchor=north] at (-1,-.07/0.703125) {$-1$};
\draw[gray] (1,-.035/0.703125)--(1,.035/0.703125);
\node[anchor=north] at (1,-.07/0.703125) {$1$};
\begin{scope}
\clip (-3,-3.2) rectangle (3,3.2);
\draw[blue,thick,solid] plot coordinates {(-3.0000000,-2.6234637) (-2.9900000,-2.6188946) (-2.9800000,-2.6143118) (-2.9700000,-2.6097151) (-2.9600000,-2.6051046) (-2.9500000,-2.6004800) (-2.9400000,-2.5958413) (-2.9300000,-2.5911885) (-2.9200000,-2.5865214) (-2.9100000,-2.5818400) (-2.9000000,-2.5771441) (-2.8900000,-2.5724338) (-2.8800000,-2.5677089) (-2.8700000,-2.5629694) (-2.8600000,-2.5582151) (-2.8500000,-2.5534459) (-2.8400000,-2.5486619) (-2.8300000,-2.5438628) (-2.8200000,-2.5390486) (-2.8100000,-2.5342193) (-2.8000000,-2.5293746) (-2.7900000,-2.5245146) (-2.7800000,-2.5196392) (-2.7700000,-2.5147481) (-2.7600000,-2.5098415) (-2.7500000,-2.5049191) (-2.7400000,-2.4999808) (-2.7300000,-2.4950266) (-2.7200000,-2.4900564) (-2.7100000,-2.4850701) (-2.7000000,-2.4800676) (-2.6900000,-2.4750487) (-2.6800000,-2.4700134) (-2.6700000,-2.4649616) (-2.6600000,-2.4598932) (-2.6500000,-2.4548081) (-2.6400000,-2.4497061) (-2.6300000,-2.4445872) (-2.6200000,-2.4394513) (-2.6100000,-2.4342982) (-2.6000000,-2.4291279) (-2.5900000,-2.4239402) (-2.5800000,-2.4187351) (-2.5700000,-2.4135124) (-2.5600000,-2.4082720) (-2.5500000,-2.4030138) (-2.5400000,-2.3977378) (-2.5300000,-2.3924437) (-2.5200000,-2.3871315) (-2.5100000,-2.3818010) (-2.5000000,-2.3764522) (-2.4900000,-2.3710849) (-2.4800000,-2.3656990) (-2.4700000,-2.3602944) (-2.4600000,-2.3548710) (-2.4500000,-2.3494286) (-2.4400000,-2.3439672) (-2.4300000,-2.3384865) (-2.4200000,-2.3329865) (-2.4100000,-2.3274671) (-2.4000000,-2.3219281) (-2.3900000,-2.3163694) (-2.3800000,-2.3107909) (-2.3700000,-2.3051924) (-2.3600000,-2.2995738) (-2.3500000,-2.2939350) (-2.3400000,-2.2882759) (-2.3300000,-2.2825962) (-2.3200000,-2.2768960) (-2.3100000,-2.2711750) (-2.3000000,-2.2654331) (-2.2900000,-2.2596701) (-2.2800000,-2.2538860) (-2.2700000,-2.2480806) (-2.2600000,-2.2422537) (-2.2500000,-2.2364052) (-2.2400000,-2.2305350) (-2.2300000,-2.2246428) (-2.2200000,-2.2187287) (-2.2100000,-2.2127923) (-2.2000000,-2.2068336) (-2.1900000,-2.2008524) (-2.1800000,-2.1948486) (-2.1700000,-2.1888220) (-2.1600000,-2.1827724) (-2.1500000,-2.1766997) (-2.1400000,-2.1706038) (-2.1300000,-2.1644844) (-2.1200000,-2.1583414) (-2.1100000,-2.1521747) (-2.1000000,-2.1459841) (-2.0900000,-2.1397694) (-2.0800000,-2.1335305) (-2.0700000,-2.1272671) (-2.0600000,-2.1209792) (-2.0500000,-2.1146665) (-2.0400000,-2.1083289) (-2.0300000,-2.1019662) (-2.0200000,-2.0955782) (-2.0100000,-2.0891648) (-2.0000000,-2.0827257) (-1.9900000,-2.0762609) (-1.9800000,-2.0697700) (-1.9700000,-2.0632530) (-1.9600000,-2.0567096) (-1.9500000,-2.0501397) (-1.9400000,-2.0435430) (-1.9300000,-2.0369195) (-1.9200000,-2.0302688) (-1.9100000,-2.0235908) (-1.9000000,-2.0168853) (-1.8900000,-2.0101521) (-1.8800000,-2.0033911) (-1.8700000,-1.9966019) (-1.8600000,-1.9897844) (-1.8500000,-1.9829385) (-1.8400000,-1.9760639) (-1.8300000,-1.9691603) (-1.8200000,-1.9622276) (-1.8100000,-1.9552657) (-1.8000000,-1.9482742) (-1.7900000,-1.9412529) (-1.7800000,-1.9342017) (-1.7700000,-1.9271203) (-1.7600000,-1.9200085) (-1.7500000,-1.9128661) (-1.7400000,-1.9056929) (-1.7300000,-1.8984886) (-1.7200000,-1.8912531) (-1.7100000,-1.8839860) (-1.7000000,-1.8766872) (-1.6900000,-1.8693564) (-1.6800000,-1.8619934) (-1.6700000,-1.8545980) (-1.6600000,-1.8471699) (-1.6500000,-1.8397090) (-1.6400000,-1.8322148) (-1.6300000,-1.8246873) (-1.6200000,-1.8171262) (-1.6100000,-1.8095312) (-1.6000000,-1.8019021) (-1.5900000,-1.7942386) (-1.5800000,-1.7865405) (-1.5700000,-1.7788075) (-1.5600000,-1.7710394) (-1.5500000,-1.7632359) (-1.5400000,-1.7553968) (-1.5300000,-1.7475218) (-1.5200000,-1.7396107) (-1.5100000,-1.7316632) (-1.5000000,-1.7236790) (-1.4900000,-1.7156578) (-1.4800000,-1.7075995) (-1.4700000,-1.6995037) (-1.4600000,-1.6913701) (-1.4500000,-1.6831986) (-1.4400000,-1.6749887) (-1.4300000,-1.6667403) (-1.4200000,-1.6584531) (-1.4100000,-1.6501268) (-1.4000000,-1.6417611) (-1.3900000,-1.6333557) (-1.3800000,-1.6249104) (-1.3700000,-1.6164248) (-1.3600000,-1.6078988) (-1.3500000,-1.5993320) (-1.3400000,-1.5907241) (-1.3300000,-1.5820748) (-1.3200000,-1.5733839) (-1.3100000,-1.5646511) (-1.3000000,-1.5558760) (-1.2900000,-1.5470584) (-1.2800000,-1.5381981) (-1.2700000,-1.5292946) (-1.2600000,-1.5203478) (-1.2500000,-1.5113572) (-1.2400000,-1.5023228) (-1.2300000,-1.4932440) (-1.2200000,-1.4841208) (-1.2100000,-1.4749526) (-1.2000000,-1.4657394) (-1.1900000,-1.4564807) (-1.1800000,-1.4471764) (-1.1700000,-1.4378260) (-1.1600000,-1.4284293) (-1.1500000,-1.4189860) (-1.1400000,-1.4094959) (-1.1300000,-1.3999586) (-1.1200000,-1.3903738) (-1.1100000,-1.3807413) (-1.1000000,-1.3710608) (-1.0900000,-1.3613320) (-1.0800000,-1.3515546) (-1.0700000,-1.3417283) (-1.0600000,-1.3318529) (-1.0500000,-1.3219281) (-1.0400000,-1.3119536) (-1.0300000,-1.3019291) (-1.0200000,-1.2918544) (-1.0100000,-1.2817292) (-1.0000000,-1.2715533) (-0.9900000,-1.2613264) (-0.9800000,-1.2510482) (-0.9700000,-1.2407185) (-0.9600000,-1.2303370) (-0.9500000,-1.2199035) (-0.9400000,-1.2094179) (-0.9300000,-1.1988797) (-0.9200000,-1.1882889) (-0.9100000,-1.1776452) (-0.9000000,-1.1669483) (-0.8900000,-1.1561982) (-0.8800000,-1.1453945) (-0.8700000,-1.1345371) (-0.8600000,-1.1236258) (-0.8500000,-1.1126604) (-0.8400000,-1.1016408) (-0.8300000,-1.0905668) (-0.8200000,-1.0794382) (-0.8100000,-1.0682550) (-0.8000000,-1.0570169) (-0.7900000,-1.0457238) (-0.7800000,-1.0343757) (-0.7700000,-1.0229724) (-0.7600000,-1.0115139) (-0.7500000,-1.0000000) (-0.7400000,-0.9884307) (-0.7300000,-0.9768060) (-0.7200000,-0.9651258) (-0.7100000,-0.9533901) (-0.7000000,-0.9415988) (-0.6900000,-0.9297520) (-0.6800000,-0.9178498) (-0.6700000,-0.9058920) (-0.6600000,-0.8938788) (-0.6500000,-0.8818103) (-0.6400000,-0.8696865) (-0.6300000,-0.8575075) (-0.6200000,-0.8452735) (-0.6100000,-0.8329845) (-0.6000000,-0.8206409) (-0.5900000,-0.8082426) (-0.5800000,-0.7957900) (-0.5700000,-0.7832831) (-0.5600000,-0.7707224) (-0.5500000,-0.7581080) (-0.5400000,-0.7454403) (-0.5300000,-0.7327194) (-0.5200000,-0.7199459) (-0.5100000,-0.7071199) (-0.5000000,-0.6942419) (-0.4900000,-0.6813123) (-0.4800000,-0.6683315) (-0.4700000,-0.6553000) (-0.4600000,-0.6422182) (-0.4500000,-0.6290867) (-0.4400000,-0.6159059) (-0.4300000,-0.6026765) (-0.4200000,-0.5893989) (-0.4100000,-0.5760739) (-0.4000000,-0.5627020) (-0.3900000,-0.5492840) (-0.3800000,-0.5358204) (-0.3700000,-0.5223120) (-0.3600000,-0.5087597) (-0.3500000,-0.4951640) (-0.3400000,-0.4815259) (-0.3300000,-0.4678462) (-0.3200000,-0.4541257) (-0.3100000,-0.4403654) (-0.3000000,-0.4265660) (-0.2900000,-0.4127287) (-0.2800000,-0.3988542) (-0.2700000,-0.3849437) (-0.2600000,-0.3709982) (-0.2500000,-0.3570186) (-0.2400000,-0.3430062) (-0.2300000,-0.3289618) (-0.2200000,-0.3148868) (-0.2100000,-0.3007822) (-0.2000000,-0.2866492) (-0.1900000,-0.2724890) (-0.1800000,-0.2583029) (-0.1700000,-0.2440919) (-0.1600000,-0.2298575) (-0.1500000,-0.2156008) (-0.1400000,-0.2013233) (-0.1300000,-0.1870261) (-0.1200000,-0.1727106) (-0.1100000,-0.1583781) (-0.1000000,-0.1440301) (-0.0900000,-0.1296679) (-0.0800000,-0.1152928) (-0.0700000,-0.1009064) (-0.0600000,-0.0865098) (-0.0500000,-0.0721047) (-0.0400000,-0.0576924) (-0.0300000,-0.0432744) (-0.0200000,-0.0288520) (-0.0100000,-0.0144267) (0.0000000,0.0000000) (0.0100000,0.0144267) (0.0200000,0.0288520) (0.0300000,0.0432744) (0.0400000,0.0576924) (0.0500000,0.0721047) (0.0600000,0.0865098) (0.0700000,0.1009064) (0.0800000,0.1152928) (0.0900000,0.1296679) (0.1000000,0.1440301) (0.1100000,0.1583781) (0.1200000,0.1727106) (0.1300000,0.1870261) (0.1400000,0.2013233) (0.1500000,0.2156008) (0.1600000,0.2298575) (0.1700000,0.2440919) (0.1800000,0.2583029) (0.1900000,0.2724890) (0.2000000,0.2866492) (0.2100000,0.3007822) (0.2200000,0.3148868) (0.2300000,0.3289618) (0.2400000,0.3430062) (0.2500000,0.3570186) (0.2600000,0.3709982) (0.2700000,0.3849437) (0.2800000,0.3988542) (0.2900000,0.4127287) (0.3000000,0.4265660) (0.3100000,0.4403654) (0.3200000,0.4541257) (0.3300000,0.4678462) (0.3400000,0.4815259) (0.3500000,0.4951640) (0.3600000,0.5087597) (0.3700000,0.5223120) (0.3800000,0.5358204) (0.3900000,0.5492840) (0.4000000,0.5627020) (0.4100000,0.5760739) (0.4200000,0.5893989) (0.4300000,0.6026765) (0.4400000,0.6159059) (0.4500000,0.6290867) (0.4600000,0.6422182) (0.4700000,0.6553000) (0.4800000,0.6683315) (0.4900000,0.6813123) (0.5000000,0.6942419) (0.5100000,0.7071199) (0.5200000,0.7199459) (0.5300000,0.7327194) (0.5400000,0.7454403) (0.5500000,0.7581080) (0.5600000,0.7707224) (0.5700000,0.7832831) (0.5800000,0.7957900) (0.5900000,0.8082426) (0.6000000,0.8206409) (0.6100000,0.8329845) (0.6200000,0.8452735) (0.6300000,0.8575075) (0.6400000,0.8696865) (0.6500000,0.8818103) (0.6600000,0.8938788) (0.6700000,0.9058920) (0.6800000,0.9178498) (0.6900000,0.9297520) (0.7000000,0.9415988) (0.7100000,0.9533901) (0.7200000,0.9651258) (0.7300000,0.9768060) (0.7400000,0.9884307) (0.7500000,1.0000000) (0.7600000,1.0115139) (0.7700000,1.0229724) (0.7800000,1.0343757) (0.7900000,1.0457238) (0.8000000,1.0570169) (0.8100000,1.0682550) (0.8200000,1.0794382) (0.8300000,1.0905668) (0.8400000,1.1016408) (0.8500000,1.1126604) (0.8600000,1.1236258) (0.8700000,1.1345371) (0.8800000,1.1453945) (0.8900000,1.1561982) (0.9000000,1.1669483) (0.9100000,1.1776452) (0.9200000,1.1882889) (0.9300000,1.1988797) (0.9400000,1.2094179) (0.9500000,1.2199035) (0.9600000,1.2303370) (0.9700000,1.2407185) (0.9800000,1.2510482) (0.9900000,1.2613264) (1.0000000,1.2715533) (1.0100000,1.2817292) (1.0200000,1.2918544) (1.0300000,1.3019291) (1.0400000,1.3119536) (1.0500000,1.3219281) (1.0600000,1.3318529) (1.0700000,1.3417283) (1.0800000,1.3515546) (1.0900000,1.3613320) (1.1000000,1.3710608) (1.1100000,1.3807413) (1.1200000,1.3903738) (1.1300000,1.3999586) (1.1400000,1.4094959) (1.1500000,1.4189860) (1.1600000,1.4284293) (1.1700000,1.4378260) (1.1800000,1.4471764) (1.1900000,1.4564807) (1.2000000,1.4657394) (1.2100000,1.4749526) (1.2200000,1.4841208) (1.2300000,1.4932440) (1.2400000,1.5023228) (1.2500000,1.5113572) (1.2600000,1.5203478) (1.2700000,1.5292946) (1.2800000,1.5381981) (1.2900000,1.5470584) (1.3000000,1.5558760) (1.3100000,1.5646511) (1.3200000,1.5733839) (1.3300000,1.5820748) (1.3400000,1.5907241) (1.3500000,1.5993320) (1.3600000,1.6078988) (1.3700000,1.6164248) (1.3800000,1.6249104) (1.3900000,1.6333557) (1.4000000,1.6417611) (1.4100000,1.6501268) (1.4200000,1.6584531) (1.4300000,1.6667403) (1.4400000,1.6749887) (1.4500000,1.6831986) (1.4600000,1.6913701) (1.4700000,1.6995037) (1.4800000,1.7075995) (1.4900000,1.7156578) (1.5000000,1.7236790) (1.5100000,1.7316632) (1.5200000,1.7396107) (1.5300000,1.7475218) (1.5400000,1.7553968) (1.5500000,1.7632359) (1.5600000,1.7710394) (1.5700000,1.7788075) (1.5800000,1.7865405) (1.5900000,1.7942386) (1.6000000,1.8019021) (1.6100000,1.8095312) (1.6200000,1.8171262) (1.6300000,1.8246873) (1.6400000,1.8322148) (1.6500000,1.8397090) (1.6600000,1.8471699) (1.6700000,1.8545980) (1.6800000,1.8619934) (1.6900000,1.8693564) (1.7000000,1.8766872) (1.7100000,1.8839860) (1.7200000,1.8912531) (1.7300000,1.8984886) (1.7400000,1.9056929) (1.7500000,1.9128661) (1.7600000,1.9200085) (1.7700000,1.9271203) (1.7800000,1.9342017) (1.7900000,1.9412529) (1.8000000,1.9482742) (1.8100000,1.9552657) (1.8200000,1.9622276) (1.8300000,1.9691603) (1.8400000,1.9760639) (1.8500000,1.9829385) (1.8600000,1.9897844) (1.8700000,1.9966019) (1.8800000,2.0033911) (1.8900000,2.0101521) (1.9000000,2.0168853) (1.9100000,2.0235908) (1.9200000,2.0302688) (1.9300000,2.0369195) (1.9400000,2.0435430) (1.9500000,2.0501397) (1.9600000,2.0567096) (1.9700000,2.0632530) (1.9800000,2.0697700) (1.9900000,2.0762609) (2.0000000,2.0827257) (2.0100000,2.0891648) (2.0200000,2.0955782) (2.0300000,2.1019662) (2.0400000,2.1083289) (2.0500000,2.1146665) (2.0600000,2.1209792) (2.0700000,2.1272671) (2.0800000,2.1335305) (2.0900000,2.1397694) (2.1000000,2.1459841) (2.1100000,2.1521747) (2.1200000,2.1583414) (2.1300000,2.1644844) (2.1400000,2.1706038) (2.1500000,2.1766997) (2.1600000,2.1827724) (2.1700000,2.1888220) (2.1800000,2.1948486) (2.1900000,2.2008524) (2.2000000,2.2068336) (2.2100000,2.2127923) (2.2200000,2.2187287) (2.2300000,2.2246428) (2.2400000,2.2305350) (2.2500000,2.2364052) (2.2600000,2.2422537) (2.2700000,2.2480806) (2.2800000,2.2538860) (2.2900000,2.2596701) (2.3000000,2.2654331) (2.3100000,2.2711750) (2.3200000,2.2768960) (2.3300000,2.2825962) (2.3400000,2.2882759) (2.3500000,2.2939350) (2.3600000,2.2995738) (2.3700000,2.3051924) (2.3800000,2.3107909) (2.3900000,2.3163694) (2.4000000,2.3219281) (2.4100000,2.3274671) (2.4200000,2.3329865) (2.4300000,2.3384865) (2.4400000,2.3439672) (2.4500000,2.3494286) (2.4600000,2.3548710) (2.4700000,2.3602944) (2.4800000,2.3656990) (2.4900000,2.3710849) (2.5000000,2.3764522) (2.5100000,2.3818010) (2.5200000,2.3871315) (2.5300000,2.3924437) (2.5400000,2.3977378) (2.5500000,2.4030138) (2.5600000,2.4082720) (2.5700000,2.4135124) (2.5800000,2.4187351) (2.5900000,2.4239402) (2.6000000,2.4291279) (2.6100000,2.4342982) (2.6200000,2.4394513) (2.6300000,2.4445872) (2.6400000,2.4497061) (2.6500000,2.4548081) (2.6600000,2.4598932) (2.6700000,2.4649616) (2.6800000,2.4700134) (2.6900000,2.4750487) (2.7000000,2.4800676) (2.7100000,2.4850701) (2.7200000,2.4900564) (2.7300000,2.4950266) (2.7400000,2.4999808) (2.7500000,2.5049191) (2.7600000,2.5098415) (2.7700000,2.5147481) (2.7800000,2.5196392) (2.7900000,2.5245146) (2.8000000,2.5293746) (2.8100000,2.5342193) (2.8200000,2.5390486) (2.8300000,2.5438628) (2.8400000,2.5486619) (2.8500000,2.5534459) (2.8600000,2.5582151) (2.8700000,2.5629694) (2.8800000,2.5677089) (2.8900000,2.5724338) (2.9000000,2.5771441) (2.9100000,2.5818400) (2.9200000,2.5865214) (2.9300000,2.5911885) (2.9400000,2.5958413) (2.9500000,2.6004800) (2.9600000,2.6051046) (2.9700000,2.6097151) (2.9800000,2.6143118) (2.9900000,2.6188946) (3.0000000,2.6234637)};
\draw[orange!85!black,thick,dashed] plot coordinates {(-3.0000000,2.6234637) (-2.9900000,2.6188946) (-2.9800000,2.6143118) (-2.9700000,2.6097151) (-2.9600000,2.6051046) (-2.9500000,2.6004800) (-2.9400000,2.5958413) (-2.9300000,2.5911885) (-2.9200000,2.5865214) (-2.9100000,2.5818400) (-2.9000000,2.5771441) (-2.8900000,2.5724338) (-2.8800000,2.5677089) (-2.8700000,2.5629694) (-2.8600000,2.5582151) (-2.8500000,2.5534459) (-2.8400000,2.5486619) (-2.8300000,2.5438628) (-2.8200000,2.5390486) (-2.8100000,2.5342193) (-2.8000000,2.5293746) (-2.7900000,2.5245146) (-2.7800000,2.5196392) (-2.7700000,2.5147481) (-2.7600000,2.5098415) (-2.7500000,2.5049191) (-2.7400000,2.4999808) (-2.7300000,2.4950266) (-2.7200000,2.4900564) (-2.7100000,2.4850701) (-2.7000000,2.4800676) (-2.6900000,2.4750487) (-2.6800000,2.4700134) (-2.6700000,2.4649616) (-2.6600000,2.4598932) (-2.6500000,2.4548081) (-2.6400000,2.4497061) (-2.6300000,2.4445872) (-2.6200000,2.4394513) (-2.6100000,2.4342982) (-2.6000000,2.4291279) (-2.5900000,2.4239402) (-2.5800000,2.4187351) (-2.5700000,2.4135124) (-2.5600000,2.4082720) (-2.5500000,2.4030138) (-2.5400000,2.3977378) (-2.5300000,2.3924437) (-2.5200000,2.3871315) (-2.5100000,2.3818010) (-2.5000000,2.3764522) (-2.4900000,2.3710849) (-2.4800000,2.3656990) (-2.4700000,2.3602944) (-2.4600000,2.3548710) (-2.4500000,2.3494286) (-2.4400000,2.3439672) (-2.4300000,2.3384865) (-2.4200000,2.3329865) (-2.4100000,2.3274671) (-2.4000000,2.3219281) (-2.3900000,2.3163694) (-2.3800000,2.3107909) (-2.3700000,2.3051924) (-2.3600000,2.2995738) (-2.3500000,2.2939350) (-2.3400000,2.2882759) (-2.3300000,2.2825962) (-2.3200000,2.2768960) (-2.3100000,2.2711750) (-2.3000000,2.2654331) (-2.2900000,2.2596701) (-2.2800000,2.2538860) (-2.2700000,2.2480806) (-2.2600000,2.2422537) (-2.2500000,2.2364052) (-2.2400000,2.2305350) (-2.2300000,2.2246428) (-2.2200000,2.2187287) (-2.2100000,2.2127923) (-2.2000000,2.2068336) (-2.1900000,2.2008524) (-2.1800000,2.1948486) (-2.1700000,2.1888220) (-2.1600000,2.1827724) (-2.1500000,2.1766997) (-2.1400000,2.1706038) (-2.1300000,2.1644844) (-2.1200000,2.1583414) (-2.1100000,2.1521747) (-2.1000000,2.1459841) (-2.0900000,2.1397694) (-2.0800000,2.1335305) (-2.0700000,2.1272671) (-2.0600000,2.1209792) (-2.0500000,2.1146665) (-2.0400000,2.1083289) (-2.0300000,2.1019662) (-2.0200000,2.0955782) (-2.0100000,2.0891648) (-2.0000000,2.0827257) (-1.9900000,2.0762609) (-1.9800000,2.0697700) (-1.9700000,2.0632530) (-1.9600000,2.0567096) (-1.9500000,2.0501397) (-1.9400000,2.0435430) (-1.9300000,2.0369195) (-1.9200000,2.0302688) (-1.9100000,2.0235908) (-1.9000000,2.0168853) (-1.8900000,2.0101521) (-1.8800000,2.0033911) (-1.8700000,1.9966019) (-1.8600000,1.9897844) (-1.8500000,1.9829385) (-1.8400000,1.9760639) (-1.8300000,1.9691603) (-1.8200000,1.9622276) (-1.8100000,1.9552657) (-1.8000000,1.9482742) (-1.7900000,1.9412529) (-1.7800000,1.9342017) (-1.7700000,1.9271203) (-1.7600000,1.9200085) (-1.7500000,1.9128661) (-1.7400000,1.9056929) (-1.7300000,1.8984886) (-1.7200000,1.8912531) (-1.7100000,1.8839860) (-1.7000000,1.8766872) (-1.6900000,1.8693564) (-1.6800000,1.8619934) (-1.6700000,1.8545980) (-1.6600000,1.8471699) (-1.6500000,1.8397090) (-1.6400000,1.8322148) (-1.6300000,1.8246873) (-1.6200000,1.8171262) (-1.6100000,1.8095312) (-1.6000000,1.8019021) (-1.5900000,1.7942386) (-1.5800000,1.7865405) (-1.5700000,1.7788075) (-1.5600000,1.7710394) (-1.5500000,1.7632359) (-1.5400000,1.7553968) (-1.5300000,1.7475218) (-1.5200000,1.7396107) (-1.5100000,1.7316632) (-1.5000000,1.7236790) (-1.4900000,1.7156578) (-1.4800000,1.7075995) (-1.4700000,1.6995037) (-1.4600000,1.6913701) (-1.4500000,1.6831986) (-1.4400000,1.6749887) (-1.4300000,1.6667403) (-1.4200000,1.6584531) (-1.4100000,1.6501268) (-1.4000000,1.6417611) (-1.3900000,1.6333557) (-1.3800000,1.6249104) (-1.3700000,1.6164248) (-1.3600000,1.6078988) (-1.3500000,1.5993320) (-1.3400000,1.5907241) (-1.3300000,1.5820748) (-1.3200000,1.5733839) (-1.3100000,1.5646511) (-1.3000000,1.5558760) (-1.2900000,1.5470584) (-1.2800000,1.5381981) (-1.2700000,1.5292946) (-1.2600000,1.5203478) (-1.2500000,1.5113572) (-1.2400000,1.5023228) (-1.2300000,1.4932440) (-1.2200000,1.4841208) (-1.2100000,1.4749526) (-1.2000000,1.4657394) (-1.1900000,1.4564807) (-1.1800000,1.4471764) (-1.1700000,1.4378260) (-1.1600000,1.4284293) (-1.1500000,1.4189860) (-1.1400000,1.4094959) (-1.1300000,1.3999586) (-1.1200000,1.3903738) (-1.1100000,1.3807413) (-1.1000000,1.3710608) (-1.0900000,1.3613320) (-1.0800000,1.3515546) (-1.0700000,1.3417283) (-1.0600000,1.3318529) (-1.0500000,1.3219281) (-1.0400000,1.3119536) (-1.0300000,1.3019291) (-1.0200000,1.2918544) (-1.0100000,1.2817292) (-1.0000000,1.2715533) (-0.9900000,1.2613264) (-0.9800000,1.2510482) (-0.9700000,1.2407185) (-0.9600000,1.2303370) (-0.9500000,1.2199035) (-0.9400000,1.2094179) (-0.9300000,1.1988797) (-0.9200000,1.1882889) (-0.9100000,1.1776452) (-0.9000000,1.1669483) (-0.8900000,1.1561982) (-0.8800000,1.1453945) (-0.8700000,1.1345371) (-0.8600000,1.1236258) (-0.8500000,1.1126604) (-0.8400000,1.1016408) (-0.8300000,1.0905668) (-0.8200000,1.0794382) (-0.8100000,1.0682550) (-0.8000000,1.0570169) (-0.7900000,1.0457238) (-0.7800000,1.0343757) (-0.7700000,1.0229724) (-0.7600000,1.0115139) (-0.7500000,1.0000000) (-0.7400000,0.9884307) (-0.7300000,0.9768060) (-0.7200000,0.9651258) (-0.7100000,0.9533901) (-0.7000000,0.9415988) (-0.6900000,0.9297520) (-0.6800000,0.9178498) (-0.6700000,0.9058920) (-0.6600000,0.8938788) (-0.6500000,0.8818103) (-0.6400000,0.8696865) (-0.6300000,0.8575075) (-0.6200000,0.8452735) (-0.6100000,0.8329845) (-0.6000000,0.8206409) (-0.5900000,0.8082426) (-0.5800000,0.7957900) (-0.5700000,0.7832831) (-0.5600000,0.7707224) (-0.5500000,0.7581080) (-0.5400000,0.7454403) (-0.5300000,0.7327194) (-0.5200000,0.7199459) (-0.5100000,0.7071199) (-0.5000000,0.6942419) (-0.4900000,0.6813123) (-0.4800000,0.6683315) (-0.4700000,0.6553000) (-0.4600000,0.6422182) (-0.4500000,0.6290867) (-0.4400000,0.6159059) (-0.4300000,0.6026765) (-0.4200000,0.5893989) (-0.4100000,0.5760739) (-0.4000000,0.5627020) (-0.3900000,0.5492840) (-0.3800000,0.5358204) (-0.3700000,0.5223120) (-0.3600000,0.5087597) (-0.3500000,0.4951640) (-0.3400000,0.4815259) (-0.3300000,0.4678462) (-0.3200000,0.4541257) (-0.3100000,0.4403654) (-0.3000000,0.4265660) (-0.2900000,0.4127287) (-0.2800000,0.3988542) (-0.2700000,0.3849437) (-0.2600000,0.3709982) (-0.2500000,0.3570186) (-0.2400000,0.3430062) (-0.2300000,0.3289618) (-0.2200000,0.3148868) (-0.2100000,0.3007822) (-0.2000000,0.2866492) (-0.1900000,0.2724890) (-0.1800000,0.2583029) (-0.1700000,0.2440919) (-0.1600000,0.2298575) (-0.1500000,0.2156008) (-0.1400000,0.2013233) (-0.1300000,0.1870261) (-0.1200000,0.1727106) (-0.1100000,0.1583781) (-0.1000000,0.1440301) (-0.0900000,0.1296679) (-0.0800000,0.1152928) (-0.0700000,0.1009064) (-0.0600000,0.0865098) (-0.0500000,0.0721047) (-0.0400000,0.0576924) (-0.0300000,0.0432744) (-0.0200000,0.0288520) (-0.0100000,0.0144267) (0.0000000,-0.0000000) (0.0100000,-0.0144267) (0.0200000,-0.0288520) (0.0300000,-0.0432744) (0.0400000,-0.0576924) (0.0500000,-0.0721047) (0.0600000,-0.0865098) (0.0700000,-0.1009064) (0.0800000,-0.1152928) (0.0900000,-0.1296679) (0.1000000,-0.1440301) (0.1100000,-0.1583781) (0.1200000,-0.1727106) (0.1300000,-0.1870261) (0.1400000,-0.2013233) (0.1500000,-0.2156008) (0.1600000,-0.2298575) (0.1700000,-0.2440919) (0.1800000,-0.2583029) (0.1900000,-0.2724890) (0.2000000,-0.2866492) (0.2100000,-0.3007822) (0.2200000,-0.3148868) (0.2300000,-0.3289618) (0.2400000,-0.3430062) (0.2500000,-0.3570186) (0.2600000,-0.3709982) (0.2700000,-0.3849437) (0.2800000,-0.3988542) (0.2900000,-0.4127287) (0.3000000,-0.4265660) (0.3100000,-0.4403654) (0.3200000,-0.4541257) (0.3300000,-0.4678462) (0.3400000,-0.4815259) (0.3500000,-0.4951640) (0.3600000,-0.5087597) (0.3700000,-0.5223120) (0.3800000,-0.5358204) (0.3900000,-0.5492840) (0.4000000,-0.5627020) (0.4100000,-0.5760739) (0.4200000,-0.5893989) (0.4300000,-0.6026765) (0.4400000,-0.6159059) (0.4500000,-0.6290867) (0.4600000,-0.6422182) (0.4700000,-0.6553000) (0.4800000,-0.6683315) (0.4900000,-0.6813123) (0.5000000,-0.6942419) (0.5100000,-0.7071199) (0.5200000,-0.7199459) (0.5300000,-0.7327194) (0.5400000,-0.7454403) (0.5500000,-0.7581080) (0.5600000,-0.7707224) (0.5700000,-0.7832831) (0.5800000,-0.7957900) (0.5900000,-0.8082426) (0.6000000,-0.8206409) (0.6100000,-0.8329845) (0.6200000,-0.8452735) (0.6300000,-0.8575075) (0.6400000,-0.8696865) (0.6500000,-0.8818103) (0.6600000,-0.8938788) (0.6700000,-0.9058920) (0.6800000,-0.9178498) (0.6900000,-0.9297520) (0.7000000,-0.9415988) (0.7100000,-0.9533901) (0.7200000,-0.9651258) (0.7300000,-0.9768060) (0.7400000,-0.9884307) (0.7500000,-1.0000000) (0.7600000,-1.0115139) (0.7700000,-1.0229724) (0.7800000,-1.0343757) (0.7900000,-1.0457238) (0.8000000,-1.0570169) (0.8100000,-1.0682550) (0.8200000,-1.0794382) (0.8300000,-1.0905668) (0.8400000,-1.1016408) (0.8500000,-1.1126604) (0.8600000,-1.1236258) (0.8700000,-1.1345371) (0.8800000,-1.1453945) (0.8900000,-1.1561982) (0.9000000,-1.1669483) (0.9100000,-1.1776452) (0.9200000,-1.1882889) (0.9300000,-1.1988797) (0.9400000,-1.2094179) (0.9500000,-1.2199035) (0.9600000,-1.2303370) (0.9700000,-1.2407185) (0.9800000,-1.2510482) (0.9900000,-1.2613264) (1.0000000,-1.2715533) (1.0100000,-1.2817292) (1.0200000,-1.2918544) (1.0300000,-1.3019291) (1.0400000,-1.3119536) (1.0500000,-1.3219281) (1.0600000,-1.3318529) (1.0700000,-1.3417283) (1.0800000,-1.3515546) (1.0900000,-1.3613320) (1.1000000,-1.3710608) (1.1100000,-1.3807413) (1.1200000,-1.3903738) (1.1300000,-1.3999586) (1.1400000,-1.4094959) (1.1500000,-1.4189860) (1.1600000,-1.4284293) (1.1700000,-1.4378260) (1.1800000,-1.4471764) (1.1900000,-1.4564807) (1.2000000,-1.4657394) (1.2100000,-1.4749526) (1.2200000,-1.4841208) (1.2300000,-1.4932440) (1.2400000,-1.5023228) (1.2500000,-1.5113572) (1.2600000,-1.5203478) (1.2700000,-1.5292946) (1.2800000,-1.5381981) (1.2900000,-1.5470584) (1.3000000,-1.5558760) (1.3100000,-1.5646511) (1.3200000,-1.5733839) (1.3300000,-1.5820748) (1.3400000,-1.5907241) (1.3500000,-1.5993320) (1.3600000,-1.6078988) (1.3700000,-1.6164248) (1.3800000,-1.6249104) (1.3900000,-1.6333557) (1.4000000,-1.6417611) (1.4100000,-1.6501268) (1.4200000,-1.6584531) (1.4300000,-1.6667403) (1.4400000,-1.6749887) (1.4500000,-1.6831986) (1.4600000,-1.6913701) (1.4700000,-1.6995037) (1.4800000,-1.7075995) (1.4900000,-1.7156578) (1.5000000,-1.7236790) (1.5100000,-1.7316632) (1.5200000,-1.7396107) (1.5300000,-1.7475218) (1.5400000,-1.7553968) (1.5500000,-1.7632359) (1.5600000,-1.7710394) (1.5700000,-1.7788075) (1.5800000,-1.7865405) (1.5900000,-1.7942386) (1.6000000,-1.8019021) (1.6100000,-1.8095312) (1.6200000,-1.8171262) (1.6300000,-1.8246873) (1.6400000,-1.8322148) (1.6500000,-1.8397090) (1.6600000,-1.8471699) (1.6700000,-1.8545980) (1.6800000,-1.8619934) (1.6900000,-1.8693564) (1.7000000,-1.8766872) (1.7100000,-1.8839860) (1.7200000,-1.8912531) (1.7300000,-1.8984886) (1.7400000,-1.9056929) (1.7500000,-1.9128661) (1.7600000,-1.9200085) (1.7700000,-1.9271203) (1.7800000,-1.9342017) (1.7900000,-1.9412529) (1.8000000,-1.9482742) (1.8100000,-1.9552657) (1.8200000,-1.9622276) (1.8300000,-1.9691603) (1.8400000,-1.9760639) (1.8500000,-1.9829385) (1.8600000,-1.9897844) (1.8700000,-1.9966019) (1.8800000,-2.0033911) (1.8900000,-2.0101521) (1.9000000,-2.0168853) (1.9100000,-2.0235908) (1.9200000,-2.0302688) (1.9300000,-2.0369195) (1.9400000,-2.0435430) (1.9500000,-2.0501397) (1.9600000,-2.0567096) (1.9700000,-2.0632530) (1.9800000,-2.0697700) (1.9900000,-2.0762609) (2.0000000,-2.0827257) (2.0100000,-2.0891648) (2.0200000,-2.0955782) (2.0300000,-2.1019662) (2.0400000,-2.1083289) (2.0500000,-2.1146665) (2.0600000,-2.1209792) (2.0700000,-2.1272671) (2.0800000,-2.1335305) (2.0900000,-2.1397694) (2.1000000,-2.1459841) (2.1100000,-2.1521747) (2.1200000,-2.1583414) (2.1300000,-2.1644844) (2.1400000,-2.1706038) (2.1500000,-2.1766997) (2.1600000,-2.1827724) (2.1700000,-2.1888220) (2.1800000,-2.1948486) (2.1900000,-2.2008524) (2.2000000,-2.2068336) (2.2100000,-2.2127923) (2.2200000,-2.2187287) (2.2300000,-2.2246428) (2.2400000,-2.2305350) (2.2500000,-2.2364052) (2.2600000,-2.2422537) (2.2700000,-2.2480806) (2.2800000,-2.2538860) (2.2900000,-2.2596701) (2.3000000,-2.2654331) (2.3100000,-2.2711750) (2.3200000,-2.2768960) (2.3300000,-2.2825962) (2.3400000,-2.2882759) (2.3500000,-2.2939350) (2.3600000,-2.2995738) (2.3700000,-2.3051924) (2.3800000,-2.3107909) (2.3900000,-2.3163694) (2.4000000,-2.3219281) (2.4100000,-2.3274671) (2.4200000,-2.3329865) (2.4300000,-2.3384865) (2.4400000,-2.3439672) (2.4500000,-2.3494286) (2.4600000,-2.3548710) (2.4700000,-2.3602944) (2.4800000,-2.3656990) (2.4900000,-2.3710849) (2.5000000,-2.3764522) (2.5100000,-2.3818010) (2.5200000,-2.3871315) (2.5300000,-2.3924437) (2.5400000,-2.3977378) (2.5500000,-2.4030138) (2.5600000,-2.4082720) (2.5700000,-2.4135124) (2.5800000,-2.4187351) (2.5900000,-2.4239402) (2.6000000,-2.4291279) (2.6100000,-2.4342982) (2.6200000,-2.4394513) (2.6300000,-2.4445872) (2.6400000,-2.4497061) (2.6500000,-2.4548081) (2.6600000,-2.4598932) (2.6700000,-2.4649616) (2.6800000,-2.4700134) (2.6900000,-2.4750487) (2.7000000,-2.4800676) (2.7100000,-2.4850701) (2.7200000,-2.4900564) (2.7300000,-2.4950266) (2.7400000,-2.4999808) (2.7500000,-2.5049191) (2.7600000,-2.5098415) (2.7700000,-2.5147481) (2.7800000,-2.5196392) (2.7900000,-2.5245146) (2.8000000,-2.5293746) (2.8100000,-2.5342193) (2.8200000,-2.5390486) (2.8300000,-2.5438628) (2.8400000,-2.5486619) (2.8500000,-2.5534459) (2.8600000,-2.5582151) (2.8700000,-2.5629694) (2.8800000,-2.5677089) (2.8900000,-2.5724338) (2.9000000,-2.5771441) (2.9100000,-2.5818400) (2.9200000,-2.5865214) (2.9300000,-2.5911885) (2.9400000,-2.5958413) (2.9500000,-2.6004800) (2.9600000,-2.6051046) (2.9700000,-2.6097151) (2.9800000,-2.6143118) (2.9900000,-2.6188946) (3.0000000,-2.6234637)};
\draw[violet!75!black,densely dotted] (-1,-1.2715533)--(1,1.2715533);
\fill[violet!75!black] (-1,-1.2715533) circle(2.3pt);
\fill[violet!75!black] (1,1.2715533) circle(2.3pt);
\end{scope}
\end{scope}
\node[anchor=west,align=left,text width=4.35cm] at (8.9,4.4) {蓝实线：$a=2$\\[8pt]橙虚线：$a=\frac12$\\[14pt]每一条曲线都\\关于原点对称。\\[8pt]$x=0$ 时真数为 $1$。};
\node[text width=12.8cm,align=center] at (7.1,.65) {$D=\mathbb R$；$\sqrt{x^2+1}>|x|$，故真数始终为正，$f(0)=0$。};
\end{tikzpicture}\end{center}

\textbf{先看图：} 因为 $\sqrt{x^2+1}>|x|$，$U(x)=x+\sqrt{x^2+1}>0$，所以定义域为 $\mathbb R$。换成 $-x$ 后，真数变成正数 $V(x)=\sqrt{x^2+1}-x$。

\textbf{再说明一般情形：}
\[
U(x)V(x)=1,\qquad \log_a U(-x)=\log_a\frac1{U(x)}=-\log_a U(x).
\]
底数从 $2$ 换成 $\frac12$ 时，增减方向改变，原点对称保持不变。
\end{explanationbox}
\clearpage
\begin{explanationbox}
\textbf{\textcolor{CExplanation}{根号减去 $x$，取对数后与上一式反号}}
% Formula-based samples; see F083 checks/build_tikz_diagrams.py.
\begin{center}\begin{tikzpicture}[font=\small,>=stealth,line join=round]
\draw[gray!40,rounded corners=3pt] (0,0) rectangle (14.2,8.4);
\node at (7.1,7.75) {奇7：$y=\log_a\!\bigl(\sqrt{x^2+1}-x\bigr)$};
\begin{scope}[shift={(4.400000,3.900000)},x=1.133333cm,y=0.703125cm]
\draw[->,gray] (-3,0)--(3,0) node[right] {$x$};
\draw[->,gray] (0,-3.2)--(0,3.2) node[above] {$y$};
\draw[gray] (-1,-.035/0.703125)--(-1,.035/0.703125);
\node[anchor=north] at (-1,-.07/0.703125) {$-1$};
\draw[gray] (1,-.035/0.703125)--(1,.035/0.703125);
\node[anchor=north] at (1,-.07/0.703125) {$1$};
\begin{scope}
\clip (-3,-3.2) rectangle (3,3.2);
\draw[blue,thick,solid] plot coordinates {(-3.0000000,2.6234637) (-2.9900000,2.6188946) (-2.9800000,2.6143118) (-2.9700000,2.6097151) (-2.9600000,2.6051046) (-2.9500000,2.6004800) (-2.9400000,2.5958413) (-2.9300000,2.5911885) (-2.9200000,2.5865214) (-2.9100000,2.5818400) (-2.9000000,2.5771441) (-2.8900000,2.5724338) (-2.8800000,2.5677089) (-2.8700000,2.5629694) (-2.8600000,2.5582151) (-2.8500000,2.5534459) (-2.8400000,2.5486619) (-2.8300000,2.5438628) (-2.8200000,2.5390486) (-2.8100000,2.5342193) (-2.8000000,2.5293746) (-2.7900000,2.5245146) (-2.7800000,2.5196392) (-2.7700000,2.5147481) (-2.7600000,2.5098415) (-2.7500000,2.5049191) (-2.7400000,2.4999808) (-2.7300000,2.4950266) (-2.7200000,2.4900564) (-2.7100000,2.4850701) (-2.7000000,2.4800676) (-2.6900000,2.4750487) (-2.6800000,2.4700134) (-2.6700000,2.4649616) (-2.6600000,2.4598932) (-2.6500000,2.4548081) (-2.6400000,2.4497061) (-2.6300000,2.4445872) (-2.6200000,2.4394513) (-2.6100000,2.4342982) (-2.6000000,2.4291279) (-2.5900000,2.4239402) (-2.5800000,2.4187351) (-2.5700000,2.4135124) (-2.5600000,2.4082720) (-2.5500000,2.4030138) (-2.5400000,2.3977378) (-2.5300000,2.3924437) (-2.5200000,2.3871315) (-2.5100000,2.3818010) (-2.5000000,2.3764522) (-2.4900000,2.3710849) (-2.4800000,2.3656990) (-2.4700000,2.3602944) (-2.4600000,2.3548710) (-2.4500000,2.3494286) (-2.4400000,2.3439672) (-2.4300000,2.3384865) (-2.4200000,2.3329865) (-2.4100000,2.3274671) (-2.4000000,2.3219281) (-2.3900000,2.3163694) (-2.3800000,2.3107909) (-2.3700000,2.3051924) (-2.3600000,2.2995738) (-2.3500000,2.2939350) (-2.3400000,2.2882759) (-2.3300000,2.2825962) (-2.3200000,2.2768960) (-2.3100000,2.2711750) (-2.3000000,2.2654331) (-2.2900000,2.2596701) (-2.2800000,2.2538860) (-2.2700000,2.2480806) (-2.2600000,2.2422537) (-2.2500000,2.2364052) (-2.2400000,2.2305350) (-2.2300000,2.2246428) (-2.2200000,2.2187287) (-2.2100000,2.2127923) (-2.2000000,2.2068336) (-2.1900000,2.2008524) (-2.1800000,2.1948486) (-2.1700000,2.1888220) (-2.1600000,2.1827724) (-2.1500000,2.1766997) (-2.1400000,2.1706038) (-2.1300000,2.1644844) (-2.1200000,2.1583414) (-2.1100000,2.1521747) (-2.1000000,2.1459841) (-2.0900000,2.1397694) (-2.0800000,2.1335305) (-2.0700000,2.1272671) (-2.0600000,2.1209792) (-2.0500000,2.1146665) (-2.0400000,2.1083289) (-2.0300000,2.1019662) (-2.0200000,2.0955782) (-2.0100000,2.0891648) (-2.0000000,2.0827257) (-1.9900000,2.0762609) (-1.9800000,2.0697700) (-1.9700000,2.0632530) (-1.9600000,2.0567096) (-1.9500000,2.0501397) (-1.9400000,2.0435430) (-1.9300000,2.0369195) (-1.9200000,2.0302688) (-1.9100000,2.0235908) (-1.9000000,2.0168853) (-1.8900000,2.0101521) (-1.8800000,2.0033911) (-1.8700000,1.9966019) (-1.8600000,1.9897844) (-1.8500000,1.9829385) (-1.8400000,1.9760639) (-1.8300000,1.9691603) (-1.8200000,1.9622276) (-1.8100000,1.9552657) (-1.8000000,1.9482742) (-1.7900000,1.9412529) (-1.7800000,1.9342017) (-1.7700000,1.9271203) (-1.7600000,1.9200085) (-1.7500000,1.9128661) (-1.7400000,1.9056929) (-1.7300000,1.8984886) (-1.7200000,1.8912531) (-1.7100000,1.8839860) (-1.7000000,1.8766872) (-1.6900000,1.8693564) (-1.6800000,1.8619934) (-1.6700000,1.8545980) (-1.6600000,1.8471699) (-1.6500000,1.8397090) (-1.6400000,1.8322148) (-1.6300000,1.8246873) (-1.6200000,1.8171262) (-1.6100000,1.8095312) (-1.6000000,1.8019021) (-1.5900000,1.7942386) (-1.5800000,1.7865405) (-1.5700000,1.7788075) (-1.5600000,1.7710394) (-1.5500000,1.7632359) (-1.5400000,1.7553968) (-1.5300000,1.7475218) (-1.5200000,1.7396107) (-1.5100000,1.7316632) (-1.5000000,1.7236790) (-1.4900000,1.7156578) (-1.4800000,1.7075995) (-1.4700000,1.6995037) (-1.4600000,1.6913701) (-1.4500000,1.6831986) (-1.4400000,1.6749887) (-1.4300000,1.6667403) (-1.4200000,1.6584531) (-1.4100000,1.6501268) (-1.4000000,1.6417611) (-1.3900000,1.6333557) (-1.3800000,1.6249104) (-1.3700000,1.6164248) (-1.3600000,1.6078988) (-1.3500000,1.5993320) (-1.3400000,1.5907241) (-1.3300000,1.5820748) (-1.3200000,1.5733839) (-1.3100000,1.5646511) (-1.3000000,1.5558760) (-1.2900000,1.5470584) (-1.2800000,1.5381981) (-1.2700000,1.5292946) (-1.2600000,1.5203478) (-1.2500000,1.5113572) (-1.2400000,1.5023228) (-1.2300000,1.4932440) (-1.2200000,1.4841208) (-1.2100000,1.4749526) (-1.2000000,1.4657394) (-1.1900000,1.4564807) (-1.1800000,1.4471764) (-1.1700000,1.4378260) (-1.1600000,1.4284293) (-1.1500000,1.4189860) (-1.1400000,1.4094959) (-1.1300000,1.3999586) (-1.1200000,1.3903738) (-1.1100000,1.3807413) (-1.1000000,1.3710608) (-1.0900000,1.3613320) (-1.0800000,1.3515546) (-1.0700000,1.3417283) (-1.0600000,1.3318529) (-1.0500000,1.3219281) (-1.0400000,1.3119536) (-1.0300000,1.3019291) (-1.0200000,1.2918544) (-1.0100000,1.2817292) (-1.0000000,1.2715533) (-0.9900000,1.2613264) (-0.9800000,1.2510482) (-0.9700000,1.2407185) (-0.9600000,1.2303370) (-0.9500000,1.2199035) (-0.9400000,1.2094179) (-0.9300000,1.1988797) (-0.9200000,1.1882889) (-0.9100000,1.1776452) (-0.9000000,1.1669483) (-0.8900000,1.1561982) (-0.8800000,1.1453945) (-0.8700000,1.1345371) (-0.8600000,1.1236258) (-0.8500000,1.1126604) (-0.8400000,1.1016408) (-0.8300000,1.0905668) (-0.8200000,1.0794382) (-0.8100000,1.0682550) (-0.8000000,1.0570169) (-0.7900000,1.0457238) (-0.7800000,1.0343757) (-0.7700000,1.0229724) (-0.7600000,1.0115139) (-0.7500000,1.0000000) (-0.7400000,0.9884307) (-0.7300000,0.9768060) (-0.7200000,0.9651258) (-0.7100000,0.9533901) (-0.7000000,0.9415988) (-0.6900000,0.9297520) (-0.6800000,0.9178498) (-0.6700000,0.9058920) (-0.6600000,0.8938788) (-0.6500000,0.8818103) (-0.6400000,0.8696865) (-0.6300000,0.8575075) (-0.6200000,0.8452735) (-0.6100000,0.8329845) (-0.6000000,0.8206409) (-0.5900000,0.8082426) (-0.5800000,0.7957900) (-0.5700000,0.7832831) (-0.5600000,0.7707224) (-0.5500000,0.7581080) (-0.5400000,0.7454403) (-0.5300000,0.7327194) (-0.5200000,0.7199459) (-0.5100000,0.7071199) (-0.5000000,0.6942419) (-0.4900000,0.6813123) (-0.4800000,0.6683315) (-0.4700000,0.6553000) (-0.4600000,0.6422182) (-0.4500000,0.6290867) (-0.4400000,0.6159059) (-0.4300000,0.6026765) (-0.4200000,0.5893989) (-0.4100000,0.5760739) (-0.4000000,0.5627020) (-0.3900000,0.5492840) (-0.3800000,0.5358204) (-0.3700000,0.5223120) (-0.3600000,0.5087597) (-0.3500000,0.4951640) (-0.3400000,0.4815259) (-0.3300000,0.4678462) (-0.3200000,0.4541257) (-0.3100000,0.4403654) (-0.3000000,0.4265660) (-0.2900000,0.4127287) (-0.2800000,0.3988542) (-0.2700000,0.3849437) (-0.2600000,0.3709982) (-0.2500000,0.3570186) (-0.2400000,0.3430062) (-0.2300000,0.3289618) (-0.2200000,0.3148868) (-0.2100000,0.3007822) (-0.2000000,0.2866492) (-0.1900000,0.2724890) (-0.1800000,0.2583029) (-0.1700000,0.2440919) (-0.1600000,0.2298575) (-0.1500000,0.2156008) (-0.1400000,0.2013233) (-0.1300000,0.1870261) (-0.1200000,0.1727106) (-0.1100000,0.1583781) (-0.1000000,0.1440301) (-0.0900000,0.1296679) (-0.0800000,0.1152928) (-0.0700000,0.1009064) (-0.0600000,0.0865098) (-0.0500000,0.0721047) (-0.0400000,0.0576924) (-0.0300000,0.0432744) (-0.0200000,0.0288520) (-0.0100000,0.0144267) (0.0000000,-0.0000000) (0.0100000,-0.0144267) (0.0200000,-0.0288520) (0.0300000,-0.0432744) (0.0400000,-0.0576924) (0.0500000,-0.0721047) (0.0600000,-0.0865098) (0.0700000,-0.1009064) (0.0800000,-0.1152928) (0.0900000,-0.1296679) (0.1000000,-0.1440301) (0.1100000,-0.1583781) (0.1200000,-0.1727106) (0.1300000,-0.1870261) (0.1400000,-0.2013233) (0.1500000,-0.2156008) (0.1600000,-0.2298575) (0.1700000,-0.2440919) (0.1800000,-0.2583029) (0.1900000,-0.2724890) (0.2000000,-0.2866492) (0.2100000,-0.3007822) (0.2200000,-0.3148868) (0.2300000,-0.3289618) (0.2400000,-0.3430062) (0.2500000,-0.3570186) (0.2600000,-0.3709982) (0.2700000,-0.3849437) (0.2800000,-0.3988542) (0.2900000,-0.4127287) (0.3000000,-0.4265660) (0.3100000,-0.4403654) (0.3200000,-0.4541257) (0.3300000,-0.4678462) (0.3400000,-0.4815259) (0.3500000,-0.4951640) (0.3600000,-0.5087597) (0.3700000,-0.5223120) (0.3800000,-0.5358204) (0.3900000,-0.5492840) (0.4000000,-0.5627020) (0.4100000,-0.5760739) (0.4200000,-0.5893989) (0.4300000,-0.6026765) (0.4400000,-0.6159059) (0.4500000,-0.6290867) (0.4600000,-0.6422182) (0.4700000,-0.6553000) (0.4800000,-0.6683315) (0.4900000,-0.6813123) (0.5000000,-0.6942419) (0.5100000,-0.7071199) (0.5200000,-0.7199459) (0.5300000,-0.7327194) (0.5400000,-0.7454403) (0.5500000,-0.7581080) (0.5600000,-0.7707224) (0.5700000,-0.7832831) (0.5800000,-0.7957900) (0.5900000,-0.8082426) (0.6000000,-0.8206409) (0.6100000,-0.8329845) (0.6200000,-0.8452735) (0.6300000,-0.8575075) (0.6400000,-0.8696865) (0.6500000,-0.8818103) (0.6600000,-0.8938788) (0.6700000,-0.9058920) (0.6800000,-0.9178498) (0.6900000,-0.9297520) (0.7000000,-0.9415988) (0.7100000,-0.9533901) (0.7200000,-0.9651258) (0.7300000,-0.9768060) (0.7400000,-0.9884307) (0.7500000,-1.0000000) (0.7600000,-1.0115139) (0.7700000,-1.0229724) (0.7800000,-1.0343757) (0.7900000,-1.0457238) (0.8000000,-1.0570169) (0.8100000,-1.0682550) (0.8200000,-1.0794382) (0.8300000,-1.0905668) (0.8400000,-1.1016408) (0.8500000,-1.1126604) (0.8600000,-1.1236258) (0.8700000,-1.1345371) (0.8800000,-1.1453945) (0.8900000,-1.1561982) (0.9000000,-1.1669483) (0.9100000,-1.1776452) (0.9200000,-1.1882889) (0.9300000,-1.1988797) (0.9400000,-1.2094179) (0.9500000,-1.2199035) (0.9600000,-1.2303370) (0.9700000,-1.2407185) (0.9800000,-1.2510482) (0.9900000,-1.2613264) (1.0000000,-1.2715533) (1.0100000,-1.2817292) (1.0200000,-1.2918544) (1.0300000,-1.3019291) (1.0400000,-1.3119536) (1.0500000,-1.3219281) (1.0600000,-1.3318529) (1.0700000,-1.3417283) (1.0800000,-1.3515546) (1.0900000,-1.3613320) (1.1000000,-1.3710608) (1.1100000,-1.3807413) (1.1200000,-1.3903738) (1.1300000,-1.3999586) (1.1400000,-1.4094959) (1.1500000,-1.4189860) (1.1600000,-1.4284293) (1.1700000,-1.4378260) (1.1800000,-1.4471764) (1.1900000,-1.4564807) (1.2000000,-1.4657394) (1.2100000,-1.4749526) (1.2200000,-1.4841208) (1.2300000,-1.4932440) (1.2400000,-1.5023228) (1.2500000,-1.5113572) (1.2600000,-1.5203478) (1.2700000,-1.5292946) (1.2800000,-1.5381981) (1.2900000,-1.5470584) (1.3000000,-1.5558760) (1.3100000,-1.5646511) (1.3200000,-1.5733839) (1.3300000,-1.5820748) (1.3400000,-1.5907241) (1.3500000,-1.5993320) (1.3600000,-1.6078988) (1.3700000,-1.6164248) (1.3800000,-1.6249104) (1.3900000,-1.6333557) (1.4000000,-1.6417611) (1.4100000,-1.6501268) (1.4200000,-1.6584531) (1.4300000,-1.6667403) (1.4400000,-1.6749887) (1.4500000,-1.6831986) (1.4600000,-1.6913701) (1.4700000,-1.6995037) (1.4800000,-1.7075995) (1.4900000,-1.7156578) (1.5000000,-1.7236790) (1.5100000,-1.7316632) (1.5200000,-1.7396107) (1.5300000,-1.7475218) (1.5400000,-1.7553968) (1.5500000,-1.7632359) (1.5600000,-1.7710394) (1.5700000,-1.7788075) (1.5800000,-1.7865405) (1.5900000,-1.7942386) (1.6000000,-1.8019021) (1.6100000,-1.8095312) (1.6200000,-1.8171262) (1.6300000,-1.8246873) (1.6400000,-1.8322148) (1.6500000,-1.8397090) (1.6600000,-1.8471699) (1.6700000,-1.8545980) (1.6800000,-1.8619934) (1.6900000,-1.8693564) (1.7000000,-1.8766872) (1.7100000,-1.8839860) (1.7200000,-1.8912531) (1.7300000,-1.8984886) (1.7400000,-1.9056929) (1.7500000,-1.9128661) (1.7600000,-1.9200085) (1.7700000,-1.9271203) (1.7800000,-1.9342017) (1.7900000,-1.9412529) (1.8000000,-1.9482742) (1.8100000,-1.9552657) (1.8200000,-1.9622276) (1.8300000,-1.9691603) (1.8400000,-1.9760639) (1.8500000,-1.9829385) (1.8600000,-1.9897844) (1.8700000,-1.9966019) (1.8800000,-2.0033911) (1.8900000,-2.0101521) (1.9000000,-2.0168853) (1.9100000,-2.0235908) (1.9200000,-2.0302688) (1.9300000,-2.0369195) (1.9400000,-2.0435430) (1.9500000,-2.0501397) (1.9600000,-2.0567096) (1.9700000,-2.0632530) (1.9800000,-2.0697700) (1.9900000,-2.0762609) (2.0000000,-2.0827257) (2.0100000,-2.0891648) (2.0200000,-2.0955782) (2.0300000,-2.1019662) (2.0400000,-2.1083289) (2.0500000,-2.1146665) (2.0600000,-2.1209792) (2.0700000,-2.1272671) (2.0800000,-2.1335305) (2.0900000,-2.1397694) (2.1000000,-2.1459841) (2.1100000,-2.1521747) (2.1200000,-2.1583414) (2.1300000,-2.1644844) (2.1400000,-2.1706038) (2.1500000,-2.1766997) (2.1600000,-2.1827724) (2.1700000,-2.1888220) (2.1800000,-2.1948486) (2.1900000,-2.2008524) (2.2000000,-2.2068336) (2.2100000,-2.2127923) (2.2200000,-2.2187287) (2.2300000,-2.2246428) (2.2400000,-2.2305350) (2.2500000,-2.2364052) (2.2600000,-2.2422537) (2.2700000,-2.2480806) (2.2800000,-2.2538860) (2.2900000,-2.2596701) (2.3000000,-2.2654331) (2.3100000,-2.2711750) (2.3200000,-2.2768960) (2.3300000,-2.2825962) (2.3400000,-2.2882759) (2.3500000,-2.2939350) (2.3600000,-2.2995738) (2.3700000,-2.3051924) (2.3800000,-2.3107909) (2.3900000,-2.3163694) (2.4000000,-2.3219281) (2.4100000,-2.3274671) (2.4200000,-2.3329865) (2.4300000,-2.3384865) (2.4400000,-2.3439672) (2.4500000,-2.3494286) (2.4600000,-2.3548710) (2.4700000,-2.3602944) (2.4800000,-2.3656990) (2.4900000,-2.3710849) (2.5000000,-2.3764522) (2.5100000,-2.3818010) (2.5200000,-2.3871315) (2.5300000,-2.3924437) (2.5400000,-2.3977378) (2.5500000,-2.4030138) (2.5600000,-2.4082720) (2.5700000,-2.4135124) (2.5800000,-2.4187351) (2.5900000,-2.4239402) (2.6000000,-2.4291279) (2.6100000,-2.4342982) (2.6200000,-2.4394513) (2.6300000,-2.4445872) (2.6400000,-2.4497061) (2.6500000,-2.4548081) (2.6600000,-2.4598932) (2.6700000,-2.4649616) (2.6800000,-2.4700134) (2.6900000,-2.4750487) (2.7000000,-2.4800676) (2.7100000,-2.4850701) (2.7200000,-2.4900564) (2.7300000,-2.4950266) (2.7400000,-2.4999808) (2.7500000,-2.5049191) (2.7600000,-2.5098415) (2.7700000,-2.5147481) (2.7800000,-2.5196392) (2.7900000,-2.5245146) (2.8000000,-2.5293746) (2.8100000,-2.5342193) (2.8200000,-2.5390486) (2.8300000,-2.5438628) (2.8400000,-2.5486619) (2.8500000,-2.5534459) (2.8600000,-2.5582151) (2.8700000,-2.5629694) (2.8800000,-2.5677089) (2.8900000,-2.5724338) (2.9000000,-2.5771441) (2.9100000,-2.5818400) (2.9200000,-2.5865214) (2.9300000,-2.5911885) (2.9400000,-2.5958413) (2.9500000,-2.6004800) (2.9600000,-2.6051046) (2.9700000,-2.6097151) (2.9800000,-2.6143118) (2.9900000,-2.6188946) (3.0000000,-2.6234637)};
\draw[orange!85!black,thick,dashed] plot coordinates {(-3.0000000,-2.6234637) (-2.9900000,-2.6188946) (-2.9800000,-2.6143118) (-2.9700000,-2.6097151) (-2.9600000,-2.6051046) (-2.9500000,-2.6004800) (-2.9400000,-2.5958413) (-2.9300000,-2.5911885) (-2.9200000,-2.5865214) (-2.9100000,-2.5818400) (-2.9000000,-2.5771441) (-2.8900000,-2.5724338) (-2.8800000,-2.5677089) (-2.8700000,-2.5629694) (-2.8600000,-2.5582151) (-2.8500000,-2.5534459) (-2.8400000,-2.5486619) (-2.8300000,-2.5438628) (-2.8200000,-2.5390486) (-2.8100000,-2.5342193) (-2.8000000,-2.5293746) (-2.7900000,-2.5245146) (-2.7800000,-2.5196392) (-2.7700000,-2.5147481) (-2.7600000,-2.5098415) (-2.7500000,-2.5049191) (-2.7400000,-2.4999808) (-2.7300000,-2.4950266) (-2.7200000,-2.4900564) (-2.7100000,-2.4850701) (-2.7000000,-2.4800676) (-2.6900000,-2.4750487) (-2.6800000,-2.4700134) (-2.6700000,-2.4649616) (-2.6600000,-2.4598932) (-2.6500000,-2.4548081) (-2.6400000,-2.4497061) (-2.6300000,-2.4445872) (-2.6200000,-2.4394513) (-2.6100000,-2.4342982) (-2.6000000,-2.4291279) (-2.5900000,-2.4239402) (-2.5800000,-2.4187351) (-2.5700000,-2.4135124) (-2.5600000,-2.4082720) (-2.5500000,-2.4030138) (-2.5400000,-2.3977378) (-2.5300000,-2.3924437) (-2.5200000,-2.3871315) (-2.5100000,-2.3818010) (-2.5000000,-2.3764522) (-2.4900000,-2.3710849) (-2.4800000,-2.3656990) (-2.4700000,-2.3602944) (-2.4600000,-2.3548710) (-2.4500000,-2.3494286) (-2.4400000,-2.3439672) (-2.4300000,-2.3384865) (-2.4200000,-2.3329865) (-2.4100000,-2.3274671) (-2.4000000,-2.3219281) (-2.3900000,-2.3163694) (-2.3800000,-2.3107909) (-2.3700000,-2.3051924) (-2.3600000,-2.2995738) (-2.3500000,-2.2939350) (-2.3400000,-2.2882759) (-2.3300000,-2.2825962) (-2.3200000,-2.2768960) (-2.3100000,-2.2711750) (-2.3000000,-2.2654331) (-2.2900000,-2.2596701) (-2.2800000,-2.2538860) (-2.2700000,-2.2480806) (-2.2600000,-2.2422537) (-2.2500000,-2.2364052) (-2.2400000,-2.2305350) (-2.2300000,-2.2246428) (-2.2200000,-2.2187287) (-2.2100000,-2.2127923) (-2.2000000,-2.2068336) (-2.1900000,-2.2008524) (-2.1800000,-2.1948486) (-2.1700000,-2.1888220) (-2.1600000,-2.1827724) (-2.1500000,-2.1766997) (-2.1400000,-2.1706038) (-2.1300000,-2.1644844) (-2.1200000,-2.1583414) (-2.1100000,-2.1521747) (-2.1000000,-2.1459841) (-2.0900000,-2.1397694) (-2.0800000,-2.1335305) (-2.0700000,-2.1272671) (-2.0600000,-2.1209792) (-2.0500000,-2.1146665) (-2.0400000,-2.1083289) (-2.0300000,-2.1019662) (-2.0200000,-2.0955782) (-2.0100000,-2.0891648) (-2.0000000,-2.0827257) (-1.9900000,-2.0762609) (-1.9800000,-2.0697700) (-1.9700000,-2.0632530) (-1.9600000,-2.0567096) (-1.9500000,-2.0501397) (-1.9400000,-2.0435430) (-1.9300000,-2.0369195) (-1.9200000,-2.0302688) (-1.9100000,-2.0235908) (-1.9000000,-2.0168853) (-1.8900000,-2.0101521) (-1.8800000,-2.0033911) (-1.8700000,-1.9966019) (-1.8600000,-1.9897844) (-1.8500000,-1.9829385) (-1.8400000,-1.9760639) (-1.8300000,-1.9691603) (-1.8200000,-1.9622276) (-1.8100000,-1.9552657) (-1.8000000,-1.9482742) (-1.7900000,-1.9412529) (-1.7800000,-1.9342017) (-1.7700000,-1.9271203) (-1.7600000,-1.9200085) (-1.7500000,-1.9128661) (-1.7400000,-1.9056929) (-1.7300000,-1.8984886) (-1.7200000,-1.8912531) (-1.7100000,-1.8839860) (-1.7000000,-1.8766872) (-1.6900000,-1.8693564) (-1.6800000,-1.8619934) (-1.6700000,-1.8545980) (-1.6600000,-1.8471699) (-1.6500000,-1.8397090) (-1.6400000,-1.8322148) (-1.6300000,-1.8246873) (-1.6200000,-1.8171262) (-1.6100000,-1.8095312) (-1.6000000,-1.8019021) (-1.5900000,-1.7942386) (-1.5800000,-1.7865405) (-1.5700000,-1.7788075) (-1.5600000,-1.7710394) (-1.5500000,-1.7632359) (-1.5400000,-1.7553968) (-1.5300000,-1.7475218) (-1.5200000,-1.7396107) (-1.5100000,-1.7316632) (-1.5000000,-1.7236790) (-1.4900000,-1.7156578) (-1.4800000,-1.7075995) (-1.4700000,-1.6995037) (-1.4600000,-1.6913701) (-1.4500000,-1.6831986) (-1.4400000,-1.6749887) (-1.4300000,-1.6667403) (-1.4200000,-1.6584531) (-1.4100000,-1.6501268) (-1.4000000,-1.6417611) (-1.3900000,-1.6333557) (-1.3800000,-1.6249104) (-1.3700000,-1.6164248) (-1.3600000,-1.6078988) (-1.3500000,-1.5993320) (-1.3400000,-1.5907241) (-1.3300000,-1.5820748) (-1.3200000,-1.5733839) (-1.3100000,-1.5646511) (-1.3000000,-1.5558760) (-1.2900000,-1.5470584) (-1.2800000,-1.5381981) (-1.2700000,-1.5292946) (-1.2600000,-1.5203478) (-1.2500000,-1.5113572) (-1.2400000,-1.5023228) (-1.2300000,-1.4932440) (-1.2200000,-1.4841208) (-1.2100000,-1.4749526) (-1.2000000,-1.4657394) (-1.1900000,-1.4564807) (-1.1800000,-1.4471764) (-1.1700000,-1.4378260) (-1.1600000,-1.4284293) (-1.1500000,-1.4189860) (-1.1400000,-1.4094959) (-1.1300000,-1.3999586) (-1.1200000,-1.3903738) (-1.1100000,-1.3807413) (-1.1000000,-1.3710608) (-1.0900000,-1.3613320) (-1.0800000,-1.3515546) (-1.0700000,-1.3417283) (-1.0600000,-1.3318529) (-1.0500000,-1.3219281) (-1.0400000,-1.3119536) (-1.0300000,-1.3019291) (-1.0200000,-1.2918544) (-1.0100000,-1.2817292) (-1.0000000,-1.2715533) (-0.9900000,-1.2613264) (-0.9800000,-1.2510482) (-0.9700000,-1.2407185) (-0.9600000,-1.2303370) (-0.9500000,-1.2199035) (-0.9400000,-1.2094179) (-0.9300000,-1.1988797) (-0.9200000,-1.1882889) (-0.9100000,-1.1776452) (-0.9000000,-1.1669483) (-0.8900000,-1.1561982) (-0.8800000,-1.1453945) (-0.8700000,-1.1345371) (-0.8600000,-1.1236258) (-0.8500000,-1.1126604) (-0.8400000,-1.1016408) (-0.8300000,-1.0905668) (-0.8200000,-1.0794382) (-0.8100000,-1.0682550) (-0.8000000,-1.0570169) (-0.7900000,-1.0457238) (-0.7800000,-1.0343757) (-0.7700000,-1.0229724) (-0.7600000,-1.0115139) (-0.7500000,-1.0000000) (-0.7400000,-0.9884307) (-0.7300000,-0.9768060) (-0.7200000,-0.9651258) (-0.7100000,-0.9533901) (-0.7000000,-0.9415988) (-0.6900000,-0.9297520) (-0.6800000,-0.9178498) (-0.6700000,-0.9058920) (-0.6600000,-0.8938788) (-0.6500000,-0.8818103) (-0.6400000,-0.8696865) (-0.6300000,-0.8575075) (-0.6200000,-0.8452735) (-0.6100000,-0.8329845) (-0.6000000,-0.8206409) (-0.5900000,-0.8082426) (-0.5800000,-0.7957900) (-0.5700000,-0.7832831) (-0.5600000,-0.7707224) (-0.5500000,-0.7581080) (-0.5400000,-0.7454403) (-0.5300000,-0.7327194) (-0.5200000,-0.7199459) (-0.5100000,-0.7071199) (-0.5000000,-0.6942419) (-0.4900000,-0.6813123) (-0.4800000,-0.6683315) (-0.4700000,-0.6553000) (-0.4600000,-0.6422182) (-0.4500000,-0.6290867) (-0.4400000,-0.6159059) (-0.4300000,-0.6026765) (-0.4200000,-0.5893989) (-0.4100000,-0.5760739) (-0.4000000,-0.5627020) (-0.3900000,-0.5492840) (-0.3800000,-0.5358204) (-0.3700000,-0.5223120) (-0.3600000,-0.5087597) (-0.3500000,-0.4951640) (-0.3400000,-0.4815259) (-0.3300000,-0.4678462) (-0.3200000,-0.4541257) (-0.3100000,-0.4403654) (-0.3000000,-0.4265660) (-0.2900000,-0.4127287) (-0.2800000,-0.3988542) (-0.2700000,-0.3849437) (-0.2600000,-0.3709982) (-0.2500000,-0.3570186) (-0.2400000,-0.3430062) (-0.2300000,-0.3289618) (-0.2200000,-0.3148868) (-0.2100000,-0.3007822) (-0.2000000,-0.2866492) (-0.1900000,-0.2724890) (-0.1800000,-0.2583029) (-0.1700000,-0.2440919) (-0.1600000,-0.2298575) (-0.1500000,-0.2156008) (-0.1400000,-0.2013233) (-0.1300000,-0.1870261) (-0.1200000,-0.1727106) (-0.1100000,-0.1583781) (-0.1000000,-0.1440301) (-0.0900000,-0.1296679) (-0.0800000,-0.1152928) (-0.0700000,-0.1009064) (-0.0600000,-0.0865098) (-0.0500000,-0.0721047) (-0.0400000,-0.0576924) (-0.0300000,-0.0432744) (-0.0200000,-0.0288520) (-0.0100000,-0.0144267) (0.0000000,0.0000000) (0.0100000,0.0144267) (0.0200000,0.0288520) (0.0300000,0.0432744) (0.0400000,0.0576924) (0.0500000,0.0721047) (0.0600000,0.0865098) (0.0700000,0.1009064) (0.0800000,0.1152928) (0.0900000,0.1296679) (0.1000000,0.1440301) (0.1100000,0.1583781) (0.1200000,0.1727106) (0.1300000,0.1870261) (0.1400000,0.2013233) (0.1500000,0.2156008) (0.1600000,0.2298575) (0.1700000,0.2440919) (0.1800000,0.2583029) (0.1900000,0.2724890) (0.2000000,0.2866492) (0.2100000,0.3007822) (0.2200000,0.3148868) (0.2300000,0.3289618) (0.2400000,0.3430062) (0.2500000,0.3570186) (0.2600000,0.3709982) (0.2700000,0.3849437) (0.2800000,0.3988542) (0.2900000,0.4127287) (0.3000000,0.4265660) (0.3100000,0.4403654) (0.3200000,0.4541257) (0.3300000,0.4678462) (0.3400000,0.4815259) (0.3500000,0.4951640) (0.3600000,0.5087597) (0.3700000,0.5223120) (0.3800000,0.5358204) (0.3900000,0.5492840) (0.4000000,0.5627020) (0.4100000,0.5760739) (0.4200000,0.5893989) (0.4300000,0.6026765) (0.4400000,0.6159059) (0.4500000,0.6290867) (0.4600000,0.6422182) (0.4700000,0.6553000) (0.4800000,0.6683315) (0.4900000,0.6813123) (0.5000000,0.6942419) (0.5100000,0.7071199) (0.5200000,0.7199459) (0.5300000,0.7327194) (0.5400000,0.7454403) (0.5500000,0.7581080) (0.5600000,0.7707224) (0.5700000,0.7832831) (0.5800000,0.7957900) (0.5900000,0.8082426) (0.6000000,0.8206409) (0.6100000,0.8329845) (0.6200000,0.8452735) (0.6300000,0.8575075) (0.6400000,0.8696865) (0.6500000,0.8818103) (0.6600000,0.8938788) (0.6700000,0.9058920) (0.6800000,0.9178498) (0.6900000,0.9297520) (0.7000000,0.9415988) (0.7100000,0.9533901) (0.7200000,0.9651258) (0.7300000,0.9768060) (0.7400000,0.9884307) (0.7500000,1.0000000) (0.7600000,1.0115139) (0.7700000,1.0229724) (0.7800000,1.0343757) (0.7900000,1.0457238) (0.8000000,1.0570169) (0.8100000,1.0682550) (0.8200000,1.0794382) (0.8300000,1.0905668) (0.8400000,1.1016408) (0.8500000,1.1126604) (0.8600000,1.1236258) (0.8700000,1.1345371) (0.8800000,1.1453945) (0.8900000,1.1561982) (0.9000000,1.1669483) (0.9100000,1.1776452) (0.9200000,1.1882889) (0.9300000,1.1988797) (0.9400000,1.2094179) (0.9500000,1.2199035) (0.9600000,1.2303370) (0.9700000,1.2407185) (0.9800000,1.2510482) (0.9900000,1.2613264) (1.0000000,1.2715533) (1.0100000,1.2817292) (1.0200000,1.2918544) (1.0300000,1.3019291) (1.0400000,1.3119536) (1.0500000,1.3219281) (1.0600000,1.3318529) (1.0700000,1.3417283) (1.0800000,1.3515546) (1.0900000,1.3613320) (1.1000000,1.3710608) (1.1100000,1.3807413) (1.1200000,1.3903738) (1.1300000,1.3999586) (1.1400000,1.4094959) (1.1500000,1.4189860) (1.1600000,1.4284293) (1.1700000,1.4378260) (1.1800000,1.4471764) (1.1900000,1.4564807) (1.2000000,1.4657394) (1.2100000,1.4749526) (1.2200000,1.4841208) (1.2300000,1.4932440) (1.2400000,1.5023228) (1.2500000,1.5113572) (1.2600000,1.5203478) (1.2700000,1.5292946) (1.2800000,1.5381981) (1.2900000,1.5470584) (1.3000000,1.5558760) (1.3100000,1.5646511) (1.3200000,1.5733839) (1.3300000,1.5820748) (1.3400000,1.5907241) (1.3500000,1.5993320) (1.3600000,1.6078988) (1.3700000,1.6164248) (1.3800000,1.6249104) (1.3900000,1.6333557) (1.4000000,1.6417611) (1.4100000,1.6501268) (1.4200000,1.6584531) (1.4300000,1.6667403) (1.4400000,1.6749887) (1.4500000,1.6831986) (1.4600000,1.6913701) (1.4700000,1.6995037) (1.4800000,1.7075995) (1.4900000,1.7156578) (1.5000000,1.7236790) (1.5100000,1.7316632) (1.5200000,1.7396107) (1.5300000,1.7475218) (1.5400000,1.7553968) (1.5500000,1.7632359) (1.5600000,1.7710394) (1.5700000,1.7788075) (1.5800000,1.7865405) (1.5900000,1.7942386) (1.6000000,1.8019021) (1.6100000,1.8095312) (1.6200000,1.8171262) (1.6300000,1.8246873) (1.6400000,1.8322148) (1.6500000,1.8397090) (1.6600000,1.8471699) (1.6700000,1.8545980) (1.6800000,1.8619934) (1.6900000,1.8693564) (1.7000000,1.8766872) (1.7100000,1.8839860) (1.7200000,1.8912531) (1.7300000,1.8984886) (1.7400000,1.9056929) (1.7500000,1.9128661) (1.7600000,1.9200085) (1.7700000,1.9271203) (1.7800000,1.9342017) (1.7900000,1.9412529) (1.8000000,1.9482742) (1.8100000,1.9552657) (1.8200000,1.9622276) (1.8300000,1.9691603) (1.8400000,1.9760639) (1.8500000,1.9829385) (1.8600000,1.9897844) (1.8700000,1.9966019) (1.8800000,2.0033911) (1.8900000,2.0101521) (1.9000000,2.0168853) (1.9100000,2.0235908) (1.9200000,2.0302688) (1.9300000,2.0369195) (1.9400000,2.0435430) (1.9500000,2.0501397) (1.9600000,2.0567096) (1.9700000,2.0632530) (1.9800000,2.0697700) (1.9900000,2.0762609) (2.0000000,2.0827257) (2.0100000,2.0891648) (2.0200000,2.0955782) (2.0300000,2.1019662) (2.0400000,2.1083289) (2.0500000,2.1146665) (2.0600000,2.1209792) (2.0700000,2.1272671) (2.0800000,2.1335305) (2.0900000,2.1397694) (2.1000000,2.1459841) (2.1100000,2.1521747) (2.1200000,2.1583414) (2.1300000,2.1644844) (2.1400000,2.1706038) (2.1500000,2.1766997) (2.1600000,2.1827724) (2.1700000,2.1888220) (2.1800000,2.1948486) (2.1900000,2.2008524) (2.2000000,2.2068336) (2.2100000,2.2127923) (2.2200000,2.2187287) (2.2300000,2.2246428) (2.2400000,2.2305350) (2.2500000,2.2364052) (2.2600000,2.2422537) (2.2700000,2.2480806) (2.2800000,2.2538860) (2.2900000,2.2596701) (2.3000000,2.2654331) (2.3100000,2.2711750) (2.3200000,2.2768960) (2.3300000,2.2825962) (2.3400000,2.2882759) (2.3500000,2.2939350) (2.3600000,2.2995738) (2.3700000,2.3051924) (2.3800000,2.3107909) (2.3900000,2.3163694) (2.4000000,2.3219281) (2.4100000,2.3274671) (2.4200000,2.3329865) (2.4300000,2.3384865) (2.4400000,2.3439672) (2.4500000,2.3494286) (2.4600000,2.3548710) (2.4700000,2.3602944) (2.4800000,2.3656990) (2.4900000,2.3710849) (2.5000000,2.3764522) (2.5100000,2.3818010) (2.5200000,2.3871315) (2.5300000,2.3924437) (2.5400000,2.3977378) (2.5500000,2.4030138) (2.5600000,2.4082720) (2.5700000,2.4135124) (2.5800000,2.4187351) (2.5900000,2.4239402) (2.6000000,2.4291279) (2.6100000,2.4342982) (2.6200000,2.4394513) (2.6300000,2.4445872) (2.6400000,2.4497061) (2.6500000,2.4548081) (2.6600000,2.4598932) (2.6700000,2.4649616) (2.6800000,2.4700134) (2.6900000,2.4750487) (2.7000000,2.4800676) (2.7100000,2.4850701) (2.7200000,2.4900564) (2.7300000,2.4950266) (2.7400000,2.4999808) (2.7500000,2.5049191) (2.7600000,2.5098415) (2.7700000,2.5147481) (2.7800000,2.5196392) (2.7900000,2.5245146) (2.8000000,2.5293746) (2.8100000,2.5342193) (2.8200000,2.5390486) (2.8300000,2.5438628) (2.8400000,2.5486619) (2.8500000,2.5534459) (2.8600000,2.5582151) (2.8700000,2.5629694) (2.8800000,2.5677089) (2.8900000,2.5724338) (2.9000000,2.5771441) (2.9100000,2.5818400) (2.9200000,2.5865214) (2.9300000,2.5911885) (2.9400000,2.5958413) (2.9500000,2.6004800) (2.9600000,2.6051046) (2.9700000,2.6097151) (2.9800000,2.6143118) (2.9900000,2.6188946) (3.0000000,2.6234637)};
\draw[violet!75!black,densely dotted] (-1,1.2715533)--(1,-1.2715533);
\fill[violet!75!black] (-1,1.2715533) circle(2.3pt);
\fill[violet!75!black] (1,-1.2715533) circle(2.3pt);
\end{scope}
\end{scope}
\node[anchor=west,align=left,text width=4.35cm] at (8.9,4.4) {蓝实线：$a=2$\\[8pt]橙虚线：$a=\frac12$\\[14pt]同一底数下，本式\\与奇6上下互为相反数。\\每条曲线也原点对称。};
\node[text width=12.8cm,align=center] at (7.1,.65) {$D=\mathbb R$；两个根式真数互为倒数，取对数后互为相反数。};
\end{tikzpicture}\end{center}

\textbf{先看图：} 本式的真数 $V(x)$ 也在所有实数上为正。与奇6使用相同底数时，$V=1/U$，所以函数值恰好互为相反数。

\textbf{再说明一般情形：}
\[
\log_a(\sqrt{x^2+1}-x)=-\log_a(x+\sqrt{x^2+1}).
\]
它仍是奇函数，且 $f(0)=0$。两个根式真数本身都是正数，不能说它们互为相反数。
\end{explanationbox}
\clearpage
\begin{explanationbox}
\textbf{\textcolor{CExplanation}{指数互换，和保持不变}}
% Formula-based samples; see F083 checks/build_tikz_diagrams.py.
\begin{center}\begin{tikzpicture}[font=\small,>=stealth,line join=round]
\draw[gray!40,rounded corners=3pt] (0,0) rectangle (14.2,8.4);
\node at (7.1,7.75) {偶1：$y=a^x+a^{-x}$};
\begin{scope}[shift={(4.400000,1.650000)},x=1.700000cm,y=0.529412cm]
\draw[->,gray] (-2,0)--(2,0) node[right] {$x$};
\draw[->,gray] (0,0)--(0,8.5) node[above] {$y$};
\draw[gray] (-1,-.035/0.529412)--(-1,.035/0.529412);
\node[anchor=north] at (-1,-.07/0.529412) {$-1$};
\draw[gray] (1,-.035/0.529412)--(1,.035/0.529412);
\node[anchor=north] at (1,-.07/0.529412) {$1$};
\begin{scope}
\clip (-2,0) rectangle (2,8.5);
\draw[blue,thick,solid] plot coordinates {(-2.0000000,4.2500000) (-1.9933333,4.2327166) (-1.9866667,4.2155237) (-1.9800000,4.1984207) (-1.9733333,4.1814074) (-1.9666667,4.1644833) (-1.9600000,4.1476482) (-1.9533333,4.1309017) (-1.9466667,4.1142434) (-1.9400000,4.0976729) (-1.9333333,4.0811899) (-1.9266667,4.0647941) (-1.9200000,4.0484851) (-1.9133333,4.0322625) (-1.9066667,4.0161260) (-1.9000000,4.0000753) (-1.8933333,3.9841100) (-1.8866667,3.9682298) (-1.8800000,3.9524343) (-1.8733333,3.9367232) (-1.8666667,3.9210962) (-1.8600000,3.9055529) (-1.8533333,3.8900930) (-1.8466667,3.8747162) (-1.8400000,3.8594221) (-1.8333333,3.8442104) (-1.8266667,3.8290808) (-1.8200000,3.8140330) (-1.8133333,3.7990666) (-1.8066667,3.7841813) (-1.8000000,3.7693768) (-1.7933333,3.7546529) (-1.7866667,3.7400091) (-1.7800000,3.7254451) (-1.7733333,3.7109608) (-1.7666667,3.6965556) (-1.7600000,3.6822294) (-1.7533333,3.6679818) (-1.7466667,3.6538126) (-1.7400000,3.6397214) (-1.7333333,3.6257078) (-1.7266667,3.6117718) (-1.7200000,3.5979128) (-1.7133333,3.5841307) (-1.7066667,3.5704251) (-1.7000000,3.5567957) (-1.6933333,3.5432423) (-1.6866667,3.5297645) (-1.6800000,3.5163621) (-1.6733333,3.5030349) (-1.6666667,3.4897824) (-1.6600000,3.4766044) (-1.6533333,3.4635007) (-1.6466667,3.4504709) (-1.6400000,3.4375148) (-1.6333333,3.4246321) (-1.6266667,3.4118225) (-1.6200000,3.3990858) (-1.6133333,3.3864217) (-1.6066667,3.3738299) (-1.6000000,3.3613101) (-1.5933333,3.3488621) (-1.5866667,3.3364856) (-1.5800000,3.3241804) (-1.5733333,3.3119461) (-1.5666667,3.2997826) (-1.5600000,3.2876895) (-1.5533333,3.2756666) (-1.5466667,3.2637137) (-1.5400000,3.2518305) (-1.5333333,3.2400167) (-1.5266667,3.2282721) (-1.5200000,3.2165964) (-1.5133333,3.2049894) (-1.5066667,3.1934509) (-1.5000000,3.1819805) (-1.4933333,3.1705781) (-1.4866667,3.1592434) (-1.4800000,3.1479761) (-1.4733333,3.1367761) (-1.4666667,3.1256431) (-1.4600000,3.1145768) (-1.4533333,3.1035770) (-1.4466667,3.0926434) (-1.4400000,3.0817760) (-1.4333333,3.0709743) (-1.4266667,3.0602382) (-1.4200000,3.0495674) (-1.4133333,3.0389618) (-1.4066667,3.0284210) (-1.4000000,3.0179450) (-1.3933333,3.0075333) (-1.3866667,2.9971859) (-1.3800000,2.9869025) (-1.3733333,2.9766829) (-1.3666667,2.9665268) (-1.3600000,2.9564341) (-1.3533333,2.9464045) (-1.3466667,2.9364378) (-1.3400000,2.9265338) (-1.3333333,2.9166924) (-1.3266667,2.9069132) (-1.3200000,2.8971960) (-1.3133333,2.8875408) (-1.3066667,2.8779472) (-1.3000000,2.8684150) (-1.2933333,2.8589441) (-1.2866667,2.8495343) (-1.2800000,2.8401853) (-1.2733333,2.8308969) (-1.2666667,2.8216690) (-1.2600000,2.8125014) (-1.2533333,2.8033938) (-1.2466667,2.7943460) (-1.2400000,2.7853580) (-1.2333333,2.7764294) (-1.2266667,2.7675601) (-1.2200000,2.7587499) (-1.2133333,2.7499986) (-1.2066667,2.7413060) (-1.2000000,2.7326720) (-1.1933333,2.7240963) (-1.1866667,2.7155788) (-1.1800000,2.7071193) (-1.1733333,2.6987175) (-1.1666667,2.6903735) (-1.1600000,2.6820868) (-1.1533333,2.6738574) (-1.1466667,2.6656852) (-1.1400000,2.6575698) (-1.1333333,2.6495112) (-1.1266667,2.6415092) (-1.1200000,2.6335636) (-1.1133333,2.6256742) (-1.1066667,2.6178408) (-1.1000000,2.6100634) (-1.0933333,2.6023417) (-1.0866667,2.5946756) (-1.0800000,2.5870649) (-1.0733333,2.5795094) (-1.0666667,2.5720090) (-1.0600000,2.5645636) (-1.0533333,2.5571729) (-1.0466667,2.5498368) (-1.0400000,2.5425551) (-1.0333333,2.5353278) (-1.0266667,2.5281545) (-1.0200000,2.5210353) (-1.0133333,2.5139699) (-1.0066667,2.5069582) (-1.0000000,2.5000000) (-0.9933333,2.4930952) (-0.9866667,2.4862436) (-0.9800000,2.4794451) (-0.9733333,2.4726996) (-0.9666667,2.4660069) (-0.9600000,2.4593668) (-0.9533333,2.4527792) (-0.9466667,2.4462441) (-0.9400000,2.4397611) (-0.9333333,2.4333303) (-0.9266667,2.4269514) (-0.9200000,2.4206243) (-0.9133333,2.4143489) (-0.9066667,2.4081251) (-0.9000000,2.4019527) (-0.8933333,2.3958316) (-0.8866667,2.3897617) (-0.8800000,2.3837427) (-0.8733333,2.3777747) (-0.8666667,2.3718575) (-0.8600000,2.3659909) (-0.8533333,2.3601748) (-0.8466667,2.3544091) (-0.8400000,2.3486937) (-0.8333333,2.3430285) (-0.8266667,2.3374132) (-0.8200000,2.3318479) (-0.8133333,2.3263324) (-0.8066667,2.3208666) (-0.8000000,2.3154503) (-0.7933333,2.3100835) (-0.7866667,2.3047660) (-0.7800000,2.2994977) (-0.7733333,2.2942785) (-0.7666667,2.2891083) (-0.7600000,2.2839870) (-0.7533333,2.2789144) (-0.7466667,2.2738905) (-0.7400000,2.2689152) (-0.7333333,2.2639883) (-0.7266667,2.2591098) (-0.7200000,2.2542795) (-0.7133333,2.2494973) (-0.7066667,2.2447632) (-0.7000000,2.2400770) (-0.6933333,2.2354386) (-0.6866667,2.2308480) (-0.6800000,2.2263050) (-0.6733333,2.2218096) (-0.6666667,2.2173616) (-0.6600000,2.2129609) (-0.6533333,2.2086075) (-0.6466667,2.2043013) (-0.6400000,2.2000421) (-0.6333333,2.1958299) (-0.6266667,2.1916646) (-0.6200000,2.1875461) (-0.6133333,2.1834743) (-0.6066667,2.1794492) (-0.6000000,2.1754705) (-0.5933333,2.1715383) (-0.5866667,2.1676525) (-0.5800000,2.1638130) (-0.5733333,2.1600197) (-0.5666667,2.1562725) (-0.5600000,2.1525714) (-0.5533333,2.1489162) (-0.5466667,2.1453069) (-0.5400000,2.1417434) (-0.5333333,2.1382257) (-0.5266667,2.1347536) (-0.5200000,2.1313271) (-0.5133333,2.1279461) (-0.5066667,2.1246105) (-0.5000000,2.1213203) (-0.4933333,2.1180755) (-0.4866667,2.1148758) (-0.4800000,2.1117213) (-0.4733333,2.1086119) (-0.4666667,2.1055475) (-0.4600000,2.1025281) (-0.4533333,2.0995536) (-0.4466667,2.0966239) (-0.4400000,2.0937389) (-0.4333333,2.0908987) (-0.4266667,2.0881032) (-0.4200000,2.0853522) (-0.4133333,2.0826457) (-0.4066667,2.0799838) (-0.4000000,2.0773662) (-0.3933333,2.0747930) (-0.3866667,2.0722641) (-0.3800000,2.0697794) (-0.3733333,2.0673390) (-0.3666667,2.0649427) (-0.3600000,2.0625905) (-0.3533333,2.0602823) (-0.3466667,2.0580181) (-0.3400000,2.0557979) (-0.3333333,2.0536216) (-0.3266667,2.0514891) (-0.3200000,2.0494004) (-0.3133333,2.0473555) (-0.3066667,2.0453543) (-0.3000000,2.0433968) (-0.2933333,2.0414829) (-0.2866667,2.0396126) (-0.2800000,2.0377859) (-0.2733333,2.0360027) (-0.2666667,2.0342629) (-0.2600000,2.0325666) (-0.2533333,2.0309137) (-0.2466667,2.0293042) (-0.2400000,2.0277380) (-0.2333333,2.0262151) (-0.2266667,2.0247354) (-0.2200000,2.0232990) (-0.2133333,2.0219058) (-0.2066667,2.0205558) (-0.2000000,2.0192489) (-0.1933333,2.0179852) (-0.1866667,2.0167645) (-0.1800000,2.0155869) (-0.1733333,2.0144523) (-0.1666667,2.0133608) (-0.1600000,2.0123122) (-0.1533333,2.0113066) (-0.1466667,2.0103440) (-0.1400000,2.0094243) (-0.1333333,2.0085475) (-0.1266667,2.0077136) (-0.1200000,2.0069225) (-0.1133333,2.0061743) (-0.1066667,2.0054690) (-0.1000000,2.0048065) (-0.0933333,2.0041867) (-0.0866667,2.0036098) (-0.0800000,2.0030757) (-0.0733333,2.0025843) (-0.0666667,2.0021357) (-0.0600000,2.0017299) (-0.0533333,2.0013668) (-0.0466667,2.0010464) (-0.0400000,2.0007688) (-0.0333333,2.0005339) (-0.0266667,2.0003417) (-0.0200000,2.0001922) (-0.0133333,2.0000854) (-0.0066667,2.0000214) (0.0000000,2.0000000) (0.0066667,2.0000214) (0.0133333,2.0000854) (0.0200000,2.0001922) (0.0266667,2.0003417) (0.0333333,2.0005339) (0.0400000,2.0007688) (0.0466667,2.0010464) (0.0533333,2.0013668) (0.0600000,2.0017299) (0.0666667,2.0021357) (0.0733333,2.0025843) (0.0800000,2.0030757) (0.0866667,2.0036098) (0.0933333,2.0041867) (0.1000000,2.0048065) (0.1066667,2.0054690) (0.1133333,2.0061743) (0.1200000,2.0069225) (0.1266667,2.0077136) (0.1333333,2.0085475) (0.1400000,2.0094243) (0.1466667,2.0103440) (0.1533333,2.0113066) (0.1600000,2.0123122) (0.1666667,2.0133608) (0.1733333,2.0144523) (0.1800000,2.0155869) (0.1866667,2.0167645) (0.1933333,2.0179852) (0.2000000,2.0192489) (0.2066667,2.0205558) (0.2133333,2.0219058) (0.2200000,2.0232990) (0.2266667,2.0247354) (0.2333333,2.0262151) (0.2400000,2.0277380) (0.2466667,2.0293042) (0.2533333,2.0309137) (0.2600000,2.0325666) (0.2666667,2.0342629) (0.2733333,2.0360027) (0.2800000,2.0377859) (0.2866667,2.0396126) (0.2933333,2.0414829) (0.3000000,2.0433968) (0.3066667,2.0453543) (0.3133333,2.0473555) (0.3200000,2.0494004) (0.3266667,2.0514891) (0.3333333,2.0536216) (0.3400000,2.0557979) (0.3466667,2.0580181) (0.3533333,2.0602823) (0.3600000,2.0625905) (0.3666667,2.0649427) (0.3733333,2.0673390) (0.3800000,2.0697794) (0.3866667,2.0722641) (0.3933333,2.0747930) (0.4000000,2.0773662) (0.4066667,2.0799838) (0.4133333,2.0826457) (0.4200000,2.0853522) (0.4266667,2.0881032) (0.4333333,2.0908987) (0.4400000,2.0937389) (0.4466667,2.0966239) (0.4533333,2.0995536) (0.4600000,2.1025281) (0.4666667,2.1055475) (0.4733333,2.1086119) (0.4800000,2.1117213) (0.4866667,2.1148758) (0.4933333,2.1180755) (0.5000000,2.1213203) (0.5066667,2.1246105) (0.5133333,2.1279461) (0.5200000,2.1313271) (0.5266667,2.1347536) (0.5333333,2.1382257) (0.5400000,2.1417434) (0.5466667,2.1453069) (0.5533333,2.1489162) (0.5600000,2.1525714) (0.5666667,2.1562725) (0.5733333,2.1600197) (0.5800000,2.1638130) (0.5866667,2.1676525) (0.5933333,2.1715383) (0.6000000,2.1754705) (0.6066667,2.1794492) (0.6133333,2.1834743) (0.6200000,2.1875461) (0.6266667,2.1916646) (0.6333333,2.1958299) (0.6400000,2.2000421) (0.6466667,2.2043013) (0.6533333,2.2086075) (0.6600000,2.2129609) (0.6666667,2.2173616) (0.6733333,2.2218096) (0.6800000,2.2263050) (0.6866667,2.2308480) (0.6933333,2.2354386) (0.7000000,2.2400770) (0.7066667,2.2447632) (0.7133333,2.2494973) (0.7200000,2.2542795) (0.7266667,2.2591098) (0.7333333,2.2639883) (0.7400000,2.2689152) (0.7466667,2.2738905) (0.7533333,2.2789144) (0.7600000,2.2839870) (0.7666667,2.2891083) (0.7733333,2.2942785) (0.7800000,2.2994977) (0.7866667,2.3047660) (0.7933333,2.3100835) (0.8000000,2.3154503) (0.8066667,2.3208666) (0.8133333,2.3263324) (0.8200000,2.3318479) (0.8266667,2.3374132) (0.8333333,2.3430285) (0.8400000,2.3486937) (0.8466667,2.3544091) (0.8533333,2.3601748) (0.8600000,2.3659909) (0.8666667,2.3718575) (0.8733333,2.3777747) (0.8800000,2.3837427) (0.8866667,2.3897617) (0.8933333,2.3958316) (0.9000000,2.4019527) (0.9066667,2.4081251) (0.9133333,2.4143489) (0.9200000,2.4206243) (0.9266667,2.4269514) (0.9333333,2.4333303) (0.9400000,2.4397611) (0.9466667,2.4462441) (0.9533333,2.4527792) (0.9600000,2.4593668) (0.9666667,2.4660069) (0.9733333,2.4726996) (0.9800000,2.4794451) (0.9866667,2.4862436) (0.9933333,2.4930952) (1.0000000,2.5000000) (1.0066667,2.5069582) (1.0133333,2.5139699) (1.0200000,2.5210353) (1.0266667,2.5281545) (1.0333333,2.5353278) (1.0400000,2.5425551) (1.0466667,2.5498368) (1.0533333,2.5571729) (1.0600000,2.5645636) (1.0666667,2.5720090) (1.0733333,2.5795094) (1.0800000,2.5870649) (1.0866667,2.5946756) (1.0933333,2.6023417) (1.1000000,2.6100634) (1.1066667,2.6178408) (1.1133333,2.6256742) (1.1200000,2.6335636) (1.1266667,2.6415092) (1.1333333,2.6495112) (1.1400000,2.6575698) (1.1466667,2.6656852) (1.1533333,2.6738574) (1.1600000,2.6820868) (1.1666667,2.6903735) (1.1733333,2.6987175) (1.1800000,2.7071193) (1.1866667,2.7155788) (1.1933333,2.7240963) (1.2000000,2.7326720) (1.2066667,2.7413060) (1.2133333,2.7499986) (1.2200000,2.7587499) (1.2266667,2.7675601) (1.2333333,2.7764294) (1.2400000,2.7853580) (1.2466667,2.7943460) (1.2533333,2.8033938) (1.2600000,2.8125014) (1.2666667,2.8216690) (1.2733333,2.8308969) (1.2800000,2.8401853) (1.2866667,2.8495343) (1.2933333,2.8589441) (1.3000000,2.8684150) (1.3066667,2.8779472) (1.3133333,2.8875408) (1.3200000,2.8971960) (1.3266667,2.9069132) (1.3333333,2.9166924) (1.3400000,2.9265338) (1.3466667,2.9364378) (1.3533333,2.9464045) (1.3600000,2.9564341) (1.3666667,2.9665268) (1.3733333,2.9766829) (1.3800000,2.9869025) (1.3866667,2.9971859) (1.3933333,3.0075333) (1.4000000,3.0179450) (1.4066667,3.0284210) (1.4133333,3.0389618) (1.4200000,3.0495674) (1.4266667,3.0602382) (1.4333333,3.0709743) (1.4400000,3.0817760) (1.4466667,3.0926434) (1.4533333,3.1035770) (1.4600000,3.1145768) (1.4666667,3.1256431) (1.4733333,3.1367761) (1.4800000,3.1479761) (1.4866667,3.1592434) (1.4933333,3.1705781) (1.5000000,3.1819805) (1.5066667,3.1934509) (1.5133333,3.2049894) (1.5200000,3.2165964) (1.5266667,3.2282721) (1.5333333,3.2400167) (1.5400000,3.2518305) (1.5466667,3.2637137) (1.5533333,3.2756666) (1.5600000,3.2876895) (1.5666667,3.2997826) (1.5733333,3.3119461) (1.5800000,3.3241804) (1.5866667,3.3364856) (1.5933333,3.3488621) (1.6000000,3.3613101) (1.6066667,3.3738299) (1.6133333,3.3864217) (1.6200000,3.3990858) (1.6266667,3.4118225) (1.6333333,3.4246321) (1.6400000,3.4375148) (1.6466667,3.4504709) (1.6533333,3.4635007) (1.6600000,3.4766044) (1.6666667,3.4897824) (1.6733333,3.5030349) (1.6800000,3.5163621) (1.6866667,3.5297645) (1.6933333,3.5432423) (1.7000000,3.5567957) (1.7066667,3.5704251) (1.7133333,3.5841307) (1.7200000,3.5979128) (1.7266667,3.6117718) (1.7333333,3.6257078) (1.7400000,3.6397214) (1.7466667,3.6538126) (1.7533333,3.6679818) (1.7600000,3.6822294) (1.7666667,3.6965556) (1.7733333,3.7109608) (1.7800000,3.7254451) (1.7866667,3.7400091) (1.7933333,3.7546529) (1.8000000,3.7693768) (1.8066667,3.7841813) (1.8133333,3.7990666) (1.8200000,3.8140330) (1.8266667,3.8290808) (1.8333333,3.8442104) (1.8400000,3.8594221) (1.8466667,3.8747162) (1.8533333,3.8900930) (1.8600000,3.9055529) (1.8666667,3.9210962) (1.8733333,3.9367232) (1.8800000,3.9524343) (1.8866667,3.9682298) (1.8933333,3.9841100) (1.9000000,4.0000753) (1.9066667,4.0161260) (1.9133333,4.0322625) (1.9200000,4.0484851) (1.9266667,4.0647941) (1.9333333,4.0811899) (1.9400000,4.0976729) (1.9466667,4.1142434) (1.9533333,4.1309017) (1.9600000,4.1476482) (1.9666667,4.1644833) (1.9733333,4.1814074) (1.9800000,4.1984207) (1.9866667,4.2155237) (1.9933333,4.2327166) (2.0000000,4.2500000)};
\draw[orange!85!black,thick,dashed] plot coordinates {(-1.9333333,8.4839473) (-1.9266667,8.4237885) (-1.9200000,8.3640815) (-1.9133333,8.3048231) (-1.9066667,8.2460103) (-1.9000000,8.1876398) (-1.8933333,8.1297085) (-1.8866667,8.0722133) (-1.8800000,8.0151512) (-1.8733333,7.9585189) (-1.8666667,7.9023136) (-1.8600000,7.8465322) (-1.8533333,7.7911717) (-1.8466667,7.7362291) (-1.8400000,7.6817016) (-1.8333333,7.6275860) (-1.8266667,7.5738797) (-1.8200000,7.5205796) (-1.8133333,7.4676830) (-1.8066667,7.4151869) (-1.8000000,7.3630886) (-1.7933333,7.3113853) (-1.7866667,7.2600742) (-1.7800000,7.2091525) (-1.7733333,7.1586175) (-1.7666667,7.1084666) (-1.7600000,7.0586969) (-1.7533333,7.0093059) (-1.7466667,6.9602910) (-1.7400000,6.9116493) (-1.7333333,6.8633785) (-1.7266667,6.8154758) (-1.7200000,6.7679386) (-1.7133333,6.7207646) (-1.7066667,6.6739511) (-1.7000000,6.6274955) (-1.6933333,6.5813955) (-1.6866667,6.5356485) (-1.6800000,6.4902522) (-1.6733333,6.4452039) (-1.6666667,6.4005014) (-1.6600000,6.3561423) (-1.6533333,6.3121241) (-1.6466667,6.2684445) (-1.6400000,6.2251011) (-1.6333333,6.1820917) (-1.6266667,6.1394139) (-1.6200000,6.0970654) (-1.6133333,6.0550440) (-1.6066667,6.0133474) (-1.6000000,5.9719734) (-1.5933333,5.9309197) (-1.5866667,5.8901842) (-1.5800000,5.8497647) (-1.5733333,5.8096589) (-1.5666667,5.7698648) (-1.5600000,5.7303802) (-1.5533333,5.6912030) (-1.5466667,5.6523310) (-1.5400000,5.6137623) (-1.5333333,5.5754947) (-1.5266667,5.5375262) (-1.5200000,5.4998547) (-1.5133333,5.4624783) (-1.5066667,5.4253949) (-1.5000000,5.3886025) (-1.4933333,5.3520992) (-1.4866667,5.3158830) (-1.4800000,5.2799519) (-1.4733333,5.2443041) (-1.4666667,5.2089375) (-1.4600000,5.1738504) (-1.4533333,5.1390409) (-1.4466667,5.1045070) (-1.4400000,5.0702469) (-1.4333333,5.0362588) (-1.4266667,5.0025409) (-1.4200000,4.9690913) (-1.4133333,4.9359082) (-1.4066667,4.9029900) (-1.4000000,4.8703347) (-1.3933333,4.8379407) (-1.3866667,4.8058063) (-1.3800000,4.7739296) (-1.3733333,4.7423090) (-1.3666667,4.7109428) (-1.3600000,4.6798293) (-1.3533333,4.6489668) (-1.3466667,4.6183537) (-1.3400000,4.5879884) (-1.3333333,4.5578691) (-1.3266667,4.5279944) (-1.3200000,4.4983626) (-1.3133333,4.4689720) (-1.3066667,4.4398212) (-1.3000000,4.4109085) (-1.2933333,4.3822325) (-1.2866667,4.3537915) (-1.2800000,4.3255841) (-1.2733333,4.2976087) (-1.2666667,4.2698639) (-1.2600000,4.2423481) (-1.2533333,4.2150598) (-1.2466667,4.1879977) (-1.2400000,4.1611602) (-1.2333333,4.1345460) (-1.2266667,4.1081535) (-1.2200000,4.0819814) (-1.2133333,4.0560283) (-1.2066667,4.0302927) (-1.2000000,4.0047733) (-1.1933333,3.9794688) (-1.1866667,3.9543777) (-1.1800000,3.9294988) (-1.1733333,3.9048306) (-1.1666667,3.8803719) (-1.1600000,3.8561214) (-1.1533333,3.8320777) (-1.1466667,3.8082396) (-1.1400000,3.7846057) (-1.1333333,3.7611749) (-1.1266667,3.7379458) (-1.1200000,3.7149172) (-1.1133333,3.6920880) (-1.1066667,3.6694567) (-1.1000000,3.6470223) (-1.0933333,3.6247836) (-1.0866667,3.6027393) (-1.0800000,3.5808882) (-1.0733333,3.5592292) (-1.0666667,3.5377612) (-1.0600000,3.5164829) (-1.0533333,3.4953933) (-1.0466667,3.4744912) (-1.0400000,3.4537754) (-1.0333333,3.4332449) (-1.0266667,3.4128986) (-1.0200000,3.3927354) (-1.0133333,3.3727541) (-1.0066667,3.3529538) (-1.0000000,3.3333333) (-0.9933333,3.3138917) (-0.9866667,3.2946278) (-0.9800000,3.2755406) (-0.9733333,3.2566292) (-0.9666667,3.2378924) (-0.9600000,3.2193294) (-0.9533333,3.2009390) (-0.9466667,3.1827203) (-0.9400000,3.1646724) (-0.9333333,3.1467942) (-0.9266667,3.1290848) (-0.9200000,3.1115433) (-0.9133333,3.0941687) (-0.9066667,3.0769600) (-0.9000000,3.0599164) (-0.8933333,3.0430370) (-0.8866667,3.0263208) (-0.8800000,3.0097669) (-0.8733333,2.9933745) (-0.8666667,2.9771427) (-0.8600000,2.9610705) (-0.8533333,2.9451572) (-0.8466667,2.9294019) (-0.8400000,2.9138037) (-0.8333333,2.8983619) (-0.8266667,2.8830755) (-0.8200000,2.8679437) (-0.8133333,2.8529658) (-0.8066667,2.8381409) (-0.8000000,2.8234683) (-0.7933333,2.8089472) (-0.7866667,2.7945767) (-0.7800000,2.7803561) (-0.7733333,2.7662847) (-0.7666667,2.7523617) (-0.7600000,2.7385863) (-0.7533333,2.7249578) (-0.7466667,2.7114755) (-0.7400000,2.6981386) (-0.7333333,2.6849465) (-0.7266667,2.6718984) (-0.7200000,2.6589936) (-0.7133333,2.6462314) (-0.7066667,2.6336113) (-0.7000000,2.6211323) (-0.6933333,2.6087940) (-0.6866667,2.5965957) (-0.6800000,2.5845366) (-0.6733333,2.5726161) (-0.6666667,2.5608337) (-0.6600000,2.5491886) (-0.6533333,2.5376803) (-0.6466667,2.5263081) (-0.6400000,2.5150714) (-0.6333333,2.5039696) (-0.6266667,2.4930022) (-0.6200000,2.4821685) (-0.6133333,2.4714679) (-0.6066667,2.4608999) (-0.6000000,2.4504639) (-0.5933333,2.4401594) (-0.5866667,2.4299857) (-0.5800000,2.4199424) (-0.5733333,2.4100290) (-0.5666667,2.4002448) (-0.5600000,2.3905893) (-0.5533333,2.3810621) (-0.5466667,2.3716626) (-0.5400000,2.3623904) (-0.5333333,2.3532448) (-0.5266667,2.3442255) (-0.5200000,2.3353320) (-0.5133333,2.3265637) (-0.5066667,2.3179202) (-0.5000000,2.3094011) (-0.4933333,2.3010058) (-0.4866667,2.2927340) (-0.4800000,2.2845852) (-0.4733333,2.2765589) (-0.4666667,2.2686547) (-0.4600000,2.2608722) (-0.4533333,2.2532110) (-0.4466667,2.2456707) (-0.4400000,2.2382509) (-0.4333333,2.2309511) (-0.4266667,2.2237709) (-0.4200000,2.2167101) (-0.4133333,2.2097682) (-0.4066667,2.2029448) (-0.4000000,2.1962396) (-0.3933333,2.1896522) (-0.3866667,2.1831822) (-0.3800000,2.1768294) (-0.3733333,2.1705933) (-0.3666667,2.1644737) (-0.3600000,2.1584702) (-0.3533333,2.1525825) (-0.3466667,2.1468102) (-0.3400000,2.1411531) (-0.3333333,2.1356108) (-0.3266667,2.1301832) (-0.3200000,2.1248697) (-0.3133333,2.1196703) (-0.3066667,2.1145846) (-0.3000000,2.1096123) (-0.2933333,2.1047531) (-0.2866667,2.1000069) (-0.2800000,2.0953733) (-0.2733333,2.0908521) (-0.2666667,2.0864431) (-0.2600000,2.0821460) (-0.2533333,2.0779606) (-0.2466667,2.0738866) (-0.2400000,2.0699240) (-0.2333333,2.0660723) (-0.2266667,2.0623315) (-0.2200000,2.0587013) (-0.2133333,2.0551815) (-0.2066667,2.0517720) (-0.2000000,2.0484725) (-0.1933333,2.0452829) (-0.1866667,2.0422031) (-0.1800000,2.0392327) (-0.1733333,2.0363718) (-0.1666667,2.0336201) (-0.1600000,2.0309775) (-0.1533333,2.0284439) (-0.1466667,2.0260190) (-0.1400000,2.0237029) (-0.1333333,2.0214953) (-0.1266667,2.0193961) (-0.1200000,2.0174053) (-0.1133333,2.0155226) (-0.1066667,2.0137481) (-0.1000000,2.0120816) (-0.0933333,2.0105231) (-0.0866667,2.0090724) (-0.0800000,2.0077294) (-0.0733333,2.0064942) (-0.0666667,2.0053666) (-0.0600000,2.0043466) (-0.0533333,2.0034341) (-0.0466667,2.0026290) (-0.0400000,2.0019314) (-0.0333333,2.0013412) (-0.0266667,2.0008583) (-0.0200000,2.0004828) (-0.0133333,2.0002146) (-0.0066667,2.0000536) (0.0000000,2.0000000) (0.0066667,2.0000536) (0.0133333,2.0002146) (0.0200000,2.0004828) (0.0266667,2.0008583) (0.0333333,2.0013412) (0.0400000,2.0019314) (0.0466667,2.0026290) (0.0533333,2.0034341) (0.0600000,2.0043466) (0.0666667,2.0053666) (0.0733333,2.0064942) (0.0800000,2.0077294) (0.0866667,2.0090724) (0.0933333,2.0105231) (0.1000000,2.0120816) (0.1066667,2.0137481) (0.1133333,2.0155226) (0.1200000,2.0174053) (0.1266667,2.0193961) (0.1333333,2.0214953) (0.1400000,2.0237029) (0.1466667,2.0260190) (0.1533333,2.0284439) (0.1600000,2.0309775) (0.1666667,2.0336201) (0.1733333,2.0363718) (0.1800000,2.0392327) (0.1866667,2.0422031) (0.1933333,2.0452829) (0.2000000,2.0484725) (0.2066667,2.0517720) (0.2133333,2.0551815) (0.2200000,2.0587013) (0.2266667,2.0623315) (0.2333333,2.0660723) (0.2400000,2.0699240) (0.2466667,2.0738866) (0.2533333,2.0779606) (0.2600000,2.0821460) (0.2666667,2.0864431) (0.2733333,2.0908521) (0.2800000,2.0953733) (0.2866667,2.1000069) (0.2933333,2.1047531) (0.3000000,2.1096123) (0.3066667,2.1145846) (0.3133333,2.1196703) (0.3200000,2.1248697) (0.3266667,2.1301832) (0.3333333,2.1356108) (0.3400000,2.1411531) (0.3466667,2.1468102) (0.3533333,2.1525825) (0.3600000,2.1584702) (0.3666667,2.1644737) (0.3733333,2.1705933) (0.3800000,2.1768294) (0.3866667,2.1831822) (0.3933333,2.1896522) (0.4000000,2.1962396) (0.4066667,2.2029448) (0.4133333,2.2097682) (0.4200000,2.2167101) (0.4266667,2.2237709) (0.4333333,2.2309511) (0.4400000,2.2382509) (0.4466667,2.2456707) (0.4533333,2.2532110) (0.4600000,2.2608722) (0.4666667,2.2686547) (0.4733333,2.2765589) (0.4800000,2.2845852) (0.4866667,2.2927340) (0.4933333,2.3010058) (0.5000000,2.3094011) (0.5066667,2.3179202) (0.5133333,2.3265637) (0.5200000,2.3353320) (0.5266667,2.3442255) (0.5333333,2.3532448) (0.5400000,2.3623904) (0.5466667,2.3716626) (0.5533333,2.3810621) (0.5600000,2.3905893) (0.5666667,2.4002448) (0.5733333,2.4100290) (0.5800000,2.4199424) (0.5866667,2.4299857) (0.5933333,2.4401594) (0.6000000,2.4504639) (0.6066667,2.4608999) (0.6133333,2.4714679) (0.6200000,2.4821685) (0.6266667,2.4930022) (0.6333333,2.5039696) (0.6400000,2.5150714) (0.6466667,2.5263081) (0.6533333,2.5376803) (0.6600000,2.5491886) (0.6666667,2.5608337) (0.6733333,2.5726161) (0.6800000,2.5845366) (0.6866667,2.5965957) (0.6933333,2.6087940) (0.7000000,2.6211323) (0.7066667,2.6336113) (0.7133333,2.6462314) (0.7200000,2.6589936) (0.7266667,2.6718984) (0.7333333,2.6849465) (0.7400000,2.6981386) (0.7466667,2.7114755) (0.7533333,2.7249578) (0.7600000,2.7385863) (0.7666667,2.7523617) (0.7733333,2.7662847) (0.7800000,2.7803561) (0.7866667,2.7945767) (0.7933333,2.8089472) (0.8000000,2.8234683) (0.8066667,2.8381409) (0.8133333,2.8529658) (0.8200000,2.8679437) (0.8266667,2.8830755) (0.8333333,2.8983619) (0.8400000,2.9138037) (0.8466667,2.9294019) (0.8533333,2.9451572) (0.8600000,2.9610705) (0.8666667,2.9771427) (0.8733333,2.9933745) (0.8800000,3.0097669) (0.8866667,3.0263208) (0.8933333,3.0430370) (0.9000000,3.0599164) (0.9066667,3.0769600) (0.9133333,3.0941687) (0.9200000,3.1115433) (0.9266667,3.1290848) (0.9333333,3.1467942) (0.9400000,3.1646724) (0.9466667,3.1827203) (0.9533333,3.2009390) (0.9600000,3.2193294) (0.9666667,3.2378924) (0.9733333,3.2566292) (0.9800000,3.2755406) (0.9866667,3.2946278) (0.9933333,3.3138917) (1.0000000,3.3333333) (1.0066667,3.3529538) (1.0133333,3.3727541) (1.0200000,3.3927354) (1.0266667,3.4128986) (1.0333333,3.4332449) (1.0400000,3.4537754) (1.0466667,3.4744912) (1.0533333,3.4953933) (1.0600000,3.5164829) (1.0666667,3.5377612) (1.0733333,3.5592292) (1.0800000,3.5808882) (1.0866667,3.6027393) (1.0933333,3.6247836) (1.1000000,3.6470223) (1.1066667,3.6694567) (1.1133333,3.6920880) (1.1200000,3.7149172) (1.1266667,3.7379458) (1.1333333,3.7611749) (1.1400000,3.7846057) (1.1466667,3.8082396) (1.1533333,3.8320777) (1.1600000,3.8561214) (1.1666667,3.8803719) (1.1733333,3.9048306) (1.1800000,3.9294988) (1.1866667,3.9543777) (1.1933333,3.9794688) (1.2000000,4.0047733) (1.2066667,4.0302927) (1.2133333,4.0560283) (1.2200000,4.0819814) (1.2266667,4.1081535) (1.2333333,4.1345460) (1.2400000,4.1611602) (1.2466667,4.1879977) (1.2533333,4.2150598) (1.2600000,4.2423481) (1.2666667,4.2698639) (1.2733333,4.2976087) (1.2800000,4.3255841) (1.2866667,4.3537915) (1.2933333,4.3822325) (1.3000000,4.4109085) (1.3066667,4.4398212) (1.3133333,4.4689720) (1.3200000,4.4983626) (1.3266667,4.5279944) (1.3333333,4.5578691) (1.3400000,4.5879884) (1.3466667,4.6183537) (1.3533333,4.6489668) (1.3600000,4.6798293) (1.3666667,4.7109428) (1.3733333,4.7423090) (1.3800000,4.7739296) (1.3866667,4.8058063) (1.3933333,4.8379407) (1.4000000,4.8703347) (1.4066667,4.9029900) (1.4133333,4.9359082) (1.4200000,4.9690913) (1.4266667,5.0025409) (1.4333333,5.0362588) (1.4400000,5.0702469) (1.4466667,5.1045070) (1.4533333,5.1390409) (1.4600000,5.1738504) (1.4666667,5.2089375) (1.4733333,5.2443041) (1.4800000,5.2799519) (1.4866667,5.3158830) (1.4933333,5.3520992) (1.5000000,5.3886025) (1.5066667,5.4253949) (1.5133333,5.4624783) (1.5200000,5.4998547) (1.5266667,5.5375262) (1.5333333,5.5754947) (1.5400000,5.6137623) (1.5466667,5.6523310) (1.5533333,5.6912030) (1.5600000,5.7303802) (1.5666667,5.7698648) (1.5733333,5.8096589) (1.5800000,5.8497647) (1.5866667,5.8901842) (1.5933333,5.9309197) (1.6000000,5.9719734) (1.6066667,6.0133474) (1.6133333,6.0550440) (1.6200000,6.0970654) (1.6266667,6.1394139) (1.6333333,6.1820917) (1.6400000,6.2251011) (1.6466667,6.2684445) (1.6533333,6.3121241) (1.6600000,6.3561423) (1.6666667,6.4005014) (1.6733333,6.4452039) (1.6800000,6.4902522) (1.6866667,6.5356485) (1.6933333,6.5813955) (1.7000000,6.6274955) (1.7066667,6.6739511) (1.7133333,6.7207646) (1.7200000,6.7679386) (1.7266667,6.8154758) (1.7333333,6.8633785) (1.7400000,6.9116493) (1.7466667,6.9602910) (1.7533333,7.0093059) (1.7600000,7.0586969) (1.7666667,7.1084666) (1.7733333,7.1586175) (1.7800000,7.2091525) (1.7866667,7.2600742) (1.7933333,7.3113853) (1.8000000,7.3630886) (1.8066667,7.4151869) (1.8133333,7.4676830) (1.8200000,7.5205796) (1.8266667,7.5738797) (1.8333333,7.6275860) (1.8400000,7.6817016) (1.8466667,7.7362291) (1.8533333,7.7911717) (1.8600000,7.8465322) (1.8666667,7.9023136) (1.8733333,7.9585189) (1.8800000,8.0151512) (1.8866667,8.0722133) (1.8933333,8.1297085) (1.9000000,8.1876398) (1.9066667,8.2460103) (1.9133333,8.3048231) (1.9200000,8.3640815) (1.9266667,8.4237885) (1.9333333,8.4839473)};
\draw[violet!75!black,densely dotted] (-1,2.5000000)--(1,2.5000000);
\fill[violet!75!black] (-1,2.5000000) circle(2.3pt);
\fill[violet!75!black] (1,2.5000000) circle(2.3pt);
\end{scope}
\end{scope}
\node[anchor=west,align=left,text width=4.35cm] at (8.9,4.4) {蓝实线：$a=2$\\[8pt]橙虚线：$a=3$\\[14pt]左右同高，关于 $y$ 轴对称。\\[8pt]两图最低点都是 $(0,2)$。};
\node[text width=12.8cm,align=center] at (7.1,.65) {$D=\mathbb R$；$f(x)\ge2$。把底数换成 $1/a$，整条曲线不变。};
\end{tikzpicture}\end{center}

\textbf{先看图：} 图中 $a=2$ 与 $a=3$ 的曲线都左右同高，最低点均为 $(0,2)$。把底数换为其倒数，只是交换 $a^x$ 与 $a^{-x}$，曲线完全不变，所以也覆盖 $0<a<1$ 的情况。

\textbf{再说明一般情形：}
\[
f(-x)=a^{-x}+a^x=f(x),\qquad a^x+a^{-x}\ge2.
\]
偶函数的图像关于 $y$ 轴对称；取等号时 $a^x=1$，即 $x=0$。
\end{explanationbox}
\clearpage
\begin{explanationbox}
\textbf{\textcolor{CExplanation}{绝对值把右半图复制到左边}}
% Formula-based samples; see F083 checks/build_tikz_diagrams.py.
\begin{center}\begin{tikzpicture}[font=\small,>=stealth,line join=round]
\draw[gray!40,rounded corners=3pt] (0,0) rectangle (14.2,8.4);
\node at (7.1,7.75) {偶2：$y=a^{|x|}$};
\begin{scope}[shift={(4.400000,1.650000)},x=1.360000cm,y=0.750000cm]
\draw[->,gray] (-2.5,0)--(2.5,0) node[right] {$x$};
\draw[->,gray] (0,0)--(0,6) node[above] {$y$};
\draw[gray] (-1,-.035/0.750000)--(-1,.035/0.750000);
\node[anchor=north] at (-1,-.07/0.750000) {$-1$};
\draw[gray] (1,-.035/0.750000)--(1,.035/0.750000);
\node[anchor=north] at (1,-.07/0.750000) {$1$};
\begin{scope}
\clip (-2.5,0) rectangle (2.5,6);
\draw[blue,thick,solid] plot coordinates {(-2.5000000,5.6568542) (-2.4916667,5.6242732) (-2.4833333,5.5918797) (-2.4750000,5.5596729) (-2.4666667,5.5276515) (-2.4583333,5.4958146) (-2.4500000,5.4641610) (-2.4416667,5.4326898) (-2.4333333,5.4013998) (-2.4250000,5.3702900) (-2.4166667,5.3393594) (-2.4083333,5.3086070) (-2.4000000,5.2780316) (-2.3916667,5.2476324) (-2.3833333,5.2174083) (-2.3750000,5.1873582) (-2.3666667,5.1574812) (-2.3583333,5.1277763) (-2.3500000,5.0982425) (-2.3416667,5.0688788) (-2.3333333,5.0396842) (-2.3250000,5.0106578) (-2.3166667,4.9817985) (-2.3083333,4.9531054) (-2.3000000,4.9245777) (-2.2916667,4.8962142) (-2.2833333,4.8680141) (-2.2750000,4.8399764) (-2.2666667,4.8121001) (-2.2583333,4.7843845) (-2.2500000,4.7568285) (-2.2416667,4.7294311) (-2.2333333,4.7021916) (-2.2250000,4.6751090) (-2.2166667,4.6481823) (-2.2083333,4.6214108) (-2.2000000,4.5947934) (-2.1916667,4.5683294) (-2.1833333,4.5420177) (-2.1750000,4.5158576) (-2.1666667,4.4898482) (-2.1583333,4.4639886) (-2.1500000,4.4382779) (-2.1416667,4.4127153) (-2.1333333,4.3872999) (-2.1250000,4.3620309) (-2.1166667,4.3369075) (-2.1083333,4.3119287) (-2.1000000,4.2870939) (-2.0916667,4.2624020) (-2.0833333,4.2378524) (-2.0750000,4.2134441) (-2.0666667,4.1891765) (-2.0583333,4.1650486) (-2.0500000,4.1410597) (-2.0416667,4.1172089) (-2.0333333,4.0934956) (-2.0250000,4.0699188) (-2.0166667,4.0464778) (-2.0083333,4.0231718) (-2.0000000,4.0000000) (-1.9916667,3.9769617) (-1.9833333,3.9540561) (-1.9750000,3.9312824) (-1.9666667,3.9086399) (-1.9583333,3.8861278) (-1.9500000,3.8637453) (-1.9416667,3.8414918) (-1.9333333,3.8193664) (-1.9250000,3.7973685) (-1.9166667,3.7754973) (-1.9083333,3.7537520) (-1.9000000,3.7321320) (-1.8916667,3.7106365) (-1.8833333,3.6892648) (-1.8750000,3.6680162) (-1.8666667,3.6468900) (-1.8583333,3.6258854) (-1.8500000,3.6050019) (-1.8416667,3.5842386) (-1.8333333,3.5635949) (-1.8250000,3.5430701) (-1.8166667,3.5226635) (-1.8083333,3.5023744) (-1.8000000,3.4822023) (-1.7916667,3.4621462) (-1.7833333,3.4422057) (-1.7750000,3.4223801) (-1.7666667,3.4026686) (-1.7583333,3.3830707) (-1.7500000,3.3635857) (-1.7416667,3.3442128) (-1.7333333,3.3249516) (-1.7250000,3.3058013) (-1.7166667,3.2867613) (-1.7083333,3.2678309) (-1.7000000,3.2490096) (-1.6916667,3.2302967) (-1.6833333,3.2116915) (-1.6750000,3.1931935) (-1.6666667,3.1748021) (-1.6583333,3.1565166) (-1.6500000,3.1383364) (-1.6416667,3.1202609) (-1.6333333,3.1022895) (-1.6250000,3.0844217) (-1.6166667,3.0666567) (-1.6083333,3.0489940) (-1.6000000,3.0314331) (-1.5916667,3.0139734) (-1.5833333,2.9966142) (-1.5750000,2.9793549) (-1.5666667,2.9621951) (-1.5583333,2.9451341) (-1.5500000,2.9281714) (-1.5416667,2.9113064) (-1.5333333,2.8945385) (-1.5250000,2.8778672) (-1.5166667,2.8612919) (-1.5083333,2.8448120) (-1.5000000,2.8284271) (-1.4916667,2.8121366) (-1.4833333,2.7959399) (-1.4750000,2.7798364) (-1.4666667,2.7638258) (-1.4583333,2.7479073) (-1.4500000,2.7320805) (-1.4416667,2.7163449) (-1.4333333,2.7006999) (-1.4250000,2.6851450) (-1.4166667,2.6696797) (-1.4083333,2.6543035) (-1.4000000,2.6390158) (-1.3916667,2.6238162) (-1.3833333,2.6087041) (-1.3750000,2.5936791) (-1.3666667,2.5787406) (-1.3583333,2.5638882) (-1.3500000,2.5491213) (-1.3416667,2.5344394) (-1.3333333,2.5198421) (-1.3250000,2.5053289) (-1.3166667,2.4908992) (-1.3083333,2.4765527) (-1.3000000,2.4622888) (-1.2916667,2.4481071) (-1.2833333,2.4340070) (-1.2750000,2.4199882) (-1.2666667,2.4060501) (-1.2583333,2.3921922) (-1.2500000,2.3784142) (-1.2416667,2.3647156) (-1.2333333,2.3510958) (-1.2250000,2.3375545) (-1.2166667,2.3240912) (-1.2083333,2.3107054) (-1.2000000,2.2973967) (-1.1916667,2.2841647) (-1.1833333,2.2710089) (-1.1750000,2.2579288) (-1.1666667,2.2449241) (-1.1583333,2.2319943) (-1.1500000,2.2191389) (-1.1416667,2.2063576) (-1.1333333,2.1936500) (-1.1250000,2.1810155) (-1.1166667,2.1684537) (-1.1083333,2.1559644) (-1.1000000,2.1435469) (-1.0916667,2.1312010) (-1.0833333,2.1189262) (-1.0750000,2.1067221) (-1.0666667,2.0945882) (-1.0583333,2.0825243) (-1.0500000,2.0705298) (-1.0416667,2.0586045) (-1.0333333,2.0467478) (-1.0250000,2.0349594) (-1.0166667,2.0232389) (-1.0083333,2.0115859) (-1.0000000,2.0000000) (-0.9916667,1.9884808) (-0.9833333,1.9770280) (-0.9750000,1.9656412) (-0.9666667,1.9543199) (-0.9583333,1.9430639) (-0.9500000,1.9318727) (-0.9416667,1.9207459) (-0.9333333,1.9096832) (-0.9250000,1.8986842) (-0.9166667,1.8877486) (-0.9083333,1.8768760) (-0.9000000,1.8660660) (-0.8916667,1.8553182) (-0.8833333,1.8446324) (-0.8750000,1.8340081) (-0.8666667,1.8234450) (-0.8583333,1.8129427) (-0.8500000,1.8025009) (-0.8416667,1.7921193) (-0.8333333,1.7817974) (-0.8250000,1.7715350) (-0.8166667,1.7613317) (-0.8083333,1.7511872) (-0.8000000,1.7411011) (-0.7916667,1.7310731) (-0.7833333,1.7211029) (-0.7750000,1.7111901) (-0.7666667,1.7013343) (-0.7583333,1.6915354) (-0.7500000,1.6817928) (-0.7416667,1.6721064) (-0.7333333,1.6624758) (-0.7250000,1.6529006) (-0.7166667,1.6433806) (-0.7083333,1.6339155) (-0.7000000,1.6245048) (-0.6916667,1.6151483) (-0.6833333,1.6058458) (-0.6750000,1.5965968) (-0.6666667,1.5874011) (-0.6583333,1.5782583) (-0.6500000,1.5691682) (-0.6416667,1.5601305) (-0.6333333,1.5511448) (-0.6250000,1.5422108) (-0.6166667,1.5333283) (-0.6083333,1.5244970) (-0.6000000,1.5157166) (-0.5916667,1.5069867) (-0.5833333,1.4983071) (-0.5750000,1.4896775) (-0.5666667,1.4810976) (-0.5583333,1.4725671) (-0.5500000,1.4640857) (-0.5416667,1.4556532) (-0.5333333,1.4472692) (-0.5250000,1.4389336) (-0.5166667,1.4306459) (-0.5083333,1.4224060) (-0.5000000,1.4142136) (-0.4916667,1.4060683) (-0.4833333,1.3979699) (-0.4750000,1.3899182) (-0.4666667,1.3819129) (-0.4583333,1.3739536) (-0.4500000,1.3660403) (-0.4416667,1.3581724) (-0.4333333,1.3503499) (-0.4250000,1.3425725) (-0.4166667,1.3348399) (-0.4083333,1.3271517) (-0.4000000,1.3195079) (-0.3916667,1.3119081) (-0.3833333,1.3043521) (-0.3750000,1.2968396) (-0.3666667,1.2893703) (-0.3583333,1.2819441) (-0.3500000,1.2745606) (-0.3416667,1.2672197) (-0.3333333,1.2599210) (-0.3250000,1.2526644) (-0.3166667,1.2454496) (-0.3083333,1.2382764) (-0.3000000,1.2311444) (-0.2916667,1.2240535) (-0.2833333,1.2170035) (-0.2750000,1.2099941) (-0.2666667,1.2030250) (-0.2583333,1.1960961) (-0.2500000,1.1892071) (-0.2416667,1.1823578) (-0.2333333,1.1755479) (-0.2250000,1.1687772) (-0.2166667,1.1620456) (-0.2083333,1.1553527) (-0.2000000,1.1486984) (-0.1916667,1.1420823) (-0.1833333,1.1355044) (-0.1750000,1.1289644) (-0.1666667,1.1224620) (-0.1583333,1.1159971) (-0.1500000,1.1095695) (-0.1416667,1.1031788) (-0.1333333,1.0968250) (-0.1250000,1.0905077) (-0.1166667,1.0842269) (-0.1083333,1.0779822) (-0.1000000,1.0717735) (-0.0916667,1.0656005) (-0.0833333,1.0594631) (-0.0750000,1.0533610) (-0.0666667,1.0472941) (-0.0583333,1.0412622) (-0.0500000,1.0352649) (-0.0416667,1.0293022) (-0.0333333,1.0233739) (-0.0250000,1.0174797) (-0.0166667,1.0116194) (-0.0083333,1.0057929) (0.0000000,1.0000000) (0.0083333,1.0057929) (0.0166667,1.0116194) (0.0250000,1.0174797) (0.0333333,1.0233739) (0.0416667,1.0293022) (0.0500000,1.0352649) (0.0583333,1.0412622) (0.0666667,1.0472941) (0.0750000,1.0533610) (0.0833333,1.0594631) (0.0916667,1.0656005) (0.1000000,1.0717735) (0.1083333,1.0779822) (0.1166667,1.0842269) (0.1250000,1.0905077) (0.1333333,1.0968250) (0.1416667,1.1031788) (0.1500000,1.1095695) (0.1583333,1.1159971) (0.1666667,1.1224620) (0.1750000,1.1289644) (0.1833333,1.1355044) (0.1916667,1.1420823) (0.2000000,1.1486984) (0.2083333,1.1553527) (0.2166667,1.1620456) (0.2250000,1.1687772) (0.2333333,1.1755479) (0.2416667,1.1823578) (0.2500000,1.1892071) (0.2583333,1.1960961) (0.2666667,1.2030250) (0.2750000,1.2099941) (0.2833333,1.2170035) (0.2916667,1.2240535) (0.3000000,1.2311444) (0.3083333,1.2382764) (0.3166667,1.2454496) (0.3250000,1.2526644) (0.3333333,1.2599210) (0.3416667,1.2672197) (0.3500000,1.2745606) (0.3583333,1.2819441) (0.3666667,1.2893703) (0.3750000,1.2968396) (0.3833333,1.3043521) (0.3916667,1.3119081) (0.4000000,1.3195079) (0.4083333,1.3271517) (0.4166667,1.3348399) (0.4250000,1.3425725) (0.4333333,1.3503499) (0.4416667,1.3581724) (0.4500000,1.3660403) (0.4583333,1.3739536) (0.4666667,1.3819129) (0.4750000,1.3899182) (0.4833333,1.3979699) (0.4916667,1.4060683) (0.5000000,1.4142136) (0.5083333,1.4224060) (0.5166667,1.4306459) (0.5250000,1.4389336) (0.5333333,1.4472692) (0.5416667,1.4556532) (0.5500000,1.4640857) (0.5583333,1.4725671) (0.5666667,1.4810976) (0.5750000,1.4896775) (0.5833333,1.4983071) (0.5916667,1.5069867) (0.6000000,1.5157166) (0.6083333,1.5244970) (0.6166667,1.5333283) (0.6250000,1.5422108) (0.6333333,1.5511448) (0.6416667,1.5601305) (0.6500000,1.5691682) (0.6583333,1.5782583) (0.6666667,1.5874011) (0.6750000,1.5965968) (0.6833333,1.6058458) (0.6916667,1.6151483) (0.7000000,1.6245048) (0.7083333,1.6339155) (0.7166667,1.6433806) (0.7250000,1.6529006) (0.7333333,1.6624758) (0.7416667,1.6721064) (0.7500000,1.6817928) (0.7583333,1.6915354) (0.7666667,1.7013343) (0.7750000,1.7111901) (0.7833333,1.7211029) (0.7916667,1.7310731) (0.8000000,1.7411011) (0.8083333,1.7511872) (0.8166667,1.7613317) (0.8250000,1.7715350) (0.8333333,1.7817974) (0.8416667,1.7921193) (0.8500000,1.8025009) (0.8583333,1.8129427) (0.8666667,1.8234450) (0.8750000,1.8340081) (0.8833333,1.8446324) (0.8916667,1.8553182) (0.9000000,1.8660660) (0.9083333,1.8768760) (0.9166667,1.8877486) (0.9250000,1.8986842) (0.9333333,1.9096832) (0.9416667,1.9207459) (0.9500000,1.9318727) (0.9583333,1.9430639) (0.9666667,1.9543199) (0.9750000,1.9656412) (0.9833333,1.9770280) (0.9916667,1.9884808) (1.0000000,2.0000000) (1.0083333,2.0115859) (1.0166667,2.0232389) (1.0250000,2.0349594) (1.0333333,2.0467478) (1.0416667,2.0586045) (1.0500000,2.0705298) (1.0583333,2.0825243) (1.0666667,2.0945882) (1.0750000,2.1067221) (1.0833333,2.1189262) (1.0916667,2.1312010) (1.1000000,2.1435469) (1.1083333,2.1559644) (1.1166667,2.1684537) (1.1250000,2.1810155) (1.1333333,2.1936500) (1.1416667,2.2063576) (1.1500000,2.2191389) (1.1583333,2.2319943) (1.1666667,2.2449241) (1.1750000,2.2579288) (1.1833333,2.2710089) (1.1916667,2.2841647) (1.2000000,2.2973967) (1.2083333,2.3107054) (1.2166667,2.3240912) (1.2250000,2.3375545) (1.2333333,2.3510958) (1.2416667,2.3647156) (1.2500000,2.3784142) (1.2583333,2.3921922) (1.2666667,2.4060501) (1.2750000,2.4199882) (1.2833333,2.4340070) (1.2916667,2.4481071) (1.3000000,2.4622888) (1.3083333,2.4765527) (1.3166667,2.4908992) (1.3250000,2.5053289) (1.3333333,2.5198421) (1.3416667,2.5344394) (1.3500000,2.5491213) (1.3583333,2.5638882) (1.3666667,2.5787406) (1.3750000,2.5936791) (1.3833333,2.6087041) (1.3916667,2.6238162) (1.4000000,2.6390158) (1.4083333,2.6543035) (1.4166667,2.6696797) (1.4250000,2.6851450) (1.4333333,2.7006999) (1.4416667,2.7163449) (1.4500000,2.7320805) (1.4583333,2.7479073) (1.4666667,2.7638258) (1.4750000,2.7798364) (1.4833333,2.7959399) (1.4916667,2.8121366) (1.5000000,2.8284271) (1.5083333,2.8448120) (1.5166667,2.8612919) (1.5250000,2.8778672) (1.5333333,2.8945385) (1.5416667,2.9113064) (1.5500000,2.9281714) (1.5583333,2.9451341) (1.5666667,2.9621951) (1.5750000,2.9793549) (1.5833333,2.9966142) (1.5916667,3.0139734) (1.6000000,3.0314331) (1.6083333,3.0489940) (1.6166667,3.0666567) (1.6250000,3.0844217) (1.6333333,3.1022895) (1.6416667,3.1202609) (1.6500000,3.1383364) (1.6583333,3.1565166) (1.6666667,3.1748021) (1.6750000,3.1931935) (1.6833333,3.2116915) (1.6916667,3.2302967) (1.7000000,3.2490096) (1.7083333,3.2678309) (1.7166667,3.2867613) (1.7250000,3.3058013) (1.7333333,3.3249516) (1.7416667,3.3442128) (1.7500000,3.3635857) (1.7583333,3.3830707) (1.7666667,3.4026686) (1.7750000,3.4223801) (1.7833333,3.4422057) (1.7916667,3.4621462) (1.8000000,3.4822023) (1.8083333,3.5023744) (1.8166667,3.5226635) (1.8250000,3.5430701) (1.8333333,3.5635949) (1.8416667,3.5842386) (1.8500000,3.6050019) (1.8583333,3.6258854) (1.8666667,3.6468900) (1.8750000,3.6680162) (1.8833333,3.6892648) (1.8916667,3.7106365) (1.9000000,3.7321320) (1.9083333,3.7537520) (1.9166667,3.7754973) (1.9250000,3.7973685) (1.9333333,3.8193664) (1.9416667,3.8414918) (1.9500000,3.8637453) (1.9583333,3.8861278) (1.9666667,3.9086399) (1.9750000,3.9312824) (1.9833333,3.9540561) (1.9916667,3.9769617) (2.0000000,4.0000000) (2.0083333,4.0231718) (2.0166667,4.0464778) (2.0250000,4.0699188) (2.0333333,4.0934956) (2.0416667,4.1172089) (2.0500000,4.1410597) (2.0583333,4.1650486) (2.0666667,4.1891765) (2.0750000,4.2134441) (2.0833333,4.2378524) (2.0916667,4.2624020) (2.1000000,4.2870939) (2.1083333,4.3119287) (2.1166667,4.3369075) (2.1250000,4.3620309) (2.1333333,4.3872999) (2.1416667,4.4127153) (2.1500000,4.4382779) (2.1583333,4.4639886) (2.1666667,4.4898482) (2.1750000,4.5158576) (2.1833333,4.5420177) (2.1916667,4.5683294) (2.2000000,4.5947934) (2.2083333,4.6214108) (2.2166667,4.6481823) (2.2250000,4.6751090) (2.2333333,4.7021916) (2.2416667,4.7294311) (2.2500000,4.7568285) (2.2583333,4.7843845) (2.2666667,4.8121001) (2.2750000,4.8399764) (2.2833333,4.8680141) (2.2916667,4.8962142) (2.3000000,4.9245777) (2.3083333,4.9531054) (2.3166667,4.9817985) (2.3250000,5.0106578) (2.3333333,5.0396842) (2.3416667,5.0688788) (2.3500000,5.0982425) (2.3583333,5.1277763) (2.3666667,5.1574812) (2.3750000,5.1873582) (2.3833333,5.2174083) (2.3916667,5.2476324) (2.4000000,5.2780316) (2.4083333,5.3086070) (2.4166667,5.3393594) (2.4250000,5.3702900) (2.4333333,5.4013998) (2.4416667,5.4326898) (2.4500000,5.4641610) (2.4583333,5.4958146) (2.4666667,5.5276515) (2.4750000,5.5596729) (2.4833333,5.5918797) (2.4916667,5.6242732) (2.5000000,5.6568542)};
\draw[orange!85!black,thick,dashed] plot coordinates {(-2.5000000,0.1767767) (-2.4916667,0.1778008) (-2.4833333,0.1788307) (-2.4750000,0.1798667) (-2.4666667,0.1809087) (-2.4583333,0.1819566) (-2.4500000,0.1830107) (-2.4416667,0.1840709) (-2.4333333,0.1851372) (-2.4250000,0.1862097) (-2.4166667,0.1872884) (-2.4083333,0.1883733) (-2.4000000,0.1894646) (-2.3916667,0.1905621) (-2.3833333,0.1916660) (-2.3750000,0.1927764) (-2.3666667,0.1938931) (-2.3583333,0.1950163) (-2.3500000,0.1961460) (-2.3416667,0.1972823) (-2.3333333,0.1984251) (-2.3250000,0.1995746) (-2.3166667,0.2007307) (-2.3083333,0.2018935) (-2.3000000,0.2030631) (-2.2916667,0.2042394) (-2.2833333,0.2054226) (-2.2750000,0.2066126) (-2.2666667,0.2078095) (-2.2583333,0.2090133) (-2.2500000,0.2102241) (-2.2416667,0.2114419) (-2.2333333,0.2126668) (-2.2250000,0.2138988) (-2.2166667,0.2151379) (-2.2083333,0.2163841) (-2.2000000,0.2176376) (-2.1916667,0.2188984) (-2.1833333,0.2201665) (-2.1750000,0.2214419) (-2.1666667,0.2227247) (-2.1583333,0.2240149) (-2.1500000,0.2253126) (-2.1416667,0.2266178) (-2.1333333,0.2279306) (-2.1250000,0.2292510) (-2.1166667,0.2305790) (-2.1083333,0.2319148) (-2.1000000,0.2332582) (-2.0916667,0.2346095) (-2.0833333,0.2359686) (-2.0750000,0.2373355) (-2.0666667,0.2387104) (-2.0583333,0.2400932) (-2.0500000,0.2414841) (-2.0416667,0.2428830) (-2.0333333,0.2442900) (-2.0250000,0.2457051) (-2.0166667,0.2471285) (-2.0083333,0.2485601) (-2.0000000,0.2500000) (-1.9916667,0.2514482) (-1.9833333,0.2529049) (-1.9750000,0.2543699) (-1.9666667,0.2558435) (-1.9583333,0.2573256) (-1.9500000,0.2588162) (-1.9416667,0.2603155) (-1.9333333,0.2618235) (-1.9250000,0.2633403) (-1.9166667,0.2648658) (-1.9083333,0.2664001) (-1.9000000,0.2679434) (-1.8916667,0.2694955) (-1.8833333,0.2710567) (-1.8750000,0.2726269) (-1.8666667,0.2742062) (-1.8583333,0.2757947) (-1.8500000,0.2773924) (-1.8416667,0.2789993) (-1.8333333,0.2806155) (-1.8250000,0.2822411) (-1.8166667,0.2838761) (-1.8083333,0.2855206) (-1.8000000,0.2871746) (-1.7916667,0.2888382) (-1.7833333,0.2905114) (-1.7750000,0.2921943) (-1.7666667,0.2938870) (-1.7583333,0.2955894) (-1.7500000,0.2973018) (-1.7416667,0.2990240) (-1.7333333,0.3007563) (-1.7250000,0.3024985) (-1.7166667,0.3042509) (-1.7083333,0.3060134) (-1.7000000,0.3077861) (-1.6916667,0.3095691) (-1.6833333,0.3113624) (-1.6750000,0.3131661) (-1.6666667,0.3149803) (-1.6583333,0.3168049) (-1.6500000,0.3186402) (-1.6416667,0.3204860) (-1.6333333,0.3223426) (-1.6250000,0.3242099) (-1.6166667,0.3260880) (-1.6083333,0.3279770) (-1.6000000,0.3298770) (-1.5916667,0.3317879) (-1.5833333,0.3337100) (-1.5750000,0.3356431) (-1.5666667,0.3375875) (-1.5583333,0.3395431) (-1.5500000,0.3415101) (-1.5416667,0.3434884) (-1.5333333,0.3454782) (-1.5250000,0.3474796) (-1.5166667,0.3494925) (-1.5083333,0.3515171) (-1.5000000,0.3535534) (-1.4916667,0.3556015) (-1.4833333,0.3576615) (-1.4750000,0.3597334) (-1.4666667,0.3618173) (-1.4583333,0.3639133) (-1.4500000,0.3660214) (-1.4416667,0.3681418) (-1.4333333,0.3702744) (-1.4250000,0.3724194) (-1.4166667,0.3745768) (-1.4083333,0.3767467) (-1.4000000,0.3789291) (-1.3916667,0.3811243) (-1.3833333,0.3833321) (-1.3750000,0.3855527) (-1.3666667,0.3877862) (-1.3583333,0.3900326) (-1.3500000,0.3922920) (-1.3416667,0.3945646) (-1.3333333,0.3968503) (-1.3250000,0.3991492) (-1.3166667,0.4014614) (-1.3083333,0.4037871) (-1.3000000,0.4061262) (-1.2916667,0.4084789) (-1.2833333,0.4108452) (-1.2750000,0.4132252) (-1.2666667,0.4156189) (-1.2583333,0.4180266) (-1.2500000,0.4204482) (-1.2416667,0.4228838) (-1.2333333,0.4253336) (-1.2250000,0.4277975) (-1.2166667,0.4302757) (-1.2083333,0.4327683) (-1.2000000,0.4352753) (-1.1916667,0.4377968) (-1.1833333,0.4403329) (-1.1750000,0.4428838) (-1.1666667,0.4454494) (-1.1583333,0.4480298) (-1.1500000,0.4506252) (-1.1416667,0.4532357) (-1.1333333,0.4558612) (-1.1250000,0.4585020) (-1.1166667,0.4611581) (-1.1083333,0.4638296) (-1.1000000,0.4665165) (-1.0916667,0.4692190) (-1.0833333,0.4719372) (-1.0750000,0.4746711) (-1.0666667,0.4774208) (-1.0583333,0.4801865) (-1.0500000,0.4829682) (-1.0416667,0.4857660) (-1.0333333,0.4885800) (-1.0250000,0.4914103) (-1.0166667,0.4942570) (-1.0083333,0.4971202) (-1.0000000,0.5000000) (-0.9916667,0.5028965) (-0.9833333,0.5058097) (-0.9750000,0.5087398) (-0.9666667,0.5116869) (-0.9583333,0.5146511) (-0.9500000,0.5176325) (-0.9416667,0.5206311) (-0.9333333,0.5236471) (-0.9250000,0.5266805) (-0.9166667,0.5297315) (-0.9083333,0.5328003) (-0.9000000,0.5358867) (-0.8916667,0.5389911) (-0.8833333,0.5421134) (-0.8750000,0.5452539) (-0.8666667,0.5484125) (-0.8583333,0.5515894) (-0.8500000,0.5547847) (-0.8416667,0.5579986) (-0.8333333,0.5612310) (-0.8250000,0.5644822) (-0.8166667,0.5677522) (-0.8083333,0.5710412) (-0.8000000,0.5743492) (-0.7916667,0.5776763) (-0.7833333,0.5810228) (-0.7750000,0.5843886) (-0.7666667,0.5877740) (-0.7583333,0.5911789) (-0.7500000,0.5946036) (-0.7416667,0.5980481) (-0.7333333,0.6015125) (-0.7250000,0.6049970) (-0.7166667,0.6085018) (-0.7083333,0.6120268) (-0.7000000,0.6155722) (-0.6916667,0.6191382) (-0.6833333,0.6227248) (-0.6750000,0.6263322) (-0.6666667,0.6299605) (-0.6583333,0.6336098) (-0.6500000,0.6372803) (-0.6416667,0.6409720) (-0.6333333,0.6446852) (-0.6250000,0.6484198) (-0.6166667,0.6521760) (-0.6083333,0.6559541) (-0.6000000,0.6597540) (-0.5916667,0.6635759) (-0.5833333,0.6674199) (-0.5750000,0.6712863) (-0.5666667,0.6751750) (-0.5583333,0.6790862) (-0.5500000,0.6830201) (-0.5416667,0.6869768) (-0.5333333,0.6909564) (-0.5250000,0.6949591) (-0.5166667,0.6989850) (-0.5083333,0.7030341) (-0.5000000,0.7071068) (-0.4916667,0.7112030) (-0.4833333,0.7153230) (-0.4750000,0.7194668) (-0.4666667,0.7236346) (-0.4583333,0.7278266) (-0.4500000,0.7320428) (-0.4416667,0.7362835) (-0.4333333,0.7405488) (-0.4250000,0.7448387) (-0.4166667,0.7491535) (-0.4083333,0.7534933) (-0.4000000,0.7578583) (-0.3916667,0.7622485) (-0.3833333,0.7666642) (-0.3750000,0.7711054) (-0.3666667,0.7755724) (-0.3583333,0.7800652) (-0.3500000,0.7845841) (-0.3416667,0.7891291) (-0.3333333,0.7937005) (-0.3250000,0.7982984) (-0.3166667,0.8029229) (-0.3083333,0.8075742) (-0.3000000,0.8122524) (-0.2916667,0.8169577) (-0.2833333,0.8216903) (-0.2750000,0.8264503) (-0.2666667,0.8312379) (-0.2583333,0.8360532) (-0.2500000,0.8408964) (-0.2416667,0.8457677) (-0.2333333,0.8506672) (-0.2250000,0.8555950) (-0.2166667,0.8605514) (-0.2083333,0.8655366) (-0.2000000,0.8705506) (-0.1916667,0.8755936) (-0.1833333,0.8806659) (-0.1750000,0.8857675) (-0.1666667,0.8908987) (-0.1583333,0.8960596) (-0.1500000,0.9012505) (-0.1416667,0.9064714) (-0.1333333,0.9117225) (-0.1250000,0.9170040) (-0.1166667,0.9223162) (-0.1083333,0.9276591) (-0.1000000,0.9330330) (-0.0916667,0.9384380) (-0.0833333,0.9438743) (-0.0750000,0.9493421) (-0.0666667,0.9548416) (-0.0583333,0.9603729) (-0.0500000,0.9659363) (-0.0416667,0.9715319) (-0.0333333,0.9771600) (-0.0250000,0.9828206) (-0.0166667,0.9885140) (-0.0083333,0.9942404) (0.0000000,1.0000000) (0.0083333,0.9942404) (0.0166667,0.9885140) (0.0250000,0.9828206) (0.0333333,0.9771600) (0.0416667,0.9715319) (0.0500000,0.9659363) (0.0583333,0.9603729) (0.0666667,0.9548416) (0.0750000,0.9493421) (0.0833333,0.9438743) (0.0916667,0.9384380) (0.1000000,0.9330330) (0.1083333,0.9276591) (0.1166667,0.9223162) (0.1250000,0.9170040) (0.1333333,0.9117225) (0.1416667,0.9064714) (0.1500000,0.9012505) (0.1583333,0.8960596) (0.1666667,0.8908987) (0.1750000,0.8857675) (0.1833333,0.8806659) (0.1916667,0.8755936) (0.2000000,0.8705506) (0.2083333,0.8655366) (0.2166667,0.8605514) (0.2250000,0.8555950) (0.2333333,0.8506672) (0.2416667,0.8457677) (0.2500000,0.8408964) (0.2583333,0.8360532) (0.2666667,0.8312379) (0.2750000,0.8264503) (0.2833333,0.8216903) (0.2916667,0.8169577) (0.3000000,0.8122524) (0.3083333,0.8075742) (0.3166667,0.8029229) (0.3250000,0.7982984) (0.3333333,0.7937005) (0.3416667,0.7891291) (0.3500000,0.7845841) (0.3583333,0.7800652) (0.3666667,0.7755724) (0.3750000,0.7711054) (0.3833333,0.7666642) (0.3916667,0.7622485) (0.4000000,0.7578583) (0.4083333,0.7534933) (0.4166667,0.7491535) (0.4250000,0.7448387) (0.4333333,0.7405488) (0.4416667,0.7362835) (0.4500000,0.7320428) (0.4583333,0.7278266) (0.4666667,0.7236346) (0.4750000,0.7194668) (0.4833333,0.7153230) (0.4916667,0.7112030) (0.5000000,0.7071068) (0.5083333,0.7030341) (0.5166667,0.6989850) (0.5250000,0.6949591) (0.5333333,0.6909564) (0.5416667,0.6869768) (0.5500000,0.6830201) (0.5583333,0.6790862) (0.5666667,0.6751750) (0.5750000,0.6712863) (0.5833333,0.6674199) (0.5916667,0.6635759) (0.6000000,0.6597540) (0.6083333,0.6559541) (0.6166667,0.6521760) (0.6250000,0.6484198) (0.6333333,0.6446852) (0.6416667,0.6409720) (0.6500000,0.6372803) (0.6583333,0.6336098) (0.6666667,0.6299605) (0.6750000,0.6263322) (0.6833333,0.6227248) (0.6916667,0.6191382) (0.7000000,0.6155722) (0.7083333,0.6120268) (0.7166667,0.6085018) (0.7250000,0.6049970) (0.7333333,0.6015125) (0.7416667,0.5980481) (0.7500000,0.5946036) (0.7583333,0.5911789) (0.7666667,0.5877740) (0.7750000,0.5843886) (0.7833333,0.5810228) (0.7916667,0.5776763) (0.8000000,0.5743492) (0.8083333,0.5710412) (0.8166667,0.5677522) (0.8250000,0.5644822) (0.8333333,0.5612310) (0.8416667,0.5579986) (0.8500000,0.5547847) (0.8583333,0.5515894) (0.8666667,0.5484125) (0.8750000,0.5452539) (0.8833333,0.5421134) (0.8916667,0.5389911) (0.9000000,0.5358867) (0.9083333,0.5328003) (0.9166667,0.5297315) (0.9250000,0.5266805) (0.9333333,0.5236471) (0.9416667,0.5206311) (0.9500000,0.5176325) (0.9583333,0.5146511) (0.9666667,0.5116869) (0.9750000,0.5087398) (0.9833333,0.5058097) (0.9916667,0.5028965) (1.0000000,0.5000000) (1.0083333,0.4971202) (1.0166667,0.4942570) (1.0250000,0.4914103) (1.0333333,0.4885800) (1.0416667,0.4857660) (1.0500000,0.4829682) (1.0583333,0.4801865) (1.0666667,0.4774208) (1.0750000,0.4746711) (1.0833333,0.4719372) (1.0916667,0.4692190) (1.1000000,0.4665165) (1.1083333,0.4638296) (1.1166667,0.4611581) (1.1250000,0.4585020) (1.1333333,0.4558612) (1.1416667,0.4532357) (1.1500000,0.4506252) (1.1583333,0.4480298) (1.1666667,0.4454494) (1.1750000,0.4428838) (1.1833333,0.4403329) (1.1916667,0.4377968) (1.2000000,0.4352753) (1.2083333,0.4327683) (1.2166667,0.4302757) (1.2250000,0.4277975) (1.2333333,0.4253336) (1.2416667,0.4228838) (1.2500000,0.4204482) (1.2583333,0.4180266) (1.2666667,0.4156189) (1.2750000,0.4132252) (1.2833333,0.4108452) (1.2916667,0.4084789) (1.3000000,0.4061262) (1.3083333,0.4037871) (1.3166667,0.4014614) (1.3250000,0.3991492) (1.3333333,0.3968503) (1.3416667,0.3945646) (1.3500000,0.3922920) (1.3583333,0.3900326) (1.3666667,0.3877862) (1.3750000,0.3855527) (1.3833333,0.3833321) (1.3916667,0.3811243) (1.4000000,0.3789291) (1.4083333,0.3767467) (1.4166667,0.3745768) (1.4250000,0.3724194) (1.4333333,0.3702744) (1.4416667,0.3681418) (1.4500000,0.3660214) (1.4583333,0.3639133) (1.4666667,0.3618173) (1.4750000,0.3597334) (1.4833333,0.3576615) (1.4916667,0.3556015) (1.5000000,0.3535534) (1.5083333,0.3515171) (1.5166667,0.3494925) (1.5250000,0.3474796) (1.5333333,0.3454782) (1.5416667,0.3434884) (1.5500000,0.3415101) (1.5583333,0.3395431) (1.5666667,0.3375875) (1.5750000,0.3356431) (1.5833333,0.3337100) (1.5916667,0.3317879) (1.6000000,0.3298770) (1.6083333,0.3279770) (1.6166667,0.3260880) (1.6250000,0.3242099) (1.6333333,0.3223426) (1.6416667,0.3204860) (1.6500000,0.3186402) (1.6583333,0.3168049) (1.6666667,0.3149803) (1.6750000,0.3131661) (1.6833333,0.3113624) (1.6916667,0.3095691) (1.7000000,0.3077861) (1.7083333,0.3060134) (1.7166667,0.3042509) (1.7250000,0.3024985) (1.7333333,0.3007563) (1.7416667,0.2990240) (1.7500000,0.2973018) (1.7583333,0.2955894) (1.7666667,0.2938870) (1.7750000,0.2921943) (1.7833333,0.2905114) (1.7916667,0.2888382) (1.8000000,0.2871746) (1.8083333,0.2855206) (1.8166667,0.2838761) (1.8250000,0.2822411) (1.8333333,0.2806155) (1.8416667,0.2789993) (1.8500000,0.2773924) (1.8583333,0.2757947) (1.8666667,0.2742062) (1.8750000,0.2726269) (1.8833333,0.2710567) (1.8916667,0.2694955) (1.9000000,0.2679434) (1.9083333,0.2664001) (1.9166667,0.2648658) (1.9250000,0.2633403) (1.9333333,0.2618235) (1.9416667,0.2603155) (1.9500000,0.2588162) (1.9583333,0.2573256) (1.9666667,0.2558435) (1.9750000,0.2543699) (1.9833333,0.2529049) (1.9916667,0.2514482) (2.0000000,0.2500000) (2.0083333,0.2485601) (2.0166667,0.2471285) (2.0250000,0.2457051) (2.0333333,0.2442900) (2.0416667,0.2428830) (2.0500000,0.2414841) (2.0583333,0.2400932) (2.0666667,0.2387104) (2.0750000,0.2373355) (2.0833333,0.2359686) (2.0916667,0.2346095) (2.1000000,0.2332582) (2.1083333,0.2319148) (2.1166667,0.2305790) (2.1250000,0.2292510) (2.1333333,0.2279306) (2.1416667,0.2266178) (2.1500000,0.2253126) (2.1583333,0.2240149) (2.1666667,0.2227247) (2.1750000,0.2214419) (2.1833333,0.2201665) (2.1916667,0.2188984) (2.2000000,0.2176376) (2.2083333,0.2163841) (2.2166667,0.2151379) (2.2250000,0.2138988) (2.2333333,0.2126668) (2.2416667,0.2114419) (2.2500000,0.2102241) (2.2583333,0.2090133) (2.2666667,0.2078095) (2.2750000,0.2066126) (2.2833333,0.2054226) (2.2916667,0.2042394) (2.3000000,0.2030631) (2.3083333,0.2018935) (2.3166667,0.2007307) (2.3250000,0.1995746) (2.3333333,0.1984251) (2.3416667,0.1972823) (2.3500000,0.1961460) (2.3583333,0.1950163) (2.3666667,0.1938931) (2.3750000,0.1927764) (2.3833333,0.1916660) (2.3916667,0.1905621) (2.4000000,0.1894646) (2.4083333,0.1883733) (2.4166667,0.1872884) (2.4250000,0.1862097) (2.4333333,0.1851372) (2.4416667,0.1840709) (2.4500000,0.1830107) (2.4583333,0.1819566) (2.4666667,0.1809087) (2.4750000,0.1798667) (2.4833333,0.1788307) (2.4916667,0.1778008) (2.5000000,0.1767767)};
\draw[violet!75!black,densely dotted] (-1,2.0000000)--(1,2.0000000);
\fill[violet!75!black] (-1,2.0000000) circle(2.3pt);
\fill[violet!75!black] (1,2.0000000) circle(2.3pt);
\end{scope}
\end{scope}
\node[anchor=west,align=left,text width=4.35cm] at (8.9,4.4) {蓝实线：$a=2$\\[8pt]橙虚线：$a=\frac12$\\[14pt]两图都关于 $y$ 轴对称。\\[8pt]$a>1$ 时中间最低；\\$0<a<1$ 时中间最高。};
\node[text width=12.8cm,align=center] at (7.1,.65) {$D=\mathbb R$；$f(0)=1$。保留指数函数右半图，再向左翻折。};
\end{tikzpicture}\end{center}

\textbf{先看图：} 把 $x$ 与 $-x$ 先取绝对值，得到同一个非负输入，两侧高度相同。蓝线 $a=2$ 在中间最低；橙线 $a=\frac12$ 在中间最高，但两条曲线都是偶函数。

\textbf{再说明一般情形：}
\[
f(-x)=a^{|-x|}=a^{|x|}=f(x).
\]
$a>1$ 时值域为 $[1,+\infty)$；$0<a<1$ 时为 $(0,1]$。更换峰谷不改变奇偶性。
\end{explanationbox}
\clearpage
\begin{explanationbox}
\textbf{\textcolor{CExplanation}{负的平方也是偶函数}}
% Formula-based samples; see F083 checks/build_tikz_diagrams.py.
\begin{center}\begin{tikzpicture}[font=\small,>=stealth,line join=round]
\draw[gray!40,rounded corners=3pt] (0,0) rectangle (14.2,8.4);
\node at (7.1,7.75) {偶3：$y=a^{-x^2}$};
\begin{scope}[shift={(4.400000,1.650000)},x=1.700000cm,y=0.900000cm]
\draw[->,gray] (-2,0)--(2,0) node[right] {$x$};
\draw[->,gray] (0,0)--(0,5) node[above] {$y$};
\draw[gray] (-1,-.035/0.900000)--(-1,.035/0.900000);
\node[anchor=north] at (-1,-.07/0.900000) {$-1$};
\draw[gray] (1,-.035/0.900000)--(1,.035/0.900000);
\node[anchor=north] at (1,-.07/0.900000) {$1$};
\begin{scope}
\clip (-2,0) rectangle (2,5);
\draw[blue,thick,solid] plot coordinates {(-2.0000000,0.0625000) (-1.9933333,0.0636640) (-1.9866667,0.0648457) (-1.9800000,0.0660453) (-1.9733333,0.0672629) (-1.9666667,0.0684988) (-1.9600000,0.0697530) (-1.9533333,0.0710259) (-1.9466667,0.0723175) (-1.9400000,0.0736281) (-1.9333333,0.0749578) (-1.9266667,0.0763068) (-1.9200000,0.0776753) (-1.9133333,0.0790635) (-1.9066667,0.0804715) (-1.9000000,0.0818996) (-1.8933333,0.0833479) (-1.8866667,0.0848165) (-1.8800000,0.0863057) (-1.8733333,0.0878157) (-1.8666667,0.0893465) (-1.8600000,0.0908985) (-1.8533333,0.0924717) (-1.8466667,0.0940663) (-1.8400000,0.0956826) (-1.8333333,0.0973206) (-1.8266667,0.0989806) (-1.8200000,0.1006626) (-1.8133333,0.1023670) (-1.8066667,0.1040938) (-1.8000000,0.1058432) (-1.7933333,0.1076153) (-1.7866667,0.1094104) (-1.7800000,0.1112286) (-1.7733333,0.1130701) (-1.7666667,0.1149349) (-1.7600000,0.1168233) (-1.7533333,0.1187354) (-1.7466667,0.1206714) (-1.7400000,0.1226314) (-1.7333333,0.1246155) (-1.7266667,0.1266240) (-1.7200000,0.1286568) (-1.7133333,0.1307143) (-1.7066667,0.1327965) (-1.7000000,0.1349035) (-1.6933333,0.1370356) (-1.6866667,0.1391927) (-1.6800000,0.1413751) (-1.6733333,0.1435829) (-1.6666667,0.1458161) (-1.6600000,0.1480750) (-1.6533333,0.1503596) (-1.6466667,0.1526700) (-1.6400000,0.1550064) (-1.6333333,0.1573689) (-1.6266667,0.1597575) (-1.6200000,0.1621724) (-1.6133333,0.1646136) (-1.6066667,0.1670813) (-1.6000000,0.1695755) (-1.5933333,0.1720964) (-1.5866667,0.1746440) (-1.5800000,0.1772184) (-1.5733333,0.1798196) (-1.5666667,0.1824478) (-1.5600000,0.1851030) (-1.5533333,0.1877852) (-1.5466667,0.1904946) (-1.5400000,0.1932312) (-1.5333333,0.1959950) (-1.5266667,0.1987861) (-1.5200000,0.2016045) (-1.5133333,0.2044503) (-1.5066667,0.2073235) (-1.5000000,0.2102241) (-1.4933333,0.2131522) (-1.4866667,0.2161077) (-1.4800000,0.2190907) (-1.4733333,0.2221012) (-1.4666667,0.2251392) (-1.4600000,0.2282046) (-1.4533333,0.2312976) (-1.4466667,0.2344180) (-1.4400000,0.2375660) (-1.4333333,0.2407413) (-1.4266667,0.2439441) (-1.4200000,0.2471742) (-1.4133333,0.2504317) (-1.4066667,0.2537164) (-1.4000000,0.2570285) (-1.3933333,0.2603677) (-1.3866667,0.2637340) (-1.3800000,0.2671274) (-1.3733333,0.2705478) (-1.3666667,0.2739951) (-1.3600000,0.2774693) (-1.3533333,0.2809702) (-1.3466667,0.2844977) (-1.3400000,0.2880518) (-1.3333333,0.2916323) (-1.3266667,0.2952391) (-1.3200000,0.2988721) (-1.3133333,0.3025311) (-1.3066667,0.3062161) (-1.3000000,0.3099269) (-1.2933333,0.3136634) (-1.2866667,0.3174253) (-1.2800000,0.3212125) (-1.2733333,0.3250249) (-1.2666667,0.3288623) (-1.2600000,0.3327245) (-1.2533333,0.3366113) (-1.2466667,0.3405225) (-1.2400000,0.3444580) (-1.2333333,0.3484175) (-1.2266667,0.3524007) (-1.2200000,0.3564076) (-1.2133333,0.3604378) (-1.2066667,0.3644911) (-1.2000000,0.3685673) (-1.1933333,0.3726661) (-1.1866667,0.3767873) (-1.1800000,0.3809306) (-1.1733333,0.3850957) (-1.1666667,0.3892824) (-1.1600000,0.3934903) (-1.1533333,0.3977192) (-1.1466667,0.4019688) (-1.1400000,0.4062388) (-1.1333333,0.4105289) (-1.1266667,0.4148387) (-1.1200000,0.4191679) (-1.1133333,0.4235161) (-1.1066667,0.4278832) (-1.1000000,0.4322686) (-1.0933333,0.4366721) (-1.0866667,0.4410932) (-1.0800000,0.4455317) (-1.0733333,0.4499871) (-1.0666667,0.4544591) (-1.0600000,0.4589472) (-1.0533333,0.4634511) (-1.0466667,0.4679703) (-1.0400000,0.4725045) (-1.0333333,0.4770533) (-1.0266667,0.4816161) (-1.0200000,0.4861927) (-1.0133333,0.4907824) (-1.0066667,0.4953850) (-1.0000000,0.5000000) (-0.9933333,0.5046269) (-0.9866667,0.5092651) (-0.9800000,0.5139144) (-0.9733333,0.5185742) (-0.9666667,0.5232439) (-0.9600000,0.5279232) (-0.9533333,0.5326115) (-0.9466667,0.5373084) (-0.9400000,0.5420132) (-0.9333333,0.5467256) (-0.9266667,0.5514450) (-0.9200000,0.5561708) (-0.9133333,0.5609026) (-0.9066667,0.5656398) (-0.9000000,0.5703819) (-0.8933333,0.5751282) (-0.8866667,0.5798784) (-0.8800000,0.5846317) (-0.8733333,0.5893877) (-0.8666667,0.5941458) (-0.8600000,0.5989054) (-0.8533333,0.6036659) (-0.8466667,0.6084268) (-0.8400000,0.6131874) (-0.8333333,0.6179472) (-0.8266667,0.6227056) (-0.8200000,0.6274620) (-0.8133333,0.6322158) (-0.8066667,0.6369663) (-0.8000000,0.6417129) (-0.7933333,0.6464552) (-0.7866667,0.6511923) (-0.7800000,0.6559237) (-0.7733333,0.6606489) (-0.7666667,0.6653670) (-0.7600000,0.6700776) (-0.7533333,0.6747799) (-0.7466667,0.6794734) (-0.7400000,0.6841573) (-0.7333333,0.6888311) (-0.7266667,0.6934941) (-0.7200000,0.6981457) (-0.7133333,0.7027851) (-0.7066667,0.7074118) (-0.7000000,0.7120251) (-0.6933333,0.7166243) (-0.6866667,0.7212088) (-0.6800000,0.7257779) (-0.6733333,0.7303309) (-0.6666667,0.7348672) (-0.6600000,0.7393862) (-0.6533333,0.7438871) (-0.6466667,0.7483693) (-0.6400000,0.7528321) (-0.6333333,0.7572748) (-0.6266667,0.7616969) (-0.6200000,0.7660975) (-0.6133333,0.7704762) (-0.6066667,0.7748321) (-0.6000000,0.7791646) (-0.5933333,0.7834731) (-0.5866667,0.7877568) (-0.5800000,0.7920152) (-0.5733333,0.7962475) (-0.5666667,0.8004532) (-0.5600000,0.8046314) (-0.5533333,0.8087817) (-0.5466667,0.8129032) (-0.5400000,0.8169955) (-0.5333333,0.8210577) (-0.5266667,0.8250893) (-0.5200000,0.8290896) (-0.5133333,0.8330580) (-0.5066667,0.8369938) (-0.5000000,0.8408964) (-0.4933333,0.8447651) (-0.4866667,0.8485994) (-0.4800000,0.8523985) (-0.4733333,0.8561619) (-0.4666667,0.8598889) (-0.4600000,0.8635790) (-0.4533333,0.8672314) (-0.4466667,0.8708456) (-0.4400000,0.8744210) (-0.4333333,0.8779570) (-0.4266667,0.8814530) (-0.4200000,0.8849084) (-0.4133333,0.8883226) (-0.4066667,0.8916950) (-0.4000000,0.8950251) (-0.3933333,0.8983122) (-0.3866667,0.9015559) (-0.3800000,0.9047556) (-0.3733333,0.9079106) (-0.3666667,0.9110206) (-0.3600000,0.9140849) (-0.3533333,0.9171029) (-0.3466667,0.9200743) (-0.3400000,0.9229984) (-0.3333333,0.9258747) (-0.3266667,0.9287028) (-0.3200000,0.9314821) (-0.3133333,0.9342122) (-0.3066667,0.9368926) (-0.3000000,0.9395227) (-0.2933333,0.9421023) (-0.2866667,0.9446306) (-0.2800000,0.9471074) (-0.2733333,0.9495322) (-0.2666667,0.9519046) (-0.2600000,0.9542241) (-0.2533333,0.9564903) (-0.2466667,0.9587028) (-0.2400000,0.9608612) (-0.2333333,0.9629652) (-0.2266667,0.9650143) (-0.2200000,0.9670082) (-0.2133333,0.9689465) (-0.2066667,0.9708289) (-0.2000000,0.9726549) (-0.1933333,0.9744244) (-0.1866667,0.9761370) (-0.1800000,0.9777923) (-0.1733333,0.9793901) (-0.1666667,0.9809301) (-0.1600000,0.9824119) (-0.1533333,0.9838354) (-0.1466667,0.9852002) (-0.1400000,0.9865062) (-0.1333333,0.9877530) (-0.1266667,0.9889405) (-0.1200000,0.9900683) (-0.1133333,0.9911364) (-0.1066667,0.9921445) (-0.1000000,0.9930925) (-0.0933333,0.9939801) (-0.0866667,0.9948072) (-0.0800000,0.9955737) (-0.0733333,0.9962793) (-0.0666667,0.9969241) (-0.0600000,0.9975078) (-0.0533333,0.9980303) (-0.0466667,0.9984916) (-0.0400000,0.9988916) (-0.0333333,0.9992301) (-0.0266667,0.9995072) (-0.0200000,0.9997228) (-0.0133333,0.9998768) (-0.0066667,0.9999692) (0.0000000,1.0000000) (0.0066667,0.9999692) (0.0133333,0.9998768) (0.0200000,0.9997228) (0.0266667,0.9995072) (0.0333333,0.9992301) (0.0400000,0.9988916) (0.0466667,0.9984916) (0.0533333,0.9980303) (0.0600000,0.9975078) (0.0666667,0.9969241) (0.0733333,0.9962793) (0.0800000,0.9955737) (0.0866667,0.9948072) (0.0933333,0.9939801) (0.1000000,0.9930925) (0.1066667,0.9921445) (0.1133333,0.9911364) (0.1200000,0.9900683) (0.1266667,0.9889405) (0.1333333,0.9877530) (0.1400000,0.9865062) (0.1466667,0.9852002) (0.1533333,0.9838354) (0.1600000,0.9824119) (0.1666667,0.9809301) (0.1733333,0.9793901) (0.1800000,0.9777923) (0.1866667,0.9761370) (0.1933333,0.9744244) (0.2000000,0.9726549) (0.2066667,0.9708289) (0.2133333,0.9689465) (0.2200000,0.9670082) (0.2266667,0.9650143) (0.2333333,0.9629652) (0.2400000,0.9608612) (0.2466667,0.9587028) (0.2533333,0.9564903) (0.2600000,0.9542241) (0.2666667,0.9519046) (0.2733333,0.9495322) (0.2800000,0.9471074) (0.2866667,0.9446306) (0.2933333,0.9421023) (0.3000000,0.9395227) (0.3066667,0.9368926) (0.3133333,0.9342122) (0.3200000,0.9314821) (0.3266667,0.9287028) (0.3333333,0.9258747) (0.3400000,0.9229984) (0.3466667,0.9200743) (0.3533333,0.9171029) (0.3600000,0.9140849) (0.3666667,0.9110206) (0.3733333,0.9079106) (0.3800000,0.9047556) (0.3866667,0.9015559) (0.3933333,0.8983122) (0.4000000,0.8950251) (0.4066667,0.8916950) (0.4133333,0.8883226) (0.4200000,0.8849084) (0.4266667,0.8814530) (0.4333333,0.8779570) (0.4400000,0.8744210) (0.4466667,0.8708456) (0.4533333,0.8672314) (0.4600000,0.8635790) (0.4666667,0.8598889) (0.4733333,0.8561619) (0.4800000,0.8523985) (0.4866667,0.8485994) (0.4933333,0.8447651) (0.5000000,0.8408964) (0.5066667,0.8369938) (0.5133333,0.8330580) (0.5200000,0.8290896) (0.5266667,0.8250893) (0.5333333,0.8210577) (0.5400000,0.8169955) (0.5466667,0.8129032) (0.5533333,0.8087817) (0.5600000,0.8046314) (0.5666667,0.8004532) (0.5733333,0.7962475) (0.5800000,0.7920152) (0.5866667,0.7877568) (0.5933333,0.7834731) (0.6000000,0.7791646) (0.6066667,0.7748321) (0.6133333,0.7704762) (0.6200000,0.7660975) (0.6266667,0.7616969) (0.6333333,0.7572748) (0.6400000,0.7528321) (0.6466667,0.7483693) (0.6533333,0.7438871) (0.6600000,0.7393862) (0.6666667,0.7348672) (0.6733333,0.7303309) (0.6800000,0.7257779) (0.6866667,0.7212088) (0.6933333,0.7166243) (0.7000000,0.7120251) (0.7066667,0.7074118) (0.7133333,0.7027851) (0.7200000,0.6981457) (0.7266667,0.6934941) (0.7333333,0.6888311) (0.7400000,0.6841573) (0.7466667,0.6794734) (0.7533333,0.6747799) (0.7600000,0.6700776) (0.7666667,0.6653670) (0.7733333,0.6606489) (0.7800000,0.6559237) (0.7866667,0.6511923) (0.7933333,0.6464552) (0.8000000,0.6417129) (0.8066667,0.6369663) (0.8133333,0.6322158) (0.8200000,0.6274620) (0.8266667,0.6227056) (0.8333333,0.6179472) (0.8400000,0.6131874) (0.8466667,0.6084268) (0.8533333,0.6036659) (0.8600000,0.5989054) (0.8666667,0.5941458) (0.8733333,0.5893877) (0.8800000,0.5846317) (0.8866667,0.5798784) (0.8933333,0.5751282) (0.9000000,0.5703819) (0.9066667,0.5656398) (0.9133333,0.5609026) (0.9200000,0.5561708) (0.9266667,0.5514450) (0.9333333,0.5467256) (0.9400000,0.5420132) (0.9466667,0.5373084) (0.9533333,0.5326115) (0.9600000,0.5279232) (0.9666667,0.5232439) (0.9733333,0.5185742) (0.9800000,0.5139144) (0.9866667,0.5092651) (0.9933333,0.5046269) (1.0000000,0.5000000) (1.0066667,0.4953850) (1.0133333,0.4907824) (1.0200000,0.4861927) (1.0266667,0.4816161) (1.0333333,0.4770533) (1.0400000,0.4725045) (1.0466667,0.4679703) (1.0533333,0.4634511) (1.0600000,0.4589472) (1.0666667,0.4544591) (1.0733333,0.4499871) (1.0800000,0.4455317) (1.0866667,0.4410932) (1.0933333,0.4366721) (1.1000000,0.4322686) (1.1066667,0.4278832) (1.1133333,0.4235161) (1.1200000,0.4191679) (1.1266667,0.4148387) (1.1333333,0.4105289) (1.1400000,0.4062388) (1.1466667,0.4019688) (1.1533333,0.3977192) (1.1600000,0.3934903) (1.1666667,0.3892824) (1.1733333,0.3850957) (1.1800000,0.3809306) (1.1866667,0.3767873) (1.1933333,0.3726661) (1.2000000,0.3685673) (1.2066667,0.3644911) (1.2133333,0.3604378) (1.2200000,0.3564076) (1.2266667,0.3524007) (1.2333333,0.3484175) (1.2400000,0.3444580) (1.2466667,0.3405225) (1.2533333,0.3366113) (1.2600000,0.3327245) (1.2666667,0.3288623) (1.2733333,0.3250249) (1.2800000,0.3212125) (1.2866667,0.3174253) (1.2933333,0.3136634) (1.3000000,0.3099269) (1.3066667,0.3062161) (1.3133333,0.3025311) (1.3200000,0.2988721) (1.3266667,0.2952391) (1.3333333,0.2916323) (1.3400000,0.2880518) (1.3466667,0.2844977) (1.3533333,0.2809702) (1.3600000,0.2774693) (1.3666667,0.2739951) (1.3733333,0.2705478) (1.3800000,0.2671274) (1.3866667,0.2637340) (1.3933333,0.2603677) (1.4000000,0.2570285) (1.4066667,0.2537164) (1.4133333,0.2504317) (1.4200000,0.2471742) (1.4266667,0.2439441) (1.4333333,0.2407413) (1.4400000,0.2375660) (1.4466667,0.2344180) (1.4533333,0.2312976) (1.4600000,0.2282046) (1.4666667,0.2251392) (1.4733333,0.2221012) (1.4800000,0.2190907) (1.4866667,0.2161077) (1.4933333,0.2131522) (1.5000000,0.2102241) (1.5066667,0.2073235) (1.5133333,0.2044503) (1.5200000,0.2016045) (1.5266667,0.1987861) (1.5333333,0.1959950) (1.5400000,0.1932312) (1.5466667,0.1904946) (1.5533333,0.1877852) (1.5600000,0.1851030) (1.5666667,0.1824478) (1.5733333,0.1798196) (1.5800000,0.1772184) (1.5866667,0.1746440) (1.5933333,0.1720964) (1.6000000,0.1695755) (1.6066667,0.1670813) (1.6133333,0.1646136) (1.6200000,0.1621724) (1.6266667,0.1597575) (1.6333333,0.1573689) (1.6400000,0.1550064) (1.6466667,0.1526700) (1.6533333,0.1503596) (1.6600000,0.1480750) (1.6666667,0.1458161) (1.6733333,0.1435829) (1.6800000,0.1413751) (1.6866667,0.1391927) (1.6933333,0.1370356) (1.7000000,0.1349035) (1.7066667,0.1327965) (1.7133333,0.1307143) (1.7200000,0.1286568) (1.7266667,0.1266240) (1.7333333,0.1246155) (1.7400000,0.1226314) (1.7466667,0.1206714) (1.7533333,0.1187354) (1.7600000,0.1168233) (1.7666667,0.1149349) (1.7733333,0.1130701) (1.7800000,0.1112286) (1.7866667,0.1094104) (1.7933333,0.1076153) (1.8000000,0.1058432) (1.8066667,0.1040938) (1.8133333,0.1023670) (1.8200000,0.1006626) (1.8266667,0.0989806) (1.8333333,0.0973206) (1.8400000,0.0956826) (1.8466667,0.0940663) (1.8533333,0.0924717) (1.8600000,0.0908985) (1.8666667,0.0893465) (1.8733333,0.0878157) (1.8800000,0.0863057) (1.8866667,0.0848165) (1.8933333,0.0833479) (1.9000000,0.0818996) (1.9066667,0.0804715) (1.9133333,0.0790635) (1.9200000,0.0776753) (1.9266667,0.0763068) (1.9333333,0.0749578) (1.9400000,0.0736281) (1.9466667,0.0723175) (1.9533333,0.0710259) (1.9600000,0.0697530) (1.9666667,0.0684988) (1.9733333,0.0672629) (1.9800000,0.0660453) (1.9866667,0.0648457) (1.9933333,0.0636640) (2.0000000,0.0625000)};
\draw[orange!85!black,thick,dashed] plot coordinates {(-1.5200000,4.9602059) (-1.5133333,4.8911638) (-1.5066667,4.8233799) (-1.5000000,4.7568285) (-1.4933333,4.6914843) (-1.4866667,4.6273229) (-1.4800000,4.5643202) (-1.4733333,4.5024527) (-1.4666667,4.4416974) (-1.4600000,4.3820319) (-1.4533333,4.3234343) (-1.4466667,4.2658831) (-1.4400000,4.2093574) (-1.4333333,4.1538365) (-1.4266667,4.0993006) (-1.4200000,4.0457299) (-1.4133333,3.9931053) (-1.4066667,3.9414080) (-1.4000000,3.8906198) (-1.3933333,3.8407226) (-1.3866667,3.7916990) (-1.3800000,3.7435318) (-1.3733333,3.6962042) (-1.3666667,3.6496997) (-1.3600000,3.6040025) (-1.3533333,3.5590967) (-1.3466667,3.5149669) (-1.3400000,3.4715982) (-1.3333333,3.4289759) (-1.3266667,3.3870856) (-1.3200000,3.3459132) (-1.3133333,3.3054449) (-1.3066667,3.2656672) (-1.3000000,3.2265670) (-1.2933333,3.1881314) (-1.2866667,3.1503478) (-1.2800000,3.1132037) (-1.2733333,3.0766872) (-1.2666667,3.0407863) (-1.2600000,3.0054895) (-1.2533333,2.9707855) (-1.2466667,2.9366631) (-1.2400000,2.9031115) (-1.2333333,2.8701201) (-1.2266667,2.8376784) (-1.2200000,2.8057763) (-1.2133333,2.7744038) (-1.2066667,2.7435511) (-1.2000000,2.7132087) (-1.1933333,2.6833671) (-1.1866667,2.6540173) (-1.1800000,2.6251503) (-1.1733333,2.5967572) (-1.1666667,2.5688295) (-1.1600000,2.5413587) (-1.1533333,2.5143366) (-1.1466667,2.4877551) (-1.1400000,2.4616062) (-1.1333333,2.4358823) (-1.1266667,2.4105758) (-1.1200000,2.3856791) (-1.1133333,2.3611850) (-1.1066667,2.3370865) (-1.1000000,2.3133764) (-1.0933333,2.2900479) (-1.0866667,2.2670944) (-1.0800000,2.2445092) (-1.0733333,2.2222859) (-1.0666667,2.2004183) (-1.0600000,2.1789000) (-1.0533333,2.1577252) (-1.0466667,2.1368877) (-1.0400000,2.1163819) (-1.0333333,2.0962020) (-1.0266667,2.0763425) (-1.0200000,2.0567978) (-1.0133333,2.0375627) (-1.0066667,2.0186318) (-1.0000000,2.0000000) (-0.9933333,1.9816623) (-0.9866667,1.9636137) (-0.9800000,1.9458493) (-0.9733333,1.9283645) (-0.9666667,1.9111545) (-0.9600000,1.8942149) (-0.9533333,1.8775410) (-0.9466667,1.8611286) (-0.9400000,1.8449734) (-0.9333333,1.8290710) (-0.9266667,1.8134175) (-0.9200000,1.7980087) (-0.9133333,1.7828407) (-0.9066667,1.7679095) (-0.9000000,1.7532114) (-0.8933333,1.7387427) (-0.8866667,1.7244996) (-0.8800000,1.7104785) (-0.8733333,1.6966760) (-0.8666667,1.6830886) (-0.8600000,1.6697128) (-0.8533333,1.6565454) (-0.8466667,1.6435831) (-0.8400000,1.6308228) (-0.8333333,1.6182612) (-0.8266667,1.6058952) (-0.8200000,1.5937220) (-0.8133333,1.5817385) (-0.8066667,1.5699418) (-0.8000000,1.5583292) (-0.7933333,1.5468977) (-0.7866667,1.5356447) (-0.7800000,1.5245675) (-0.7733333,1.5136634) (-0.7666667,1.5029300) (-0.7600000,1.4923646) (-0.7533333,1.4819647) (-0.7466667,1.4717281) (-0.7400000,1.4616521) (-0.7333333,1.4517346) (-0.7266667,1.4419733) (-0.7200000,1.4323658) (-0.7133333,1.4229100) (-0.7066667,1.4136038) (-0.7000000,1.4044449) (-0.6933333,1.3954313) (-0.6866667,1.3865610) (-0.6800000,1.3778320) (-0.6733333,1.3692423) (-0.6666667,1.3607900) (-0.6600000,1.3524732) (-0.6533333,1.3442900) (-0.6466667,1.3362387) (-0.6400000,1.3283175) (-0.6333333,1.3205245) (-0.6266667,1.3128582) (-0.6200000,1.3053168) (-0.6133333,1.2978987) (-0.6066667,1.2906022) (-0.6000000,1.2834259) (-0.5933333,1.2763681) (-0.5866667,1.2694273) (-0.5800000,1.2626021) (-0.5733333,1.2558909) (-0.5666667,1.2492923) (-0.5600000,1.2428050) (-0.5533333,1.2364276) (-0.5466667,1.2301587) (-0.5400000,1.2239970) (-0.5333333,1.2179412) (-0.5266667,1.2119900) (-0.5200000,1.2061422) (-0.5133333,1.2003966) (-0.5066667,1.1947519) (-0.5000000,1.1892071) (-0.4933333,1.1837610) (-0.4866667,1.1784123) (-0.4800000,1.1731602) (-0.4733333,1.1680034) (-0.4666667,1.1629409) (-0.4600000,1.1579717) (-0.4533333,1.1530948) (-0.4466667,1.1483092) (-0.4400000,1.1436139) (-0.4333333,1.1390079) (-0.4266667,1.1344904) (-0.4200000,1.1300605) (-0.4133333,1.1257172) (-0.4066667,1.1214597) (-0.4000000,1.1172871) (-0.3933333,1.1131987) (-0.3866667,1.1091935) (-0.3800000,1.1052709) (-0.3733333,1.1014300) (-0.3666667,1.0976700) (-0.3600000,1.0939903) (-0.3533333,1.0903902) (-0.3466667,1.0868688) (-0.3400000,1.0834255) (-0.3333333,1.0800597) (-0.3266667,1.0767707) (-0.3200000,1.0735579) (-0.3133333,1.0704206) (-0.3066667,1.0673582) (-0.3000000,1.0643702) (-0.2933333,1.0614559) (-0.2866667,1.0586148) (-0.2800000,1.0558464) (-0.2733333,1.0531501) (-0.2666667,1.0505254) (-0.2600000,1.0479719) (-0.2533333,1.0454889) (-0.2466667,1.0430761) (-0.2400000,1.0407330) (-0.2333333,1.0384591) (-0.2266667,1.0362541) (-0.2200000,1.0341174) (-0.2133333,1.0320487) (-0.2066667,1.0300477) (-0.2000000,1.0281138) (-0.1933333,1.0262468) (-0.1866667,1.0244464) (-0.1800000,1.0227120) (-0.1733333,1.0210436) (-0.1666667,1.0194406) (-0.1600000,1.0179029) (-0.1533333,1.0164302) (-0.1466667,1.0150221) (-0.1400000,1.0136784) (-0.1333333,1.0123989) (-0.1266667,1.0111832) (-0.1200000,1.0100313) (-0.1133333,1.0089428) (-0.1066667,1.0079177) (-0.1000000,1.0069556) (-0.0933333,1.0060563) (-0.0866667,1.0052199) (-0.0800000,1.0044460) (-0.0733333,1.0037345) (-0.0666667,1.0030854) (-0.0600000,1.0024984) (-0.0533333,1.0019736) (-0.0466667,1.0015107) (-0.0400000,1.0011097) (-0.0333333,1.0007705) (-0.0266667,1.0004930) (-0.0200000,1.0002773) (-0.0133333,1.0001232) (-0.0066667,1.0000308) (0.0000000,1.0000000) (0.0066667,1.0000308) (0.0133333,1.0001232) (0.0200000,1.0002773) (0.0266667,1.0004930) (0.0333333,1.0007705) (0.0400000,1.0011097) (0.0466667,1.0015107) (0.0533333,1.0019736) (0.0600000,1.0024984) (0.0666667,1.0030854) (0.0733333,1.0037345) (0.0800000,1.0044460) (0.0866667,1.0052199) (0.0933333,1.0060563) (0.1000000,1.0069556) (0.1066667,1.0079177) (0.1133333,1.0089428) (0.1200000,1.0100313) (0.1266667,1.0111832) (0.1333333,1.0123989) (0.1400000,1.0136784) (0.1466667,1.0150221) (0.1533333,1.0164302) (0.1600000,1.0179029) (0.1666667,1.0194406) (0.1733333,1.0210436) (0.1800000,1.0227120) (0.1866667,1.0244464) (0.1933333,1.0262468) (0.2000000,1.0281138) (0.2066667,1.0300477) (0.2133333,1.0320487) (0.2200000,1.0341174) (0.2266667,1.0362541) (0.2333333,1.0384591) (0.2400000,1.0407330) (0.2466667,1.0430761) (0.2533333,1.0454889) (0.2600000,1.0479719) (0.2666667,1.0505254) (0.2733333,1.0531501) (0.2800000,1.0558464) (0.2866667,1.0586148) (0.2933333,1.0614559) (0.3000000,1.0643702) (0.3066667,1.0673582) (0.3133333,1.0704206) (0.3200000,1.0735579) (0.3266667,1.0767707) (0.3333333,1.0800597) (0.3400000,1.0834255) (0.3466667,1.0868688) (0.3533333,1.0903902) (0.3600000,1.0939903) (0.3666667,1.0976700) (0.3733333,1.1014300) (0.3800000,1.1052709) (0.3866667,1.1091935) (0.3933333,1.1131987) (0.4000000,1.1172871) (0.4066667,1.1214597) (0.4133333,1.1257172) (0.4200000,1.1300605) (0.4266667,1.1344904) (0.4333333,1.1390079) (0.4400000,1.1436139) (0.4466667,1.1483092) (0.4533333,1.1530948) (0.4600000,1.1579717) (0.4666667,1.1629409) (0.4733333,1.1680034) (0.4800000,1.1731602) (0.4866667,1.1784123) (0.4933333,1.1837610) (0.5000000,1.1892071) (0.5066667,1.1947519) (0.5133333,1.2003966) (0.5200000,1.2061422) (0.5266667,1.2119900) (0.5333333,1.2179412) (0.5400000,1.2239970) (0.5466667,1.2301587) (0.5533333,1.2364276) (0.5600000,1.2428050) (0.5666667,1.2492923) (0.5733333,1.2558909) (0.5800000,1.2626021) (0.5866667,1.2694273) (0.5933333,1.2763681) (0.6000000,1.2834259) (0.6066667,1.2906022) (0.6133333,1.2978987) (0.6200000,1.3053168) (0.6266667,1.3128582) (0.6333333,1.3205245) (0.6400000,1.3283175) (0.6466667,1.3362387) (0.6533333,1.3442900) (0.6600000,1.3524732) (0.6666667,1.3607900) (0.6733333,1.3692423) (0.6800000,1.3778320) (0.6866667,1.3865610) (0.6933333,1.3954313) (0.7000000,1.4044449) (0.7066667,1.4136038) (0.7133333,1.4229100) (0.7200000,1.4323658) (0.7266667,1.4419733) (0.7333333,1.4517346) (0.7400000,1.4616521) (0.7466667,1.4717281) (0.7533333,1.4819647) (0.7600000,1.4923646) (0.7666667,1.5029300) (0.7733333,1.5136634) (0.7800000,1.5245675) (0.7866667,1.5356447) (0.7933333,1.5468977) (0.8000000,1.5583292) (0.8066667,1.5699418) (0.8133333,1.5817385) (0.8200000,1.5937220) (0.8266667,1.6058952) (0.8333333,1.6182612) (0.8400000,1.6308228) (0.8466667,1.6435831) (0.8533333,1.6565454) (0.8600000,1.6697128) (0.8666667,1.6830886) (0.8733333,1.6966760) (0.8800000,1.7104785) (0.8866667,1.7244996) (0.8933333,1.7387427) (0.9000000,1.7532114) (0.9066667,1.7679095) (0.9133333,1.7828407) (0.9200000,1.7980087) (0.9266667,1.8134175) (0.9333333,1.8290710) (0.9400000,1.8449734) (0.9466667,1.8611286) (0.9533333,1.8775410) (0.9600000,1.8942149) (0.9666667,1.9111545) (0.9733333,1.9283645) (0.9800000,1.9458493) (0.9866667,1.9636137) (0.9933333,1.9816623) (1.0000000,2.0000000) (1.0066667,2.0186318) (1.0133333,2.0375627) (1.0200000,2.0567978) (1.0266667,2.0763425) (1.0333333,2.0962020) (1.0400000,2.1163819) (1.0466667,2.1368877) (1.0533333,2.1577252) (1.0600000,2.1789000) (1.0666667,2.2004183) (1.0733333,2.2222859) (1.0800000,2.2445092) (1.0866667,2.2670944) (1.0933333,2.2900479) (1.1000000,2.3133764) (1.1066667,2.3370865) (1.1133333,2.3611850) (1.1200000,2.3856791) (1.1266667,2.4105758) (1.1333333,2.4358823) (1.1400000,2.4616062) (1.1466667,2.4877551) (1.1533333,2.5143366) (1.1600000,2.5413587) (1.1666667,2.5688295) (1.1733333,2.5967572) (1.1800000,2.6251503) (1.1866667,2.6540173) (1.1933333,2.6833671) (1.2000000,2.7132087) (1.2066667,2.7435511) (1.2133333,2.7744038) (1.2200000,2.8057763) (1.2266667,2.8376784) (1.2333333,2.8701201) (1.2400000,2.9031115) (1.2466667,2.9366631) (1.2533333,2.9707855) (1.2600000,3.0054895) (1.2666667,3.0407863) (1.2733333,3.0766872) (1.2800000,3.1132037) (1.2866667,3.1503478) (1.2933333,3.1881314) (1.3000000,3.2265670) (1.3066667,3.2656672) (1.3133333,3.3054449) (1.3200000,3.3459132) (1.3266667,3.3870856) (1.3333333,3.4289759) (1.3400000,3.4715982) (1.3466667,3.5149669) (1.3533333,3.5590967) (1.3600000,3.6040025) (1.3666667,3.6496997) (1.3733333,3.6962042) (1.3800000,3.7435318) (1.3866667,3.7916990) (1.3933333,3.8407226) (1.4000000,3.8906198) (1.4066667,3.9414080) (1.4133333,3.9931053) (1.4200000,4.0457299) (1.4266667,4.0993006) (1.4333333,4.1538365) (1.4400000,4.2093574) (1.4466667,4.2658831) (1.4533333,4.3234343) (1.4600000,4.3820319) (1.4666667,4.4416974) (1.4733333,4.5024527) (1.4800000,4.5643202) (1.4866667,4.6273229) (1.4933333,4.6914843) (1.5000000,4.7568285) (1.5066667,4.8233799) (1.5133333,4.8911638) (1.5200000,4.9602059)};
\draw[violet!75!black,densely dotted] (-1,0.5000000)--(1,0.5000000);
\fill[violet!75!black] (-1,0.5000000) circle(2.3pt);
\fill[violet!75!black] (1,0.5000000) circle(2.3pt);
\end{scope}
\end{scope}
\node[anchor=west,align=left,text width=4.35cm] at (8.9,4.4) {蓝实线：$a=2$\\[8pt]橙虚线：$a=\frac12$\\[14pt]正负输入给出同一指数。\\[8pt]底数改变峰谷形状，\\不改变左右对称。};
\node[text width=12.8cm,align=center] at (7.1,.65) {$D=\mathbb R$；$f(0)=1$。$a>1$ 时 $0<f\le1$；$0<a<1$ 时 $f\ge1$。};
\end{tikzpicture}\end{center}

\textbf{先看图：} 因为 $-(-x)^2=-x^2$，正负输入会先得到同一个指数。图中一条曲线中间高、另一条中间低，但两条曲线都左右同高。

\textbf{再说明一般情形：}
\[
f(-x)=a^{-(-x)^2}=a^{-x^2}=f(x).
\]
$a>1$ 时值域为 $(0,1]$；$0<a<1$ 时为 $[1,+\infty)$。不要只根据曲线是峰还是谷判断奇偶。
\end{explanationbox}
\clearpage
\begin{explanationbox}
\textbf{\textcolor{CExplanation}{真数左右相同，但零不能取}}
% Formula-based samples; see F083 checks/build_tikz_diagrams.py.
\begin{center}\begin{tikzpicture}[font=\small,>=stealth,line join=round]
\draw[gray!40,rounded corners=3pt] (0,0) rectangle (14.2,8.4);
\node at (7.1,7.75) {偶4：$y=\log_a|x|$};
\begin{scope}[shift={(4.400000,3.900000)},x=0.850000cm,y=0.750000cm]
\draw[->,gray] (-4,0)--(4,0) node[right] {$x$};
\draw[->,gray] (0,-3)--(0,3) node[above] {$y$};
\draw[gray] (-2,-.035/0.750000)--(-2,.035/0.750000);
\node[anchor=north] at (-2,-.07/0.750000) {$-2$};
\draw[gray] (2,-.035/0.750000)--(2,.035/0.750000);
\node[anchor=north] at (2,-.07/0.750000) {$2$};
\begin{scope}
\clip (-4,-3) rectangle (4,3);
\draw[blue,thick,solid] plot coordinates {(-4.0000000,2.0000000) (-3.9933333,1.9975935) (-3.9866667,1.9951830) (-3.9800000,1.9927684) (-3.9733333,1.9903498) (-3.9666667,1.9879272) (-3.9600000,1.9855004) (-3.9533333,1.9830696) (-3.9466667,1.9806347) (-3.9400000,1.9781956) (-3.9333333,1.9757525) (-3.9266667,1.9733051) (-3.9200000,1.9708537) (-3.9133333,1.9683980) (-3.9066667,1.9659382) (-3.9000000,1.9634741) (-3.8933333,1.9610059) (-3.8866667,1.9585334) (-3.8800000,1.9560567) (-3.8733333,1.9535757) (-3.8666667,1.9510904) (-3.8600000,1.9486008) (-3.8533333,1.9461070) (-3.8466667,1.9436088) (-3.8400000,1.9411063) (-3.8333333,1.9385995) (-3.8266667,1.9360882) (-3.8200000,1.9335726) (-3.8133333,1.9310526) (-3.8066667,1.9285282) (-3.8000000,1.9259994) (-3.7933333,1.9234662) (-3.7866667,1.9209284) (-3.7800000,1.9183862) (-3.7733333,1.9158396) (-3.7666667,1.9132884) (-3.7600000,1.9107327) (-3.7533333,1.9081724) (-3.7466667,1.9056076) (-3.7400000,1.9030383) (-3.7333333,1.9004643) (-3.7266667,1.8978858) (-3.7200000,1.8953026) (-3.7133333,1.8927148) (-3.7066667,1.8901224) (-3.7000000,1.8875253) (-3.6933333,1.8849235) (-3.6866667,1.8823170) (-3.6800000,1.8797058) (-3.6733333,1.8770898) (-3.6666667,1.8744691) (-3.6600000,1.8718436) (-3.6533333,1.8692134) (-3.6466667,1.8665783) (-3.6400000,1.8639385) (-3.6333333,1.8612937) (-3.6266667,1.8586442) (-3.6200000,1.8559897) (-3.6133333,1.8533304) (-3.6066667,1.8506661) (-3.6000000,1.8479969) (-3.5933333,1.8453228) (-3.5866667,1.8426437) (-3.5800000,1.8399596) (-3.5733333,1.8372705) (-3.5666667,1.8345764) (-3.5600000,1.8318772) (-3.5533333,1.8291730) (-3.5466667,1.8264637) (-3.5400000,1.8237494) (-3.5333333,1.8210299) (-3.5266667,1.8183052) (-3.5200000,1.8155754) (-3.5133333,1.8128405) (-3.5066667,1.8101003) (-3.5000000,1.8073549) (-3.4933333,1.8046043) (-3.4866667,1.8018484) (-3.4800000,1.7990873) (-3.4733333,1.7963209) (-3.4666667,1.7935491) (-3.4600000,1.7907720) (-3.4533333,1.7879896) (-3.4466667,1.7852018) (-3.4400000,1.7824086) (-3.4333333,1.7796099) (-3.4266667,1.7768059) (-3.4200000,1.7739963) (-3.4133333,1.7711813) (-3.4066667,1.7683608) (-3.4000000,1.7655347) (-3.3933333,1.7627032) (-3.3866667,1.7598660) (-3.3800000,1.7570232) (-3.3733333,1.7541749) (-3.3666667,1.7513209) (-3.3600000,1.7484612) (-3.3533333,1.7455959) (-3.3466667,1.7427249) (-3.3400000,1.7398481) (-3.3333333,1.7369656) (-3.3266667,1.7340773) (-3.3200000,1.7311832) (-3.3133333,1.7282834) (-3.3066667,1.7253776) (-3.3000000,1.7224660) (-3.2933333,1.7195485) (-3.2866667,1.7166251) (-3.2800000,1.7136958) (-3.2733333,1.7107605) (-3.2666667,1.7078192) (-3.2600000,1.7048720) (-3.2533333,1.7019186) (-3.2466667,1.6989593) (-3.2400000,1.6959938) (-3.2333333,1.6930222) (-3.2266667,1.6900445) (-3.2200000,1.6870607) (-3.2133333,1.6840706) (-3.2066667,1.6810744) (-3.2000000,1.6780719) (-3.1933333,1.6750632) (-3.1866667,1.6720481) (-3.1800000,1.6690268) (-3.1733333,1.6659991) (-3.1666667,1.6629650) (-3.1600000,1.6599246) (-3.1533333,1.6568777) (-3.1466667,1.6538244) (-3.1400000,1.6507646) (-3.1333333,1.6476983) (-3.1266667,1.6446254) (-3.1200000,1.6415460) (-3.1133333,1.6384600) (-3.1066667,1.6353675) (-3.1000000,1.6322682) (-3.0933333,1.6291623) (-3.0866667,1.6260497) (-3.0800000,1.6229304) (-3.0733333,1.6198042) (-3.0666667,1.6166714) (-3.0600000,1.6135317) (-3.0533333,1.6103851) (-3.0466667,1.6072317) (-3.0400000,1.6040713) (-3.0333333,1.6009040) (-3.0266667,1.5977298) (-3.0200000,1.5945485) (-3.0133333,1.5913603) (-3.0066667,1.5881649) (-3.0000000,1.5849625) (-2.9933333,1.5817529) (-2.9866667,1.5785362) (-2.9800000,1.5753123) (-2.9733333,1.5720812) (-2.9666667,1.5688428) (-2.9600000,1.5655972) (-2.9533333,1.5623442) (-2.9466667,1.5590839) (-2.9400000,1.5558162) (-2.9333333,1.5525410) (-2.9266667,1.5492584) (-2.9200000,1.5459684) (-2.9133333,1.5426708) (-2.9066667,1.5393656) (-2.9000000,1.5360529) (-2.8933333,1.5327325) (-2.8866667,1.5294045) (-2.8800000,1.5260688) (-2.8733333,1.5227254) (-2.8666667,1.5193742) (-2.8600000,1.5160151) (-2.8533333,1.5126483) (-2.8466667,1.5092736) (-2.8400000,1.5058909) (-2.8333333,1.5025003) (-2.8266667,1.4991018) (-2.8200000,1.4956952) (-2.8133333,1.4922805) (-2.8066667,1.4888577) (-2.8000000,1.4854268) (-2.7933333,1.4819877) (-2.7866667,1.4785404) (-2.7800000,1.4750849) (-2.7733333,1.4716210) (-2.7666667,1.4681488) (-2.7600000,1.4646683) (-2.7533333,1.4611793) (-2.7466667,1.4576818) (-2.7400000,1.4541759) (-2.7333333,1.4506614) (-2.7266667,1.4471383) (-2.7200000,1.4436067) (-2.7133333,1.4400663) (-2.7066667,1.4365172) (-2.7000000,1.4329594) (-2.6933333,1.4293928) (-2.6866667,1.4258173) (-2.6800000,1.4222330) (-2.6733333,1.4186397) (-2.6666667,1.4150375) (-2.6600000,1.4114262) (-2.6533333,1.4078059) (-2.6466667,1.4041765) (-2.6400000,1.4005379) (-2.6333333,1.3968902) (-2.6266667,1.3932331) (-2.6200000,1.3895668) (-2.6133333,1.3858912) (-2.6066667,1.3822061) (-2.6000000,1.3785116) (-2.5933333,1.3748077) (-2.5866667,1.3710942) (-2.5800000,1.3673711) (-2.5733333,1.3636383) (-2.5666667,1.3598959) (-2.5600000,1.3561438) (-2.5533333,1.3523819) (-2.5466667,1.3486101) (-2.5400000,1.3448285) (-2.5333333,1.3410369) (-2.5266667,1.3372353) (-2.5200000,1.3334237) (-2.5133333,1.3296020) (-2.5066667,1.3257702) (-2.5000000,1.3219281) (-2.4933333,1.3180758) (-2.4866667,1.3142131) (-2.4800000,1.3103401) (-2.4733333,1.3064567) (-2.4666667,1.3025628) (-2.4600000,1.2986583) (-2.4533333,1.2947433) (-2.4466667,1.2908176) (-2.4400000,1.2868811) (-2.4333333,1.2829340) (-2.4266667,1.2789759) (-2.4200000,1.2750070) (-2.4133333,1.2710272) (-2.4066667,1.2670363) (-2.4000000,1.2630344) (-2.3933333,1.2590213) (-2.3866667,1.2549971) (-2.3800000,1.2509616) (-2.3733333,1.2469147) (-2.3666667,1.2428565) (-2.3600000,1.2387869) (-2.3533333,1.2347057) (-2.3466667,1.2306129) (-2.3400000,1.2265085) (-2.3333333,1.2223924) (-2.3266667,1.2182645) (-2.3200000,1.2141248) (-2.3133333,1.2099732) (-2.3066667,1.2058095) (-2.3000000,1.2016339) (-2.2933333,1.1974461) (-2.2866667,1.1932461) (-2.2800000,1.1890338) (-2.2733333,1.1848092) (-2.2666667,1.1805722) (-2.2600000,1.1763228) (-2.2533333,1.1720607) (-2.2466667,1.1677861) (-2.2400000,1.1634987) (-2.2333333,1.1591986) (-2.2266667,1.1548856) (-2.2200000,1.1505597) (-2.2133333,1.1462207) (-2.2066667,1.1418687) (-2.2000000,1.1375035) (-2.1933333,1.1331251) (-2.1866667,1.1287333) (-2.1800000,1.1243281) (-2.1733333,1.1199095) (-2.1666667,1.1154772) (-2.1600000,1.1110313) (-2.1533333,1.1065717) (-2.1466667,1.1020982) (-2.1400000,1.0976108) (-2.1333333,1.0931094) (-2.1266667,1.0885939) (-2.1200000,1.0840643) (-2.1133333,1.0795203) (-2.1066667,1.0749621) (-2.1000000,1.0703893) (-2.0933333,1.0658021) (-2.0866667,1.0612002) (-2.0800000,1.0565835) (-2.0733333,1.0519521) (-2.0666667,1.0473057) (-2.0600000,1.0426443) (-2.0533333,1.0379679) (-2.0466667,1.0332762) (-2.0400000,1.0285692) (-2.0333333,1.0238467) (-2.0266667,1.0191088) (-2.0200000,1.0143553) (-2.0133333,1.0095860) (-2.0066667,1.0048010) (-2.0000000,1.0000000) (-1.9933333,0.9951830) (-1.9866667,0.9903498) (-1.9800000,0.9855004) (-1.9733333,0.9806347) (-1.9666667,0.9757525) (-1.9600000,0.9708537) (-1.9533333,0.9659382) (-1.9466667,0.9610059) (-1.9400000,0.9560567) (-1.9333333,0.9510904) (-1.9266667,0.9461070) (-1.9200000,0.9411063) (-1.9133333,0.9360882) (-1.9066667,0.9310526) (-1.9000000,0.9259994) (-1.8933333,0.9209284) (-1.8866667,0.9158396) (-1.8800000,0.9107327) (-1.8733333,0.9056076) (-1.8666667,0.9004643) (-1.8600000,0.8953026) (-1.8533333,0.8901224) (-1.8466667,0.8849235) (-1.8400000,0.8797058) (-1.8333333,0.8744691) (-1.8266667,0.8692134) (-1.8200000,0.8639385) (-1.8133333,0.8586442) (-1.8066667,0.8533304) (-1.8000000,0.8479969) (-1.7933333,0.8426437) (-1.7866667,0.8372705) (-1.7800000,0.8318772) (-1.7733333,0.8264637) (-1.7666667,0.8210299) (-1.7600000,0.8155754) (-1.7533333,0.8101003) (-1.7466667,0.8046043) (-1.7400000,0.7990873) (-1.7333333,0.7935491) (-1.7266667,0.7879896) (-1.7200000,0.7824086) (-1.7133333,0.7768059) (-1.7066667,0.7711813) (-1.7000000,0.7655347) (-1.6933333,0.7598660) (-1.6866667,0.7541749) (-1.6800000,0.7484612) (-1.6733333,0.7427249) (-1.6666667,0.7369656) (-1.6600000,0.7311832) (-1.6533333,0.7253776) (-1.6466667,0.7195485) (-1.6400000,0.7136958) (-1.6333333,0.7078192) (-1.6266667,0.7019186) (-1.6200000,0.6959938) (-1.6133333,0.6900445) (-1.6066667,0.6840706) (-1.6000000,0.6780719) (-1.5933333,0.6720481) (-1.5866667,0.6659991) (-1.5800000,0.6599246) (-1.5733333,0.6538244) (-1.5666667,0.6476983) (-1.5600000,0.6415460) (-1.5533333,0.6353675) (-1.5466667,0.6291623) (-1.5400000,0.6229304) (-1.5333333,0.6166714) (-1.5266667,0.6103851) (-1.5200000,0.6040713) (-1.5133333,0.5977298) (-1.5066667,0.5913603) (-1.5000000,0.5849625) (-1.4933333,0.5785362) (-1.4866667,0.5720812) (-1.4800000,0.5655972) (-1.4733333,0.5590839) (-1.4666667,0.5525410) (-1.4600000,0.5459684) (-1.4533333,0.5393656) (-1.4466667,0.5327325) (-1.4400000,0.5260688) (-1.4333333,0.5193742) (-1.4266667,0.5126483) (-1.4200000,0.5058909) (-1.4133333,0.4991018) (-1.4066667,0.4922805) (-1.4000000,0.4854268) (-1.3933333,0.4785404) (-1.3866667,0.4716210) (-1.3800000,0.4646683) (-1.3733333,0.4576818) (-1.3666667,0.4506614) (-1.3600000,0.4436067) (-1.3533333,0.4365172) (-1.3466667,0.4293928) (-1.3400000,0.4222330) (-1.3333333,0.4150375) (-1.3266667,0.4078059) (-1.3200000,0.4005379) (-1.3133333,0.3932331) (-1.3066667,0.3858912) (-1.3000000,0.3785116) (-1.2933333,0.3710942) (-1.2866667,0.3636383) (-1.2800000,0.3561438) (-1.2733333,0.3486101) (-1.2666667,0.3410369) (-1.2600000,0.3334237) (-1.2533333,0.3257702) (-1.2466667,0.3180758) (-1.2400000,0.3103401) (-1.2333333,0.3025628) (-1.2266667,0.2947433) (-1.2200000,0.2868811) (-1.2133333,0.2789759) (-1.2066667,0.2710272) (-1.2000000,0.2630344) (-1.1933333,0.2549971) (-1.1866667,0.2469147) (-1.1800000,0.2387869) (-1.1733333,0.2306129) (-1.1666667,0.2223924) (-1.1600000,0.2141248) (-1.1533333,0.2058095) (-1.1466667,0.1974461) (-1.1400000,0.1890338) (-1.1333333,0.1805722) (-1.1266667,0.1720607) (-1.1200000,0.1634987) (-1.1133333,0.1548856) (-1.1066667,0.1462207) (-1.1000000,0.1375035) (-1.0933333,0.1287333) (-1.0866667,0.1199095) (-1.0800000,0.1110313) (-1.0733333,0.1020982) (-1.0666667,0.0931094) (-1.0600000,0.0840643) (-1.0533333,0.0749621) (-1.0466667,0.0658021) (-1.0400000,0.0565835) (-1.0333333,0.0473057) (-1.0266667,0.0379679) (-1.0200000,0.0285692) (-1.0133333,0.0191088) (-1.0066667,0.0095860) (-1.0000000,0.0000000) (-0.9933333,-0.0096502) (-0.9866667,-0.0193653) (-0.9800000,-0.0291463) (-0.9733333,-0.0389941) (-0.9666667,-0.0489096) (-0.9600000,-0.0588937) (-0.9533333,-0.0689474) (-0.9466667,-0.0790716) (-0.9400000,-0.0892673) (-0.9333333,-0.0995357) (-0.9266667,-0.1098776) (-0.9200000,-0.1202942) (-0.9133333,-0.1307866) (-0.9066667,-0.1413558) (-0.9000000,-0.1520031) (-0.8933333,-0.1627295) (-0.8866667,-0.1735363) (-0.8800000,-0.1844246) (-0.8733333,-0.1953957) (-0.8666667,-0.2064509) (-0.8600000,-0.2175914) (-0.8533333,-0.2288187) (-0.8466667,-0.2401340) (-0.8400000,-0.2515388) (-0.8333333,-0.2630344) (-0.8266667,-0.2746224) (-0.8200000,-0.2863042) (-0.8133333,-0.2980814) (-0.8066667,-0.3099555) (-0.8000000,-0.3219281) (-0.7933333,-0.3340009) (-0.7866667,-0.3461756) (-0.7800000,-0.3584540) (-0.7733333,-0.3708377) (-0.7666667,-0.3833286) (-0.7600000,-0.3959287) (-0.7533333,-0.4086397) (-0.7466667,-0.4214638) (-0.7400000,-0.4344028) (-0.7333333,-0.4474590) (-0.7266667,-0.4606344) (-0.7200000,-0.4739312) (-0.7133333,-0.4873517) (-0.7066667,-0.5008982) (-0.7000000,-0.5145732) (-0.6933333,-0.5283790) (-0.6866667,-0.5423182) (-0.6800000,-0.5563933) (-0.6733333,-0.5706072) (-0.6666667,-0.5849625) (-0.6600000,-0.5994621) (-0.6533333,-0.6141088) (-0.6466667,-0.6289058) (-0.6400000,-0.6438562) (-0.6333333,-0.6589631) (-0.6266667,-0.6742298) (-0.6200000,-0.6896599) (-0.6133333,-0.7052567) (-0.6066667,-0.7210241) (-0.6000000,-0.7369656) (-0.5933333,-0.7530853) (-0.5866667,-0.7693871) (-0.5800000,-0.7858752) (-0.5733333,-0.8025539) (-0.5666667,-0.8194278) (-0.5600000,-0.8365013) (-0.5533333,-0.8537793) (-0.5466667,-0.8712667) (-0.5400000,-0.8889687) (-0.5333333,-0.9068906) (-0.5266667,-0.9250379) (-0.5200000,-0.9434165) (-0.5133333,-0.9620321) (-0.5066667,-0.9808912) (-0.5000000,-1.0000000) (-0.4933333,-1.0193653) (-0.4866667,-1.0389941) (-0.4800000,-1.0588937) (-0.4733333,-1.0790716) (-0.4666667,-1.0995357) (-0.4600000,-1.1202942) (-0.4533333,-1.1413558) (-0.4466667,-1.1627295) (-0.4400000,-1.1844246) (-0.4333333,-1.2064509) (-0.4266667,-1.2288187) (-0.4200000,-1.2515388) (-0.4133333,-1.2746224) (-0.4066667,-1.2980814) (-0.4000000,-1.3219281) (-0.3933333,-1.3461756) (-0.3866667,-1.3708377) (-0.3800000,-1.3959287) (-0.3733333,-1.4214638) (-0.3666667,-1.4474590) (-0.3600000,-1.4739312) (-0.3533333,-1.5008982) (-0.3466667,-1.5283790) (-0.3400000,-1.5563933) (-0.3333333,-1.5849625) (-0.3266667,-1.6141088) (-0.3200000,-1.6438562) (-0.3133333,-1.6742298) (-0.3066667,-1.7052567) (-0.3000000,-1.7369656) (-0.2933333,-1.7693871) (-0.2866667,-1.8025539) (-0.2800000,-1.8365013) (-0.2733333,-1.8712667) (-0.2666667,-1.9068906) (-0.2600000,-1.9434165) (-0.2533333,-1.9808912) (-0.2466667,-2.0193653) (-0.2400000,-2.0588937) (-0.2333333,-2.0995357) (-0.2266667,-2.1413558) (-0.2200000,-2.1844246) (-0.2133333,-2.2288187) (-0.2066667,-2.2746224) (-0.2000000,-2.3219281) (-0.1933333,-2.3708377) (-0.1866667,-2.4214638) (-0.1800000,-2.4739312) (-0.1733333,-2.5283790) (-0.1666667,-2.5849625) (-0.1600000,-2.6438562) (-0.1533333,-2.7052567) (-0.1466667,-2.7693871) (-0.1400000,-2.8365013) (-0.1333333,-2.9068906) (-0.1266667,-2.9808912)};
\draw[blue,thick,solid] plot coordinates {(0.1266667,-2.9808912) (0.1333333,-2.9068906) (0.1400000,-2.8365013) (0.1466667,-2.7693871) (0.1533333,-2.7052567) (0.1600000,-2.6438562) (0.1666667,-2.5849625) (0.1733333,-2.5283790) (0.1800000,-2.4739312) (0.1866667,-2.4214638) (0.1933333,-2.3708377) (0.2000000,-2.3219281) (0.2066667,-2.2746224) (0.2133333,-2.2288187) (0.2200000,-2.1844246) (0.2266667,-2.1413558) (0.2333333,-2.0995357) (0.2400000,-2.0588937) (0.2466667,-2.0193653) (0.2533333,-1.9808912) (0.2600000,-1.9434165) (0.2666667,-1.9068906) (0.2733333,-1.8712667) (0.2800000,-1.8365013) (0.2866667,-1.8025539) (0.2933333,-1.7693871) (0.3000000,-1.7369656) (0.3066667,-1.7052567) (0.3133333,-1.6742298) (0.3200000,-1.6438562) (0.3266667,-1.6141088) (0.3333333,-1.5849625) (0.3400000,-1.5563933) (0.3466667,-1.5283790) (0.3533333,-1.5008982) (0.3600000,-1.4739312) (0.3666667,-1.4474590) (0.3733333,-1.4214638) (0.3800000,-1.3959287) (0.3866667,-1.3708377) (0.3933333,-1.3461756) (0.4000000,-1.3219281) (0.4066667,-1.2980814) (0.4133333,-1.2746224) (0.4200000,-1.2515388) (0.4266667,-1.2288187) (0.4333333,-1.2064509) (0.4400000,-1.1844246) (0.4466667,-1.1627295) (0.4533333,-1.1413558) (0.4600000,-1.1202942) (0.4666667,-1.0995357) (0.4733333,-1.0790716) (0.4800000,-1.0588937) (0.4866667,-1.0389941) (0.4933333,-1.0193653) (0.5000000,-1.0000000) (0.5066667,-0.9808912) (0.5133333,-0.9620321) (0.5200000,-0.9434165) (0.5266667,-0.9250379) (0.5333333,-0.9068906) (0.5400000,-0.8889687) (0.5466667,-0.8712667) (0.5533333,-0.8537793) (0.5600000,-0.8365013) (0.5666667,-0.8194278) (0.5733333,-0.8025539) (0.5800000,-0.7858752) (0.5866667,-0.7693871) (0.5933333,-0.7530853) (0.6000000,-0.7369656) (0.6066667,-0.7210241) (0.6133333,-0.7052567) (0.6200000,-0.6896599) (0.6266667,-0.6742298) (0.6333333,-0.6589631) (0.6400000,-0.6438562) (0.6466667,-0.6289058) (0.6533333,-0.6141088) (0.6600000,-0.5994621) (0.6666667,-0.5849625) (0.6733333,-0.5706072) (0.6800000,-0.5563933) (0.6866667,-0.5423182) (0.6933333,-0.5283790) (0.7000000,-0.5145732) (0.7066667,-0.5008982) (0.7133333,-0.4873517) (0.7200000,-0.4739312) (0.7266667,-0.4606344) (0.7333333,-0.4474590) (0.7400000,-0.4344028) (0.7466667,-0.4214638) (0.7533333,-0.4086397) (0.7600000,-0.3959287) (0.7666667,-0.3833286) (0.7733333,-0.3708377) (0.7800000,-0.3584540) (0.7866667,-0.3461756) (0.7933333,-0.3340009) (0.8000000,-0.3219281) (0.8066667,-0.3099555) (0.8133333,-0.2980814) (0.8200000,-0.2863042) (0.8266667,-0.2746224) (0.8333333,-0.2630344) (0.8400000,-0.2515388) (0.8466667,-0.2401340) (0.8533333,-0.2288187) (0.8600000,-0.2175914) (0.8666667,-0.2064509) (0.8733333,-0.1953957) (0.8800000,-0.1844246) (0.8866667,-0.1735363) (0.8933333,-0.1627295) (0.9000000,-0.1520031) (0.9066667,-0.1413558) (0.9133333,-0.1307866) (0.9200000,-0.1202942) (0.9266667,-0.1098776) (0.9333333,-0.0995357) (0.9400000,-0.0892673) (0.9466667,-0.0790716) (0.9533333,-0.0689474) (0.9600000,-0.0588937) (0.9666667,-0.0489096) (0.9733333,-0.0389941) (0.9800000,-0.0291463) (0.9866667,-0.0193653) (0.9933333,-0.0096502) (1.0000000,0.0000000) (1.0066667,0.0095860) (1.0133333,0.0191088) (1.0200000,0.0285692) (1.0266667,0.0379679) (1.0333333,0.0473057) (1.0400000,0.0565835) (1.0466667,0.0658021) (1.0533333,0.0749621) (1.0600000,0.0840643) (1.0666667,0.0931094) (1.0733333,0.1020982) (1.0800000,0.1110313) (1.0866667,0.1199095) (1.0933333,0.1287333) (1.1000000,0.1375035) (1.1066667,0.1462207) (1.1133333,0.1548856) (1.1200000,0.1634987) (1.1266667,0.1720607) (1.1333333,0.1805722) (1.1400000,0.1890338) (1.1466667,0.1974461) (1.1533333,0.2058095) (1.1600000,0.2141248) (1.1666667,0.2223924) (1.1733333,0.2306129) (1.1800000,0.2387869) (1.1866667,0.2469147) (1.1933333,0.2549971) (1.2000000,0.2630344) (1.2066667,0.2710272) (1.2133333,0.2789759) (1.2200000,0.2868811) (1.2266667,0.2947433) (1.2333333,0.3025628) (1.2400000,0.3103401) (1.2466667,0.3180758) (1.2533333,0.3257702) (1.2600000,0.3334237) (1.2666667,0.3410369) (1.2733333,0.3486101) (1.2800000,0.3561438) (1.2866667,0.3636383) (1.2933333,0.3710942) (1.3000000,0.3785116) (1.3066667,0.3858912) (1.3133333,0.3932331) (1.3200000,0.4005379) (1.3266667,0.4078059) (1.3333333,0.4150375) (1.3400000,0.4222330) (1.3466667,0.4293928) (1.3533333,0.4365172) (1.3600000,0.4436067) (1.3666667,0.4506614) (1.3733333,0.4576818) (1.3800000,0.4646683) (1.3866667,0.4716210) (1.3933333,0.4785404) (1.4000000,0.4854268) (1.4066667,0.4922805) (1.4133333,0.4991018) (1.4200000,0.5058909) (1.4266667,0.5126483) (1.4333333,0.5193742) (1.4400000,0.5260688) (1.4466667,0.5327325) (1.4533333,0.5393656) (1.4600000,0.5459684) (1.4666667,0.5525410) (1.4733333,0.5590839) (1.4800000,0.5655972) (1.4866667,0.5720812) (1.4933333,0.5785362) (1.5000000,0.5849625) (1.5066667,0.5913603) (1.5133333,0.5977298) (1.5200000,0.6040713) (1.5266667,0.6103851) (1.5333333,0.6166714) (1.5400000,0.6229304) (1.5466667,0.6291623) (1.5533333,0.6353675) (1.5600000,0.6415460) (1.5666667,0.6476983) (1.5733333,0.6538244) (1.5800000,0.6599246) (1.5866667,0.6659991) (1.5933333,0.6720481) (1.6000000,0.6780719) (1.6066667,0.6840706) (1.6133333,0.6900445) (1.6200000,0.6959938) (1.6266667,0.7019186) (1.6333333,0.7078192) (1.6400000,0.7136958) (1.6466667,0.7195485) (1.6533333,0.7253776) (1.6600000,0.7311832) (1.6666667,0.7369656) (1.6733333,0.7427249) (1.6800000,0.7484612) (1.6866667,0.7541749) (1.6933333,0.7598660) (1.7000000,0.7655347) (1.7066667,0.7711813) (1.7133333,0.7768059) (1.7200000,0.7824086) (1.7266667,0.7879896) (1.7333333,0.7935491) (1.7400000,0.7990873) (1.7466667,0.8046043) (1.7533333,0.8101003) (1.7600000,0.8155754) (1.7666667,0.8210299) (1.7733333,0.8264637) (1.7800000,0.8318772) (1.7866667,0.8372705) (1.7933333,0.8426437) (1.8000000,0.8479969) (1.8066667,0.8533304) (1.8133333,0.8586442) (1.8200000,0.8639385) (1.8266667,0.8692134) (1.8333333,0.8744691) (1.8400000,0.8797058) (1.8466667,0.8849235) (1.8533333,0.8901224) (1.8600000,0.8953026) (1.8666667,0.9004643) (1.8733333,0.9056076) (1.8800000,0.9107327) (1.8866667,0.9158396) (1.8933333,0.9209284) (1.9000000,0.9259994) (1.9066667,0.9310526) (1.9133333,0.9360882) (1.9200000,0.9411063) (1.9266667,0.9461070) (1.9333333,0.9510904) (1.9400000,0.9560567) (1.9466667,0.9610059) (1.9533333,0.9659382) (1.9600000,0.9708537) (1.9666667,0.9757525) (1.9733333,0.9806347) (1.9800000,0.9855004) (1.9866667,0.9903498) (1.9933333,0.9951830) (2.0000000,1.0000000) (2.0066667,1.0048010) (2.0133333,1.0095860) (2.0200000,1.0143553) (2.0266667,1.0191088) (2.0333333,1.0238467) (2.0400000,1.0285692) (2.0466667,1.0332762) (2.0533333,1.0379679) (2.0600000,1.0426443) (2.0666667,1.0473057) (2.0733333,1.0519521) (2.0800000,1.0565835) (2.0866667,1.0612002) (2.0933333,1.0658021) (2.1000000,1.0703893) (2.1066667,1.0749621) (2.1133333,1.0795203) (2.1200000,1.0840643) (2.1266667,1.0885939) (2.1333333,1.0931094) (2.1400000,1.0976108) (2.1466667,1.1020982) (2.1533333,1.1065717) (2.1600000,1.1110313) (2.1666667,1.1154772) (2.1733333,1.1199095) (2.1800000,1.1243281) (2.1866667,1.1287333) (2.1933333,1.1331251) (2.2000000,1.1375035) (2.2066667,1.1418687) (2.2133333,1.1462207) (2.2200000,1.1505597) (2.2266667,1.1548856) (2.2333333,1.1591986) (2.2400000,1.1634987) (2.2466667,1.1677861) (2.2533333,1.1720607) (2.2600000,1.1763228) (2.2666667,1.1805722) (2.2733333,1.1848092) (2.2800000,1.1890338) (2.2866667,1.1932461) (2.2933333,1.1974461) (2.3000000,1.2016339) (2.3066667,1.2058095) (2.3133333,1.2099732) (2.3200000,1.2141248) (2.3266667,1.2182645) (2.3333333,1.2223924) (2.3400000,1.2265085) (2.3466667,1.2306129) (2.3533333,1.2347057) (2.3600000,1.2387869) (2.3666667,1.2428565) (2.3733333,1.2469147) (2.3800000,1.2509616) (2.3866667,1.2549971) (2.3933333,1.2590213) (2.4000000,1.2630344) (2.4066667,1.2670363) (2.4133333,1.2710272) (2.4200000,1.2750070) (2.4266667,1.2789759) (2.4333333,1.2829340) (2.4400000,1.2868811) (2.4466667,1.2908176) (2.4533333,1.2947433) (2.4600000,1.2986583) (2.4666667,1.3025628) (2.4733333,1.3064567) (2.4800000,1.3103401) (2.4866667,1.3142131) (2.4933333,1.3180758) (2.5000000,1.3219281) (2.5066667,1.3257702) (2.5133333,1.3296020) (2.5200000,1.3334237) (2.5266667,1.3372353) (2.5333333,1.3410369) (2.5400000,1.3448285) (2.5466667,1.3486101) (2.5533333,1.3523819) (2.5600000,1.3561438) (2.5666667,1.3598959) (2.5733333,1.3636383) (2.5800000,1.3673711) (2.5866667,1.3710942) (2.5933333,1.3748077) (2.6000000,1.3785116) (2.6066667,1.3822061) (2.6133333,1.3858912) (2.6200000,1.3895668) (2.6266667,1.3932331) (2.6333333,1.3968902) (2.6400000,1.4005379) (2.6466667,1.4041765) (2.6533333,1.4078059) (2.6600000,1.4114262) (2.6666667,1.4150375) (2.6733333,1.4186397) (2.6800000,1.4222330) (2.6866667,1.4258173) (2.6933333,1.4293928) (2.7000000,1.4329594) (2.7066667,1.4365172) (2.7133333,1.4400663) (2.7200000,1.4436067) (2.7266667,1.4471383) (2.7333333,1.4506614) (2.7400000,1.4541759) (2.7466667,1.4576818) (2.7533333,1.4611793) (2.7600000,1.4646683) (2.7666667,1.4681488) (2.7733333,1.4716210) (2.7800000,1.4750849) (2.7866667,1.4785404) (2.7933333,1.4819877) (2.8000000,1.4854268) (2.8066667,1.4888577) (2.8133333,1.4922805) (2.8200000,1.4956952) (2.8266667,1.4991018) (2.8333333,1.5025003) (2.8400000,1.5058909) (2.8466667,1.5092736) (2.8533333,1.5126483) (2.8600000,1.5160151) (2.8666667,1.5193742) (2.8733333,1.5227254) (2.8800000,1.5260688) (2.8866667,1.5294045) (2.8933333,1.5327325) (2.9000000,1.5360529) (2.9066667,1.5393656) (2.9133333,1.5426708) (2.9200000,1.5459684) (2.9266667,1.5492584) (2.9333333,1.5525410) (2.9400000,1.5558162) (2.9466667,1.5590839) (2.9533333,1.5623442) (2.9600000,1.5655972) (2.9666667,1.5688428) (2.9733333,1.5720812) (2.9800000,1.5753123) (2.9866667,1.5785362) (2.9933333,1.5817529) (3.0000000,1.5849625) (3.0066667,1.5881649) (3.0133333,1.5913603) (3.0200000,1.5945485) (3.0266667,1.5977298) (3.0333333,1.6009040) (3.0400000,1.6040713) (3.0466667,1.6072317) (3.0533333,1.6103851) (3.0600000,1.6135317) (3.0666667,1.6166714) (3.0733333,1.6198042) (3.0800000,1.6229304) (3.0866667,1.6260497) (3.0933333,1.6291623) (3.1000000,1.6322682) (3.1066667,1.6353675) (3.1133333,1.6384600) (3.1200000,1.6415460) (3.1266667,1.6446254) (3.1333333,1.6476983) (3.1400000,1.6507646) (3.1466667,1.6538244) (3.1533333,1.6568777) (3.1600000,1.6599246) (3.1666667,1.6629650) (3.1733333,1.6659991) (3.1800000,1.6690268) (3.1866667,1.6720481) (3.1933333,1.6750632) (3.2000000,1.6780719) (3.2066667,1.6810744) (3.2133333,1.6840706) (3.2200000,1.6870607) (3.2266667,1.6900445) (3.2333333,1.6930222) (3.2400000,1.6959938) (3.2466667,1.6989593) (3.2533333,1.7019186) (3.2600000,1.7048720) (3.2666667,1.7078192) (3.2733333,1.7107605) (3.2800000,1.7136958) (3.2866667,1.7166251) (3.2933333,1.7195485) (3.3000000,1.7224660) (3.3066667,1.7253776) (3.3133333,1.7282834) (3.3200000,1.7311832) (3.3266667,1.7340773) (3.3333333,1.7369656) (3.3400000,1.7398481) (3.3466667,1.7427249) (3.3533333,1.7455959) (3.3600000,1.7484612) (3.3666667,1.7513209) (3.3733333,1.7541749) (3.3800000,1.7570232) (3.3866667,1.7598660) (3.3933333,1.7627032) (3.4000000,1.7655347) (3.4066667,1.7683608) (3.4133333,1.7711813) (3.4200000,1.7739963) (3.4266667,1.7768059) (3.4333333,1.7796099) (3.4400000,1.7824086) (3.4466667,1.7852018) (3.4533333,1.7879896) (3.4600000,1.7907720) (3.4666667,1.7935491) (3.4733333,1.7963209) (3.4800000,1.7990873) (3.4866667,1.8018484) (3.4933333,1.8046043) (3.5000000,1.8073549) (3.5066667,1.8101003) (3.5133333,1.8128405) (3.5200000,1.8155754) (3.5266667,1.8183052) (3.5333333,1.8210299) (3.5400000,1.8237494) (3.5466667,1.8264637) (3.5533333,1.8291730) (3.5600000,1.8318772) (3.5666667,1.8345764) (3.5733333,1.8372705) (3.5800000,1.8399596) (3.5866667,1.8426437) (3.5933333,1.8453228) (3.6000000,1.8479969) (3.6066667,1.8506661) (3.6133333,1.8533304) (3.6200000,1.8559897) (3.6266667,1.8586442) (3.6333333,1.8612937) (3.6400000,1.8639385) (3.6466667,1.8665783) (3.6533333,1.8692134) (3.6600000,1.8718436) (3.6666667,1.8744691) (3.6733333,1.8770898) (3.6800000,1.8797058) (3.6866667,1.8823170) (3.6933333,1.8849235) (3.7000000,1.8875253) (3.7066667,1.8901224) (3.7133333,1.8927148) (3.7200000,1.8953026) (3.7266667,1.8978858) (3.7333333,1.9004643) (3.7400000,1.9030383) (3.7466667,1.9056076) (3.7533333,1.9081724) (3.7600000,1.9107327) (3.7666667,1.9132884) (3.7733333,1.9158396) (3.7800000,1.9183862) (3.7866667,1.9209284) (3.7933333,1.9234662) (3.8000000,1.9259994) (3.8066667,1.9285282) (3.8133333,1.9310526) (3.8200000,1.9335726) (3.8266667,1.9360882) (3.8333333,1.9385995) (3.8400000,1.9411063) (3.8466667,1.9436088) (3.8533333,1.9461070) (3.8600000,1.9486008) (3.8666667,1.9510904) (3.8733333,1.9535757) (3.8800000,1.9560567) (3.8866667,1.9585334) (3.8933333,1.9610059) (3.9000000,1.9634741) (3.9066667,1.9659382) (3.9133333,1.9683980) (3.9200000,1.9708537) (3.9266667,1.9733051) (3.9333333,1.9757525) (3.9400000,1.9781956) (3.9466667,1.9806347) (3.9533333,1.9830696) (3.9600000,1.9855004) (3.9666667,1.9879272) (3.9733333,1.9903498) (3.9800000,1.9927684) (3.9866667,1.9951830) (3.9933333,1.9975935) (4.0000000,2.0000000)};
\draw[orange!85!black,thick,dashed] plot coordinates {(-4.0000000,-2.0000000) (-3.9933333,-1.9975935) (-3.9866667,-1.9951830) (-3.9800000,-1.9927684) (-3.9733333,-1.9903498) (-3.9666667,-1.9879272) (-3.9600000,-1.9855004) (-3.9533333,-1.9830696) (-3.9466667,-1.9806347) (-3.9400000,-1.9781956) (-3.9333333,-1.9757525) (-3.9266667,-1.9733051) (-3.9200000,-1.9708537) (-3.9133333,-1.9683980) (-3.9066667,-1.9659382) (-3.9000000,-1.9634741) (-3.8933333,-1.9610059) (-3.8866667,-1.9585334) (-3.8800000,-1.9560567) (-3.8733333,-1.9535757) (-3.8666667,-1.9510904) (-3.8600000,-1.9486008) (-3.8533333,-1.9461070) (-3.8466667,-1.9436088) (-3.8400000,-1.9411063) (-3.8333333,-1.9385995) (-3.8266667,-1.9360882) (-3.8200000,-1.9335726) (-3.8133333,-1.9310526) (-3.8066667,-1.9285282) (-3.8000000,-1.9259994) (-3.7933333,-1.9234662) (-3.7866667,-1.9209284) (-3.7800000,-1.9183862) (-3.7733333,-1.9158396) (-3.7666667,-1.9132884) (-3.7600000,-1.9107327) (-3.7533333,-1.9081724) (-3.7466667,-1.9056076) (-3.7400000,-1.9030383) (-3.7333333,-1.9004643) (-3.7266667,-1.8978858) (-3.7200000,-1.8953026) (-3.7133333,-1.8927148) (-3.7066667,-1.8901224) (-3.7000000,-1.8875253) (-3.6933333,-1.8849235) (-3.6866667,-1.8823170) (-3.6800000,-1.8797058) (-3.6733333,-1.8770898) (-3.6666667,-1.8744691) (-3.6600000,-1.8718436) (-3.6533333,-1.8692134) (-3.6466667,-1.8665783) (-3.6400000,-1.8639385) (-3.6333333,-1.8612937) (-3.6266667,-1.8586442) (-3.6200000,-1.8559897) (-3.6133333,-1.8533304) (-3.6066667,-1.8506661) (-3.6000000,-1.8479969) (-3.5933333,-1.8453228) (-3.5866667,-1.8426437) (-3.5800000,-1.8399596) (-3.5733333,-1.8372705) (-3.5666667,-1.8345764) (-3.5600000,-1.8318772) (-3.5533333,-1.8291730) (-3.5466667,-1.8264637) (-3.5400000,-1.8237494) (-3.5333333,-1.8210299) (-3.5266667,-1.8183052) (-3.5200000,-1.8155754) (-3.5133333,-1.8128405) (-3.5066667,-1.8101003) (-3.5000000,-1.8073549) (-3.4933333,-1.8046043) (-3.4866667,-1.8018484) (-3.4800000,-1.7990873) (-3.4733333,-1.7963209) (-3.4666667,-1.7935491) (-3.4600000,-1.7907720) (-3.4533333,-1.7879896) (-3.4466667,-1.7852018) (-3.4400000,-1.7824086) (-3.4333333,-1.7796099) (-3.4266667,-1.7768059) (-3.4200000,-1.7739963) (-3.4133333,-1.7711813) (-3.4066667,-1.7683608) (-3.4000000,-1.7655347) (-3.3933333,-1.7627032) (-3.3866667,-1.7598660) (-3.3800000,-1.7570232) (-3.3733333,-1.7541749) (-3.3666667,-1.7513209) (-3.3600000,-1.7484612) (-3.3533333,-1.7455959) (-3.3466667,-1.7427249) (-3.3400000,-1.7398481) (-3.3333333,-1.7369656) (-3.3266667,-1.7340773) (-3.3200000,-1.7311832) (-3.3133333,-1.7282834) (-3.3066667,-1.7253776) (-3.3000000,-1.7224660) (-3.2933333,-1.7195485) (-3.2866667,-1.7166251) (-3.2800000,-1.7136958) (-3.2733333,-1.7107605) (-3.2666667,-1.7078192) (-3.2600000,-1.7048720) (-3.2533333,-1.7019186) (-3.2466667,-1.6989593) (-3.2400000,-1.6959938) (-3.2333333,-1.6930222) (-3.2266667,-1.6900445) (-3.2200000,-1.6870607) (-3.2133333,-1.6840706) (-3.2066667,-1.6810744) (-3.2000000,-1.6780719) (-3.1933333,-1.6750632) (-3.1866667,-1.6720481) (-3.1800000,-1.6690268) (-3.1733333,-1.6659991) (-3.1666667,-1.6629650) (-3.1600000,-1.6599246) (-3.1533333,-1.6568777) (-3.1466667,-1.6538244) (-3.1400000,-1.6507646) (-3.1333333,-1.6476983) (-3.1266667,-1.6446254) (-3.1200000,-1.6415460) (-3.1133333,-1.6384600) (-3.1066667,-1.6353675) (-3.1000000,-1.6322682) (-3.0933333,-1.6291623) (-3.0866667,-1.6260497) (-3.0800000,-1.6229304) (-3.0733333,-1.6198042) (-3.0666667,-1.6166714) (-3.0600000,-1.6135317) (-3.0533333,-1.6103851) (-3.0466667,-1.6072317) (-3.0400000,-1.6040713) (-3.0333333,-1.6009040) (-3.0266667,-1.5977298) (-3.0200000,-1.5945485) (-3.0133333,-1.5913603) (-3.0066667,-1.5881649) (-3.0000000,-1.5849625) (-2.9933333,-1.5817529) (-2.9866667,-1.5785362) (-2.9800000,-1.5753123) (-2.9733333,-1.5720812) (-2.9666667,-1.5688428) (-2.9600000,-1.5655972) (-2.9533333,-1.5623442) (-2.9466667,-1.5590839) (-2.9400000,-1.5558162) (-2.9333333,-1.5525410) (-2.9266667,-1.5492584) (-2.9200000,-1.5459684) (-2.9133333,-1.5426708) (-2.9066667,-1.5393656) (-2.9000000,-1.5360529) (-2.8933333,-1.5327325) (-2.8866667,-1.5294045) (-2.8800000,-1.5260688) (-2.8733333,-1.5227254) (-2.8666667,-1.5193742) (-2.8600000,-1.5160151) (-2.8533333,-1.5126483) (-2.8466667,-1.5092736) (-2.8400000,-1.5058909) (-2.8333333,-1.5025003) (-2.8266667,-1.4991018) (-2.8200000,-1.4956952) (-2.8133333,-1.4922805) (-2.8066667,-1.4888577) (-2.8000000,-1.4854268) (-2.7933333,-1.4819877) (-2.7866667,-1.4785404) (-2.7800000,-1.4750849) (-2.7733333,-1.4716210) (-2.7666667,-1.4681488) (-2.7600000,-1.4646683) (-2.7533333,-1.4611793) (-2.7466667,-1.4576818) (-2.7400000,-1.4541759) (-2.7333333,-1.4506614) (-2.7266667,-1.4471383) (-2.7200000,-1.4436067) (-2.7133333,-1.4400663) (-2.7066667,-1.4365172) (-2.7000000,-1.4329594) (-2.6933333,-1.4293928) (-2.6866667,-1.4258173) (-2.6800000,-1.4222330) (-2.6733333,-1.4186397) (-2.6666667,-1.4150375) (-2.6600000,-1.4114262) (-2.6533333,-1.4078059) (-2.6466667,-1.4041765) (-2.6400000,-1.4005379) (-2.6333333,-1.3968902) (-2.6266667,-1.3932331) (-2.6200000,-1.3895668) (-2.6133333,-1.3858912) (-2.6066667,-1.3822061) (-2.6000000,-1.3785116) (-2.5933333,-1.3748077) (-2.5866667,-1.3710942) (-2.5800000,-1.3673711) (-2.5733333,-1.3636383) (-2.5666667,-1.3598959) (-2.5600000,-1.3561438) (-2.5533333,-1.3523819) (-2.5466667,-1.3486101) (-2.5400000,-1.3448285) (-2.5333333,-1.3410369) (-2.5266667,-1.3372353) (-2.5200000,-1.3334237) (-2.5133333,-1.3296020) (-2.5066667,-1.3257702) (-2.5000000,-1.3219281) (-2.4933333,-1.3180758) (-2.4866667,-1.3142131) (-2.4800000,-1.3103401) (-2.4733333,-1.3064567) (-2.4666667,-1.3025628) (-2.4600000,-1.2986583) (-2.4533333,-1.2947433) (-2.4466667,-1.2908176) (-2.4400000,-1.2868811) (-2.4333333,-1.2829340) (-2.4266667,-1.2789759) (-2.4200000,-1.2750070) (-2.4133333,-1.2710272) (-2.4066667,-1.2670363) (-2.4000000,-1.2630344) (-2.3933333,-1.2590213) (-2.3866667,-1.2549971) (-2.3800000,-1.2509616) (-2.3733333,-1.2469147) (-2.3666667,-1.2428565) (-2.3600000,-1.2387869) (-2.3533333,-1.2347057) (-2.3466667,-1.2306129) (-2.3400000,-1.2265085) (-2.3333333,-1.2223924) (-2.3266667,-1.2182645) (-2.3200000,-1.2141248) (-2.3133333,-1.2099732) (-2.3066667,-1.2058095) (-2.3000000,-1.2016339) (-2.2933333,-1.1974461) (-2.2866667,-1.1932461) (-2.2800000,-1.1890338) (-2.2733333,-1.1848092) (-2.2666667,-1.1805722) (-2.2600000,-1.1763228) (-2.2533333,-1.1720607) (-2.2466667,-1.1677861) (-2.2400000,-1.1634987) (-2.2333333,-1.1591986) (-2.2266667,-1.1548856) (-2.2200000,-1.1505597) (-2.2133333,-1.1462207) (-2.2066667,-1.1418687) (-2.2000000,-1.1375035) (-2.1933333,-1.1331251) (-2.1866667,-1.1287333) (-2.1800000,-1.1243281) (-2.1733333,-1.1199095) (-2.1666667,-1.1154772) (-2.1600000,-1.1110313) (-2.1533333,-1.1065717) (-2.1466667,-1.1020982) (-2.1400000,-1.0976108) (-2.1333333,-1.0931094) (-2.1266667,-1.0885939) (-2.1200000,-1.0840643) (-2.1133333,-1.0795203) (-2.1066667,-1.0749621) (-2.1000000,-1.0703893) (-2.0933333,-1.0658021) (-2.0866667,-1.0612002) (-2.0800000,-1.0565835) (-2.0733333,-1.0519521) (-2.0666667,-1.0473057) (-2.0600000,-1.0426443) (-2.0533333,-1.0379679) (-2.0466667,-1.0332762) (-2.0400000,-1.0285692) (-2.0333333,-1.0238467) (-2.0266667,-1.0191088) (-2.0200000,-1.0143553) (-2.0133333,-1.0095860) (-2.0066667,-1.0048010) (-2.0000000,-1.0000000) (-1.9933333,-0.9951830) (-1.9866667,-0.9903498) (-1.9800000,-0.9855004) (-1.9733333,-0.9806347) (-1.9666667,-0.9757525) (-1.9600000,-0.9708537) (-1.9533333,-0.9659382) (-1.9466667,-0.9610059) (-1.9400000,-0.9560567) (-1.9333333,-0.9510904) (-1.9266667,-0.9461070) (-1.9200000,-0.9411063) (-1.9133333,-0.9360882) (-1.9066667,-0.9310526) (-1.9000000,-0.9259994) (-1.8933333,-0.9209284) (-1.8866667,-0.9158396) (-1.8800000,-0.9107327) (-1.8733333,-0.9056076) (-1.8666667,-0.9004643) (-1.8600000,-0.8953026) (-1.8533333,-0.8901224) (-1.8466667,-0.8849235) (-1.8400000,-0.8797058) (-1.8333333,-0.8744691) (-1.8266667,-0.8692134) (-1.8200000,-0.8639385) (-1.8133333,-0.8586442) (-1.8066667,-0.8533304) (-1.8000000,-0.8479969) (-1.7933333,-0.8426437) (-1.7866667,-0.8372705) (-1.7800000,-0.8318772) (-1.7733333,-0.8264637) (-1.7666667,-0.8210299) (-1.7600000,-0.8155754) (-1.7533333,-0.8101003) (-1.7466667,-0.8046043) (-1.7400000,-0.7990873) (-1.7333333,-0.7935491) (-1.7266667,-0.7879896) (-1.7200000,-0.7824086) (-1.7133333,-0.7768059) (-1.7066667,-0.7711813) (-1.7000000,-0.7655347) (-1.6933333,-0.7598660) (-1.6866667,-0.7541749) (-1.6800000,-0.7484612) (-1.6733333,-0.7427249) (-1.6666667,-0.7369656) (-1.6600000,-0.7311832) (-1.6533333,-0.7253776) (-1.6466667,-0.7195485) (-1.6400000,-0.7136958) (-1.6333333,-0.7078192) (-1.6266667,-0.7019186) (-1.6200000,-0.6959938) (-1.6133333,-0.6900445) (-1.6066667,-0.6840706) (-1.6000000,-0.6780719) (-1.5933333,-0.6720481) (-1.5866667,-0.6659991) (-1.5800000,-0.6599246) (-1.5733333,-0.6538244) (-1.5666667,-0.6476983) (-1.5600000,-0.6415460) (-1.5533333,-0.6353675) (-1.5466667,-0.6291623) (-1.5400000,-0.6229304) (-1.5333333,-0.6166714) (-1.5266667,-0.6103851) (-1.5200000,-0.6040713) (-1.5133333,-0.5977298) (-1.5066667,-0.5913603) (-1.5000000,-0.5849625) (-1.4933333,-0.5785362) (-1.4866667,-0.5720812) (-1.4800000,-0.5655972) (-1.4733333,-0.5590839) (-1.4666667,-0.5525410) (-1.4600000,-0.5459684) (-1.4533333,-0.5393656) (-1.4466667,-0.5327325) (-1.4400000,-0.5260688) (-1.4333333,-0.5193742) (-1.4266667,-0.5126483) (-1.4200000,-0.5058909) (-1.4133333,-0.4991018) (-1.4066667,-0.4922805) (-1.4000000,-0.4854268) (-1.3933333,-0.4785404) (-1.3866667,-0.4716210) (-1.3800000,-0.4646683) (-1.3733333,-0.4576818) (-1.3666667,-0.4506614) (-1.3600000,-0.4436067) (-1.3533333,-0.4365172) (-1.3466667,-0.4293928) (-1.3400000,-0.4222330) (-1.3333333,-0.4150375) (-1.3266667,-0.4078059) (-1.3200000,-0.4005379) (-1.3133333,-0.3932331) (-1.3066667,-0.3858912) (-1.3000000,-0.3785116) (-1.2933333,-0.3710942) (-1.2866667,-0.3636383) (-1.2800000,-0.3561438) (-1.2733333,-0.3486101) (-1.2666667,-0.3410369) (-1.2600000,-0.3334237) (-1.2533333,-0.3257702) (-1.2466667,-0.3180758) (-1.2400000,-0.3103401) (-1.2333333,-0.3025628) (-1.2266667,-0.2947433) (-1.2200000,-0.2868811) (-1.2133333,-0.2789759) (-1.2066667,-0.2710272) (-1.2000000,-0.2630344) (-1.1933333,-0.2549971) (-1.1866667,-0.2469147) (-1.1800000,-0.2387869) (-1.1733333,-0.2306129) (-1.1666667,-0.2223924) (-1.1600000,-0.2141248) (-1.1533333,-0.2058095) (-1.1466667,-0.1974461) (-1.1400000,-0.1890338) (-1.1333333,-0.1805722) (-1.1266667,-0.1720607) (-1.1200000,-0.1634987) (-1.1133333,-0.1548856) (-1.1066667,-0.1462207) (-1.1000000,-0.1375035) (-1.0933333,-0.1287333) (-1.0866667,-0.1199095) (-1.0800000,-0.1110313) (-1.0733333,-0.1020982) (-1.0666667,-0.0931094) (-1.0600000,-0.0840643) (-1.0533333,-0.0749621) (-1.0466667,-0.0658021) (-1.0400000,-0.0565835) (-1.0333333,-0.0473057) (-1.0266667,-0.0379679) (-1.0200000,-0.0285692) (-1.0133333,-0.0191088) (-1.0066667,-0.0095860) (-1.0000000,-0.0000000) (-0.9933333,0.0096502) (-0.9866667,0.0193653) (-0.9800000,0.0291463) (-0.9733333,0.0389941) (-0.9666667,0.0489096) (-0.9600000,0.0588937) (-0.9533333,0.0689474) (-0.9466667,0.0790716) (-0.9400000,0.0892673) (-0.9333333,0.0995357) (-0.9266667,0.1098776) (-0.9200000,0.1202942) (-0.9133333,0.1307866) (-0.9066667,0.1413558) (-0.9000000,0.1520031) (-0.8933333,0.1627295) (-0.8866667,0.1735363) (-0.8800000,0.1844246) (-0.8733333,0.1953957) (-0.8666667,0.2064509) (-0.8600000,0.2175914) (-0.8533333,0.2288187) (-0.8466667,0.2401340) (-0.8400000,0.2515388) (-0.8333333,0.2630344) (-0.8266667,0.2746224) (-0.8200000,0.2863042) (-0.8133333,0.2980814) (-0.8066667,0.3099555) (-0.8000000,0.3219281) (-0.7933333,0.3340009) (-0.7866667,0.3461756) (-0.7800000,0.3584540) (-0.7733333,0.3708377) (-0.7666667,0.3833286) (-0.7600000,0.3959287) (-0.7533333,0.4086397) (-0.7466667,0.4214638) (-0.7400000,0.4344028) (-0.7333333,0.4474590) (-0.7266667,0.4606344) (-0.7200000,0.4739312) (-0.7133333,0.4873517) (-0.7066667,0.5008982) (-0.7000000,0.5145732) (-0.6933333,0.5283790) (-0.6866667,0.5423182) (-0.6800000,0.5563933) (-0.6733333,0.5706072) (-0.6666667,0.5849625) (-0.6600000,0.5994621) (-0.6533333,0.6141088) (-0.6466667,0.6289058) (-0.6400000,0.6438562) (-0.6333333,0.6589631) (-0.6266667,0.6742298) (-0.6200000,0.6896599) (-0.6133333,0.7052567) (-0.6066667,0.7210241) (-0.6000000,0.7369656) (-0.5933333,0.7530853) (-0.5866667,0.7693871) (-0.5800000,0.7858752) (-0.5733333,0.8025539) (-0.5666667,0.8194278) (-0.5600000,0.8365013) (-0.5533333,0.8537793) (-0.5466667,0.8712667) (-0.5400000,0.8889687) (-0.5333333,0.9068906) (-0.5266667,0.9250379) (-0.5200000,0.9434165) (-0.5133333,0.9620321) (-0.5066667,0.9808912) (-0.5000000,1.0000000) (-0.4933333,1.0193653) (-0.4866667,1.0389941) (-0.4800000,1.0588937) (-0.4733333,1.0790716) (-0.4666667,1.0995357) (-0.4600000,1.1202942) (-0.4533333,1.1413558) (-0.4466667,1.1627295) (-0.4400000,1.1844246) (-0.4333333,1.2064509) (-0.4266667,1.2288187) (-0.4200000,1.2515388) (-0.4133333,1.2746224) (-0.4066667,1.2980814) (-0.4000000,1.3219281) (-0.3933333,1.3461756) (-0.3866667,1.3708377) (-0.3800000,1.3959287) (-0.3733333,1.4214638) (-0.3666667,1.4474590) (-0.3600000,1.4739312) (-0.3533333,1.5008982) (-0.3466667,1.5283790) (-0.3400000,1.5563933) (-0.3333333,1.5849625) (-0.3266667,1.6141088) (-0.3200000,1.6438562) (-0.3133333,1.6742298) (-0.3066667,1.7052567) (-0.3000000,1.7369656) (-0.2933333,1.7693871) (-0.2866667,1.8025539) (-0.2800000,1.8365013) (-0.2733333,1.8712667) (-0.2666667,1.9068906) (-0.2600000,1.9434165) (-0.2533333,1.9808912) (-0.2466667,2.0193653) (-0.2400000,2.0588937) (-0.2333333,2.0995357) (-0.2266667,2.1413558) (-0.2200000,2.1844246) (-0.2133333,2.2288187) (-0.2066667,2.2746224) (-0.2000000,2.3219281) (-0.1933333,2.3708377) (-0.1866667,2.4214638) (-0.1800000,2.4739312) (-0.1733333,2.5283790) (-0.1666667,2.5849625) (-0.1600000,2.6438562) (-0.1533333,2.7052567) (-0.1466667,2.7693871) (-0.1400000,2.8365013) (-0.1333333,2.9068906) (-0.1266667,2.9808912)};
\draw[orange!85!black,thick,dashed] plot coordinates {(0.1266667,2.9808912) (0.1333333,2.9068906) (0.1400000,2.8365013) (0.1466667,2.7693871) (0.1533333,2.7052567) (0.1600000,2.6438562) (0.1666667,2.5849625) (0.1733333,2.5283790) (0.1800000,2.4739312) (0.1866667,2.4214638) (0.1933333,2.3708377) (0.2000000,2.3219281) (0.2066667,2.2746224) (0.2133333,2.2288187) (0.2200000,2.1844246) (0.2266667,2.1413558) (0.2333333,2.0995357) (0.2400000,2.0588937) (0.2466667,2.0193653) (0.2533333,1.9808912) (0.2600000,1.9434165) (0.2666667,1.9068906) (0.2733333,1.8712667) (0.2800000,1.8365013) (0.2866667,1.8025539) (0.2933333,1.7693871) (0.3000000,1.7369656) (0.3066667,1.7052567) (0.3133333,1.6742298) (0.3200000,1.6438562) (0.3266667,1.6141088) (0.3333333,1.5849625) (0.3400000,1.5563933) (0.3466667,1.5283790) (0.3533333,1.5008982) (0.3600000,1.4739312) (0.3666667,1.4474590) (0.3733333,1.4214638) (0.3800000,1.3959287) (0.3866667,1.3708377) (0.3933333,1.3461756) (0.4000000,1.3219281) (0.4066667,1.2980814) (0.4133333,1.2746224) (0.4200000,1.2515388) (0.4266667,1.2288187) (0.4333333,1.2064509) (0.4400000,1.1844246) (0.4466667,1.1627295) (0.4533333,1.1413558) (0.4600000,1.1202942) (0.4666667,1.0995357) (0.4733333,1.0790716) (0.4800000,1.0588937) (0.4866667,1.0389941) (0.4933333,1.0193653) (0.5000000,1.0000000) (0.5066667,0.9808912) (0.5133333,0.9620321) (0.5200000,0.9434165) (0.5266667,0.9250379) (0.5333333,0.9068906) (0.5400000,0.8889687) (0.5466667,0.8712667) (0.5533333,0.8537793) (0.5600000,0.8365013) (0.5666667,0.8194278) (0.5733333,0.8025539) (0.5800000,0.7858752) (0.5866667,0.7693871) (0.5933333,0.7530853) (0.6000000,0.7369656) (0.6066667,0.7210241) (0.6133333,0.7052567) (0.6200000,0.6896599) (0.6266667,0.6742298) (0.6333333,0.6589631) (0.6400000,0.6438562) (0.6466667,0.6289058) (0.6533333,0.6141088) (0.6600000,0.5994621) (0.6666667,0.5849625) (0.6733333,0.5706072) (0.6800000,0.5563933) (0.6866667,0.5423182) (0.6933333,0.5283790) (0.7000000,0.5145732) (0.7066667,0.5008982) (0.7133333,0.4873517) (0.7200000,0.4739312) (0.7266667,0.4606344) (0.7333333,0.4474590) (0.7400000,0.4344028) (0.7466667,0.4214638) (0.7533333,0.4086397) (0.7600000,0.3959287) (0.7666667,0.3833286) (0.7733333,0.3708377) (0.7800000,0.3584540) (0.7866667,0.3461756) (0.7933333,0.3340009) (0.8000000,0.3219281) (0.8066667,0.3099555) (0.8133333,0.2980814) (0.8200000,0.2863042) (0.8266667,0.2746224) (0.8333333,0.2630344) (0.8400000,0.2515388) (0.8466667,0.2401340) (0.8533333,0.2288187) (0.8600000,0.2175914) (0.8666667,0.2064509) (0.8733333,0.1953957) (0.8800000,0.1844246) (0.8866667,0.1735363) (0.8933333,0.1627295) (0.9000000,0.1520031) (0.9066667,0.1413558) (0.9133333,0.1307866) (0.9200000,0.1202942) (0.9266667,0.1098776) (0.9333333,0.0995357) (0.9400000,0.0892673) (0.9466667,0.0790716) (0.9533333,0.0689474) (0.9600000,0.0588937) (0.9666667,0.0489096) (0.9733333,0.0389941) (0.9800000,0.0291463) (0.9866667,0.0193653) (0.9933333,0.0096502) (1.0000000,-0.0000000) (1.0066667,-0.0095860) (1.0133333,-0.0191088) (1.0200000,-0.0285692) (1.0266667,-0.0379679) (1.0333333,-0.0473057) (1.0400000,-0.0565835) (1.0466667,-0.0658021) (1.0533333,-0.0749621) (1.0600000,-0.0840643) (1.0666667,-0.0931094) (1.0733333,-0.1020982) (1.0800000,-0.1110313) (1.0866667,-0.1199095) (1.0933333,-0.1287333) (1.1000000,-0.1375035) (1.1066667,-0.1462207) (1.1133333,-0.1548856) (1.1200000,-0.1634987) (1.1266667,-0.1720607) (1.1333333,-0.1805722) (1.1400000,-0.1890338) (1.1466667,-0.1974461) (1.1533333,-0.2058095) (1.1600000,-0.2141248) (1.1666667,-0.2223924) (1.1733333,-0.2306129) (1.1800000,-0.2387869) (1.1866667,-0.2469147) (1.1933333,-0.2549971) (1.2000000,-0.2630344) (1.2066667,-0.2710272) (1.2133333,-0.2789759) (1.2200000,-0.2868811) (1.2266667,-0.2947433) (1.2333333,-0.3025628) (1.2400000,-0.3103401) (1.2466667,-0.3180758) (1.2533333,-0.3257702) (1.2600000,-0.3334237) (1.2666667,-0.3410369) (1.2733333,-0.3486101) (1.2800000,-0.3561438) (1.2866667,-0.3636383) (1.2933333,-0.3710942) (1.3000000,-0.3785116) (1.3066667,-0.3858912) (1.3133333,-0.3932331) (1.3200000,-0.4005379) (1.3266667,-0.4078059) (1.3333333,-0.4150375) (1.3400000,-0.4222330) (1.3466667,-0.4293928) (1.3533333,-0.4365172) (1.3600000,-0.4436067) (1.3666667,-0.4506614) (1.3733333,-0.4576818) (1.3800000,-0.4646683) (1.3866667,-0.4716210) (1.3933333,-0.4785404) (1.4000000,-0.4854268) (1.4066667,-0.4922805) (1.4133333,-0.4991018) (1.4200000,-0.5058909) (1.4266667,-0.5126483) (1.4333333,-0.5193742) (1.4400000,-0.5260688) (1.4466667,-0.5327325) (1.4533333,-0.5393656) (1.4600000,-0.5459684) (1.4666667,-0.5525410) (1.4733333,-0.5590839) (1.4800000,-0.5655972) (1.4866667,-0.5720812) (1.4933333,-0.5785362) (1.5000000,-0.5849625) (1.5066667,-0.5913603) (1.5133333,-0.5977298) (1.5200000,-0.6040713) (1.5266667,-0.6103851) (1.5333333,-0.6166714) (1.5400000,-0.6229304) (1.5466667,-0.6291623) (1.5533333,-0.6353675) (1.5600000,-0.6415460) (1.5666667,-0.6476983) (1.5733333,-0.6538244) (1.5800000,-0.6599246) (1.5866667,-0.6659991) (1.5933333,-0.6720481) (1.6000000,-0.6780719) (1.6066667,-0.6840706) (1.6133333,-0.6900445) (1.6200000,-0.6959938) (1.6266667,-0.7019186) (1.6333333,-0.7078192) (1.6400000,-0.7136958) (1.6466667,-0.7195485) (1.6533333,-0.7253776) (1.6600000,-0.7311832) (1.6666667,-0.7369656) (1.6733333,-0.7427249) (1.6800000,-0.7484612) (1.6866667,-0.7541749) (1.6933333,-0.7598660) (1.7000000,-0.7655347) (1.7066667,-0.7711813) (1.7133333,-0.7768059) (1.7200000,-0.7824086) (1.7266667,-0.7879896) (1.7333333,-0.7935491) (1.7400000,-0.7990873) (1.7466667,-0.8046043) (1.7533333,-0.8101003) (1.7600000,-0.8155754) (1.7666667,-0.8210299) (1.7733333,-0.8264637) (1.7800000,-0.8318772) (1.7866667,-0.8372705) (1.7933333,-0.8426437) (1.8000000,-0.8479969) (1.8066667,-0.8533304) (1.8133333,-0.8586442) (1.8200000,-0.8639385) (1.8266667,-0.8692134) (1.8333333,-0.8744691) (1.8400000,-0.8797058) (1.8466667,-0.8849235) (1.8533333,-0.8901224) (1.8600000,-0.8953026) (1.8666667,-0.9004643) (1.8733333,-0.9056076) (1.8800000,-0.9107327) (1.8866667,-0.9158396) (1.8933333,-0.9209284) (1.9000000,-0.9259994) (1.9066667,-0.9310526) (1.9133333,-0.9360882) (1.9200000,-0.9411063) (1.9266667,-0.9461070) (1.9333333,-0.9510904) (1.9400000,-0.9560567) (1.9466667,-0.9610059) (1.9533333,-0.9659382) (1.9600000,-0.9708537) (1.9666667,-0.9757525) (1.9733333,-0.9806347) (1.9800000,-0.9855004) (1.9866667,-0.9903498) (1.9933333,-0.9951830) (2.0000000,-1.0000000) (2.0066667,-1.0048010) (2.0133333,-1.0095860) (2.0200000,-1.0143553) (2.0266667,-1.0191088) (2.0333333,-1.0238467) (2.0400000,-1.0285692) (2.0466667,-1.0332762) (2.0533333,-1.0379679) (2.0600000,-1.0426443) (2.0666667,-1.0473057) (2.0733333,-1.0519521) (2.0800000,-1.0565835) (2.0866667,-1.0612002) (2.0933333,-1.0658021) (2.1000000,-1.0703893) (2.1066667,-1.0749621) (2.1133333,-1.0795203) (2.1200000,-1.0840643) (2.1266667,-1.0885939) (2.1333333,-1.0931094) (2.1400000,-1.0976108) (2.1466667,-1.1020982) (2.1533333,-1.1065717) (2.1600000,-1.1110313) (2.1666667,-1.1154772) (2.1733333,-1.1199095) (2.1800000,-1.1243281) (2.1866667,-1.1287333) (2.1933333,-1.1331251) (2.2000000,-1.1375035) (2.2066667,-1.1418687) (2.2133333,-1.1462207) (2.2200000,-1.1505597) (2.2266667,-1.1548856) (2.2333333,-1.1591986) (2.2400000,-1.1634987) (2.2466667,-1.1677861) (2.2533333,-1.1720607) (2.2600000,-1.1763228) (2.2666667,-1.1805722) (2.2733333,-1.1848092) (2.2800000,-1.1890338) (2.2866667,-1.1932461) (2.2933333,-1.1974461) (2.3000000,-1.2016339) (2.3066667,-1.2058095) (2.3133333,-1.2099732) (2.3200000,-1.2141248) (2.3266667,-1.2182645) (2.3333333,-1.2223924) (2.3400000,-1.2265085) (2.3466667,-1.2306129) (2.3533333,-1.2347057) (2.3600000,-1.2387869) (2.3666667,-1.2428565) (2.3733333,-1.2469147) (2.3800000,-1.2509616) (2.3866667,-1.2549971) (2.3933333,-1.2590213) (2.4000000,-1.2630344) (2.4066667,-1.2670363) (2.4133333,-1.2710272) (2.4200000,-1.2750070) (2.4266667,-1.2789759) (2.4333333,-1.2829340) (2.4400000,-1.2868811) (2.4466667,-1.2908176) (2.4533333,-1.2947433) (2.4600000,-1.2986583) (2.4666667,-1.3025628) (2.4733333,-1.3064567) (2.4800000,-1.3103401) (2.4866667,-1.3142131) (2.4933333,-1.3180758) (2.5000000,-1.3219281) (2.5066667,-1.3257702) (2.5133333,-1.3296020) (2.5200000,-1.3334237) (2.5266667,-1.3372353) (2.5333333,-1.3410369) (2.5400000,-1.3448285) (2.5466667,-1.3486101) (2.5533333,-1.3523819) (2.5600000,-1.3561438) (2.5666667,-1.3598959) (2.5733333,-1.3636383) (2.5800000,-1.3673711) (2.5866667,-1.3710942) (2.5933333,-1.3748077) (2.6000000,-1.3785116) (2.6066667,-1.3822061) (2.6133333,-1.3858912) (2.6200000,-1.3895668) (2.6266667,-1.3932331) (2.6333333,-1.3968902) (2.6400000,-1.4005379) (2.6466667,-1.4041765) (2.6533333,-1.4078059) (2.6600000,-1.4114262) (2.6666667,-1.4150375) (2.6733333,-1.4186397) (2.6800000,-1.4222330) (2.6866667,-1.4258173) (2.6933333,-1.4293928) (2.7000000,-1.4329594) (2.7066667,-1.4365172) (2.7133333,-1.4400663) (2.7200000,-1.4436067) (2.7266667,-1.4471383) (2.7333333,-1.4506614) (2.7400000,-1.4541759) (2.7466667,-1.4576818) (2.7533333,-1.4611793) (2.7600000,-1.4646683) (2.7666667,-1.4681488) (2.7733333,-1.4716210) (2.7800000,-1.4750849) (2.7866667,-1.4785404) (2.7933333,-1.4819877) (2.8000000,-1.4854268) (2.8066667,-1.4888577) (2.8133333,-1.4922805) (2.8200000,-1.4956952) (2.8266667,-1.4991018) (2.8333333,-1.5025003) (2.8400000,-1.5058909) (2.8466667,-1.5092736) (2.8533333,-1.5126483) (2.8600000,-1.5160151) (2.8666667,-1.5193742) (2.8733333,-1.5227254) (2.8800000,-1.5260688) (2.8866667,-1.5294045) (2.8933333,-1.5327325) (2.9000000,-1.5360529) (2.9066667,-1.5393656) (2.9133333,-1.5426708) (2.9200000,-1.5459684) (2.9266667,-1.5492584) (2.9333333,-1.5525410) (2.9400000,-1.5558162) (2.9466667,-1.5590839) (2.9533333,-1.5623442) (2.9600000,-1.5655972) (2.9666667,-1.5688428) (2.9733333,-1.5720812) (2.9800000,-1.5753123) (2.9866667,-1.5785362) (2.9933333,-1.5817529) (3.0000000,-1.5849625) (3.0066667,-1.5881649) (3.0133333,-1.5913603) (3.0200000,-1.5945485) (3.0266667,-1.5977298) (3.0333333,-1.6009040) (3.0400000,-1.6040713) (3.0466667,-1.6072317) (3.0533333,-1.6103851) (3.0600000,-1.6135317) (3.0666667,-1.6166714) (3.0733333,-1.6198042) (3.0800000,-1.6229304) (3.0866667,-1.6260497) (3.0933333,-1.6291623) (3.1000000,-1.6322682) (3.1066667,-1.6353675) (3.1133333,-1.6384600) (3.1200000,-1.6415460) (3.1266667,-1.6446254) (3.1333333,-1.6476983) (3.1400000,-1.6507646) (3.1466667,-1.6538244) (3.1533333,-1.6568777) (3.1600000,-1.6599246) (3.1666667,-1.6629650) (3.1733333,-1.6659991) (3.1800000,-1.6690268) (3.1866667,-1.6720481) (3.1933333,-1.6750632) (3.2000000,-1.6780719) (3.2066667,-1.6810744) (3.2133333,-1.6840706) (3.2200000,-1.6870607) (3.2266667,-1.6900445) (3.2333333,-1.6930222) (3.2400000,-1.6959938) (3.2466667,-1.6989593) (3.2533333,-1.7019186) (3.2600000,-1.7048720) (3.2666667,-1.7078192) (3.2733333,-1.7107605) (3.2800000,-1.7136958) (3.2866667,-1.7166251) (3.2933333,-1.7195485) (3.3000000,-1.7224660) (3.3066667,-1.7253776) (3.3133333,-1.7282834) (3.3200000,-1.7311832) (3.3266667,-1.7340773) (3.3333333,-1.7369656) (3.3400000,-1.7398481) (3.3466667,-1.7427249) (3.3533333,-1.7455959) (3.3600000,-1.7484612) (3.3666667,-1.7513209) (3.3733333,-1.7541749) (3.3800000,-1.7570232) (3.3866667,-1.7598660) (3.3933333,-1.7627032) (3.4000000,-1.7655347) (3.4066667,-1.7683608) (3.4133333,-1.7711813) (3.4200000,-1.7739963) (3.4266667,-1.7768059) (3.4333333,-1.7796099) (3.4400000,-1.7824086) (3.4466667,-1.7852018) (3.4533333,-1.7879896) (3.4600000,-1.7907720) (3.4666667,-1.7935491) (3.4733333,-1.7963209) (3.4800000,-1.7990873) (3.4866667,-1.8018484) (3.4933333,-1.8046043) (3.5000000,-1.8073549) (3.5066667,-1.8101003) (3.5133333,-1.8128405) (3.5200000,-1.8155754) (3.5266667,-1.8183052) (3.5333333,-1.8210299) (3.5400000,-1.8237494) (3.5466667,-1.8264637) (3.5533333,-1.8291730) (3.5600000,-1.8318772) (3.5666667,-1.8345764) (3.5733333,-1.8372705) (3.5800000,-1.8399596) (3.5866667,-1.8426437) (3.5933333,-1.8453228) (3.6000000,-1.8479969) (3.6066667,-1.8506661) (3.6133333,-1.8533304) (3.6200000,-1.8559897) (3.6266667,-1.8586442) (3.6333333,-1.8612937) (3.6400000,-1.8639385) (3.6466667,-1.8665783) (3.6533333,-1.8692134) (3.6600000,-1.8718436) (3.6666667,-1.8744691) (3.6733333,-1.8770898) (3.6800000,-1.8797058) (3.6866667,-1.8823170) (3.6933333,-1.8849235) (3.7000000,-1.8875253) (3.7066667,-1.8901224) (3.7133333,-1.8927148) (3.7200000,-1.8953026) (3.7266667,-1.8978858) (3.7333333,-1.9004643) (3.7400000,-1.9030383) (3.7466667,-1.9056076) (3.7533333,-1.9081724) (3.7600000,-1.9107327) (3.7666667,-1.9132884) (3.7733333,-1.9158396) (3.7800000,-1.9183862) (3.7866667,-1.9209284) (3.7933333,-1.9234662) (3.8000000,-1.9259994) (3.8066667,-1.9285282) (3.8133333,-1.9310526) (3.8200000,-1.9335726) (3.8266667,-1.9360882) (3.8333333,-1.9385995) (3.8400000,-1.9411063) (3.8466667,-1.9436088) (3.8533333,-1.9461070) (3.8600000,-1.9486008) (3.8666667,-1.9510904) (3.8733333,-1.9535757) (3.8800000,-1.9560567) (3.8866667,-1.9585334) (3.8933333,-1.9610059) (3.9000000,-1.9634741) (3.9066667,-1.9659382) (3.9133333,-1.9683980) (3.9200000,-1.9708537) (3.9266667,-1.9733051) (3.9333333,-1.9757525) (3.9400000,-1.9781956) (3.9466667,-1.9806347) (3.9533333,-1.9830696) (3.9600000,-1.9855004) (3.9666667,-1.9879272) (3.9733333,-1.9903498) (3.9800000,-1.9927684) (3.9866667,-1.9951830) (3.9933333,-1.9975935) (4.0000000,-2.0000000)};
\draw[violet!75!black,densely dotted] (-2,1.0000000)--(2,1.0000000);
\fill[violet!75!black] (-2,1.0000000) circle(2.3pt);
\fill[violet!75!black] (2,1.0000000) circle(2.3pt);
\end{scope}
\end{scope}
\node[anchor=west,align=left,text width=4.35cm] at (8.9,4.4) {蓝实线：$a=2$\\[8pt]橙虚线：$a=\frac12$\\[14pt]每条曲线左右同高。\\[8pt]经过 $(-1,0)$、$(1,0)$；\\$x=0$ 不能代入。};
\node[text width=12.8cm,align=center] at (7.1,.65) {$D=\mathbb R\setminus\{0\}$；纵轴是渐近线，定义域仍然左右成对。};
\end{tikzpicture}\end{center}

\textbf{先看图：} 真数要求 $|x|>0$，即 $x\ne0$；定义域虽然断成两段，仍然左右成对。每条曲线在 $x$ 与 $-x$ 处同高，所以关于纵轴对称。

\textbf{再说明一般情形：}
\[
f(-x)=\log_a|-x|=\log_a|x|=f(x)\qquad(x\ne0).
\]
它是偶函数；与 $\log_a x$ 相比，多出的绝对值使负数输入也合法，但并没有让零合法。
\end{explanationbox}
\clearpage
\begin{explanationbox}
\textbf{\textcolor{CExplanation}{奇函数乘奇函数，空点仍要保留}}
% Formula-based samples; see F083 checks/build_tikz_diagrams.py.
\begin{center}\begin{tikzpicture}[font=\small,>=stealth,line join=round]
\draw[gray!40,rounded corners=3pt] (0,0) rectangle (14.2,8.4);
\node at (7.1,7.75) {偶5：$y=x\left(\dfrac12+\dfrac1{2^x-1}\right)$};
\begin{scope}[shift={(4.400000,1.650000)},x=0.850000cm,y=1.551724cm]
\draw[->,gray] (-4,0)--(4,0) node[right] {$x$};
\draw[->,gray] (0,0)--(0,2.9) node[above] {$y$};
\draw[gray] (-2,-.035/1.551724)--(-2,.035/1.551724);
\node[anchor=north] at (-2,-.07/1.551724) {$-2$};
\draw[gray] (2,-.035/1.551724)--(2,.035/1.551724);
\node[anchor=north] at (2,-.07/1.551724) {$2$};
\begin{scope}
\clip (-4,0) rectangle (4,2.9);
\draw[blue,thick,solid] plot coordinates {(-4.0000000,2.2666667) (-3.9933333,2.2642046) (-3.9866667,2.2617449) (-3.9800000,2.2592879) (-3.9733333,2.2568333) (-3.9666667,2.2543812) (-3.9600000,2.2519317) (-3.9533333,2.2494847) (-3.9466667,2.2470403) (-3.9400000,2.2445984) (-3.9333333,2.2421591) (-3.9266667,2.2397223) (-3.9200000,2.2372881) (-3.9133333,2.2348564) (-3.9066667,2.2324274) (-3.9000000,2.2300009) (-3.8933333,2.2275770) (-3.8866667,2.2251557) (-3.8800000,2.2227370) (-3.8733333,2.2203209) (-3.8666667,2.2179074) (-3.8600000,2.2154965) (-3.8533333,2.2130882) (-3.8466667,2.2106826) (-3.8400000,2.2082796) (-3.8333333,2.2058793) (-3.8266667,2.2034815) (-3.8200000,2.2010865) (-3.8133333,2.1986941) (-3.8066667,2.1963043) (-3.8000000,2.1939173) (-3.7933333,2.1915328) (-3.7866667,2.1891511) (-3.7800000,2.1867721) (-3.7733333,2.1843957) (-3.7666667,2.1820221) (-3.7600000,2.1796511) (-3.7533333,2.1772829) (-3.7466667,2.1749174) (-3.7400000,2.1725546) (-3.7333333,2.1701945) (-3.7266667,2.1678371) (-3.7200000,2.1654825) (-3.7133333,2.1631306) (-3.7066667,2.1607815) (-3.7000000,2.1584352) (-3.6933333,2.1560916) (-3.6866667,2.1537507) (-3.6800000,2.1514126) (-3.6733333,2.1490774) (-3.6666667,2.1467449) (-3.6600000,2.1444151) (-3.6533333,2.1420882) (-3.6466667,2.1397641) (-3.6400000,2.1374428) (-3.6333333,2.1351243) (-3.6266667,2.1328086) (-3.6200000,2.1304958) (-3.6133333,2.1281857) (-3.6066667,2.1258785) (-3.6000000,2.1235742) (-3.5933333,2.1212727) (-3.5866667,2.1189741) (-3.5800000,2.1166783) (-3.5733333,2.1143854) (-3.5666667,2.1120953) (-3.5600000,2.1098082) (-3.5533333,2.1075239) (-3.5466667,2.1052425) (-3.5400000,2.1029640) (-3.5333333,2.1006884) (-3.5266667,2.0984157) (-3.5200000,2.0961459) (-3.5133333,2.0938791) (-3.5066667,2.0916152) (-3.5000000,2.0893542) (-3.4933333,2.0870961) (-3.4866667,2.0848410) (-3.4800000,2.0825888) (-3.4733333,2.0803396) (-3.4666667,2.0780934) (-3.4600000,2.0758501) (-3.4533333,2.0736098) (-3.4466667,2.0713724) (-3.4400000,2.0691381) (-3.4333333,2.0669067) (-3.4266667,2.0646784) (-3.4200000,2.0624530) (-3.4133333,2.0602307) (-3.4066667,2.0580114) (-3.4000000,2.0557950) (-3.3933333,2.0535818) (-3.3866667,2.0513715) (-3.3800000,2.0491643) (-3.3733333,2.0469601) (-3.3666667,2.0447590) (-3.3600000,2.0425610) (-3.3533333,2.0403660) (-3.3466667,2.0381740) (-3.3400000,2.0359852) (-3.3333333,2.0337994) (-3.3266667,2.0316167) (-3.3200000,2.0294371) (-3.3133333,2.0272606) (-3.3066667,2.0250872) (-3.3000000,2.0229169) (-3.2933333,2.0207498) (-3.2866667,2.0185857) (-3.2800000,2.0164248) (-3.2733333,2.0142671) (-3.2666667,2.0121124) (-3.2600000,2.0099609) (-3.2533333,2.0078126) (-3.2466667,2.0056674) (-3.2400000,2.0035254) (-3.2333333,2.0013865) (-3.2266667,1.9992509) (-3.2200000,1.9971184) (-3.2133333,1.9949891) (-3.2066667,1.9928630) (-3.2000000,1.9907401) (-3.1933333,1.9886204) (-3.1866667,1.9865039) (-3.1800000,1.9843907) (-3.1733333,1.9822806) (-3.1666667,1.9801738) (-3.1600000,1.9780703) (-3.1533333,1.9759699) (-3.1466667,1.9738729) (-3.1400000,1.9717790) (-3.1333333,1.9696885) (-3.1266667,1.9676012) (-3.1200000,1.9655172) (-3.1133333,1.9634364) (-3.1066667,1.9613590) (-3.1000000,1.9592848) (-3.0933333,1.9572139) (-3.0866667,1.9551464) (-3.0800000,1.9530821) (-3.0733333,1.9510212) (-3.0666667,1.9489636) (-3.0600000,1.9469093) (-3.0533333,1.9448583) (-3.0466667,1.9428107) (-3.0400000,1.9407665) (-3.0333333,1.9387255) (-3.0266667,1.9366880) (-3.0200000,1.9346538) (-3.0133333,1.9326229) (-3.0066667,1.9305955) (-3.0000000,1.9285714) (-2.9933333,1.9265507) (-2.9866667,1.9245335) (-2.9800000,1.9225196) (-2.9733333,1.9205091) (-2.9666667,1.9185020) (-2.9600000,1.9164984) (-2.9533333,1.9144981) (-2.9466667,1.9125013) (-2.9400000,1.9105080) (-2.9333333,1.9085181) (-2.9266667,1.9065316) (-2.9200000,1.9045486) (-2.9133333,1.9025690) (-2.9066667,1.9005929) (-2.9000000,1.8986203) (-2.8933333,1.8966512) (-2.8866667,1.8946855) (-2.8800000,1.8927233) (-2.8733333,1.8907647) (-2.8666667,1.8888095) (-2.8600000,1.8868578) (-2.8533333,1.8849097) (-2.8466667,1.8829651) (-2.8400000,1.8810240) (-2.8333333,1.8790864) (-2.8266667,1.8771524) (-2.8200000,1.8752219) (-2.8133333,1.8732949) (-2.8066667,1.8713715) (-2.8000000,1.8694517) (-2.7933333,1.8675355) (-2.7866667,1.8656228) (-2.7800000,1.8637137) (-2.7733333,1.8618082) (-2.7666667,1.8599063) (-2.7600000,1.8580080) (-2.7533333,1.8561132) (-2.7466667,1.8542221) (-2.7400000,1.8523346) (-2.7333333,1.8504508) (-2.7266667,1.8485705) (-2.7200000,1.8466939) (-2.7133333,1.8448210) (-2.7066667,1.8429516) (-2.7000000,1.8410860) (-2.6933333,1.8392239) (-2.6866667,1.8373656) (-2.6800000,1.8355109) (-2.6733333,1.8336599) (-2.6666667,1.8318126) (-2.6600000,1.8299689) (-2.6533333,1.8281290) (-2.6466667,1.8262927) (-2.6400000,1.8244602) (-2.6333333,1.8226314) (-2.6266667,1.8208063) (-2.6200000,1.8189849) (-2.6133333,1.8171672) (-2.6066667,1.8153533) (-2.6000000,1.8135431) (-2.5933333,1.8117367) (-2.5866667,1.8099340) (-2.5800000,1.8081350) (-2.5733333,1.8063399) (-2.5666667,1.8045485) (-2.5600000,1.8027608) (-2.5533333,1.8009770) (-2.5466667,1.7991969) (-2.5400000,1.7974207) (-2.5333333,1.7956482) (-2.5266667,1.7938796) (-2.5200000,1.7921147) (-2.5133333,1.7903537) (-2.5066667,1.7885965) (-2.5000000,1.7868431) (-2.4933333,1.7850935) (-2.4866667,1.7833478) (-2.4800000,1.7816060) (-2.4733333,1.7798680) (-2.4666667,1.7781338) (-2.4600000,1.7764035) (-2.4533333,1.7746771) (-2.4466667,1.7729545) (-2.4400000,1.7712359) (-2.4333333,1.7695211) (-2.4266667,1.7678102) (-2.4200000,1.7661032) (-2.4133333,1.7644001) (-2.4066667,1.7627009) (-2.4000000,1.7610057) (-2.3933333,1.7593143) (-2.3866667,1.7576269) (-2.3800000,1.7559434) (-2.3733333,1.7542639) (-2.3666667,1.7525883) (-2.3600000,1.7509166) (-2.3533333,1.7492489) (-2.3466667,1.7475852) (-2.3400000,1.7459254) (-2.3333333,1.7442696) (-2.3266667,1.7426177) (-2.3200000,1.7409699) (-2.3133333,1.7393260) (-2.3066667,1.7376862) (-2.3000000,1.7360503) (-2.2933333,1.7344185) (-2.2866667,1.7327906) (-2.2800000,1.7311668) (-2.2733333,1.7295470) (-2.2666667,1.7279312) (-2.2600000,1.7263195) (-2.2533333,1.7247118) (-2.2466667,1.7231081) (-2.2400000,1.7215085) (-2.2333333,1.7199129) (-2.2266667,1.7183215) (-2.2200000,1.7167340) (-2.2133333,1.7151507) (-2.2066667,1.7135714) (-2.2000000,1.7119962) (-2.1933333,1.7104251) (-2.1866667,1.7088581) (-2.1800000,1.7072952) (-2.1733333,1.7057364) (-2.1666667,1.7041817) (-2.1600000,1.7026312) (-2.1533333,1.7010847) (-2.1466667,1.6995424) (-2.1400000,1.6980042) (-2.1333333,1.6964702) (-2.1266667,1.6949403) (-2.1200000,1.6934145) (-2.1133333,1.6918929) (-2.1066667,1.6903755) (-2.1000000,1.6888622) (-2.0933333,1.6873531) (-2.0866667,1.6858482) (-2.0800000,1.6843474) (-2.0733333,1.6828509) (-2.0666667,1.6813585) (-2.0600000,1.6798703) (-2.0533333,1.6783864) (-2.0466667,1.6769066) (-2.0400000,1.6754311) (-2.0333333,1.6739598) (-2.0266667,1.6724927) (-2.0200000,1.6710298) (-2.0133333,1.6695712) (-2.0066667,1.6681168) (-2.0000000,1.6666667) (-1.9933333,1.6652208) (-1.9866667,1.6637791) (-1.9800000,1.6623418) (-1.9733333,1.6609087) (-1.9666667,1.6594798) (-1.9600000,1.6580553) (-1.9533333,1.6566350) (-1.9466667,1.6552190) (-1.9400000,1.6538073) (-1.9333333,1.6523999) (-1.9266667,1.6509968) (-1.9200000,1.6495981) (-1.9133333,1.6482036) (-1.9066667,1.6468134) (-1.9000000,1.6454276) (-1.8933333,1.6440461) (-1.8866667,1.6426689) (-1.8800000,1.6412961) (-1.8733333,1.6399276) (-1.8666667,1.6385635) (-1.8600000,1.6372037) (-1.8533333,1.6358483) (-1.8466667,1.6344973) (-1.8400000,1.6331506) (-1.8333333,1.6318083) (-1.8266667,1.6304703) (-1.8200000,1.6291368) (-1.8133333,1.6278076) (-1.8066667,1.6264829) (-1.8000000,1.6251625) (-1.7933333,1.6238465) (-1.7866667,1.6225350) (-1.7800000,1.6212279) (-1.7733333,1.6199251) (-1.7666667,1.6186268) (-1.7600000,1.6173330) (-1.7533333,1.6160435) (-1.7466667,1.6147586) (-1.7400000,1.6134780) (-1.7333333,1.6122019) (-1.7266667,1.6109303) (-1.7200000,1.6096631) (-1.7133333,1.6084003) (-1.7066667,1.6071421) (-1.7000000,1.6058883) (-1.6933333,1.6046390) (-1.6866667,1.6033941) (-1.6800000,1.6021538) (-1.6733333,1.6009179) (-1.6666667,1.5996866) (-1.6600000,1.5984597) (-1.6533333,1.5972374) (-1.6466667,1.5960195) (-1.6400000,1.5948062) (-1.6333333,1.5935974) (-1.6266667,1.5923931) (-1.6200000,1.5911934) (-1.6133333,1.5899981) (-1.6066667,1.5888074) (-1.6000000,1.5876213) (-1.5933333,1.5864397) (-1.5866667,1.5852627) (-1.5800000,1.5840902) (-1.5733333,1.5829222) (-1.5666667,1.5817589) (-1.5600000,1.5806001) (-1.5533333,1.5794458) (-1.5466667,1.5782962) (-1.5400000,1.5771511) (-1.5333333,1.5760107) (-1.5266667,1.5748748) (-1.5200000,1.5737435) (-1.5133333,1.5726168) (-1.5066667,1.5714947) (-1.5000000,1.5703772) (-1.4933333,1.5692644) (-1.4866667,1.5681561) (-1.4800000,1.5670525) (-1.4733333,1.5659535) (-1.4666667,1.5648592) (-1.4600000,1.5637694) (-1.4533333,1.5626843) (-1.4466667,1.5616039) (-1.4400000,1.5605281) (-1.4333333,1.5594569) (-1.4266667,1.5583905) (-1.4200000,1.5573286) (-1.4133333,1.5562715) (-1.4066667,1.5552190) (-1.4000000,1.5541711) (-1.3933333,1.5531280) (-1.3866667,1.5520895) (-1.3800000,1.5510557) (-1.3733333,1.5500266) (-1.3666667,1.5490022) (-1.3600000,1.5479825) (-1.3533333,1.5469675) (-1.3466667,1.5459572) (-1.3400000,1.5449517) (-1.3333333,1.5439508) (-1.3266667,1.5429546) (-1.3200000,1.5419632) (-1.3133333,1.5409765) (-1.3066667,1.5399945) (-1.3000000,1.5390173) (-1.2933333,1.5380447) (-1.2866667,1.5370770) (-1.2800000,1.5361140) (-1.2733333,1.5351557) (-1.2666667,1.5342022) (-1.2600000,1.5332534) (-1.2533333,1.5323094) (-1.2466667,1.5313701) (-1.2400000,1.5304357) (-1.2333333,1.5295060) (-1.2266667,1.5285810) (-1.2200000,1.5276609) (-1.2133333,1.5267455) (-1.2066667,1.5258349) (-1.2000000,1.5249291) (-1.1933333,1.5240281) (-1.1866667,1.5231319) (-1.1800000,1.5222405) (-1.1733333,1.5213539) (-1.1666667,1.5204721) (-1.1600000,1.5195951) (-1.1533333,1.5187230) (-1.1466667,1.5178556) (-1.1400000,1.5169931) (-1.1333333,1.5161354) (-1.1266667,1.5152825) (-1.1200000,1.5144345) (-1.1133333,1.5135913) (-1.1066667,1.5127529) (-1.1000000,1.5119194) (-1.0933333,1.5110908) (-1.0866667,1.5102669) (-1.0800000,1.5094480) (-1.0733333,1.5086338) (-1.0666667,1.5078246) (-1.0600000,1.5070202) (-1.0533333,1.5062207) (-1.0466667,1.5054260) (-1.0400000,1.5046362) (-1.0333333,1.5038513) (-1.0266667,1.5030713) (-1.0200000,1.5022961) (-1.0133333,1.5015259) (-1.0066667,1.5007605) (-1.0000000,1.5000000) (-0.9933333,1.4992444) (-0.9866667,1.4984937) (-0.9800000,1.4977479) (-0.9733333,1.4970070) (-0.9666667,1.4962711) (-0.9600000,1.4955400) (-0.9533333,1.4948138) (-0.9466667,1.4940926) (-0.9400000,1.4933763) (-0.9333333,1.4926649) (-0.9266667,1.4919584) (-0.9200000,1.4912568) (-0.9133333,1.4905602) (-0.9066667,1.4898685) (-0.9000000,1.4891818) (-0.8933333,1.4885000) (-0.8866667,1.4878231) (-0.8800000,1.4871512) (-0.8733333,1.4864842) (-0.8666667,1.4858222) (-0.8600000,1.4851651) (-0.8533333,1.4845130) (-0.8466667,1.4838659) (-0.8400000,1.4832237) (-0.8333333,1.4825864) (-0.8266667,1.4819542) (-0.8200000,1.4813269) (-0.8133333,1.4807045) (-0.8066667,1.4800872) (-0.8000000,1.4794748) (-0.7933333,1.4788674) (-0.7866667,1.4782650) (-0.7800000,1.4776676) (-0.7733333,1.4770751) (-0.7666667,1.4764877) (-0.7600000,1.4759052) (-0.7533333,1.4753278) (-0.7466667,1.4747553) (-0.7400000,1.4741878) (-0.7333333,1.4736254) (-0.7266667,1.4730679) (-0.7200000,1.4725154) (-0.7133333,1.4719680) (-0.7066667,1.4714255) (-0.7000000,1.4708881) (-0.6933333,1.4703557) (-0.6866667,1.4698283) (-0.6800000,1.4693059) (-0.6733333,1.4687886) (-0.6666667,1.4682763) (-0.6600000,1.4677690) (-0.6533333,1.4672667) (-0.6466667,1.4667694) (-0.6400000,1.4662772) (-0.6333333,1.4657900) (-0.6266667,1.4653079) (-0.6200000,1.4648308) (-0.6133333,1.4643587) (-0.6066667,1.4638917) (-0.6000000,1.4634298) (-0.5933333,1.4629728) (-0.5866667,1.4625210) (-0.5800000,1.4620741) (-0.5733333,1.4616323) (-0.5666667,1.4611956) (-0.5600000,1.4607640) (-0.5533333,1.4603374) (-0.5466667,1.4599158) (-0.5400000,1.4594993) (-0.5333333,1.4590879) (-0.5266667,1.4586815) (-0.5200000,1.4582802) (-0.5133333,1.4578840) (-0.5066667,1.4574929) (-0.5000000,1.4571068) (-0.4933333,1.4567258) (-0.4866667,1.4563498) (-0.4800000,1.4559790) (-0.4733333,1.4556132) (-0.4666667,1.4552525) (-0.4600000,1.4548969) (-0.4533333,1.4545463) (-0.4466667,1.4542009) (-0.4400000,1.4538605) (-0.4333333,1.4535252) (-0.4266667,1.4531950) (-0.4200000,1.4528699) (-0.4133333,1.4525499) (-0.4066667,1.4522350) (-0.4000000,1.4519252) (-0.3933333,1.4516205) (-0.3866667,1.4513208) (-0.3800000,1.4510263) (-0.3733333,1.4507368) (-0.3666667,1.4504525) (-0.3600000,1.4501733) (-0.3533333,1.4498991) (-0.3466667,1.4496301) (-0.3400000,1.4493662) (-0.3333333,1.4491074) (-0.3266667,1.4488537) (-0.3200000,1.4486051) (-0.3133333,1.4483616) (-0.3066667,1.4481232) (-0.3000000,1.4478899) (-0.2933333,1.4476617) (-0.2866667,1.4474387) (-0.2800000,1.4472208) (-0.2733333,1.4470079) (-0.2666667,1.4468002) (-0.2600000,1.4465977) (-0.2533333,1.4464002) (-0.2466667,1.4462078) (-0.2400000,1.4460206) (-0.2333333,1.4458385) (-0.2266667,1.4456615) (-0.2200000,1.4454897) (-0.2133333,1.4453229) (-0.2066667,1.4451613) (-0.2000000,1.4450048) (-0.1933333,1.4448534) (-0.1866667,1.4447072) (-0.1800000,1.4445661) (-0.1733333,1.4444301) (-0.1666667,1.4442992) (-0.1600000,1.4441735) (-0.1533333,1.4440528) (-0.1466667,1.4439374) (-0.1400000,1.4438270) (-0.1333333,1.4437218) (-0.1266667,1.4436217) (-0.1200000,1.4435267) (-0.1133333,1.4434369) (-0.1066667,1.4433522) (-0.1000000,1.4432726) (-0.0933333,1.4431982) (-0.0866667,1.4431289) (-0.0800000,1.4430647) (-0.0733333,1.4430057) (-0.0666667,1.4429518) (-0.0600000,1.4429030) (-0.0533333,1.4428593) (-0.0466667,1.4428208) (-0.0400000,1.4427875) (-0.0333333,1.4427592) (-0.0266667,1.4427361) (-0.0200000,1.4427181) (-0.0133333,1.4427053) (-0.0066667,1.4426976)};
\draw[blue,thick,solid] plot coordinates {(0.0066667,1.4426976) (0.0133333,1.4427053) (0.0200000,1.4427181) (0.0266667,1.4427361) (0.0333333,1.4427592) (0.0400000,1.4427875) (0.0466667,1.4428208) (0.0533333,1.4428593) (0.0600000,1.4429030) (0.0666667,1.4429518) (0.0733333,1.4430057) (0.0800000,1.4430647) (0.0866667,1.4431289) (0.0933333,1.4431982) (0.1000000,1.4432726) (0.1066667,1.4433522) (0.1133333,1.4434369) (0.1200000,1.4435267) (0.1266667,1.4436217) (0.1333333,1.4437218) (0.1400000,1.4438270) (0.1466667,1.4439374) (0.1533333,1.4440528) (0.1600000,1.4441735) (0.1666667,1.4442992) (0.1733333,1.4444301) (0.1800000,1.4445661) (0.1866667,1.4447072) (0.1933333,1.4448534) (0.2000000,1.4450048) (0.2066667,1.4451613) (0.2133333,1.4453229) (0.2200000,1.4454897) (0.2266667,1.4456615) (0.2333333,1.4458385) (0.2400000,1.4460206) (0.2466667,1.4462078) (0.2533333,1.4464002) (0.2600000,1.4465977) (0.2666667,1.4468002) (0.2733333,1.4470079) (0.2800000,1.4472208) (0.2866667,1.4474387) (0.2933333,1.4476617) (0.3000000,1.4478899) (0.3066667,1.4481232) (0.3133333,1.4483616) (0.3200000,1.4486051) (0.3266667,1.4488537) (0.3333333,1.4491074) (0.3400000,1.4493662) (0.3466667,1.4496301) (0.3533333,1.4498991) (0.3600000,1.4501733) (0.3666667,1.4504525) (0.3733333,1.4507368) (0.3800000,1.4510263) (0.3866667,1.4513208) (0.3933333,1.4516205) (0.4000000,1.4519252) (0.4066667,1.4522350) (0.4133333,1.4525499) (0.4200000,1.4528699) (0.4266667,1.4531950) (0.4333333,1.4535252) (0.4400000,1.4538605) (0.4466667,1.4542009) (0.4533333,1.4545463) (0.4600000,1.4548969) (0.4666667,1.4552525) (0.4733333,1.4556132) (0.4800000,1.4559790) (0.4866667,1.4563498) (0.4933333,1.4567258) (0.5000000,1.4571068) (0.5066667,1.4574929) (0.5133333,1.4578840) (0.5200000,1.4582802) (0.5266667,1.4586815) (0.5333333,1.4590879) (0.5400000,1.4594993) (0.5466667,1.4599158) (0.5533333,1.4603374) (0.5600000,1.4607640) (0.5666667,1.4611956) (0.5733333,1.4616323) (0.5800000,1.4620741) (0.5866667,1.4625210) (0.5933333,1.4629728) (0.6000000,1.4634298) (0.6066667,1.4638917) (0.6133333,1.4643587) (0.6200000,1.4648308) (0.6266667,1.4653079) (0.6333333,1.4657900) (0.6400000,1.4662772) (0.6466667,1.4667694) (0.6533333,1.4672667) (0.6600000,1.4677690) (0.6666667,1.4682763) (0.6733333,1.4687886) (0.6800000,1.4693059) (0.6866667,1.4698283) (0.6933333,1.4703557) (0.7000000,1.4708881) (0.7066667,1.4714255) (0.7133333,1.4719680) (0.7200000,1.4725154) (0.7266667,1.4730679) (0.7333333,1.4736254) (0.7400000,1.4741878) (0.7466667,1.4747553) (0.7533333,1.4753278) (0.7600000,1.4759052) (0.7666667,1.4764877) (0.7733333,1.4770751) (0.7800000,1.4776676) (0.7866667,1.4782650) (0.7933333,1.4788674) (0.8000000,1.4794748) (0.8066667,1.4800872) (0.8133333,1.4807045) (0.8200000,1.4813269) (0.8266667,1.4819542) (0.8333333,1.4825864) (0.8400000,1.4832237) (0.8466667,1.4838659) (0.8533333,1.4845130) (0.8600000,1.4851651) (0.8666667,1.4858222) (0.8733333,1.4864842) (0.8800000,1.4871512) (0.8866667,1.4878231) (0.8933333,1.4885000) (0.9000000,1.4891818) (0.9066667,1.4898685) (0.9133333,1.4905602) (0.9200000,1.4912568) (0.9266667,1.4919584) (0.9333333,1.4926649) (0.9400000,1.4933763) (0.9466667,1.4940926) (0.9533333,1.4948138) (0.9600000,1.4955400) (0.9666667,1.4962711) (0.9733333,1.4970070) (0.9800000,1.4977479) (0.9866667,1.4984937) (0.9933333,1.4992444) (1.0000000,1.5000000) (1.0066667,1.5007605) (1.0133333,1.5015259) (1.0200000,1.5022961) (1.0266667,1.5030713) (1.0333333,1.5038513) (1.0400000,1.5046362) (1.0466667,1.5054260) (1.0533333,1.5062207) (1.0600000,1.5070202) (1.0666667,1.5078246) (1.0733333,1.5086338) (1.0800000,1.5094480) (1.0866667,1.5102669) (1.0933333,1.5110908) (1.1000000,1.5119194) (1.1066667,1.5127529) (1.1133333,1.5135913) (1.1200000,1.5144345) (1.1266667,1.5152825) (1.1333333,1.5161354) (1.1400000,1.5169931) (1.1466667,1.5178556) (1.1533333,1.5187230) (1.1600000,1.5195951) (1.1666667,1.5204721) (1.1733333,1.5213539) (1.1800000,1.5222405) (1.1866667,1.5231319) (1.1933333,1.5240281) (1.2000000,1.5249291) (1.2066667,1.5258349) (1.2133333,1.5267455) (1.2200000,1.5276609) (1.2266667,1.5285810) (1.2333333,1.5295060) (1.2400000,1.5304357) (1.2466667,1.5313701) (1.2533333,1.5323094) (1.2600000,1.5332534) (1.2666667,1.5342022) (1.2733333,1.5351557) (1.2800000,1.5361140) (1.2866667,1.5370770) (1.2933333,1.5380447) (1.3000000,1.5390173) (1.3066667,1.5399945) (1.3133333,1.5409765) (1.3200000,1.5419632) (1.3266667,1.5429546) (1.3333333,1.5439508) (1.3400000,1.5449517) (1.3466667,1.5459572) (1.3533333,1.5469675) (1.3600000,1.5479825) (1.3666667,1.5490022) (1.3733333,1.5500266) (1.3800000,1.5510557) (1.3866667,1.5520895) (1.3933333,1.5531280) (1.4000000,1.5541711) (1.4066667,1.5552190) (1.4133333,1.5562715) (1.4200000,1.5573286) (1.4266667,1.5583905) (1.4333333,1.5594569) (1.4400000,1.5605281) (1.4466667,1.5616039) (1.4533333,1.5626843) (1.4600000,1.5637694) (1.4666667,1.5648592) (1.4733333,1.5659535) (1.4800000,1.5670525) (1.4866667,1.5681561) (1.4933333,1.5692644) (1.5000000,1.5703772) (1.5066667,1.5714947) (1.5133333,1.5726168) (1.5200000,1.5737435) (1.5266667,1.5748748) (1.5333333,1.5760107) (1.5400000,1.5771511) (1.5466667,1.5782962) (1.5533333,1.5794458) (1.5600000,1.5806001) (1.5666667,1.5817589) (1.5733333,1.5829222) (1.5800000,1.5840902) (1.5866667,1.5852627) (1.5933333,1.5864397) (1.6000000,1.5876213) (1.6066667,1.5888074) (1.6133333,1.5899981) (1.6200000,1.5911934) (1.6266667,1.5923931) (1.6333333,1.5935974) (1.6400000,1.5948062) (1.6466667,1.5960195) (1.6533333,1.5972374) (1.6600000,1.5984597) (1.6666667,1.5996866) (1.6733333,1.6009179) (1.6800000,1.6021538) (1.6866667,1.6033941) (1.6933333,1.6046390) (1.7000000,1.6058883) (1.7066667,1.6071421) (1.7133333,1.6084003) (1.7200000,1.6096631) (1.7266667,1.6109303) (1.7333333,1.6122019) (1.7400000,1.6134780) (1.7466667,1.6147586) (1.7533333,1.6160435) (1.7600000,1.6173330) (1.7666667,1.6186268) (1.7733333,1.6199251) (1.7800000,1.6212279) (1.7866667,1.6225350) (1.7933333,1.6238465) (1.8000000,1.6251625) (1.8066667,1.6264829) (1.8133333,1.6278076) (1.8200000,1.6291368) (1.8266667,1.6304703) (1.8333333,1.6318083) (1.8400000,1.6331506) (1.8466667,1.6344973) (1.8533333,1.6358483) (1.8600000,1.6372037) (1.8666667,1.6385635) (1.8733333,1.6399276) (1.8800000,1.6412961) (1.8866667,1.6426689) (1.8933333,1.6440461) (1.9000000,1.6454276) (1.9066667,1.6468134) (1.9133333,1.6482036) (1.9200000,1.6495981) (1.9266667,1.6509968) (1.9333333,1.6523999) (1.9400000,1.6538073) (1.9466667,1.6552190) (1.9533333,1.6566350) (1.9600000,1.6580553) (1.9666667,1.6594798) (1.9733333,1.6609087) (1.9800000,1.6623418) (1.9866667,1.6637791) (1.9933333,1.6652208) (2.0000000,1.6666667) (2.0066667,1.6681168) (2.0133333,1.6695712) (2.0200000,1.6710298) (2.0266667,1.6724927) (2.0333333,1.6739598) (2.0400000,1.6754311) (2.0466667,1.6769066) (2.0533333,1.6783864) (2.0600000,1.6798703) (2.0666667,1.6813585) (2.0733333,1.6828509) (2.0800000,1.6843474) (2.0866667,1.6858482) (2.0933333,1.6873531) (2.1000000,1.6888622) (2.1066667,1.6903755) (2.1133333,1.6918929) (2.1200000,1.6934145) (2.1266667,1.6949403) (2.1333333,1.6964702) (2.1400000,1.6980042) (2.1466667,1.6995424) (2.1533333,1.7010847) (2.1600000,1.7026312) (2.1666667,1.7041817) (2.1733333,1.7057364) (2.1800000,1.7072952) (2.1866667,1.7088581) (2.1933333,1.7104251) (2.2000000,1.7119962) (2.2066667,1.7135714) (2.2133333,1.7151507) (2.2200000,1.7167340) (2.2266667,1.7183215) (2.2333333,1.7199129) (2.2400000,1.7215085) (2.2466667,1.7231081) (2.2533333,1.7247118) (2.2600000,1.7263195) (2.2666667,1.7279312) (2.2733333,1.7295470) (2.2800000,1.7311668) (2.2866667,1.7327906) (2.2933333,1.7344185) (2.3000000,1.7360503) (2.3066667,1.7376862) (2.3133333,1.7393260) (2.3200000,1.7409699) (2.3266667,1.7426177) (2.3333333,1.7442696) (2.3400000,1.7459254) (2.3466667,1.7475852) (2.3533333,1.7492489) (2.3600000,1.7509166) (2.3666667,1.7525883) (2.3733333,1.7542639) (2.3800000,1.7559434) (2.3866667,1.7576269) (2.3933333,1.7593143) (2.4000000,1.7610057) (2.4066667,1.7627009) (2.4133333,1.7644001) (2.4200000,1.7661032) (2.4266667,1.7678102) (2.4333333,1.7695211) (2.4400000,1.7712359) (2.4466667,1.7729545) (2.4533333,1.7746771) (2.4600000,1.7764035) (2.4666667,1.7781338) (2.4733333,1.7798680) (2.4800000,1.7816060) (2.4866667,1.7833478) (2.4933333,1.7850935) (2.5000000,1.7868431) (2.5066667,1.7885965) (2.5133333,1.7903537) (2.5200000,1.7921147) (2.5266667,1.7938796) (2.5333333,1.7956482) (2.5400000,1.7974207) (2.5466667,1.7991969) (2.5533333,1.8009770) (2.5600000,1.8027608) (2.5666667,1.8045485) (2.5733333,1.8063399) (2.5800000,1.8081350) (2.5866667,1.8099340) (2.5933333,1.8117367) (2.6000000,1.8135431) (2.6066667,1.8153533) (2.6133333,1.8171672) (2.6200000,1.8189849) (2.6266667,1.8208063) (2.6333333,1.8226314) (2.6400000,1.8244602) (2.6466667,1.8262927) (2.6533333,1.8281290) (2.6600000,1.8299689) (2.6666667,1.8318126) (2.6733333,1.8336599) (2.6800000,1.8355109) (2.6866667,1.8373656) (2.6933333,1.8392239) (2.7000000,1.8410860) (2.7066667,1.8429516) (2.7133333,1.8448210) (2.7200000,1.8466939) (2.7266667,1.8485705) (2.7333333,1.8504508) (2.7400000,1.8523346) (2.7466667,1.8542221) (2.7533333,1.8561132) (2.7600000,1.8580080) (2.7666667,1.8599063) (2.7733333,1.8618082) (2.7800000,1.8637137) (2.7866667,1.8656228) (2.7933333,1.8675355) (2.8000000,1.8694517) (2.8066667,1.8713715) (2.8133333,1.8732949) (2.8200000,1.8752219) (2.8266667,1.8771524) (2.8333333,1.8790864) (2.8400000,1.8810240) (2.8466667,1.8829651) (2.8533333,1.8849097) (2.8600000,1.8868578) (2.8666667,1.8888095) (2.8733333,1.8907647) (2.8800000,1.8927233) (2.8866667,1.8946855) (2.8933333,1.8966512) (2.9000000,1.8986203) (2.9066667,1.9005929) (2.9133333,1.9025690) (2.9200000,1.9045486) (2.9266667,1.9065316) (2.9333333,1.9085181) (2.9400000,1.9105080) (2.9466667,1.9125013) (2.9533333,1.9144981) (2.9600000,1.9164984) (2.9666667,1.9185020) (2.9733333,1.9205091) (2.9800000,1.9225196) (2.9866667,1.9245335) (2.9933333,1.9265507) (3.0000000,1.9285714) (3.0066667,1.9305955) (3.0133333,1.9326229) (3.0200000,1.9346538) (3.0266667,1.9366880) (3.0333333,1.9387255) (3.0400000,1.9407665) (3.0466667,1.9428107) (3.0533333,1.9448583) (3.0600000,1.9469093) (3.0666667,1.9489636) (3.0733333,1.9510212) (3.0800000,1.9530821) (3.0866667,1.9551464) (3.0933333,1.9572139) (3.1000000,1.9592848) (3.1066667,1.9613590) (3.1133333,1.9634364) (3.1200000,1.9655172) (3.1266667,1.9676012) (3.1333333,1.9696885) (3.1400000,1.9717790) (3.1466667,1.9738729) (3.1533333,1.9759699) (3.1600000,1.9780703) (3.1666667,1.9801738) (3.1733333,1.9822806) (3.1800000,1.9843907) (3.1866667,1.9865039) (3.1933333,1.9886204) (3.2000000,1.9907401) (3.2066667,1.9928630) (3.2133333,1.9949891) (3.2200000,1.9971184) (3.2266667,1.9992509) (3.2333333,2.0013865) (3.2400000,2.0035254) (3.2466667,2.0056674) (3.2533333,2.0078126) (3.2600000,2.0099609) (3.2666667,2.0121124) (3.2733333,2.0142671) (3.2800000,2.0164248) (3.2866667,2.0185857) (3.2933333,2.0207498) (3.3000000,2.0229169) (3.3066667,2.0250872) (3.3133333,2.0272606) (3.3200000,2.0294371) (3.3266667,2.0316167) (3.3333333,2.0337994) (3.3400000,2.0359852) (3.3466667,2.0381740) (3.3533333,2.0403660) (3.3600000,2.0425610) (3.3666667,2.0447590) (3.3733333,2.0469601) (3.3800000,2.0491643) (3.3866667,2.0513715) (3.3933333,2.0535818) (3.4000000,2.0557950) (3.4066667,2.0580114) (3.4133333,2.0602307) (3.4200000,2.0624530) (3.4266667,2.0646784) (3.4333333,2.0669067) (3.4400000,2.0691381) (3.4466667,2.0713724) (3.4533333,2.0736098) (3.4600000,2.0758501) (3.4666667,2.0780934) (3.4733333,2.0803396) (3.4800000,2.0825888) (3.4866667,2.0848410) (3.4933333,2.0870961) (3.5000000,2.0893542) (3.5066667,2.0916152) (3.5133333,2.0938791) (3.5200000,2.0961459) (3.5266667,2.0984157) (3.5333333,2.1006884) (3.5400000,2.1029640) (3.5466667,2.1052425) (3.5533333,2.1075239) (3.5600000,2.1098082) (3.5666667,2.1120953) (3.5733333,2.1143854) (3.5800000,2.1166783) (3.5866667,2.1189741) (3.5933333,2.1212727) (3.6000000,2.1235742) (3.6066667,2.1258785) (3.6133333,2.1281857) (3.6200000,2.1304958) (3.6266667,2.1328086) (3.6333333,2.1351243) (3.6400000,2.1374428) (3.6466667,2.1397641) (3.6533333,2.1420882) (3.6600000,2.1444151) (3.6666667,2.1467449) (3.6733333,2.1490774) (3.6800000,2.1514126) (3.6866667,2.1537507) (3.6933333,2.1560916) (3.7000000,2.1584352) (3.7066667,2.1607815) (3.7133333,2.1631306) (3.7200000,2.1654825) (3.7266667,2.1678371) (3.7333333,2.1701945) (3.7400000,2.1725546) (3.7466667,2.1749174) (3.7533333,2.1772829) (3.7600000,2.1796511) (3.7666667,2.1820221) (3.7733333,2.1843957) (3.7800000,2.1867721) (3.7866667,2.1891511) (3.7933333,2.1915328) (3.8000000,2.1939173) (3.8066667,2.1963043) (3.8133333,2.1986941) (3.8200000,2.2010865) (3.8266667,2.2034815) (3.8333333,2.2058793) (3.8400000,2.2082796) (3.8466667,2.2106826) (3.8533333,2.2130882) (3.8600000,2.2154965) (3.8666667,2.2179074) (3.8733333,2.2203209) (3.8800000,2.2227370) (3.8866667,2.2251557) (3.8933333,2.2275770) (3.9000000,2.2300009) (3.9066667,2.2324274) (3.9133333,2.2348564) (3.9200000,2.2372881) (3.9266667,2.2397223) (3.9333333,2.2421591) (3.9400000,2.2445984) (3.9466667,2.2470403) (3.9533333,2.2494847) (3.9600000,2.2519317) (3.9666667,2.2543812) (3.9733333,2.2568333) (3.9800000,2.2592879) (3.9866667,2.2617449) (3.9933333,2.2642046) (4.0000000,2.2666667)};
\draw[violet!75!black,densely dotted] (-2,1.6666667)--(2,1.6666667);
\fill[violet!75!black] (-2,1.6666667) circle(2.3pt);
\fill[violet!75!black] (2,1.6666667) circle(2.3pt);
\draw[blue,fill=white,thick] (0,1.44269504) circle(2.7pt);
\end{scope}
\end{scope}
\node[anchor=west,align=left,text width=4.35cm] at (8.9,4.4) {没有可变底数参数。\\[14pt]$f(-2)=f(2)=\frac53$。\\[14pt]左右同高，是偶函数；\\中间空点不能补入。};
\node[text width=12.8cm,align=center] at (7.1,.65) {$D=\mathbb R\setminus\{0\}$；括号内是奇函数，乘上奇函数 $x$ 后为偶函数。};
\end{tikzpicture}\end{center}

\textbf{先看图：} 先把括号通分：$H(x)=\dfrac{2^x+1}{2(2^x-1)}$。由奇3可知 $H$ 在 $x\ne0$ 上为奇函数；再乘上奇函数 $x$，两个负号抵消。

\textbf{再说明一般情形：}
\[
f(-x)=(-x)H(-x)=(-x)(-H(x))=f(x).
\]
图中 $f(-2)=f(2)=\frac53$。中间空点表示 $0$ 不在定义域，不能因为外面乘了 $x$ 就补为零。
\end{explanationbox}
